\documentclass[book]{magnolia}

\magtex{tex_driver={pdftex},
        tex_packages={float,caption,titling,epigraph,minitoc,slashbox,tabularx,cancel,pgfplots,xypic,nicefrac},
        tex_pstricks={pstricks,pst-plot}}
\magfiche{document_nom={Cours de Sup},
          auteur_nom={François Fayard},
          auteur_mail={fayard.prof@gmail.com}}
\magcours{cours_matiere={maths},
          cours_niveau={mpsi},
          cours_chapitre_numero={11},
          cours_chapitre={Dérivation}}
% \magmisenpage{misenpage_presentation={tikzvelvia},
%           misenpage_format={presentation},
%           misenpage_nbcolonnes={1},
%           misenpage_preuve={non},
%           misenpage_sol={non}}
\magmisenpage{}
\maglieudiff{}
\magprocess

\usepackage[scale=3]{ccicons}
\title{{\Huge\bf Cours de Mathématiques}\\\vspace{1cm}
%       \textbf{\Huge 2022--2023}\\
       \textbf{\Huge Maths Sup -- Tome III}\\\vspace{1cm}
      %  \textsc{L. Bouyge-Bono, F. Fayard, V. Lambert}\\\vspace{1cm}
       \textsc{F. Fayard}\\\vspace{1cm}
       \includegraphics[width=8cm]{../../Commun/Images/lazos-bde-2024.png}\\%\vspace{1cm}
       \includegraphics[width=8cm]{../../Commun/Images/lazos.png}}
% \author{{\sc François Fayard}, {\sc Victor Lambert}}
% \date{3 Août 2020 -- 1a20aa}

% \pretitle{%
%   \begin{center}
%   \includegraphics[width=12cm]{/Users/fayard/Desktop/lazos.png}\\[\bigskipamount]
% }
% \posttitle{\end{center}}

\mtcsettitle{minitoc}{}
\mtcsettitle{secttoc}{Table des matières}
\mtcsettitle{parttoc}{Table des matières}

\usetikzlibrary{positioning}
% \lstset{%
%   frame            = tb,    % draw frame at top and bottom of code block
%   tabsize          = 1,     % tab space width
%   numbers          = none,  % display line numbers on the left
%   framesep         = 3pt,   % expand outward
%   framerule        = 0.4pt, % expand outward 
%   commentstyle     = \color{green},      % comment color
%   keywordstyle     = \color{blue},       % keyword color
%   stringstyle      = \color{blue},    % string color
%   backgroundcolor  = \color{colorLazoBlue1Light}, % backgroundcolor color
%   showstringspaces = false,              % do not mark spaces in strings
%   basicstyle       = \ttfamily,
%   breaklines       = true
% }

\input{vc.tex}

\begin{document}

\maketitle

La version de ce document est la \textsc{\GITAbrHash}.\\

Merci à tous les élèves des lycées Janson de Sailly, du Parc et des Lazaristes pour leurs remarques et corrections. Je remercie
particulièrement \nom{Samuel Auroy}, \nom{Antonin Barbier}, \nom{Amin Belfkira}, \nom{Martin Bot}, \nom{Alexandre Brousse}, \nom{Élodie Brun}, \nom{Damien Callendrier}, \nom{Lauren Calvosa}, \nom{Sylvain Crosnier}, \nom{Enguerrand De Jaegere}, \nom{Thibaud De Valicourt}, \nom{Victor Déru}, \nom{Raphaël Des Boscs}, \nom{Grégoire Dhimoïla}, \nom{Léo Duhamel-Callot}, \nom{Sacha Evrard}, \nom{Axel Faou}, \nom{Titouan Francheteau}, \nom{Hélène Ghaleb}, \nom{Cédric Holocher},
\nom{Maxime Joubert}, \nom{Maxime Lombard}, \nom{Mira Maamri}, \nom{Gauthier Malandrin}, \nom{Pierre-Antoine Nguyen},
\nom{Hilaire Oudinot}, \nom{Eliott Pradeleix}, \nom{Yann-Ellie Ravon}, \nom{Sixtine Reynaud}, \nom{Vivien Thienot},
\nom{Carole Vacherand}, \nom{Camille Vialet}, \nom{Paul Vilars} et \nom{Antonin Villepontoux}.\\

Je tiens enfin à remercier mes anciens professeurs et
collègues qui ont eu une influence sur la rédaction de ce document~: \nom{Walter Appel}, \nom{Bruno Arsac}, \nom{Jean-Pierre Barani}, \nom{Vincent Bayle}, \nom{Christophe Bertault}, \nom{Laurence Bouyge}, \nom{Gilles Chaffard}, \nom{Alain Chillès}, \nom{Denis Choimet}, \nom{Vincent Clapiès}, \nom{Gérard Esposito}, \nom{Stéphane Gonnord}, \nom{Victor Lambert}, %
%\nom{Pierre-Louis Lions},
\nom{Frédéric Morlot}, \nom{Franz Ridde}, \nom{Emmanuel Roblet} et \nom{Alain Troesch}.\\
%, \nom{Cédric Villani}.\\
\vfill


\begin{center}
  \ccbysa\\
  \vspace{2ex}
  This work is licensed under a Creative Commons\\
  Attribution-ShareAlike 4.0 International License.\\
  \url{https://creativecommons.org/licenses/by-sa/4.0/legalcode.fr}\\
  \vspace{2ex}
  La dernière version de ce document ainsi que\\
  les sources \LaTeX{} sont disponibles à l'adresse\\
  \url{https://github.com/FayardProf/Maths-MPSI-MP2I}
  \end{center}
  \vspace{2ex}
  \begin{center}
  \textbf{Vous êtes autorisés à~:}
  \end{center}
  \vspace{2ex}
  \begin{itemize}
  \item \textbf{Partager}~: copier, distribuer et communiquer le matériel par tous les moyens et sous tous formats.
  \item \textbf{Adapter}~: remixer, transformer et créer à partir du matériel
  pour toute utilisation, y compris commerciale.
  \end{itemize}
  \vspace{2ex}
  \begin{center}
  \textbf{Selon les conditions suivantes~:}
  \end{center}
  \vspace{2ex}
  \begin{itemize}
    \item \textbf{Attribution}~: Vous devez créditer l'œuvre, intégrer un lien vers la licence et indiquer si des modifications ont été effectuées à l'œuvre. Vous devez indiquer ces informations par tous les moyens raisonnables, sans toutefois suggérer que l'offrant vous soutient ou soutient la façon dont vous avez utilisé son œuvre.
    \item \textbf{Partage dans les mêmes conditions}~: Dans le cas où vous effectuez un remix, que vous transformez, ou créez à partir du matériel composant l'œuvre originale, vous devez diffuser l'œuvre modifiée dans les même conditions, c'est à dire avec la même licence avec laquelle l'œuvre originale a été diffusée.
    \item \textbf{Pas de restrictions complémentaires}~: Vous n'êtes pas autorisé à appliquer des conditions légales ou des mesures techniques qui restreindraient légalement autrui à utiliser l'œuvre dans les conditions décrites par la licence.
    \end{itemize}

\setcounter{chapter}{10}

\tableofcontents

% \part{Sup}

\chapter{Groupes}
\setcounter{numeroexercicecours}{1}
\documentclass{magnolia}

\magtex{tex_driver={pdftex},
        tex_packages={xypic}}
\magfiche{document_nom={Cours sur les structures algebriques},
          auteur_nom={François Fayard},
          auteur_mail={fayard.prof@gmail.com}}
\magcours{cours_matiere={maths},
          cours_niveau={mpsi},
          cours_chapitre_numero={12},
          cours_chapitre={Groupes}}
\magmisenpage{}
\maglieudiff{}
\magprocess


\begin{document}

%BEGIN_BOOK
\magtoc

\section{Groupe}
\subsection{Loi de composition interne}

\begin{definition}
Soit $E$ un ensemble. On appelle loi de \emph{composition interne} toute
application $\star$ de $E\times E$ dans $E$.
\[\dspappli{\star}{E\times E}{E}{\p{x,y}}{x\star y}\]
\end{definition}

\begin{definition}
La loi $\star$ est dite
\begin{itemize}
\item \emph{associative} lorsque
  \[\forall x,y,z\in E \qsep \p{x\star y}\star z=x\star\p{y\star z}.\]
\item \emph{commutative} lorsque
  \[\forall x,y\in E \qsep x\star y=y\star x.\]
\end{itemize}
\end{definition}

\begin{exemples}
\exemple L'addition et la multiplication sont des lois de composition interne sur $\Z$,
  associatives et commutatives.
\exemple L'exponentiation est une loi de composition interne sur $\N$ qui n'est
  ni associative, ni commutative.
\exemple Si $X$ est un ensemble, la composition est une loi de composition interne
  associative sur $E\defeq\mathcal{F}(X,X)$. Elle n'est pas commutative dès que $X$ possède
  au moins deux éléments.
\exemple Le produit matriciel est une loi de composition interne associative sur
  $\mat{n}{\C}$. Elle est n'est pas commutative dès que $n\geq 2$.
\end{exemples}

\begin{remarques}
\remarque Soit $\star$ une loi \emph{associative}.
  Quels que soient $x, y, z, t\in E$, les 5 expressions
  suivantes sont égales~:
  \[(x\star y)\star(z\star t)\qsep ((x\star y)\star z)\star t,\]
  \[(x\star(y\star z))\star t\qsep
    x\star((y\star z)\star t)\qsep
    x\star(y\star(z\star t)).\]
  On admettra plus généralement que
  toute expression de $n$ éléments construite à l'aide de la loi $\star$ ne dépend
  pas de l'emplacement des parenthèses. C'est pourquoi on se permettra de les omettre.
\remarque On dit que deux éléments $x,y\in E$ \emph{commutent} lorsque $x\star y=y\star x$.
\end{remarques}
\begin{sol}
$\dspappli{\star}{\N\times\N}{\N}{(a,b)}{a^b}$ est une lci. Elle n'est pas associative :
$$(0\star 1)\star 0= 0^1 \star 0=0\star 0=0^0=1$$
$$0\star (1\star 0)=0\star 1^0=0\star 1=0^1=0$$
\end{sol}

\begin{definition}
Une partie $A$ de $E$ est dite \emph{stable} par $\star$ lorsque
\[\forall x,y\in A \qsep x\star y\in A.\]
\end{definition}

\begin{remarqueUnique}
\remarque Si $\star$ est une loi de composition interne sur $E$ et
  $A\in\parties{E}$ est stable par $\star$, alors la loi
  \[\dspappli{\star_A}{A\times A}{A}{\p{x,y}}{x\star y}\]
  est une loi de composition interne sur $A$. On continuera à la noter
  $\star$.
\end{remarqueUnique}
\begin{sol}
Elle hérite des propriétés (associativité, commutativité) de la loi ambiante.
\end{sol}

\begin{definition}
On dit que $\star$ admet un \emph{élément neutre} $e\in E$ lorsque
\[\forall x\in E \qsep x\star e=x \et e \star x=x.\]
Si tel est le cas, il est unique et on l'appelle \emph{élément neutre} de $\star$.
Lorsque la loi est notée additivement, l'élément neutre est noté $0$.
\end{definition}

\begin{preuve}
Preuve de l'unicité : $e_2=e_1\star e_2=e_1$.
\end{preuve}

\begin{exoUnique}
\exo Parmi les lois de composition interne citées plus haut, lesquelles
  admettent un élément neutre~?
  \begin{sol}
  $0$ est l'élément neutre de $(\C,+)$, $1$ de $(\C,\times)$.\\
  $0$ est l'élément neutre de $(\N,\wedge)$. $(\Z,\wedge)$, n'admet pas d'élément neutre. Sinon, $-1\wedge e=-1$. Pour la composition, c'est l'identité.
  \end{sol}
\end{exoUnique}

\medskip
Dans toute la suite de ce cours, on supposera, sauf mention explicite
du contraire, que les lois sont associatives et admettent un élément neutre.
\medskip

\begin{definition}
Soit $x\in E$. On définit par récurrence $x^n$ pour tout $n\in\N$ en posant~:
\begin{itemize}
\item $x^0\defeq e$
\item $\forall n\in\N \qsep x^{n+1}\defeq x^n\star x.$
\end{itemize}
\end{definition}

\begin{remarqueUnique}
\remarque Lorsque la loi est notée additivement, on n'utilise pas la notation $x^n$ mais plutôt la notation $n\cdot x$. On a donc~:
\begin{itemize}
\item $0\cdot x=0$
\item $\forall n\in\N \qsep \p{n+1}\cdot x=n\cdot x+x.$
\end{itemize}
\end{remarqueUnique}

\begin{proposition}
\begin{itemize}
\item Soit $x\in E$. Alors
\begin{eqnarray*}
\forall m,n\in\N, & & x^{m+n}=x^m\star x^n\\
                 & & \p{x^m}^n=x^{mn}.
\end{eqnarray*}
\item Soit $x,y\in E$ tels que $x\star y=y\star x$. Alors, pour tout $n,m\in\N$, $x^n$ et $y^m$ commutent. De plus
\[\forall n\in\N \qsep \p{x\star y}^n=x^n \star y^n.\]
\end{itemize}
\end{proposition}

\begin{preuve}
Pour les deux premiers points, on fixe $m$ et on "récurre" sur $n\in \N$.
Pour l'autre récurrence + commutativité.
\end{preuve}

\begin{remarqueUnique}
\remarque Si la loi est notée additivement, on a donc~:
  \begin{eqnarray*}
  \forall x\in E\qsep \forall m,n\in\N, & & (m+n)\cdot x=m\cdot x+n\cdot x\\
                                        & & n\cdot(m\cdot x)=(nm)\cdot x\\
  \forall x,y\in E\qsep \forall n\in\N, & & n\cdot(x+y)=n\cdot x+n\cdot y.
  \end{eqnarray*}
% \remarque Pour calculer $x^8$, on peut commencer par calculer $x\star x$, puis
%   \og multiplier\fg  6 fois ce résultat par $x$. Cette méthode nécessite 7
%   multiplications. On peut cependant faire plus rapide et se limiter à 3
%   multiplications~: il suffit de calculer $x^2=x\star x$, de multiplier le
%   résultat obtenu par lui-même pour obtenir $x^4$ puis de répéter l'opération pour obtenir $x^8$. Cette méthode se généralise facilement au calcul de $x^n$ lorsque $n$ est une puissance de 2. Mais on peut aussi calculer $x^n$ pour tout $n\in\N$ grâce à
%   l'algorithme récursif suivant.
%   \begin{itemize}
%   \item Si $n=0$, alors la réponse est $e$.
%   \item Sinon, deux cas se présentent.
%   \begin{itemize}
%   \item Si $n$ est pair, alors il existe $k\in\N$ tel que $n=2k$. On calcule
%     alors de manière récursive $x^k$, résultat qu'il suffit de multiplier par
%     lui-même pour obtenir $x^n$.
%   \item Si $n$ est impair, alors il existe $k\in\N$ tel que $n=2k+1$. On calcule
%     alors de manière récursive $x^k$, résultat qu'il suffit de multiplier par
%     lui-même, puis par $x$ pour obtenir $x^n$.
%   \end{itemize}
%   \end{itemize}
%   On peut montrer que cet algorithme, appelé \textit{méthode d'exponentiation rapide},
%   nécessite asymptotiquement ${\rm O}(\ln n)$ multiplications, contrairement à
%   l'algorithme naïf qui s'effectue asymptotiquement en ${\rm O}(n)$ multiplications. Voici une implémentation de cet algorithme en Python dans le cas où $x$ est un nombre et $\star$ est la multiplication usuelle.
% \begin{pythoncode}
% def puissance(x, n):
%     if n == 0:
%         return 1
%     else:
%         y = puissance(x, n // 2)
%         if n % 2 == 0:
%             return y * y
%         else:
%             return y * y * x
% \end{pythoncode}
\end{remarqueUnique}

%% Exponentiation rapide
%% Multiplication naïve : pour calculer x^8 il faut 7 multiplications
%% Exponentiation rapide : pour calculer x^8 il faut 3 multiplications
%% Plus généralement, pour calculer x^(2^a)
%% Méthode naïve : 2^a - 1 multiplications
%% Exponentiation rapide : a multiplications
%% Donc si n est une puissance de 2, il faut log_2 (n) multiplications
%% Ce procédé marche aussi pour les exposants qui ne sont pas des puissances
%% de 2. Pour calculer x^9, on peut utiliser 4 multiplications.
%%
%% > puiss:=proc(x,n::nonnegint)
%%     local q,r,p;
%%
%%     if n=0 then
%%       return(1);
%%     else
%%       q:=iquo(n,2);
%%       r:=irem(n,2);
%%       p:=puiss(x,q);
%%       p:=p*p;
%%       if r=0 then
%%         return(p);
%%       else
%%         return(x*p);
%%       end if;
%%     end if;
%%   end proc;
%%
%% On peut montrer que l'algorithme effectue asymptotiquement log_2 n
%% multiplications.

\begin{definition}
Soit $x\in E$. On dit que $x$ est \emph{symétrisable} pour la loi $\star$ lorsqu'il
existe $y\in E$ tel que
\[x\star y=y\star x=e.\]
Si tel est le cas, $y$ est unique et est appelé \emph{symétrique} de $x$. On l'appelle
\emph{inverse} de $x$ et on le note $x^{-1}$ lorsque la loi est notée multiplicativement.
On l'appelle \emph{opposé} de $x$ et on le note $-x$ lorsque la loi est notée
additivement.
\end{definition}

\begin{preuve}
On suppose qu'il existe $y_1$ et $y_2$ tels que $x\star y_1=y_1\star x=e=x\star y_2=y_2\star x$.
On a d'une part $(y_1\star x)\star y_2=e\star y_2=y_2$ et d'autre part $(y_1\star x)\star y_2=y_1\star (x\star y_2)=y_1\star e=y_1$ donc $y_1=y_2$.
\end{preuve}

\begin{sol}
Remarques : Si $f \in F(X,X)$ muni de la loi de composition, f est inversible si et seulement si elle est bijective. Si . n'est pas commutative, il faut vérifier les deux cotés. Par exemple $\dspappli{f}{\N}{\N}{n}{n+1}$ est inversible à gauche mais pas à droite.
\end{sol}

%% Exemples :
%% 1) L'élément neutre est toujours inversible et est égal à son propre inverse
%% 2) Si f \in F(X,X) muni de la loi de composition, f est inversible si
%%    et seulement si elle est bijective.
%%
%% Remarques:
%% 2) Si . n'est pas commutative, il faut vérifier les deux cotés
%%    Pas exemple f : N -> N , n -> n+1 est in versible à droite mais pas à
%%    gauche.
%% 3) x y^{-1} ne doit pas être noté x/y car on ne sait pas si cela signifie
%%    x y{-1} ou y^{-1} x.

\begin{proposition}
\begin{itemize}
\item Si $x$ est symétrisable, $x^{-1}$ l'est et
  \[\p{x^{-1}}^{-1}=x.\]
\item Si $x$ et $y$ sont symétrisables, $x\star y$ l'est et
  \[\p{x\star y}^{-1}=y^{-1}\star x^{-1}.\]  
\end{itemize}
\end{proposition}

\begin{definition}
Soit $x\in E$. Si $x$ est symétrisable, on étend la définition de $x^n$ en posant~:
\[\forall n\in\Z \qsep x^n\defeq
  \begin{cases}
  x^n & \text{si $n\geq 0$}\\
  \p{x^{-n}}^{-1} & \text{si $n< 0$}.
  \end{cases}\]
\end{definition}

% \begin{remarqueUnique}
% \remarque En particulier, si $x$ est symétrisable
%   \[\forall n\in\Z\qsep x^n\star x^{-n}=x^{-n}\star x^n=e.\]
% \end{remarqueUnique}

\begin{proposition}
\begin{itemize}
\item Soit $x\in E$. Si $x$ est symétrisable
\begin{eqnarray*}
\forall m,n\in\Z, & & x^{m+n}=x^m\star x^n\\
                 & & \p{x^m}^n=x^{mn}.
\end{eqnarray*}
\item Si $x,y\in E$ sont symétrisables et commutent, alors
\[\forall n\in\Z \qsep \p{x\star y}^n=x^n \star y^n.\]
\end{itemize}
\end{proposition}


\begin{preuve}
 1) On commence par remarquer que pour tout $n\in\Z$, $x^n$ et $x^{-n}$ sont inverses l'un de l'autre. Ensuite on fait une disjonction de cas :
\begin{itemize}
\item[$\bullet$] Si $m\geq 0$ et $n\geq 0$ : déjà vu
\item[$\bullet$] Si $m\geq 0$ et $n\leq 0$ :
\begin{itemize}
\item Si $m+n \geq 0$,
         $x^{m+n} x^{-n} = x^m$
\item Si $m+n \leq 0$,
         $x^{-(m+n)} x^m = x^{-n}$
\end{itemize}
\item[$\bullet$] Si $m\leq 0$ et $n\geq 0$ :
\begin{itemize}
\item Si $m+n \geq  0$,
         $x^{-m}x^{m+n} = x^n$
\item Si $m+n \leq 0$,
         $x^n x^{-{m+n}} = x^{-m}$
\end{itemize}
\item[$\bullet$] Si $m\leq 0$ et $n\leq 0$,
       $x^{-(n+m)}=x^{-n} x^{-m}$
\end{itemize}
 2) On commence par montrer que $(x^{-1})^n=(x^n)^{-1}$ pour tout $n\in\N$.
\begin{itemize}
\item[$\bullet$] Si $m$ et $n$ sont positifs, OK.
\item[$\bullet$] Si $m\leq 0$ et $n\geq 0$,
      $$(x^m)^n = ((x^{-m})^{-1})^n= ((x^{-1})^{-m})^n = ...$$
\item[$\bullet$] Si $m\geq 0$ et $n\leq 0$,
      $$(x^m)^n = ((x^m)^{-n})^{-1}= (x^{-nm})^{-1}= x^{nm} .$$
\item[$\bullet$] Si $m\leq 0$ et $n\leq 0$,
      $$(x^m)^n = ((x^m)^{-n})^{-1}= ((x^m)^{-1})^{-n}= (x^{-m})^{-n}= x^{mn}.$$
\end{itemize}

    \end{preuve}

%% Preuve
%% 1) - On commence par remarquer que pour tout n\in\Z x^n et x^{-n} sont
%%      inverses l'un de l'autre.
%%   - Ensuite on fait le cas :
%%     - Si m>=0 et n>=0 : déjà vu
%%     - Si m>=0 et n<=0
%%       - Si m+n >=0
%%         x^(m+n) x^{-n} = x^m
%%       - Si m+n <=0
%%         x^{-(m+n)} x^m = x^{-n}
%%     - Si m<=0 et n>=0
%%       - Si m+n >= 0
%%         x^{-m}x^{m+n} = x^n
%%       - Si m+n <=0
%%         x^n x^{-{m+n}} = x^{-m}
%%     - Si m<=0 et n<=0
%%       x^{-(n+m)}=x^{-n} x^{-m}
%% 2) - On commence par montrer que (x^(-1))^n=(x^n)^(-1) pour tout n\in\N
%%    - Si m et n sont positifs, OK
%%    - Si m<=0 et n>=0
%%      (x^m)^n = ((x^(-m))^(-1))^n
%%              = ((x^(-1))^(-m))^n
%%              = ...
%%    - Si m>=0 et n<=0
%%      (x^m)^n = ((x^m)^(-n))^(-1)
%%              = (x^{-nm})^{-1}
%%              = x^{nm} 
%%    - Si m<=0 et n<=0
%%      (x^m)^n = ((x^m)^{-n})^{-1}
%%              = ((x^m)^{-1})^{-n}
%%              = (x^{-m})^{-n}
%%              = x^{mn}
%%    - par récurrence sur n, puis par inverse (échanger x et y)
\begin{remarqueUnique}
\remarque Lorsque la loi est notée additivement, on a donc~:
  \begin{eqnarray*}
  \forall x\in E\qsep \forall m,n\in\Z, & & (m+n)\cdot x=m\cdot x+n\cdot x\\
                                        & & n\cdot(m\cdot x)=(nm)\cdot x\\
  \forall x,y\in E\qsep \forall n\in\Z, & & n\cdot(x+y)=n\cdot x+n\cdot y.
  \end{eqnarray*}
\end{remarqueUnique}


\begin{definition}
On dit qu'un élément $x$ de $E$ est \emph{régulier} lorsque
\begin{eqnarray*}
\forall y,z\in E \qsep x\star y=x\star z &\implique& y=z\\
                       y\star x=z\star x &\implique& y=z.
\end{eqnarray*}
\end{definition}

%% Exemples
%% 1) (\R,+) : tous les éléments sont réguliers.
%% 2) (\R,x) : tous les éléments non nuls sont réguliers.
%% 3) F(R,R),o : seules les bijections sont régulières.
%%    En effet, si f n'est pas injective, f o g_1 = f o g_2 mais g_1 \neq g_2
%ù              si f n'est pas surjective, g_1 o f = g_2 o f mais g_1 \neq g_2

% \begin{remarqueUnique}
% \remarque Si l'élément $x$ est régulier, alors, pour tout $y,z\in E$
% \[x\star y=x\star z \ssi y=z.\]
% De même
% \[y\star x=z\star x \ssi y=z.\]
% \end{remarqueUnique}

\begin{proposition}
Les éléments symétrisables sont réguliers.  
\end{proposition}

% \begin{remarqueUnique}
% \remarque Si l'élément $x$ est régulier, alors, pour tout $y,z\in E$
% \[x\star y=x\star z \ssi y=z.\]
% De même
% \[y\star x=z\star x \ssi y=z.\]
% \end{remarqueUnique}

\begin{sol} La réciproque est fausse. Par exemple dans $(\Z,\times)$, $1$ et $-1$ sont les seuls inversibles, les éléments réguliers sont les éléments non nuls.
\end{sol}
%% Remarque :
%% 1) La réciproque est fausse. Par exemples Z,x
%%    1 et -1 sont les seuls inversibles
%%    les éléments réguliers sont les éléments non nuls.


\subsection{Groupe}

\begin{definition}
Soit $G$ un ensemble muni d'une loi de composition interne $\star$. On  dit que
$\p{G,\star}$ est un \emph{groupe} lorsque
\begin{itemize}
\item $\star$ est associative
\item $\star$ admet un élément neutre
\item tout élément de $G$ est symétrisable.
\end{itemize}
Le groupe $\p{G,\star}$ est dit commutatif (ou \emph{abélien}) lorsque la loi $\star$
est commutative.
\end{definition}

\begin{remarques}
\remarque $(\C,+)$ et $(\Cs,\times)$ sont des groupes commutatifs.
\remarque Si $(G,\star)$ est un groupe et $a, b\in G$, alors
  \[\forall x\in G\qsep a\star x = b \quad\ssi\quad x=a^{-1}\star b.\]
  De même
  \[\forall x\in G\qsep x \star a = b \quad\ssi\quad x=b\star a^{-1}.\]
  % On peut ainsi résoudre facilement de nombreuses équations dans $G$.
% \remarque Si $(G,\star)$ est un groupe et $g\in G$, les applications
%   \[\dspappli{\tau_g}{G}{G}{x}{g\star x} \quad\et\quad
%     \dspappli{\sigma_g}{G}{G}{x}{x\star g}\]
%   sont des bijections de $G$ appelées respectivement translation à gauche
%   et à droite. De plus $\tau_g^{-1}=\tau_{g^{-1}}$ et $\sigma_g^{-1}=\sigma_{g^{-1}}$.
%   En particulier, la résolution d'équation du type $a\star x\star b = y$ se fait
%   simplement dans un groupe. En effet, si $a,b,y\in G$, on a
%   \begin{eqnarray*}
%   \forall x\in G\qsep
%   a\star x\star b = y
%   &\ssi& x\star b  = a^{-1} \star y\\
%   &\ssi& x  = a^{-1} \star y \star b^{-1}.
%   \end{eqnarray*}
%  \begin{sol}
%  On définit $\dspappli{\phi_g}{G}{G}{g}{x^{-1}\star g}$ et on montre que $\tau_g\circ\phi_g=Id=\phi_g\circ\tau_g$ pour voir qu'elles sont bijectives.
%  \end{sol}
\remarque Si $(G,\star)$ est un groupe fini, on appelle table de $(G,\star)$ le
  tableau à deux entrées dont les lignes et les colonnes sont 
  indexées par les éléments de $G$ et qui contient les produits $x\star y$.
  Puisque $(G,\star)$ est un groupe, chaque
  ligne et chaque colonne contient une et une seule fois chaque élément de
  $G$.
\end{remarques}

\begin{exoUnique}
\exo Montrer qu'il n'existe qu'une seule table de groupe à 3 éléments.
\end{exoUnique}

%% (Niels Abel : 1802-1829, mathématicien Norvégien)
%%
%% Exemples :
%% 1) (Z,+) (Q,+) (R,+) et (C,+) sont des groupes
%% 2) (Q^*,x) (R^*,x) (C^*,x) sont des groupes
%%
%% Remarques :
%% 1) Dans un groupe, si a\in G, les opérations
%%    g -> g a et g -> a g sont des bijections de G. On les appelle translation
%%    à droite et à gauche.
%%    Application : - détermination du groupe à 3 éléments e,a,b
%%                    par sa table.
%%                  - On peut démontrer qu'il y a deux tables pour les groupes
%%                    à 4 éléments

\begin{definition}
Soit $\p{G,\star}$ un groupe et $H$ une partie de $G$. On dit que $H$ est un
\emph{sous-groupe} de $(G,\star)$ lorsque
\begin{itemize}
\item $e\in H$
\item $\forall x,y\in H \qsep x\star y\in H$
\item $\forall x\in H \qsep x^{-1}\in H.$
\end{itemize}
Si tel est le cas, alors $\p{H,\star}$ est un groupe.
\end{definition}

\begin{preuve}
Preuve du fait qu'un sous-groupe est un groupe.
\begin{itemize}
\item[$\bullet$]$\star$ est bien une lci sur $H$ d'après (ii).
\item[$\bullet$]$\star$ est associative (hérité de $G$)
\item[$\bullet$]$\star$ admet un élément neutre d'après (i) et hérité de $G$.
\item[$\bullet$] Tout élément est symétrisable d'après (iii).
\end{itemize}
\end{preuve}

\begin{remarques}
\remarque Si $H$ est un sous-groupe de $G$, alors~:
  $\forall x\in H\qsep \forall n\in\Z \qsep x^n\in H.$
% \remarque Afin de montrer les deux derniers points, il suffit de montrer que~: $\forall x,y\in H\qsep x\star y^{-1}\in H$.
\remarque En pratique, pour montrer que $(H,\star)$ est un groupe, on le fera
  presque toujours apparaître comme sous-groupe d'un groupe connu.
  % En pratique, pour montrer que $H$ est un sous-groupe de $G$, on vérifiera le plus souvent les 3 points suivants.
  % \begin{itemize}
  % \item $H\subset G$
  % \item $e\in H$
  % \item $\forall x,y\in H\qsep x\star y^{-1}\in H$.
  % \end{itemize}
% \remarque On peut montrer que si $H$ une partie du groupe $(G,\star)$, stable par $\star$, et telle que $(H,\star)$ est un groupe, alors $H$ est un sous-groupe de $G$.
\end{remarques}

\begin{exemples}
\exemple Si $(G,\star)$ est un groupe, $G$ et $\ens{e}$ sont des sous-groupes
  de $G$. Le sous-groupe $\ens{e}$ est appelé groupe \emph{trivial}.
\exemple $\R$ et $\Z$ sont des sous-groupes de $(\C,+)$.
  De même, $\Rs$ et $\U$ sont des sous-groupes de $(\Cs,\times)$.
\end{exemples}

%% Exemples :
%% 1) Z est un sous-groupe de R pour l'addition
%% 2) nZ est un sous-groupe de Z
%% 4) L'ensemble des isométries du plan est un sous-groupe du groupe des
%%    bijections du plan. (isométrie = bijection conservant les distances)
%% 
%% Remarques :
%% 1) En pratique, pour montrer que (E,x) est un groupe, on montrera très
%%    souvent que c'est un sous-groupe d'un groupe connu.
%% 2) Il suffit de vérifier que x \in H et que x y^{-1} \in H

\begin{proposition}
Si $n\in\Ns$, $\p{\U[n],\times}$ est un groupe dont l'élément neutre est 1.
\end{proposition}

\begin{proposition}
Soit $E$ un ensemble. On note $\sigma(E)$ l'ensemble des bijections de $E$
dans $E$. Alors $\p{\sigma(E),\circ}$ est un groupe, appelé groupe des
\emph{permutations} de $E$, dont l'élément neutre est $\id_E$.
\end{proposition}

\begin{exoUnique}
\exo Montrer que l'ensemble des bijections strictement croissantes de $\R$ dans $\R$ est
  un sous-groupe de $(\sigma(\R),\circ)$.
\end{exoUnique}

\begin{proposition}
L'intersection d'une famille de sous-groupes est un sous-groupe.
\end{proposition}

\begin{remarqueUnique}
\remarque Contrairement à l'intersection, l'union de deux sous-groupes n'est
  en général pas un sous-groupe.
\end{remarqueUnique}

%% Remarque :
%% 1) C'est faux pour l'union. Par exemple 2Z et 3Z sont des sous-groupes de Z
%%    ET 2Z u 3Z n'en n'est pas un

\begin{definition}
Soit $\p{G,\star}$ un groupe et $A$ une partie de $G$. Alors, au sens de l'inclusion, il existe un plus
petit sous-groupe de $G$ contenant $A$; on l'appelle \emph{groupe engendré} par $A$ et
on le note ${\rm Gr}(A)$.
\end{definition}

\begin{preuve}
On considère $(H_i)_{i\in I}$ la famille des sous-groupes de $G$ qui contiennent $A$. On pose $\displaystyle H=\bigcap_{i\in I} H_i$. On montre ensuite successivement que $A\subset H$, $H$ est un sg de $(G,\star)$, si $K$ est un sous-groupe de $G$ contenant $A$, $H\subset K$.
\end{preuve}

\begin{remarqueUnique}
\remarque Si $(G,\star)$ est un groupe et $x$ est un élément de $G$, le
  groupe engendré par $\ens{x}$, appelé aussi groupe engendré par $x$,
  est $\ensim{x^k}{k\in\Z}$.
\end{remarqueUnique}
\begin{sol}
On pose $B=\enstq{x^k}{k\in\Z}$. On montre que $B$ est un sous-groupe de $G$. Ensuite, si $H$ est un sg de $G$ contenant $A$, il contient $x$ donc tous les $x^k, k\in \Z$ par propriété des groupes. Donc il contient $B$. Ainsi, $B$ est bien le plus petit.
\end{sol}



\begin{sol}
\begin{itemize}
$\U[n]=\set{\omega^r, r\in \Z}=\set{\omega^r, r\in \intere{0}{n-1}}$
\item[$\bullet$] Soit $k\in\Z$. On suppose que $\U[n]={\rm Gr}\p{\set{\omega^k}}=\set{\p{\omega^k}^l, l\in\Z}$. En particulier, comme $\omega \in \U[n]$, il existe $l\in\Z$ tel que $\omega^{kl}=\omega$, i.e. $\e^{i\frac{2\pi}{n}kl}=\e^{i\frac{2\pi}{n}}$ ce qui conduit à $kl\equiv 1\quad [n]$ d'où $kl+(-u)n=1$ avec $u\in \Z$ ce qui signifie d'après le théorème de Bézout que $k\wedge n=1$.
\item[$\bullet$] Réciproquement, on suppose que $k\wedge n=1$. Montrons que $\omega \in {\rm Gr}\p{\set{\omega^k}}$. Il existe $u,v \in \Z$ tel que $uk+vn=1$. Alors $$\p{\omega^k}^v=\omega^{1-vn}=\omega \p{\omega^n}^{-v}=\omega.$$ Donc $\omega \in {\rm Gr}\p{\set{\omega^k}}$. Par propriété des groupes, on en déduit que $\U[n]\subset {\rm Gr}\p{\set{\omega^k}}$. On a aussi l'autre inclusion d'où le résultat.
\end{itemize}
\end{sol}

\begin{definition}
Soit $\p{G_1,\star_1}$ et $\p{G_2,\star_2}$ deux groupes. On dit qu'une
application $\phi$ de $G_1$ dans $G_2$ est un \emph{morphisme de groupe} lorsque
\[\forall x,y\in G_1 \qsep \phi\p{x\star_1 y}=\phi(x)\star_2\phi(y).\]
Plus précisément, on dit que $\phi$ est un
\begin{itemize}
\item \emph{endomorphisme} lorsque $(G_1,\star_1)=(G_2,\star_2)$
\item \emph{isomorphisme} lorsque $\phi$ est bijective
\item \emph{automorphisme} lorsque $\phi$ est un endomorphisme et un isomorphisme.
\end{itemize}
\end{definition}

\begin{remarqueUnique}
\remarque L'application $\phi$ de $\R$ dans $\U$ qui à $\theta$ associe
  $\e^{\ii\theta}$ est un morphisme du groupe $(\R,+)$ dans le groupe $(\U,\times)$.
  L'application $\exp$ de $\R$ dans $\RPs$ est un isomorphisme du
  groupe $(\R,+)$ dans le groupe $(\RPs,\times)$. 
\end{remarqueUnique}
%% Exemples :
%% 1) R -> RPs, x -> exp(x)
%% 2) RPs -> R, x -> ln(x)
%% 3) R -> U, theta -> exp(i theta)


\begin{proposition}
Soit $\phi$ un morphisme du groupe de $\p{G_1,\star_1}$ dans $\p{G_2,\star_2}$.
Alors
\begin{eqnarray*}
& & \phi\p{e_1}=e_2\\
\forall x\in G_1, & & \phi\p{x^{-1}}=\cro{\phi(x)}^{-1} \\
\forall x\in G_1 \qsep \forall n\in\Z, & & \phi\p{x^n}=\cro{\phi(x)}^n.
\end{eqnarray*}
\end{proposition}

\begin{preuve}
Ne pas oublier d'obtenir le deuxième point en montrant inverse à gauche \underline{et} à droite.
\end{preuve}

\begin{remarqueUnique}
\remarque Si $\phi$ est un morphisme de groupe et que les lois sont notées additivement, alors
  \[\forall x\in G_1\qsep \forall n\in\Z\qsep \phi(n\cdot x)=n\cdot\phi(x).\]
\end{remarqueUnique}

\begin{exoUnique}
\exo Déterminer les endomorphismes, puis les automorphismes de $(\Z,+)$.
\begin{sol}
On montre par analyse-synthèse que ce sont les $\phi$ qui vérifient $\phi(k)=\phi(1)k, \forall k \in \Z$. Pour les automorphismes, cela impose $\phi(1)=\pm 1$. Pour montrer que ceux-ci fonctionnent, on peut montrer qu'ils sont des involutions.
\end{sol}
% \exo Quels sont les morphismes de $(\Q,+)$ dans $(\Z,+)$~?
% \begin{sol}
% On montre que le seul possible est le morphisme nul, essentiellement grâce à $\phi(kr)=k\phi(r)$ pour $k\in\Z$.
% \end{sol}
\end{exoUnique}

\begin{proposition}
Soit $\phi$ un morphisme de $\p{G_1,\star_1}$ dans $\p{G_2,\star_2}$. Alors
\begin{itemize}
\item l'image réciproque d'un sous-groupe de $G_2$ est un sous-groupe de $G_1$.
\item l'image directe d'un sous-groupe de $G_1$ est un sous-groupe de $G_2$.
\end{itemize}
\end{proposition}

\begin{definition}
Soit $\phi$ un morphisme de $\p{G_1,\star_1}$ dans $\p{G_2,\star_2}$. On
appelle \emph{noyau} de $\phi$ et on note $\ker\phi$ l'ensemble
\[\ker\phi\defeq\enstq{x\in G_1}{\phi(x)=e_2}.\]
C'est un sous-groupe de $G_1$.
\end{definition}

\begin{proposition}
Un morphisme $\phi$ de $\p{G_1,\star_1}$ dans $\p{G_2,\star_2}$ est injectif
si et seulement si
\[\ker\phi=\ens{e_1}.\]
\end{proposition}

\begin{remarqueUnique}
\remarque Pour montrer l'injectivité d'un morphisme, montrer que $\ker\phi=\ens{e_1}$ doit
  devenir un réflexe. Pour cela, il est naturel de procéder par double
  inclusion. Mais comme l'inclusion $\ens{e_1}\subset\ker\phi$
  est toujours vraie, puisque $\phi(e_1)=e_2$, il est essentiel de se concentrer sur
  l'inclusion $\ker\phi\subset\ens{e_1}$.
\end{remarqueUnique}


\begin{exoUnique}
\exo Soit $(G,\star)$ un groupe et $\phi$ l'application de $G$ dans
  $\sigma(G)$ définie par
  \[\dspappli{\phi}{G}{\sigma(G)}{x}{\dspappli{\phi(x)}{G}{G}{g}{x\star g}}\]
  Montrer que $\phi$ est bien définie et que c'est un morphisme injectif de
  groupe. En déduire que $(G,\star)$ est isomorphe à un sous-groupe du groupe
  de ses permutations.
\end{exoUnique}

\begin{proposition}
$\quad$
\begin{itemize}
\item La composée de deux morphismes de groupe est un morphisme de groupe.
\item La bijection réciproque d'un isomorphisme de groupe est un isomorphisme de
  groupe.
\end{itemize}
\end{proposition}

\begin{proposition}
Si $\p{G,\star}$ est un groupe, on note ${\rm Aut}(G)$ l'ensemble des
automorphismes de $G$. $\p{{\rm Aut}(G),\circ}$ est un groupe.
\end{proposition}

\begin{definition}
Soit $\p{G_1,\star_1}$ et $\p{G_2,\star_2}$ deux groupes. On définit la loi $\star$
sur $G_1\times G_2$ par
\[\forall \p{x_1,x_2},\p{y_1,y_2}\in G_1\times G_2 \qsep
  \p{x_1,x_2}\star\p{y_1,y_2}=\p{x_1\star_1 y_1,x_2\star_2 y_2}.\]
Alors $\p{G_1\times G_2,\star}$ est un groupe d'élément neutre $\p{e_1,e_2}$ et
\[\forall \p{x_1,x_2}\in G_1\times G_2 \qsep
  \p{x_1,x_2}^{-1}=\p{x_1^{-1},x_2^{-1}}.\]
\end{definition}

\begin{exoUnique}
\exo Montrer que $(\RPs\times\U,\times)$ est isomorphe à $(\Cs,\times)$.
\end{exoUnique}
\begin{sol}
On définit $\dspappli{\phi}{\RPs\times\U}{\Cs}{(r,u)}{ru}$. On montre que $\phi$ est bien défini, un morphisme, puis injectif, puis surjectif.
\end{sol}

\subsection{Ordre d'un élément}

\begin{proposition}
Pour tout $n\in\N$, on pose
\[n\Z\defeq\ensim{kn}{k\in\Z}.\]
C'est un sous-groupe de $(\Z,+)$.
\end{proposition}

\begin{proposition}
Une partie $H$ de $\Z$ est un sous-groupe de $(\Z,+)$ si et seulement si il
existe $n\in\N$ tel que $H=n\Z$. De plus, si tel est le cas, l'entier $n$ est unique.
\end{proposition}

\begin{preuve}
On raisonne par analyse-synthèse.
Dans l'analyse, après avoir évacué le cas nul, on s'intéresse à $A=\Ns\cap G$ qui étant non vide admet un ppe qu'on note $n$. On montre qu'alors $G=n\Z$ en utilisant la DE. La synthèse est triviale.\\

L'unicité vient du fait que si $n_1\Z=n_2\Z$ alors $n_1\in n_2\Z$ donc $n_2\mid n_1$. De même, $n_1\mid n_2$.
\end{preuve}

\begin{remarqueUnique}
\remarque Si $H$ est un sous-groupe de $(\Z,+)$ non réduit à $\ens{0}$, alors
  $H$ admet un plus petit élément strictement positif $n\in\Ns$. On a alors $H=n\Z$.
\end{remarqueUnique}

\begin{definition}
Soit $(G,\star)$ un groupe et $x\in G$.
\begin{itemize}
\item On dit que $x$ est d'\emph{ordre fini} lorsqu'il existe $n\in\Ns$ tel que $x^n=e$.
  Dans ce cas, il existe un unique $\omega\in\Ns$ tel que 
  \[\forall n\in\Z \qsep x^n=e \quad\ssi\quad \omega|n.\]
  On l'appelle \emph{ordre} de $x$. C'est le plus petit entier $n\in\Ns$ tel que
  $x^n=e$.
\item Sinon, on dit que $x$ est d'\emph{ordre infini}. On a alors
  \[\forall n\in\Z \qsep x^n=e \quad\ssi\quad n=0.\]
\end{itemize}
% Alors, l'application
% \[\dspappli{\phi}{\Z}{G}{n}{x^n}\]
% est un morphisme du groupe $(\Z,+)$ dans $(G,\star)$.
% %  dont l'image est le groupe
% % engendré par $x$.
% \begin{itemize}
% \item Si $\ker \phi=\ens{0}$, on dit que $x$ n'a pas d'ordre. Dans ce cas
% \item Sinon, il existe un unique $\omega\in\Ns$ tel que $\ker\phi=\omega\Z$.
%   On dit que $\omega$ est l'\emph{ordre} de $x$. C'est l'unique entier naturel tel que
%   %\cro{\exists k\in\Z \quad n=k\omega}\]
% \end{itemize}
\end{definition}

\begin{remarques}
\remarque Dans
  $(\Cs,\times)$, si $n\in\Ns$, $\omega\defeq\e^{\ii\frac{2\pi}{n}}$ est d'ordre
  $n$.
  \begin{sol}
  On montre que $\omega^k=1\Longleftrightarrow k\equiv 0 \quad [n] \Longleftrightarrow n\mid k$.
  \end{sol}
\remarque Dans un groupe, $e$ est l'unique élément d'ordre 1.
\remarque Dans un groupe fini, tout élément est d'ordre fini.
\remarque Soit $x\in G$ un élément d'ordre
  $\omega\in\Ns$. Alors le groupe engendré par $x$ est
  $\ens{e,x,x^2,\ldots,x^{\omega-1}}$, ces éléments étant deux à deux distincts. En particulier, l'ordre de $x$ est le
  cardinal du groupe qu'il engendre.
  \begin{sol}
  On veut montrer que $\set{x^n, n\in \Z}=\set{x^n, n\in \intere{0}{\omega-1}}$. On raisonne par double-inclusion, l'une étant triviale et l'autre s'obtient grâce à la DE.
  Reste ensuite à voir que les éléments de $\ens{e,x,x^2,\ldots,x^{\omega-1}}$ sont bien deux à deux distincts. pour cela, on en prend $2$ d'exposants $k_1$ et $k_2$. On voit assez vite que leur différence doit être multiple de $\omega$ alors qu'elle est plus petite en valeur absolue que $\omega$, donc c'est $0$.
  \end{sol}
\end{remarques}

% \begin{exoUnique}
% \exo Soit $(G,\star)$ un groupe et $x\in G$ un élément d'ordre fini $n\in\Ns$.
%   Étant donné $k\in\Z$, calculer l'ordre de $x^k$.
%   \begin{sol}
%   $\forall a \in \Z, \p{w^k}^a=1\Longleftrightarrow \omega^{ka}=1 \Longleftrightarrow n\mid ka$ (cf. prop. précédente).
%   \'Ecrivons alors $n=(k\wedge n)n'$ et $k=(k\wedge n)k'$ ce qui impose (classique maintenant) $k'\wedge n'=1$. Reprenons les équivalences précédentes :
%   $$\forall a \in \Z, \p{w^k}^a=1\Longleftrightarrow (k\wedge n)n'\mid ka=(k\wedge n)k' \Longleftrightarrow n'\mid k'a \underbrace{\Longleftrightarrow}_{Gauss} n'\mid a \Longleftrightarrow \dfrac{n}{k\wedge n}\mid a.$$ Donc l'ordre de $\omega^k$ est $\dfrac{n}{k\wedge n}$.
  
%   En particulier, l'ordre de $\omega^k$ est $n$ ssi $k\wedge n=1$.
  
%   \end{sol}
% \end{exoUnique}

\begin{theoreme}[nom={Théorème de \nom{Lagrange}}]
Soit $(G,\star)$ un groupe fini et $x$ un élément de $G$. Alors l'ordre de $x$
divise le cardinal de $G$.
\end{theoreme}

\begin{preuve}
ADMIS. C'est une conséquence du Théorème de Lagrange (HP) qui dit que l'ordre d'un sous-groupe d'un groupe fini le divise.
\medskip

Tout de même, remarquons que dans un groupe fini, tout élément admet un ordre. En effet, si tel n'était pas le cas alors $\dspappli{\phi}{\Z}{G}{k}{x^k}$ serait injective ce qui est absurde puisque $\Z$ est infini tandis que $G$ est fini.
\bigskip

\textbf{Preuve HP du théorème de Lagrange} :
On définit sur $G$ la relation binaire $R_H$ par $$x R_H y \Longleftrightarrow xy^{-1} \in H.$$
C'est une relation d'équivalence sur $G$ (à vérifier). Or
\begin{eqnarray*}
y\in Cl_H(x) &\Longleftrightarrow& y R_H x \\
&\Longleftrightarrow& yx^{-1}\in H \\
&\Longleftrightarrow& \exists h \in H \text{ tq } yx^{-1}=h \\
&\Longleftrightarrow& \exists h \in H \text{ tq } y=hx \\
&\Longleftrightarrow& y\in Hx=\set{hx, h\in H}.
\end{eqnarray*}
Donc $\overline{x}=Hx$ qui est équipotent à $H$ via $h\mapsto hx$ et a donc même cardinal que $H$, et ce peu importe la classe. Les classes ont donc toutes mêmes cardinales. Or, elles forment une partition (cf. relation d'équivalence). Disons qu'il y a $q$ classes, on a donc $|G|=q|H|$.
\end{preuve}

\begin{remarques}
\remarque Si $(G,\star)$ est un groupe fini, le cardinal de $G$ est aussi appelé
  ordre de $G$. La version faible du théorème de Lagrange nous dit donc que
  dans un groupe fini, l'ordre d'un élément divise l'ordre du groupe.
\remarque La version forte du théorème de \nom{Lagrange} dit que si $(G,\star)$ est
  un groupe fini et $H$ est un sous-groupe de $(G,\star)$, alors le
  cardinal de $H$ divise le cardinal de $G$. De cette version forte découle
  la version faible~: si $x\in G$, il suffit de remarquer que le cardinal du
  groupe $H$ engendré par $x$ est l'ordre de $x$.
\end{remarques}

\begin{exoUnique}
\exo Déterminer les sous-groupes finis de $(\U,\times)$. 
\begin{sol}
Soit $H$ un sous-groupe fini de $(\U,\times)$.  Notons $n$ son cardinal. Soit $x$ un élément de $H$, d'après le théorème précédent, l'ordre $a$ de $x$ divise $n$ donc $au=n$ et en particulier $x^n=(x^a)^u=1$ donc $x\in \U[n]$. Ainsi, on a montré $H\subset \U[n]$ et on a donc égalité par égalité des cardinaux.
\end{sol}
\end{exoUnique}



\section{Groupe symétrique}
\subsection{Groupe symétrique}

\begin{definition}
Soit $n\in\N$. On appelle \emph{groupe symétrique} et on note $\p{\gsym{n},\circ}$
l'ensemble des bijections de $\intere{1}{n}$ dans lui-même muni de la loi de
composition.
\end{definition}

\begin{remarques}
\remarque Si $\sigma\in\mathcal{F}\p{\intere{1}{n},\intere{1}{n}}$, l'application
  $\sigma$ est aussi notée
  \[\sigma=
    \begin{pmatrix}
    1 & 2 & \cdots & n\\
    \sigma\p{1} & \sigma\p{2} & \cdots & \sigma(n)
    \end{pmatrix}.\]
  Puisque $\intere{1}{n}$ est fini, $\sigma$ est bijective si et seulement si
  elle est injective ou surjective. Autrement dit, $\sigma$ est bijective si
  et seulement si l'une des deux conditions suivantes est vérifiée~:
  \begin{enumerate}
  \item Les entiers $\sigma\p{1},\ldots,\sigma(n)$ sont deux à deux distincts.
  \item $\ens{\sigma\p{1},\ldots,\sigma(n)}=\intere{1}{n}$.
  \end{enumerate}
% \remarque L'élément neutre de $\gsym{n}$ est $\id$. De plus, si
%   $\sigma\in\gsym{n}$, son symétrique pour la loi $\circ$ est sa bijection
%   réciproque $\sigma^{-1}$.
% \remarque Le groupe $\p{\gsym{2},\circ}$ ne possède que deux
%   éléments~: $\id$ et la permutation $\sigma$ définie par $\sigma\p{1}=2$ et
%   $\sigma\p{2}=1$. En particulier il est commutatif. Cependant, si $n\geq 3$,
%   le groupe $\p{\gsym{n},\circ}$ n'est pas commutatif.
\remarque Si $E$ est un ensemble fini de cardinal $n$, l'ensemble des bijections
  de $E$ muni de la loi de composition est un groupe isomorphe à
  $(\gsym{n},\circ)$.
\end{remarques}

\begin{proposition}
$\p{\gsym{n},\circ}$ est un groupe fini de cardinal $n!$.  
\end{proposition}

\begin{definition}
Soit $n\in\N$.
\begin{itemize}
\item Soit $p\in\intere{2}{n}$. On appelle \emph{cycle} de longueur $p$
  (ou p-cycle) toute
  permutation $\sigma$ tel qu'il existe $x_0,\ldots,x_{p-1}\in\intere{1}{n}$ deux à
  deux distincts tels que
  \begin{itemize}
  \item $\sigma\p{x_0}=x_1,\sigma\p{x_1}=x_2,\ldots,\sigma\p{x_{p-1}}=x_0$
  \item $\forall x\in\intere{1}{n}\setminus\ens{x_0,\ldots,x_{p-1}} \qsep
         \sigma(x)=x$
  \end{itemize}
  On note $\sigma=\p{x_0 \quad x_1 \quad \cdots \quad x_{p-1}}$.
\item On appelle \emph{transposition} tout cycle de longueur 2.
\end{itemize}
\end{definition}

\begin{remarques}
\remarque Si $n\geq 3$, $(\gsym{n},\circ)$ n'est pas commutatif.
\remarque Si $\sigma$ est une transposition, alors $\sigma^2=\id$. On en déduit que
  $\sigma^{-1}=\sigma$. 
\remarque Si $i\in\Z$ et $p\in\Ns$, on note $i\ {\rm mod}\ p$, le reste de la division
  euclidienne de $i$ par $p$. Si $\sigma=(x_0 \quad x_1 \quad \cdots \quad x_{p-1})$ est
  un $p$-cycle, on a donc
  \[\forall i\in\interefo{0}{p}\qsep \sigma(x_i)=x_{i+1\ {\rm mod}\ p}.\]
\remarque Les $p$-cycles sont des éléments d'ordre $p$.
% \remarque On appelle support d'une permutation l'ensemble des
%   $k\in\intere{1}{n}$ tels que $\sigma(k)\neq k$. Deux permutations
%   de supports disjoints commutent.
% \remarque Toute permutation s'écrit comme le produit (commutatif) de cycles
%   de supports disjoints.
% \remarque On dit que deux permutations $\sigma_1$ et $\sigma_2$ sont conjuguées
%   lorsqu'il existe une permutation $\tau$ telle que
%   $\sigma_1=\tau^{-1} \sigma_2 \tau$. Sur $\gsym{n}$, la relation
%   \og est conjuguée à \fg, est une relation d'équivalence. De plus, si
%   $\sigma_1$ et $\sigma_2$ sont conjuguées, il en est de même pour
%   $\sigma_1^k$ et $\sigma_2^k$ quel que soit $k\in\Z$.
% \remarque Deux $p$-cycles sont conjugués. En particulier, les transpositions
%   sont conjuguées entre elles.
\end{remarques}
%% Remarque :
%% 1) Toute permutation est un produit de cycles disjoints
%%    (à voir sur des exemples)
%% 2) si sigma est un p-cycle, sigma^p=id
%% 3) En particulier, l'inverse d'une transposition est elle-même

\begin{exos}
\exo Soit $\tau$ un $p$-cycle et $\sigma\in\gsym{n}$. Montrer que
  $\sigma\tau\sigma^{-1}$ est un $p$-cycle.
\exo Montrer que si $\sigma_1,\sigma_2\in\gsym{n}$ sont deux $p$-cycles, il
  existe $\sigma\in\gsym{n}$ tel que $\sigma_2=\sigma \sigma_1 \sigma^{-1}$.
  \begin{sol}
  Montrons pour commencer (pas nécessaire pour la preuve de l'exercice) que si $\sigma_1=\p{x_0 \quad x_1 \quad \cdots \quad x_{p-1}}$ alors pour toute permutation $\sigma$, $\sigma \circ \sigma_1 \circ \sigma^{-1}=\p{\sigma(x_0) \quad \sigma(x_1) \quad \cdots \quad \sigma(x_{p-1})}$. Pour cela, on pose $\gamma=\sigma \circ \sigma_1 \circ \sigma^{-1}$ et on regarde ce que vaut $\gamma(x)$ selon si $x\in \set{x_0,x_1,\cdots,x_{p-1}}$ ou non.\\
  Maintenant, si $\sigma_2$ est un autre $p$-cycle, $\sigma_2=\p{u_0 \quad u_1 \quad \cdots \quad u_{p-1}}$, il suffit de construire $\sigma$ qui envoie $x_i$ sur $u_i$ et qui se complète par une bijection de $\intere{1}{n}\setminus \set{x_0,x_1,\cdots,x_{p-1}}$ sur $\intere{1}{n}\setminus \set{u_0,u_1,\cdots,u_{p-1}}$ ce qui est possible car ce sont des ensembles de même cardinaux.
  \end{sol}
\end{exos}

% \begin{proposition}
% Soit $\sigma=\p{x_1 \quad x_2 \quad \cdots \quad x_p}$ un cycle de longueur $p$.
% Alors~:
% \[\sigma=\p{x_1 \quad x_2}\p{x_2 \quad x_3}\cdots\p{x_{p-1} \quad x_p}\]
% \end{proposition}

\subsection{Décomposition en cycles à supports disjoints}

\begin{definition}
Soit $\sigma\in\gsym{n}$. On définit la relation $\mathcal{R}$ sur
$\intere{1}{n}$ par
\[\forall x,y\in\intere{1}{n} \qsep x\mathcal{R}y \quad\ssi\quad
  \cro{\exists k\in\Z \qsep \sigma^k(x)=y}.\]
Alors $\mathcal{R}$ est une relation d'équivalence. Si $x\in\intere{1}{n}$, la
classe de $x$ est notée $\mathcal{O}(x)$ et est appelée orbite de $x$.
\end{definition}

\begin{remarques}
\remarque Les orbites étant des classes d'équivalence, elles forment une
  partition de $\intere{1}{n}$.
\remarque Si $x\in\intere{1}{n}$, alors
  $\mathcal{O}(x)=\ensim{\sigma^k(x)}{k\in\Z}$. De plus, il existe
  un plus petit $p\in\Ns$ tel que $\sigma^p(x)=x$. On a alors
  $\mathcal{O}(x)=\ens{x,\sigma(x),\ldots,\sigma^{p-1}(x)}$.
\end{remarques}


% \begin{exoUnique}
% \exo Montrer qu'une permutation est un cycle si et seulement si la
%   relation d'équivalence définie ci-dessus admet une et une seule classe
%   non réduite à un point.
% \end{exoUnique}

\begin{definition}
Soit $\sigma\in\gsym{n}$. On appelle \emph{support} de $\sigma$ et on note
$\supp(\sigma)$ l'ensemble des $x\in\intere{1}{n}$ tels que $\sigma(x)\neq x$.
\end{definition}

\begin{remarques}
\remarque Si $\sigma=(x_0 \quad x_1 \quad \cdots \quad x_{p-1})$ est un $p$-cycle,
  alors ${\rm Supp}(\sigma)=\ens{x_0,\ldots,x_{p-1}}$.
\remarque Le support de $\sigma$ est stable par $\sigma$.
\remarque Deux permutations de supports disjoints commutent. Cependant la
  réciproque est fausse.
% \remarque Soit $\sigma_1,\ldots,\sigma_m\in\gsym{n}$ une famille de
%   permutations de supports deux à deux disjoints telle que
%   $\sigma_1 \circ \cdots \circ \sigma_m=\id$. Alors, pour tout
%   $i\in\intere{1}{m}$, $\sigma_i=\id$.
\end{remarques}

\begin{theoreme}
Toute permutation s'écrit comme le produit (commutatif) de cycles à supports
disjoints. De plus, à l'ordre près, il y a unicité d'une telle décomposition.
\end{theoreme}

\begin{remarqueUnique}
\remarque Si une permutation $\sigma\in\gsym{n}$ s'écrit comme le produit
  de $m$ cycles à supports deux à deux disjoints de longueurs respectives $p_1,\ldots,p_m$, alors l'ordre de
  $\sigma$ est $\ppcm(p_1,\ldots,p_m)$.  
\end{remarqueUnique}

\begin{exos}
\exo Déterminer tous les éléments de $\gsym{3}$. Quels sont leurs ordres~?
\exo Quels sont les entiers qui sont l'ordre d'un élément de $\gsym{4}$~?
% \exo Combien de fois un mélange portant sur 6 cartes doit-il être répété
%   pour retomber à coup sûr sur l'ordre initial~?
\exo Déterminer les éléments d'ordre 2 de $\gsym{n}$~?
\end{exos}

\subsection{Signature, groupe alterné}

\begin{proposition}
Toute permutation $\sigma\in\gsym{n}$ s'écrit comme le produit d'au plus
$n-1$ transpositions.
\end{proposition}

\begin{remarqueUnique}
% \remarque Plus précisement, toute permutation $\sigma\in\intere{1}{n}$ s'écrit
%   comme le produit d'au plus $n-1$ transpositions.
\remarque Soit $\sigma=\p{x_0 \quad x_1 \quad \cdots \quad x_{p-1}}$ un cycle de
  longueur $p$. Alors
  \[\sigma=\p{x_0 \quad x_1}\p{x_1 \quad x_2}\cdots\p{x_{p-2} \quad x_{p-1}}.\]
\end{remarqueUnique}

\begin{exoUnique}
\exo Dans $\gsym{3}$, on pose $\sigma_1\defeq\p{1\quad 3}$ et
  $\sigma_2\defeq\p{1\quad 2\quad 3}$. Décomposer $\sigma_1 \sigma_2$ en
  produit de transpositions de deux manières distinctes.
\end{exoUnique}

\begin{definition}
Soit $\sigma$ une permutation et
\[\sigma=\tau_1\cdots\tau_{m} \et \sigma=\tau_1'\cdots\tau_{m'}'\]
deux décompositions de $\sigma$ en produit de transpositions.
Alors $m$ et $m'$ ont même parité; on dit que $\sigma$ est \emph{paire} lorsque ces entiers
sont pairs et que $\sigma$ est \emph{impaire} dans le cas contraire.
On définit la \emph{signature} de $\sigma$ que l'on note $\signat{\sigma}$ par
\[\signat{\sigma}\defeq
  \begin{cases}
  +1 & \text{si $\sigma$ est paire}\\
  -1 & \text{si $\sigma$ est impaire.}
  \end{cases}\]
\end{definition}

\begin{remarques}
\remarque Soit $\sigma\in\gsym{n}$ et $\sigma=\tau_1\cdots\tau_{m}$ une
  décomposition de $\sigma$ en produit de transpositions. Alors
  \[\signat{\sigma}=\p{-1}^m.\]
\remarque La signature d'un $p$-cycle est $\p{-1}^{p-1}$. En particulier, les
  transpositions sont impaires et les 3-cycles sont pairs.
\end{remarques}

% \begin{proposition}
% Soit $\sigma$ un cycle de longueur $p$. Alors~:
% \[\signat{\sigma}=\p{-1}^{p+1}\]
% En particulier, la signature d'une transposition est $-1$.
% \end{proposition}

\begin{proposition}
L'application $\signature$ de $\p{\gsym{n},\circ}$ dans $\p{\ens{-1,1},\times}$
est un morphisme de groupe.
\end{proposition}

\begin{remarqueUnique}
\remarque Si $\sigma$ est une permutation, $\sigma$ et $\sigma^{-1}$ ont la même
  signature.  
\end{remarqueUnique}

% \begin{exos}
% \exo Il n'existe que deux morphismes du groupe $\p{\gsym{n},\circ}$ dans
%   $\p{\Cs,\times}$~: le morphisme trivial qui à tout élément $\sigma$ de
%   $\gsym{n}$ associe 1, et la signature.
% \end{exos}

\begin{definition}
On note $\galt{n}$ l'ensemble des permutations paires. C'est un sous-groupe de
$(\gsym{n},\circ)$ appelé \emph{groupe symétrique alterné}.
\end{definition}

\begin{remarqueUnique}
\remarque Si $n\geq 2$, le groupe $\p{\galt{n},\circ}$ est de cardinal $n!/2$.
\end{remarqueUnique}

\begin{sol}
FAIT : Soit $\tau \in \gsym{n}$ tel que $\signature(\tau)=-1$ (par exemple une transposition), alors $\tau \galt{n}=\gsym{n}\setminus \galt{n}$.
En effet, raisonnons par double inclusion :
\begin{itemize}
\item[$\bullet$] $\subset$ : Si $\sigma \in \galt{n}$, $\signature(\tau\circ\sigma)=\ldots=-1$ donc $\tau\circ\sigma\in \gsym{n}\setminus \galt{n}$.
\item[$\bullet$] $\supset$ : Soit $\sigma\notin \galt{n}$, $\signature(\tau^{-1}\circ \sigma)=\ldots=1$ donc $\tau^{-1}\circ \sigma \in \galt{n}$ d'où $\sigma=\tau\circ(\tau^{-1}\circ \sigma) \in \tau \galt{n}$.
\end{itemize}

Il reste alors à montrer que $\abs{\tau \galt{n}}=\abs{\galt{n}}$. Pour cela définissons :
$$\dspappli{\phi}{\galt{n}}{\tau \galt{n}}{\sigma}{\tau\circ \sigma}.$$
$\phi$ est surjective par définition même de $\tau \galt{n}$ et injective car $\phi(\sigma)=\phi(\sigma')\Rightarrow \sigma=\sigma'$ par régularité de $\tau$.
D'où le résultat souhaité.

On a alors une union disjointe de deux ensembles de même cardinal dont la réunion est de cardinal $n!$, donc chacun a pour cardinal $\dfrac{n!}{2}$.
\end{sol}

% \subsection{Groupe diédral}

% \begin{definition}
% Soit $n\geq 2$. L'ensemble des similitudes du plan complexe laissant invariant
% $\U[n]$ est un groupe pour la composition appelé groupe diédral et noté
% $(D_n,\circ)$.
% \end{definition}

% \begin{proposition}
% Soit $n\geq 2$. Une application $f:\C\mapsto\C$ est un élément du groupe
% diédral si et seulement si il existe $u\in\U[n]$ tel que
% \[\cro{\forall z\in\C \quad f(z)=uz} \ou \cro{\forall z\in\C \quad
%   f(z)=u\bar{z}}\]
% En particulier, ces applications étant deux à deux distinctes, le groupe
% diédral est fini de cardinal $2n$.
% \end{proposition}

% \begin{remarques}
% \remarque Toute similitude du plan laissant invariant $\U[n]$ induit
%   une bijection de $\U[n]$ dans lui-même, donc un élément de $\gsym{n}$.
%   On construit ainsi un morphisme de groupe $\phi$ de $(D_n,\circ)$ dans
%   $(\gsym{n},\circ)$. Ce morphisme est injectif dès que $n\geq 3$.
% \end{remarques}

% \begin{exos}
% \exo Montrer que $\gsym{3}$ est isomorphe à $D_3$. Que devient
%   $\galt{3}$ par cet isomorphisme~?
% \end{exos}

% \subsection{Groupe $(\Z/n\Z,+)$}

% \begin{definition}
% Soit $n\in\Ns$ et $\mathcal{R}$ la relation d'équivalence définie sur $\Z$ par
% \[\forall a,b\in\Z \qsep a\mathcal{R}b \quad\ssi\quad a\equiv b\ [n].\]
% On appelle $\Z/n\Z$ l'ensemble des classes d'équivalence pour cette relation.
% \end{definition}

% \begin{proposition}
% Soit $n\in\Ns$. Pour tout $k\in\Z$, on note $\overline{k}$ la classe
% d'équivalence de $k$. Alors, les éléments
% $\overline{0},\overline{1},\ldots,\overline{n-1}$ sont deux à deux distincts
% et
% \[\Z/n\Z=\ens{\overline{0},\overline{1},\ldots,\overline{n-1}}.\]
% \end{proposition}

% \begin{preuve}
% On montre qu'ils sont distincts en prenant $k_1$ et $k_2$ dans $\intere{0}{n-1}$ tels que $\overline{k_1}=\overline{k_2}$ et on montre que $k_1=k_2$.\\
% Puis, on montre l'égalité par double inclusion, droite-gauche étant immédit et gauche-droite se faisant grâce à la DE.
% \end{preuve}

% \begin{definition}
% Soit $n\in\Ns$. On définit la loi de composition interne $+$ sur $\Z/n\Z$ par
% \[\forall k_1,k_2\in\Z \qsep \overline{k_1}+\overline{k_2}=\overline{k_1+k_2}\]
% Alors $(\Z/n\Z,+)$ est un groupe commutatif de cardinal $n$ et d'élément
% neutre $\overline{0}$.
% \end{definition}

% \begin{preuve}
% On montre que
% \begin{itemize}
% \item[$\bullet$] + est bien défini en montrant que la définition ne dépend pas du représentant choisi dans la classe.
% \item[$\bullet$]$(\Z/n\Z,+)$ est un groupe commutatif (associativité, commutativité, élément neutre, opposé)
% \end{itemize}
% Soient $(x,y)\in \p{\Z/n\Z}^2$. Il existe $k,k' \in \Z$ tel que $x=\overline{k}$ et $y=\overline{k'}$ et on fait tout avec ça...
% \end{preuve}

% \begin{remarqueUnique}
% \remarque Du groupe $(\Z/n\Z,+)$, on retiendra essentiellement le fait que
%   l'application
%   \[\dspappli{\phi}{\Z}{\Z/n\Z}{k}{\overline{k}}\]
%   est un morphisme surjectif de groupe et que~:
%   $\forall k_1,k_2\in\Z \qsep  \overline{k_1}=\overline{k_2} \ssi
%     k_1\equiv k_2\ [n]$.
% \end{remarqueUnique}

% \begin{sol}
% En pratique, vous pouvez faire "comme si" $\Z/n\Z=\intere{0}{n-1}$ et $\phi$ associé le reste dans la DE de $k$ par $n$.
% \end{sol}

% \begin{exoUnique}
% \exo Soit $(G,\star)$ un groupe et $x\in G$ un élément d'ordre $n\in\Ns$.
%   Montrer que l'application
%   \[\dspappli{\phi}{\Z/n\Z}{G}{\overline{k}}{x^k}\]
%   est bien définie et qu'elle réalise un isomorphisme de $\Z/n\Z$ sur le
%   groupe engendré par $x$.
%   \begin{sol}
%   On montre que c'est indépendant du représentant puis que c'est un morphisme puis que $\ker\phi=\set{0}$ puis égalité des cardinaux.
%   \end{sol}
% % $n\in\Ns$ et $\omega=e^{i\frac{2\pi}{n}}$. Montrer que
% %   l'application $\phi$ de $(\Z/n\Z,+)$ dans $(\U[n],\cdot)$ qui à
% %   $\overline{k}$ associe $\omega^k$ est bien définie et que c'est un
% %   isomorphisme de groupe.
% \end{exoUnique}
%END_BOOK

\end{document}
\section{Exercices}
\setcounter{numeroexercice}{1}
\documentclass{magnolia}

\magtex{tex_driver={pdftex}}
\magfiche{document_nom={Groupes},
          auteur_nom={Emmanuel Roblet, François Fayard},
          auteur_mail={fayard.prof@gmail.com}}
\magexos{exos_matiere={maths},
         exos_niveau={mpsi},
         exos_chapitre_numero={12},
         exos_theme={Groupes}}
\magmisenpage{}
\maglieudiff{}
\magprocess


\begin{document}
%BEGIN_BOOK
\magsection{Groupe}
\magsubsection{Loi de composition interne}


\exercice{nom={Loi de composition interne}}
\begin{questions}
\question Déterminer les propriétés de la loi de composition interne $\perp$
  sur $\Q$ définie par
  \[\forall a,b\in\Q \qsep a\perp b\defeq a+b+ab.\]
\question Faire de même avec la loi $\nabla$ définie sur  $\Rs\times\R$ par
  \[\forall \p{a,b},\p{c,d}\in\Rs\times\R \qsep
    (a,b)\nabla(c,d)\defeq\p{ac,\frac{d}{a}+bc}.\]
\end{questions}

\begin{sol}
\begin{questions}
\question Commutativité immédiate.
Associativité :
\begin{eqnarray*}
(a\perp b)\perp c&=&(a+b+ab)\perp c=a+b+ab+c+(a+b+ab)c\\
&=&a+b+c+ab+ac+bc+abc\\
&=&a+(b+c+bc)+a(b+c+bc)\\
&=&a\perp (b\perp c).
\end{eqnarray*}
Neutre : 0.
\question $\nabla$ n'est pas commutative : Prendre $(1,1)$ et $(2,0)$ par exemple. \\
$\nabla$ est en revanche associative (suffit de l'écrire).
Neutre : $(1,0)$.
\end{questions}
\end{sol}

\magsubsection{Groupe}

\exercice{nom={Transport de structure}}
\begin{questions}
\question Soit $(G,\square)$ un groupe et $H$ un ensemble tel qu'il existe une
  fonction $f:H\to G$ bijective. On définit la loi $\star$ sur $H$ par
  \[\forall x,y\in H \qsep x\star y\defeq f^{-1}\p{f(x)\square f(y)}.\]
  Montrer que $(H,\star)$ est un groupe isomorphe à $(G,\square)$. 
\question On définit la loi $\oplus$ sur $\R$ par
  \[\forall x,y\in\R \qsep x\oplus y\defeq\sqrt[3]{x^3+y^3}.\]
  Montrer que $(\R,\oplus)$ est un groupe commutatif.
\end{questions}
\begin{sol}
\begin{questions}
\question On montre aisément que $\star$ est associative sur $H$, que $f^{-1}(e_G)$ est l'élément neutre de $H$ puis que si $x\in H$, son inverse est $f^{-1}$ de l'inverse de $f(x)$. On vérifie bien tout cela à gauche et à droite. Donc $H$ est un groupe. Reste à vérifier que le $f$ considéré est alors un homomorphisme de groupes.
\question Application avec $f$ la fonction cube.
\end{questions}
\end{sol}

\exercice{nom={Sous groupe et permutations}}
Soit $E$ un ensemble et $A$ une partie de $E$. On note $\mathcal{S}(A)$
l'ensemble des permutations~$f$ de $E$ qui laissent la partie $A$ invariante,
c'est-à-dire telles que $f(A)=A$. Montrer que $\mathcal{S}(A)$ est un
sous-groupe de $(\sigma(E),\circ)$.


\exercice{nom={Sous-groupes additifs de $\R$}}
% On dit qu'une Vpartie $G$ de $\R$ est un sous-groupe de $(\R,+)$ lorsque
% \begin{itemize}
% \item $0\in G,$
% \item $\forall x,y\in G\qsep x+y\in G,$
% \item $\forall x\in G\qsep -x\in G.$
% \end{itemize}
\begin{questions}
\question Soit $a\in\RP$. Montrer que
  \[a\Z\defeq\ensim{ka}{k\in\Z}\]
  est un sous-groupe additif de $\R$.
\enonce Le but de cet exercice est de montrer que les sous-groupes additifs de $\R$ sont soit de cette forme, soit dense dans $\R$. Soit $G$ un sous-groupe de $\p{\R,+}$ distinct de $\ens{0}$. On pose $a\defeq\inf\p{G\cap\RPs}.$
\question Montrer que $a$ est bien défini.
\question Dans cette question on suppose que $a>0$.
  \begin{questions}
  \question Montrer que $a\in G$.
  \question En déduire que $G=a\Z$.
  \end{questions}
\question Montrer que si $a=0$, $G$ est dense dans $\R$.
\enonce On résume les conclusions des questions précédentes en disant que les
  sous-groupes additifs de $\R$ sont soit discrets soit denses dans $\R$.
\question 
  Soit $a$ et $b$ deux réels non nuls tels que $a/b\not\in\Q$. Montrer que
  \[a\Z+b\Z=\ensim{na+mb}{n,m\in\Z}\]
  est dense dans $\R$.
\question 
  Soit $f$ une fonction continue de $\R$ dans $\R$. On dit qu'un réel $T$ est
  une période de $f$ lorsque
  \[\forall x\in\R \qsep f\p{x+T}=f(x).\]
  \begin{questions}
  \question Montrer que l'ensemble $\mathcal{P}$ des périodes de $f$ est un
    sous-groupe additif de $\R$.
  \question On suppose que $f$ est non constante et périodique. Montrer qu'il
    existe un unique $T_0>0$ tel que $\mathcal{P}=T_0\Z$.
  \end{questions}
\end{questions}

\begin{sol}
\begin{questions}
\question Ok. Il suffit d'écrire les choses.
\question Il existe $x\neq 0$ tel que $x\in G$. $x$ ou $-x$ est donc dans $\RPs$.
\question Dans cette question on suppose que $a>0$.
  \begin{questions}
  \question Si $a\notin G$, d'après la caractérisation de la borne inférieure, il existe $b\in G$ tel que $a<b< a+a$ et il existe $c \in G$ tel que $a<c<b$. D'où $b-c \in G$ d'une part et $0<b-c<a$ ce qui fournit une contradiction car $a$ est un minorant de $G\cap\RPs$.
  \question D'après la question précédente, $a\in G$ donc par stabilité pour $+$ et parce-que $-a\in G$, $a\Z \subset G$. Soit alors $g\in G$. Par définition de la partie entière, il existe $n\in \Z$ tel que $n\leq \dfrac{g}{a}<n+1$ d'où $0\leq g-na <a$ avec $g-na \in G$ par stabilité pour $+$. Donc $g-na=0$ par définition de $a$. On a montré l'existence d'un $n\in \Z$ tel que $g=na$, d'où $g \in a\Z$. Ainsi, $G\subset a\Z$ et on a donc l'égalité.
  \end{questions}
\question Si $a=0$, soit $x,y \in \R$ tels que $x<y$. Il existe $g\in G$ tel que $0<g<y-x$ car $\inf\p{G\cap\RPs}=0$. Avec $n=\ent{\dfrac{x}{g}}+1$, on a $n-1\leq \dfrac{x}{g} <n$ donc $(n-1)g\leq x < ng$. On a donc $$x<ng=(n-1)g+g <x+(y-x)=y.$$ On a donc trouvé un élément de $G$ dans $[x,y]$.
\question 
  $a\Z+b\Z$ est un sous-groupe de $\R$. Supposons par l'absurde qu'il n'est pas dense dans $\R$, alors il est du type $c\Z$ avec $c\in \R$. Mais alors, comme $a \in a\Z+b\Z=c\Z$, il existe $u\in \Z$ tel que $a=cu$. De même, il existe $v\in \Z$ tel que $b=cv$. Alors $\dfrac{a}{b}=\dfrac{u}{v} \in \Q$, d'où la contradiction.
\question 
  
  \begin{questions}
  \question 
  \question 
  \end{questions}
\end{questions}
\end{sol}

\exercice{nom={Sous-groupe de $\U$}}
Montrer que $\cup_{n\in\Ns} \U[n]$ est un sous-groupe strict de
$(\U,\cdot)$.

\begin{sol}
On montre que c'est un sous-groupe de $(\U,\cdot)$.
\begin{itemize}
\item[$\bullet$] $G\subset \U$.
\item[$\bullet$] $G\neq \emptyset$ car $1\in G$.
\item[$\bullet$] Soit $(x,y)\in G$, montrons que $xy^{-1} \in G$. On peut fixer $(m,n)\in \N^2$ tel que $x^m=1=y^n$.
$$\p{\frac{x}{y}}^{mn}=\frac{x^{mn}}{y^{mn}}=\frac{\p{x^m}^n}{\p{y^n}^m}=1.$$
\end{itemize}
\end{sol}

\exercice{nom={Union de deux sous-groupes}}
Soit $(G,*)$ un groupe, et $H$ et $K$ deux sous-groupes de $G$. Montrer que
$H\cup K$ est un sous-groupe de $G$ si et seulement si  $H\subset K$ ou
$K\subset  H$.

\begin{sol}
Si $H\subset K$, alors $H\cup K = K$ est un sous-groupe de  $(G,\times)$. De même si $K\subset H$.
    
        Réciproquement, si $H\cup K$ est un sous-groupe de  $(G,\times)$. Supposons que $H\not\subset K$ et $K\not\subset H$. 
        Il existe donc $x \in H\setminus K$ et $y \in K\setminus H$. 
        Si $xy \in H$, alors $y = x^{-1}xy \in H$ (par stabilité de $H$ par produit et passage à l'inverse), on obtient une contradiction. De même si $xy \in K$.
        Ainsi, $H\subset K$ ou $K\subset H$.
        
        On a donc bien que 
        \begin{center}
            \fbox{\begin{minipage}{0.9\textwidth}$  H\cup K  $ est un
sous-groupe de $(G,\times)$ si et seulement si $  H\subset K  $ ou $
K\subset H $. \end{minipage}}
        \end{center}
\end{sol}

\exercice{nom={Groupes tels que $x^2=e$}}
Soit $\p{G,\star}$ un groupe tel que
\[\forall x\in G \qsep x^2=e.\]
Montrer que $G$ est commutatif.

\begin{sol}
$(xy)^2=xyxy=e$ puis multiplication par $x$ à gauche et $y$ à droite.
\end{sol}

\exercice{nom={Action d'un groupe sur lui-même}}
Soit $(G,\star)$ un groupe, dont l'élément neutre est noté $e$. Pour $a\in G$,
on définit les  applications~$\gamma_a$ et~$\tau_a$ de $G$ dans lui-même par les
formules
\[\forall g\in G \qsep \gamma_a(g)\defeq a\star g \et \tau_a(g)\defeq a\star g\star a^{-1}.\]
\begin{questions}
\question Montrer que $\gamma_a$ et $\tau_a$ sont des bijections. Montrer que
  $\tau_a$ est un automorphisme de $G$. Que dire de $\gamma_a$~?
\question Pour $a$ et $b$ éléments quelconques de $G$, calculer
  $\tau_a\circ \tau_b$. Retrouver le fait que $\tau_a$ est un automorphisme de $G$,
  et donner une autre expression de $(\tau_a)^{-1}$.
% \question On note $\text{Aut}(G)$ l'ensemble de tous les automorphismes du 
%   groupe $(G,\star )$. Montrer que $(\text{Aut}(G),\circ)$ est un groupe.
\question On note $\text{Int}(G)$ le sous ensemble de $\text{Aut}(G)$ formé
  des applications $\tau_a$ pour $a\in G$ : ce sont les automorphismes dits
  {\em intérieurs}\/ de $G$. Montrer  que l'application $a\mapsto \tau_a$ est
  un morphisme  du groupe $(G,\star )$ dans  le groupe $(\text{Aut}(G),\circ)$, et en
  déduire que $\text{Int}(G)$ est un sous-groupe de $(\text{Aut}(G),\circ)$.
\end{questions}

\exercice{nom={Groupe diédral}}
Pour tout $(a,b)\in\Cs\times\C$, on pose
\[\dspappli{\phi_{a,b}}{\C}{\C}{z}{az+b}\]
et on pose $A^+\defeq\ensim{\phi_{a,b}}{(a,b)\in\Cs\times\C}$. On pose enfin
\[\dspappli{s}{\C}{\C}{z}{\conj{z}}\]
\begin{questions}
\question
\begin{questions}
\question Montrer que $A^+$ est un groupe pour la composition.
\question On pose $A\defeq A^+ \cup A^+ s$ où $A^+ s\defeq\ensim{f\circ s}{f\in A^+}$.
  Montrer que $A$ est un sous-groupe de $\sigma(\C)$.
\end{questions}
\question Soit $n\geq 2$. On note $D_n$ l'ensemble des fonctions $f\in A$ pour
  lesquelles $f(\U[n])\subset\U[n]$.
\begin{questions}
\question Montrer que si $f\in D_n$ l'application induite à $\U[n]$ est bijective.
\question Montrer que $D_n$ est un sous-groupe de $A$ appelé
  \emph{groupe diédral de degré $n$}.
\question Montrer que $D_n$ contient $s$ et
  \[\dspappli{r}{\C}{\C}{z}{\e^{\frac{2\ii\pi}{n}}z}\]
\question Que vaut la somme des éléments de $\U[n]$~? En déduire que 0 est point fixe
  de tout élément de $D_n$.
\question Montrer que
\[D_n=\ensim{r^k s^\epsilon}{k\in\interefo{0}{n}, \ \epsilon\in\ens{0,1}}.\]
\end{questions}
\end{questions}

\magsubsection{Ordre d'un élément}




% \exercice{nom={Sous-groupes additifs de $\Z$}}
% L'objet de cet exercice est de montrer que les sous-groupes additifs de $\Z$
% sont les ensembles de la forme
% \[b\Z\defeq\ensim{nb}{n\in\Z}\]
% où $b\in\N$.
% \begin{questions}
% \question Soit $b\in\N$. Montrer que $b\Z$ est un sous-groupe de $\p{\Z,+}$.
% \question Réciproquement, soit $G$ un sous-groupe additif de $\Z$, non réduit à $\set{0}$. On pose
%   $b\defeq\min\p{G\cap\Ns}$. Justifier l'existence de $b$ et montrer que $G=b\Z$.
%   % {\it On rappelle que si $a\in\Z$ et $b\in\Ns$, il existe un unique couple
%   % $\p{q,r}\in\Z^2$ tel que $a=qb+r$ et $r\in\intere{0}{b-1}$.}
% \end{questions}
% \begin{sol}
% \begin{questions}
% \question OK.
% \question On procède par double inclusion. $b\Z\subset G$ OK. Prenons $x\in G$ et effectuons la DE de $x$ par $b$. Le reste est forcément nul, sinon, il contredirait la définition de $b$.
% \end{questions}
% \end{sol}





\exercice{nom={Théorème de Lagrange}}
Soit $(G,\star)$ un groupe fini et $H$ un sous-groupe de $G$. On souhaite montrer
que $\card(H)$ divise $\card(G)$. Pour cela, on introduit la relation binaire $\mathcal{R}$
sur $G$ par
\[\forall x,y \in G\qsep x\mathcal{R}y \ \ssi\  yx^{-1}\in H.\]
\begin{questions}
\question Montrer que $\mathcal{R}$ est une relation d'équivalence.
\question Soit $x\in G$. Montrer que
  \[\dspappli{\phi}{H}{{\rm Cl}(x)}{h}{hx}\]
  est une bijection.
\question Conclure.
\end{questions}



% \exercice{nom={Sous-groupe}}
% Soit $(G,\star)$ un groupe, et $H$ et~$K$ deux sous-groupes de $G$. On pose 
% \[HK=\enstq{x\star y}{x\in H \quad  y\in K}\] 
% Montrer que les propriétés suivantes sont équivalentes~:
% \begin{itemize}
% \item $HK$ est un sous-groupe de $(G,\star)$.
% \item $KH$ est un sous-groupe de $(G,\star)$.
% \item $HK\subset KH$
% \item $KH\subset HK$
% \end{itemize}

% \exercice{nom={Groupe produit}}
% Soit $(G,\star)$ un groupe, et $H$ et $K$ deux sous-groupes de $G$. On suppose
% que $H\cap K=\ens{e}$.
% \begin{questions}
% \question Montrer que l'application $\phi$, de $H\times K$ dans
%   \[HK\enstq{x\star y}{x\in H, y\in K}\] 
%   définie par $\phi(x,y)=x\star y$, est une bijection.
% \question[label={b}] On suppose de plus que $H$ et $K$ sont finis. Que dire du
%   cardinal de $HK$~?
% \question La conclusion de la question \qref{b} persiste-t-elle si l'on ne
%   suppose plus que  $H\cap K=\ens{e}$~?
% \end{questions}



\exercice{nom={Existence d'un élément d'ordre 2}}
Soit $\p{G,\star}$ un groupe fini de cardinal pair. Le but de cet exercice est
de montrer qu'il possède un élément d'ordre 2.
Pour cela, on considère l'ensemble
\[E\defeq\enstq{x\in G}{x^2\neq e}.\]
\begin{questions}
\question On définit la relation $\mathcal{R}$ sur $E$ par
  \[\forall x, y\in E\qsep x\mathcal{R}y\ \ssi\ y=x \ou y=x^{-1}.\]
  Montrer que $\mathcal{R}$ est une relation d'équivalence. Quel est le cardinal
  de chacune de ses classes d'équivalence.
\question Conclure.
\end{questions}
\emph{On peut généraliser ce résultat en montrant que si $p$ est un nombre premier divisant l'ordre de $G$, ce dernier possède un élément d'ordre $p$. Ce résultat est connu sous le nom de théorème de \nom{Sylow}.}

\exercice{nom={Les groupes d'ordre inférieur à 5 sont commutatifs}}
\begin{questions}
\question Soit $(G,\star)$ un groupe fini dont le cardinal $p$ est un nombre
  premier et $x\in G\setminus\ens{e}$. Montrer que
  \[G=\ensim{x^k}{k\in\intere{0}{p-1}}\]
  puis en déduire que $G$ est commutatif.
\question Montrer que les groupes finis de cardinal inférieur ou égal à 5 sont
  commutatifs.\\
  \textit{On montrera qu'il n'y a que deux tables possibles pour les groupes
  de cardinal 4.}
\question Donner un exemple de groupe de cardinal 6 non commutatif.
\end{questions}

\begin{sol}
\begin{questions}
\question 
\question Pour $n = 2; 3 \ou 5$, la première question nous assure que $G$ est commutatif.
Pour $n = 4$, s'il y a un élément d'ordre $4$ dans le groupe, celui-ci est cyclique et donc commutatif. Sinon, tous les éléments du groupe vérifient $x^2 = e$. Il est alors classique de
justifier que le groupe est commutatif.
\question 
\end{questions}
\end{sol}

% \exercice{nom={Sous-groupes de $\Z/n\Z$}}
% Soit $n\in\Ns$. Déterminer et dénombrer les sous-groupes de $(\Z/n\Z,+)$.

% \exercice{nom={Théorème chinois}}
% Soit $p,q\in\Ns$ deux entiers premiers entre eux. Montrer que l'application
% $\phi$ de $(\Z/pq\Z,+)$ dans $(\Z/p\Z\times \Z/q\Z,+)$ qui à $\overline{k}$ associe
% $(\overline{k},\overline{k})$ est bien définie et est un isomorphisme de groupe.
% \begin{sol}
% Si $n\in \N$ et $x\in \Z$, on note $cl_n(x)$ la classe de $x$ dans $\Z/n\Z$. Soit $(p,q)\in \Ns^2$ tel que $p\wedge q=1$. Soit $f$ l'application définie sur $\Z/pq\Z$, à valeurs dans $\Z/p\Z\times \Z/q\Z$ par $\alpha\mapsto (cl_p(x) ;cl_q(x))$, où $x\in \Z$ tel que $\alpha=cl_{pq}(x)$.\\
% Montrons déjà que l'application $f$ est bien définie. Soit $\alpha\in \Z/pq\Z$. Soit $x,y \in \Z$ tels que $\alpha=cl_{pq}(x) =cl_{pq}(y)$ . $x\equiv y \quad [pq]$ donc $pq$ divise $x-y$. On en déduit que $p$ divise $x-y$ (i.e. $cl_p(x) =cl_p(y)$) et $q$ divise $x-y$ (i.e. $cl_q(x) =cl_q(y)$) et donc $f$ est bien définie.\\
% Montrons  maintenant  que $f$ est  un  morphisme  d'anneaux.  Soit $\alpha,\beta \in \Z/pq\Z$.  Soit $x,y\in \Z$ tels  que $\alpha=cl_{pq}(x)$ et $\beta=cl_{pq}(y)$. On a $\alpha+\beta=cl_{pq}(x+y), \alpha \beta=cl_{pq}(xy)$ et :
% $$f(\alpha+\beta) = (cl_p(x+y) ;cl_q(x+y)) = (cl_p(x)+cl_p(y) ;cl_q(x)+cl_q(y)) = (cl_p(x) ;cl_q(x))+(cl_p(y) ;cl_q(y)) =f(\alpha)+f(\beta)$$
% $$f(\alpha \beta) = (cl_p(xy) ;cl_q(xy)) = (cl_p(x)cl_p(y) ;cl_q(x)cl_q(y)) = (cl_p(x) ;cl_p(y))(cl_q(x) ;cl_q(y)) =f(\alpha)f(\beta)$$
% $$f(cl_{pq}(1)) = (cl_p(1) ;cl_q(1))$$
% $f$ est bien un morphisme d'anneaux.\\
% Il  reste  à  montrer  que $f$ est  bijective.  Soit $\alpha \in \Z/pq\Z$ tel  que $f(\alpha)  =  (cl_p(0) ;cl_q(0))$.  Soit $x\in \Z$ tel  que $\alpha=cl_{pq}(x)$. $f(\alpha) = (cl_p(x) ;cl_q(x))$. On en déduit $cl_p(x) =cl_p(0)$ (donc $p$ divise $x$) et $cl_q(x) =cl_q(0)$ (donc $q$ divise $x$). $p$ et $q$ étant  premiers entre eux, on en déduit que $pq$ divise $x$, et donc $cl_{pq}(x) =cl_{pq}(0)$. $f$ est donc injective. On obtient alors la bijectivité par égalité des cardinaux.
% \end{sol}

\magsection{Groupe symétrique}
\magsubsection{Groupe symétrique}
\magsubsection{Décomposition en cycles à supports disjoints}
\exercice{nom={Décomposition en produit de cycles}}
Décomposer la permutation suivante en produit de cycles de support
disjoints.
$$\sigma\defeq
\begin{pmatrix}
1&2&3&4&5&6&7&8&9&10\\
9&4&3&8&7&10&1&2&5&6 
\end{pmatrix}
$$
En déduire sa signature.

\exercice{nom={Exercice}}
Déterminer l'ordre maximal d'un élément de $\gsym{10}$.

\magsubsection{Signature, groupe alterné}


\exercice{nom={Générateurs du groupe symétrique}}
\begin{questions}
\question Montrer que les transpositions $\p{1\quad i}$ (pour
  $i\in\intere{2}{n}$) engendrent le groupe symétrique $(\gsym{n}, \circ)$.
\question Montrer que les cycles de longueur $3$ engendrent $(\galt{n},\circ)$.
\end{questions}

\exercice{nom={Exercice}}
\begin{questions}
\question Montrer que pour tout $n\in\N$, il existe un morphisme injectif de
  $(\gsym{n},\circ)$ dans $(\galt{n+2},\circ)$.
\question Montrer qu'il n'existe pas de morphisme injectif de $(\gsym{4},\circ)$
  dans $(\galt{5},\circ)$.
\end{questions}
\begin{sol}
Prendre $\tau$ la transposition $(n+1,n+2)$ et $\sigma \mapsto
\tau \circ \sigma$ si $\sigma$ est impaire et $\sigma$ sinon.\\
Il n'y a pas dans $A_5$ d'élément d'ordre $4$.
\end{sol}

\exercice{nom={Définition de la signature}}
Soit $n\geq2$. Le but de cet exercice est de démontrer qu'il existe deux
et seulement deux morphismes de groupe de $(\gsym{n},\circ)$ dans
$(\Cs,\times)$.
\begin{itemize}
\item L'application $\bar{1}$ qui a toute permutation $\sigma$ associe $1$.
\item Un autre morphisme $\signature$ que l'on définira comme étant la
  signature.
\end{itemize}

\begin{questions}
\question Le but de cette partie est de montrer qu'il existe au plus un
  seul morphisme $\phi$ de $(\gsym{n},\circ)$ dans $(\Cs,\times)$
  différent de $\bar{1}$.
  Soit $\phi$ un morphisme de $(\gsym{n},\circ)$ dans $(\Cs,\times)$.
  \begin{questions}
  \question Soit $\tau$ une transposition. Montrer que $\phi(\tau)\in
    \ens{-1,1}$.
  \question En déduire que $\phi$ est à valeurs dans $\ens{-1,1}$.
  \question Soit $\tau_1$ et $\tau_2$ deux transpositions.
    \begin{questions}
    \question Montrer que $\tau_1$ et $\tau_2$ sont conjuguées, c'est-à-dire
      qu'il existe une permutation $\sigma$ telle que
      $$\tau_2 = \sigma \tau_1 \sigma^{-1}$$
    \question En déduire que $\phi(\tau_1)=\phi(\tau_2)$.
    \end{questions}
  \question Conclure
  \end{questions}
\question Le but de cette partie est de montrer l'existence d'un morphisme de
  $(\gsym{n},\circ)$ dans $(\Cs,\times)$ différent de $\bar{1}$.\\
  On dit qu'une partie $A$ de $\intere{1}{n}^2$ est une représentation
  des couples d'éléments de $\intere{1}{n}$ lorsque

  $$\forall i,j \in\intere{1}{n}
  \qsep
  \begin{cases}
  (i,j)\in A \Longrightarrow (j,i)\not\in A\\
  i\neq j \implique \cro{(i,j)\in A \quad\text{ou}\quad (j,i)\in A}
  \end{cases}
  $$
  \begin{questions}
  \question Soit $A$ une représentation des couples d'éléments de
    $\intere{1}{n}$ et $\sigma$ une permutation de
    $\intere{1}{n}$. Montrer que : 
    $$\sigma(A)\defeq\ens{ (\sigma(i),\sigma(j)) : (i,j)\in A }$$
    est une représentation des couples d'éléments de $\intere{1}{n}$.
  \question Soit $A$ une représentation des couples d'éléments de
    $\intere{1}{n}$ et $\sigma$ une permutation de
    $\intere{1}{n}$. On note $n_A$ le nombre d'inversions de $\sigma$,
    c'est-à-dire le nombre d'éléments $(i,j)$ de $A$ tels que $j-i$ et
    $\sigma(j)-\sigma(i)$ soient de signes distincts. On définit alors
    la signature de $\sigma$ par
    $$\signature(\sigma)\defeq(-1)^{n_A}$$
    \begin{questions}
    \question Montrer que la signature ne dépend pas du choix de $A$.
    \question Montrer que :
      $$\signature(\sigma)=\prod_{(i,j)\in A}
                         \frac{\sigma(j)-\sigma(i)}{j-i}$$
    \question En déduire que $\signature$ est un morphisme de groupe
    \end{questions}
  \question Montrer que $\signature$ est différent de $\bar{1}$ et conclure.
  \end{questions}
\question En déduire qu'il existe un unique morphisme de $\gsym{n}$ dans
  $\ens{-1,1}$ qui vaut $-1$ sur les transpositions. Ce morphisme est
  appelé signature.
\end{questions}


%END_BOOK


\end{document}




\chapter{Limites et continuité}
\setcounter{numeroexercicecours}{1}
\documentclass{magnolia}

\magtex{tex_driver={pdftex},
        tex_packages={epigraph,xypic}}
\magfiche{document_nom={Continuité, limites},
          auteur_nom={François Fayard},
          auteur_mail={fayard.prof@gmail.com}}
\magcours{cours_matiere={maths},
          cours_niveau={mpsi},
          cours_chapitre_numero={10},
          cours_chapitre={Continuité, limites}}
\magmisenpage{}
\maglieudiff{}
\magprocess

\begin{document}

%BEGIN_BOOK
\setlength\epigraphwidth{.75\textwidth}
\epigraph{\og Un mathématicien est une machine à transformer le café en théorèmes.\fg}{--- \textsc{Paul Erdös (1913-1996)}}
\setlength\epigraphwidth{.65\textwidth}
\epigraph{\og Si les gens ne croient pas que les mathématiques sont simples, c’est seulement parce qu’ils ne réalisent pas combien la vie est compliquée !\fg}{--- \textsc{John Von Neumann (1903--1957)}}

\magtoc

\section{Fonction numérique, topologie élémentaire}

\begin{definition}[utile=-3]
On appelle \emph{fonction numérique} toute fonction définie sur une partie $\mathcal{D}$
de $\R$, à valeurs dans $\R$ ou $\C$.
\end{definition}

\begin{remarques}
\remarque Dans la suite de ce chapitre, $\K$ désignera le corps $\R$ ou $\C$.
\remarque Il arrive que l'on définisse une fonction par son expression \og $f(x)$ \fg. C'est
  alors au lecteur de déterminer son domaine de définition, c'est-à-dire l'ensemble des
  $x\in\R$ pour lesquels \og $f(x)$ \fg a un sens.
\end{remarques}

\begin{definition}[utile=-3]
Soit $f:\mathcal{D}\to\K$ et $A$ une partie de $\mathcal{D}$. On dit que $f$
\emph{vérifie la propriété $\mathcal{P}$ sur $A$} lorsque la restriction de $f$ à $A$ vérifie
la propriété $\mathcal{P}$.
\end{definition}

% \subsection{Symétries, fonctions lipschitziennes}

% \begin{definition}[utile=-3]
% Soit $f$ une fonction dont le domaine de définition est symétrique par rapport
% à 0, c'est-à-dire tel que
% \[\forall x\in\mathcal{D} \qsep -x\in\mathcal{D}.\]
% On dit que $f$ est
% \begin{itemize}
% \item \emph{paire} lorsque
%   \[\forall x\in\mathcal{D} \qsep f\p{-x}=f(x).\]
% \item \emph{impaire} lorsque
%   \[\forall x\in\mathcal{D} \qsep f\p{-x}=-f(x).\]
% \end{itemize}
% \end{definition}

% % \begin{remarques}
% % \remarque Si $f$ est une fonction dérivable paire, alors $f'$ est impaire. De
% %   même, si $f$ est une fonction dérivable  impaire, $f'$ est paire.
% % \remarque Si $f$ est une fonction impaire (resp. paire) continue sur $\R$. Que
% %   peut-on dire de ses primitives~?
% % \end{remarques}

% %% Si f est impaire, F est paire.
% %% Si f est paire, seule la primitive s'annulant en 0 est impaire.

% \begin{definition}[utile=-3]
% Soit $T\in\R$ et $f$ une fonction dont le domaine de définition vérifie
% \[\forall x\in\mathcal{D} \qsep \forall k\in\Z \qsep x+kT\in\mathcal{D}\]
% On dit que $f$ est \emph{$T$-périodique}, ou que $T$ est une \emph{période} de $f$, lorsque
% \[\forall x\in\mathcal{D} \qsep f\p{x+T}=f(x).\]
% Lorsque $f$ admet une période non nulle, on dit que $f$ est \emph{périodique}.
% \end{definition}

% % \begin{remarques}
% % \remarque L'ensemble $G$ des périodes d'une fonction $f$ est un sous-groupe de
% %   $\p{\R,+}$. En particulier, si $T$ est une période de $f$ et si $k\in\Z$,
% %   alors $kT$ est une période de $f$.
% %   Lorsque $f$ admet une plus petite période strictement positive, notée
% %   $T_0$, ont dit que $T_0$ est la période de $f$. On montre dans ce cas
% %   que $G=T_0\Z$.
% % \remarque Si $f$ est dérivable et $T$-périodique, il en est de même pour $f'$.
% %   Cependant, il est possible que $f'$ soit $T$-périodique sans que $f$ le soit.
% % \end{remarques}

% % \begin{exos}
% % \exo Les périodes de $\tan$ sont les $k\pi$, pour $k\in\Z$. Donc $\tan$
% %   est périodique et sa période est $\pi$.
% % \end{exos}

% \begin{definition}[utile=-3]
% Soit $k\in\RP$. On dit qu'une fonction $f$ est \emph{$k$-lipschitzienne} lorsque
% \[\forall x,y\in \mathcal{D} \qsep \abs{f(x)-f(y)}\leq k\abs{x-y}.\] 
% \end{definition}

% \begin{exoUnique}
% \exo Montrer que les fonctions \og valeur absolue\fg et \og sinus\fg sont
%   1-lipschitziennes.
% % \exo La fonction valeur absolue est $1$-lipschitzienne. En effet
% %   \[\forall x,y\in\R \qsep \abs{\abs{x}-\abs{y}}\leq\abs{x-y}\]
% % \exo La fonction $\sin$ est $1$-lipschitzienne.
% \end{exoUnique}

% % \begin{remarques}
% % \remarque Si $f$ est une fonction dérivable et $k$-lipschitzienne,
% %   alors
% %   \[\forall x\in\mathcal{D} \qsep \abs{f'(x)}\leq k\]
% %   Réciproquement, nous verrons que si $f$ est dérivable sur un intervalle $I$
% %   et que
% %   \[\forall x\in I \qsep \abs{f'(x)}\leq k\]
% %   alors $f$ est $k$-lipschitzienne. Parmi les fonctions dérivable sur un
% %   intervalle $I$, les fonctions lipschitziennes sont donc les fonctions dont
% %   la dérivée est bornée.
% % \remarque L'ensemble des fonctions lipschitziennes est un sous-espace vectoriel
% %   de $\mathcal{F}\p{\R,\R}$, mais ce n'est pas une sous-algèbre. En effet, la
% %   fonction $f:x\mapsto x$ est lipschitzienne, mais $f^2$ ne l'est pas.
% % \end{remarques}




\subsection{Propriété locale}

% \begin{definition}[utile=-3]
% \begin{itemize}
% \item Soit $a\in\R$. On dit qu'une partie $\mathcal{V}$ de $\R$ est un \emph{voisinage
%   fondamental} de $a$ lorsqu'il existe $\epsilon>0$ tel que 
%   \[\mathcal{V}=\enstq{x\in\R}{\abs{x-a}\leq\epsilon} 
%     =\interf{a-\epsilon}{a+\epsilon}.\]
% \item On dit qu'une partie $\mathcal{V}$ de $\R$ est un \emph{voisinage fondamental} de $+\infty$
%   lorsqu'il existe $m\in\R$ tel que
%   \[\mathcal{V}=\interfo{m}{+\infty}.\]
% \item On dit qu'une partie $\mathcal{V}$ de $\R$ est un \emph{voisinage fondamental} de $-\infty$
%   lorsqu'il existe $M\in\R$ tel que
%   \[\mathcal{V}=\interof{-\infty}{M}.\]
% \item Soit $a\in\C$. On dit qu'une partie $\mathcal{V}$ de $\C$ est un \emph{voisinage
%   fondamental} de $a$ lorsqu'il existe $\epsilon>0$ tel que
%   \[\mathcal{V}=\enstq{z\in\C}{\abs{z-a}\leq\epsilon}.\]
% \end{itemize}
% \end{definition}

% \begin{remarques}
% \remarque Une intersection de deux voisinages fondamentaux d'un même élément est un voisinage fondamental. Plus
%   généralement, une intersection d'un nombre fini de voisinages fondamentaux est un
%   voisinage fondamental.
% \begin{sol}
% Attention, c'est faux pour une infinité de voisinages : $\displaystyle \bigcap_{n\in \Ns}\intero{a-\frac{1}{n}}{a+\frac{1}{n}}$.
% \end{sol}
% \remarque Une suite réelle (ou complexe) $\p{u_n}$ converge vers
%   $l\in\RbarouC$ si
%   et seulement si pour tout voisinage fondamental $\mathcal{V}$ de $l$, il existe un
%   voisinage fondamental $\mathcal{W}$ de $+\infty$ tel que
%   \[\forall n\in\N \qsep n\in\mathcal{W} \implique u_n\in\mathcal{V}.\]
%   \begin{sol}
%   C'est la définition de la limite pour chacun d'entre eux.
%   \end{sol}
% \remarque Une partie $A$ de $\R$ est dense dans $\R$ si et seulement si pour
%   tout $x\in\R$ et tout voisinage $\mathcal{V}$ de $x$,
%   $A\cap\mathcal{V}\neq\emptyset$.
% \remarque En seconde année, vous verrez une autre définition des voisinages
%   d'un élément $a\in\Rbar$ (ou $\C$). Les voisinages, tels qu'ils sont définis
%   dans ce cours, seront toujours des voisinages l'année prochaine, mais la
%   réciproque est fausse.
% sont des voisinages tels qu'ils seront définis l'année prochaine,
%   mais la réciproque est fausse.
% \end{remarques}

% \begin{definition}[utile=-3]
% Soit $a\in\R$. On dit qu'une partie $\mathcal{V}$ de $\R$ est~:
% \begin{itemize}
% \item un voisinage à droite de $a$ lorsqu'il existe $\epsilon>0$ tel que~:
%   \[\interof{a}{a+\epsilon} \subset \mathcal{V}\]
% \item un voisinage à gauche de $a$ lorsqu'il existe $\epsilon>0$ tel que~:
%   \[\interfo{a-\epsilon}{a} \subset \mathcal{V}\]
% \end{itemize}
% \end{definition}

% \begin{exos}
% \exo $\RPs$ est un voisinage à droite de 0.
% \end{exos}

% \begin{proposition}
% L'intersection d'un nombre fini de voisinages fondamentaux de $a\in\RbarouC$ est un voisinage fondamental.
% \end{proposition}

\begin{definition}[utile=-3]
On dit que $f:\mathcal{D}\to\K$ est \emph{définie au voisinage}
de $a\in\Rbar$ lorsque, pour tout voisinage $\mathcal{V}$ de $a$, $\mathcal{D}\cap\mathcal{V}\neq\emptyset$.
\end{definition}

\begin{proposition}[utile=-3]
Une fonction $f:\mathcal{D}\to\K$ est définie au voisinage de $a\in\Rbar$ si et seulement si il existe une
suite d'éléments de $\mathcal{D}$ qui tend vers $a$.
\end{proposition}

\begin{preuve}
On va faire le cas où $a\in \R$, les autres se traitent de la même façon.
\begin{itemize}
\item[$\bullet$] Si $f$ est définie au voisinage de $a$, alors pour tout voisinage $\mathcal{V}$ de $a$, $\mathcal{D}\cap\mathcal{V}\neq\emptyset$. En particulier, pour chaque $n\in \N$, on peut construire $u_n \in \mathcal{D}\cap \interf{a-\frac{1}{n}}{a+\frac{1}{n}}$. $(u_n)$ est une suite d'éléments de $\mathcal{D}$ qui converge vers $a$.
\item[$\bullet$] Réciproquement, s'il existe une suite $(u_n) \in \mathcal{D}^\N$ qui converge vers $a$, alors pour chaque voisinage $\mathcal{V}$ de $a$, il existe $n\in N$ tel que $\forall n \geq N$, $u_n \in \mathcal{V}$. En particulier $u_N \in \mathcal{D}\cap\mathcal{V}$ qui est donc différent du vide.
\end{itemize}
\end{preuve}

\begin{definition}[utile=-3]
On dit que $f:\mathcal{D}\to\K$ vérifie la propriété $\mathcal{P}$ \emph{au voisinage de $a\in\Rbar$}
lorsqu'il existe un voisinage $\mathcal{V}$ de $a$ tel que $f$ vérifie la propriété $\mathcal{P}$ sur $\mathcal{D}\cap\mathcal{V}$.
% \begin{itemize}
% \item On dit que $f$ est définie au voisinage de $a$ lorsque tout voisinage
%   de $a$ admet une intersection non vide avec $\mathcal{D}$. 
% \item On dit que $f$ vérifie la propriété $\mathcal{P}$ au voisinage de $a$
%   lorsqu'il existe un voisinage $\mathcal{V}$ de $a$ tel que la restriction
%   de $f$ à $\mathcal{D}\cap\mathcal{V}$ vérifie la propriété $\mathcal{P}$.
% \end{itemize}
\end{definition}

%% Remarque :
%% 1) Si f vérifie une propriété P_1  au voisinage de a et vérifie une
%%    propriété P_2 au voisinage de a, alors f vérifie la propriété P_1 et P_2
%%    au voisinage de a.
%% 2) De même, si a\in\R, on dit que f est définie au voisinage à gauche
%%    (resp. à droite) de a lorsque tout voisinage à gauche (resp. à droite) de
%%    a rencontre \mathcal{D}.
%%


\begin{remarques}
\remarque La fonction $\sin$ est croissante sur $\interf{-\pi/2}{\pi/2}$. Elle
  est donc croissante au voisinage de 0.
\remarque Une fonction $f:\mathcal{D}\to\K$ est bornée au voisinage de $a\in\R$ si et seulement si
   \[\exists \eta>0 \qsep \exists M\in\RP \qsep \forall x\in\mathcal{D}
     \qsep \abs{x-a}\leq\eta \implique \abs{f(x)}\leq M.\]
  Une fonction $f:\mathcal{D}\to\K$ est bornée au voisinage de $+\infty$ si et seulement si
   \[\exists m\in\R \qsep \exists M\in\RP \qsep \forall x\in\mathcal{D}
     \qsep x\geq m \implique \abs{f(x)}\leq M.\]
% \exo La fonction $\appli{f}{\RPs}{\R}{x}{1/x}$ est bornée au voisinage de
%   $+\infty$ mais n'est pas bornée sur son ensemble de définition.
\end{remarques}


\begin{definition}[utile=-3]
On dit qu'une propriété $\mathcal{P}$ est \emph{locale en $a\in\Rbar$} lorsque, quelles que
soient les fonctions $f:\dom\to\K$ et $g:\dom\to\K$ définies au voisinage de $a$, si il
existe un voisinage $\mathcal{V}$ de $a$ tel que
\[\forall x\in\mathcal{D}\cap\mathcal{V} \qsep g(x)=f(x)\]
alors $\mathcal{P}(f)$ est vrai si et seulement si $\mathcal{P}(g)$ est vrai.
\end{definition}


\begin{definition}
On dit qu'un élément $a\in \mathcal{D}$ est \emph{intérieur} à $\mathcal{D}$ lorsqu'il existe un voisinage
$\mathcal{V}$ de $a$ tel que $\mathcal{V}\subset\mathcal{D}$. 
\end{definition}

\begin{remarqueUnique}
\remarque Si $I$ est un intervalle, un élément $a\in I$ est intérieur à $I$ si et seulement si ce n'est pas
  une de ses extrémités.
\end{remarqueUnique}
% \begin{remarques}
% \remarque Toute personne comprenant les subtilités de cette définition gagnera
%   une sucette géante.
% \end{remarques}
%% Remarque :
%% 1) Locale à gauche, à gauche au sens large (dérivation)

% \subsection{Ouvert}

% \begin{preuve}
% \begin{itemize}
% \item[$\bullet$] On suppose que $\alpha=\sup A$. Pour chaque $n\geq 1$, $s-1/n$ ne majore pas $A$, on peut donc fixer $u_n\in A$ tel que $s-1/n<u_n(\leq s)$. Ainsi, $\forall n\in \Ns, |u_n-s|\leq 1/n$. La suite $u$ ainsi construite converge donc vers $s$.
% \item[$\bullet$] Réciproquement, si $\alpha$ est un majorant de $A$. Soit $\epsilon>0$. Il existe $N\in \N$ tel que $\forall n\geq N$, $|u_n-\alpha|\leq \epsilon$. En particulier $\alpha-\epsilon \leq u_N$. Donc $\alpha$ est bien la borne supérieure de $A$ d'après la caractérisation de la borne supérieure.
% \end{itemize}
% \end{preuve}

% \begin{remarqueUnique}
% \remarque On a bien évidemment une proposition similaire avec la borne inférieure.
% \end{remarqueUnique}


% \begin{definition}[utile=-3]
% On dit qu'une partie $O$ de $\R$ est un \emph{ouvert} lorsque c'est une réunion
% d'intervalles ouverts.
% \end{definition}

% \begin{remarqueUnique}
% \remarque Un intervalle $I$ est ouvert si et seulement si ses extremités ne sont pas des
%   éléments de $I$.
% \end{remarqueUnique}

% \begin{definition}[utile=-3]
% On dit qu'un réel $x_0$ est \emph{intérieur} à une partie $A$ de $\R$ lorsqu'il existe
% $\eta>0$ tel que $[x_0-\eta,x_0+\eta]\subset A$.
% \end{definition}

% \begin{remarques}
% \remarque Si $x_0$ est intérieur à $A$, alors $x_0\in A$. Cependant, la
%   réciproque est fausse. Par exemple $0\in\ens{0}$ mais 0 n'est pas intérieur
%   à $\ens{0}$.
% \remarque Si $A$ est un ouvert, tous les points de $A$ sont intérieurs à $A$.
% \end{remarques}

% \begin{remarques}
% \remarque Si la fonction $f$ est définie sur un intervalle $I$, $f$ est définie
%   au voisinage de tous les points de $I$ et de ses bornes.
% \end{remarques}

% \subsection{Théorème de \nom{Bolzano-Weierstrass}}




\section{Limite}

\subsection{Définition, propriétés élémentaires}

% \begin{definition}[utile=-3]
% Soit $f$ une fonction définie au voisinage de $a\in\Rbar$ et $l\in\RbarouC$. On
% dit que $f(x)$ tend vers $l$ lorsque $x$ tend vers $a$ lorsque pour tout
% voisinage $\mathcal{W}$ de $l$, il existe un voisinage $\mathcal{V}$ de $a$ tel
% que~:
% \[\forall x\in\mathcal{D}\cap\mathcal{V} \qsep f(x)\in\mathcal{W}\]
% Si tel est le cas, on note~:
% \[f(x)\tendvers{x}{a} l\]
% La propriété \og tend vers $l$ lorsque $x$ tend vers $a$ \fg est locale en $a$.
% \end{definition}

\begin{definition}[utile=-3]
Soit $f:\mathcal{D}\to\K$ une fonction définie au voisinage de $a\in\Rbar$ et $l\in\K$.
On dit que $f(x)$ \emph{tend} vers $l$ lorsque $x$ tend vers $a$ et on note
\[f(x)\tendvers{x}{a} l\]
lorsque, quel que soit le voisinage $\mathcal{W}$ de $l$, il existe un voisinage
$\mathcal{V}$ de $a$ tel que
\[\forall x\in\dom\qsep x\in\mathcal{V}\implique f(x)\in\mathcal{W}.\]
La propriété \og tend vers $l$ lorsque $x$ tend vers $a$ \fg est locale en $a$.
% lorsque pour tout
% voisinage $\mathcal{W}$ de $l$, il existe un voisinage $\mathcal{V}$ de $a$ tel
% que~:
% \[\forall x\in\mathcal{D}\cap\mathcal{V} \qsep f(x)\in\mathcal{W}\]
% Si tel est le cas, on note~:

% La propriété \og tend vers $l$ lorsque $x$ tend vers $a$ \fg est locale en $a$.


% Alors~:
% \[f(x)\tendvers{x}{a}l\]
% si et seulement si~:

\end{definition}

\begin{remarqueUnique}
\remarque En pratique, on utilisera les caractérisations suivantes~:
\begin{itemize}
  \item Pour $a\in\R$ et $l\in\K$
    \[\forall \epsilon>0 \qsep \exists \eta>0 \qsep \forall x\in\mathcal{D}
      \qsep \abs{x-a}\leq\eta \implique \abs{f(x)-l}\leq\epsilon.\]
  \item Pour $a\in\R$ et $l=-\infty$
    \[\forall M\in\R \qsep \exists \eta>0 \qsep \forall x\in\mathcal{D}
      \qsep \abs{x-a}\leq\eta \implique f(x)\leq M.\]
  \item Pour $a\in\R$ et $l=+\infty$
    \[\forall m\in\R \qsep \exists \eta>0 \qsep \forall x\in\mathcal{D}
      \qsep \abs{x-a}\leq\eta \implique f(x)\geq m.\]
  \item Pour $a=-\infty$ et $l\in\K$
    \[\forall \epsilon>0 \qsep \exists A\in\R \qsep \forall x\in\mathcal{D}
      \qsep x\leq A \implique \abs{f(x)-l}\leq\epsilon.\]
  \item Pour $a=-\infty$ et $l=-\infty$
    \[\forall M\in\R \qsep \exists A\in\R \qsep \forall x\in\mathcal{D}
      \qsep x\leq A \implique f(x)\leq M.\]
  \item Pour $a=-\infty$ et $l=+\infty$
    \[\forall m\in\R \qsep \exists A\in\R \qsep \forall x\in\mathcal{D}
      \qsep x\leq A \implique f(x)\geq m.\]
  \item Pour $a=+\infty$ et $l\in\K$
    \[\forall \epsilon>0 \qsep \exists B\in\R \qsep \forall x\in\mathcal{D}
      \qsep x\geq B \implique \abs{f(x)-l}\leq\epsilon.\]
  \item Pour $a=+\infty$ et $l=-\infty$
    \[\forall M\in\R \qsep \exists B\in\R \qsep \forall x\in\mathcal{D}
      \qsep x\geq B \implique f(x)\leq M.\]
  \item Pour $a=+\infty$ et $l=+\infty$
    \[\forall m\in\R \qsep \exists B\in\R \qsep \forall x\in\mathcal{D}
      \qsep x\geq B \implique f(x)\geq m.\]
  \end{itemize}
% \remarque On montre que $f(x)$ tend vers $l\in\RbarouC$ lorsque $x$ tend
%   vers $a\in\Rbar$ si et seulement si pour tout voisinage $\mathcal{W}$ de $l$,
%   il existe un voisinage $\mathcal{V}$ de $a$ tel que
%   \[\forall x\in\mathcal{D}\cap\mathcal{V} \qsep f(x)\in\mathcal{W}.\]
\end{remarqueUnique}

\begin{exoUnique}
\exo Soit $f:\R\to\R$ une fonction croissante telle que
  \[f(n)\tendvers{n}{+\infty}+\infty.\]
  Montrer que $f(x)$ tend vers $+\infty$ lorsque $x$ tend vers $+\infty$.
  Que dire si $f$ n'est pas croissante~?
% \exo En remarquant que $(\sin x)/x$ tend vers 1 lorsque $x$ tend vers
%   0, donner la limite, si elle existe, de
%   \[\sum_{k=0}^n \sin\frac{k}{n^2}\]
%   lorsque $n$ tend vers $+\infty$.
  %% Utiliser le fait que (sin x-x)/x^2 tend vers 0 lorsque x tend vers 0
\end{exoUnique}


\begin{proposition}[utile=-3]
Soit $f:\mathcal{D}\to\K$ une fonction définie au voisinage de $a\in\Rbar$ et $l\in\K$.
Alors
\[f(x)\tendvers{x}{a} l\]
si et seulement si, pour toute suite $\p{u_n}$ d'éléments de $\mathcal{D}$
tendant vers $a$, la suite $\p{f\p{u_n}}$ tend vers $l$.
\end{proposition}

\begin{preuve}
C'est la CARACTERISATION SEQUENTIELLE DE LA LIMITE :
On va faire la preuve avec $a\in \R$ et $\ell \in \R$.
\begin{itemize}
\item[$\bullet$] Soit $\epsilon >0$. Si $f(x)\tendvers{x}{a} \ell$, il existe $\eta>0$ tel que $\forall x \in \mathcal{D}_f, \abs{x-a}\leq \eta \Longrightarrow \abs{f(x)-\ell}\leq \epsilon$. Considérons alors $(u_n)$ une suite d'éléments de $\mathcal{D}_f$ qui tend vers $a$. Il existe donc $N\in \N$ tel que $\forall n \geq N$, $\abs{u_n-a}\leq \eta$. Mais alors, $\abs{f(u_n)-\ell}\leq \epsilon$, ce qui prouve le sens gauche-droite.
\item[$\bullet$] Réciproquement, on suppose que pour toute suite $\p{u_n}$ d'éléments de $\mathcal{D}_f$
tendant vers $a$, la suite $\p{f\p{u_n}}$ tend vers $\ell$. On raisonne par l'absurde, c'est-à-dire qu'on suppose qu'il existe $\epsilon>0$ tel que $\forall \eta>0, \exists x \in \mathcal{D}_f$, on ait $\abs{x-a}\leq \eta$ et $\abs{f(x)-\ell}>\epsilon$.\\
Soit $n\in \N$, avec $\eta=\frac{1}{n}$, il existe $x_n\in \mathcal{D}_f$ tel que $\abs{x_n-a}\leq \dfrac{1}{n}$ et $\abs{f(x_n)-\ell}>\epsilon$. $(x_n)$ converge vers $a$ donc par hypothèse, $(f(x_n))$ converge vers $\ell$, d'où $0\geq \epsilon$ par passage à la limite, ce qui fournit une contradiction.
\end{itemize}

\end{preuve}

\begin{remarques}
\remarque Si $f$ est définie en $a\in\R$ et admet une limite en $a$, cette
  limite est $f(a)$.
\exo Cette proposition est utile pour prouver qu'une fonction $f$ n'a pas
  de limite en $a$. Pour cela, il suffit de trouver deux suites
  $\p{u_n}$ et $\p{v_n}$ telles que
  \[u_n\tendvers{n}{+\infty} a \et v_n\tendvers{n}{+\infty} a.\]
  et telles que les suites de terme général $f\p{u_n}$ et $f\p{v_n}$ aient
  des limites distinctes.
%   Par exemple $x\mapsto 1/x$ n'a pas de limite en 0, et $x\mapsto \sin x$ n'a
%   pas de limite en $+\infty$.
\end{remarques}

\begin{exos}
\exo Soit $f:\R\to\R$ une fonction périodique admettant une limite finie en
  $+\infty$. Montrer que $f$ est constante.
  \begin{sol}
  $(f(x+nT))$ tend vers $\ell$, mais cette suite est aussi égale à $f(x)$ donc $\ell=f(x)$ et ce peu importe $x$.
  \end{sol}
\exo Montrer que la fonction d'expression $\sin\p{\frac{1}{x}}$ n'a pas de
  limite en 0.
  \begin{sol}
  Prendre $u_n=\dfrac{1}{\frac{\pi}{2}+2n\pi}$ et $v_n=\dfrac{1}{-\frac{\pi}{2}+2n\pi}$.
  \end{sol}
\end{exos}

\begin{proposition}[utile=-3]
Si $f$ admet une limite $l\in\RbarouC$ en $a\in\Rbar$, alors cette limite
est unique. Si tel est le cas, on écrit
\[\lim_{x\to a} f(x)=l.\]
\end{proposition}

\begin{preuve}
\'Etant donné qu'on a démontré les résultats sur les suites, on va utiliser la caractérisation séquentielle de la limite pour obtenir chacun des résultats ci-dessous. Par exemple, pour celui-ci :
On considère $a\in \Rbar$. On suppose qu'il existe $\ell_1, \ell_2$ tels que $f(x)\tendvers{x}{a} \ell_1$ et $f(x)\tendvers{x}{a} \ell_2$. Puisque $f$ est définie au voisinage de $a$, il existe une suite $(u_n)$ d'éléments de $\mathcal{D}_f$ telle que $u_n\tendvers{n}{+\infty}{a}$ donc $f(u_n)\tendvers{n}{+\infty}{\ell_1}$ et $f(u_n)\tendvers{n}{+\infty}{\ell_2}$. D'où : $\ell_1=\ell_2$ par unicité de la limite sur les suites.
\end{preuve}

\begin{proposition}[utile=-3]
Soit $f$ une fonction définie au voisinage de $a\in\Rbar$
et $l\in\K$. On suppose que
\[f(x)\tendvers{x}{a} l.\]
Alors
\[\abs{f(x)}\tendvers{x}{a} \abs{l} \et
  \conj{f}(x)\tendvers{x}{a} \conj{l}.\]
\end{proposition}

\begin{proposition}[utile=-3]
Soit $f$ une fonction complexe définie au voisinage de $a\in\Rbar$ et $l\in\C$.
Alors
\[f(x)\tendvers{x}{a} l \quad\ssi\quad \cro{\Re\cro{f(x)}\tendvers{x}{a}\Re(l) \et
  \Im\cro{f(x)}\tendvers{x}{a}\Im(l)}.\]
\end{proposition}

\begin{proposition}[utile=-3]
\begin{itemize}
\item Les théorèmes usuels portant sur les combinaisons linéaires, les produits
  et les quotients de limites de suites restent vrais pour les fonctions.
\item Soit $f$ une fonction définie au voisinage de $a\in\Rbar$ tendant vers $l_1\in\Rbar$
  en $a$ et $g$ une fonction définie au voisinage de $l_1$ tendant vers
  $l_2\in\RbarouC$ en $l_1$. Si $g\circ f$ est définie au voisinage de $a$, alors
  \[\p{g\circ f}(x) \tendvers{x}{a} l_2.\]
\end{itemize}
\end{proposition}

\begin{preuve}
Pour la composition  :
Soit $(u_n)$ une suite d'éléments de $\mathcal{D}_f$ telle que $u_n\tendvers{n}{+\infty}a$, on a $f(x)\tendvers{x}{a}\ell_1$ donc $f(u_n)\tendvers{n}{+\infty}\ell_1$. De plus, $g(x)\tendvers{x}{\ell_1}\ell_2$ donc $g(f(u_n))\tendvers{n}{+\infty}\ell_2$ et ce peu importe la suite qui tend vers $a$ donc $g(f(x))\tendvers{x}{a} \ell_2$.
\end{preuve}

\begin{remarques}
\remarque Comme pour les suites, la somme d'une fonction admettant une limite
  finie en $a$ et d'une fonction n'admettant pas de limite en $a$ n'admet pas de
  limite en $a$.
  Les autres théorèmes de ce type sont souvent faux;
  par exemple, il est possible qu'une fonction $f$ n'admette pas de limite en
  $a$ bien que $g\circ f$ admette une limite en $a$.
\remarque Soit $f$ et $g$ sont deux fonctions réelles définies sur
  $\mathcal{D}$. On définit les fonctions $\sup\p{f,g}$ et $\inf\p{f,g}$ par
  \[\forall x\in\mathcal{D} \qsep \sup\p{f,g}(x)\defeq\max\p{f(x),g(x)} \et
    \inf\p{f,g}(x)\defeq\min\p{f(x),g(x)}.\]
  Si $f$ et $g$ admettent pour limites respectives $l_f$ et $l_g\in\R$ en
  $a\in\Rbar$, alors
  \[\sup\p{f,g}(x)\tendvers{x}{a}\max\p{l_f,l_g} \et 
    \inf\p{f,g}(x)\tendvers{x}{a}\min\p{l_f,l_g}.\]
\end{remarques}

% \begin{exoUnique}
% \exo Déterminer la limite, si elle existe de
%   \[\frac{\sqrt{1+\sin x}-\sqrt{1-\sin x}}{x}\] en 0.
% \end{exoUnique}
% \begin{sol}
% On passe à la quantité conjuguée et le tout tend vers $1$.
% \end{sol}

\subsection{Limite et ordre sur $\R$}

\begin{proposition}[utile=-3]
Soit $f:\mathcal{D}\to\K$ une fonction admettant une limite finie en $a$.
Alors $f$ est bornée au voisinage de $a$.  
\end{proposition}

\begin{preuve}
On écrit la définition de la limite avec $\epsilon=1$ puis I.T.
\end{preuve}

\begin{proposition}[utile=-3]
Soit $f:\mathcal{D}\to\R$ une fonction telle que $f(x)$ tend vers $l\in\R$ lorsque
$x$ tend vers $a\in\Rbar$.
\begin{itemize}
\item Si $f$ est majorée par $M\in\R$, alors $l\leq M$.
\item Si $f$ est minorée par $m\in\R$, alors $l\geq m$.
\end{itemize}
\end{proposition}

\begin{preuve}
Cas majoré : On prend une suite $(u_n)$ d'éléments de $\mathcal{D}_f$ qui converge vers $a$. Alors $\forall n \in \N$, $f(u_n)\leq M$ d'où par passage à la limite (SUR LES SUITES), $\ell\leq M$.

\end{preuve}

%% Remarque :
%% 1) Si f est croissante sur R et f(x) -> l_1 (en  -infini) et
%%    f(x) -> l_2 (en +infini), alors l_1 <= f(x) <= l_2
%% 2) Si f est strictement croissante et f(x) -> l_1 (en  -infini) et
%%    f(x) -> l_2 (en +infini), alors l_1 < f(x) < l_2

\begin{proposition}[utile=-3]
Soit $f:\mathcal{D}\to\R$ une fonction telle que $f(x)$ tend vers $l\in\Rbar$ lorsque
$x$ tend vers $a\in\Rbar$.
\begin{itemize}
\item Si $M$ est un réel tel que $l<M$, il existe un voisinage $\mathcal{V}$
  de $a$ tel que
  \[\forall x\in\mathcal{D} \qsep x\in\mathcal{V} \implique f(x)\leq M.\]
\item Si $m$ est un réel tel que $l>m$, il existe un voisinage $\mathcal{V}$
  de $a$ tel que
  \[\forall x\in\mathcal{D} \qsep x\in\mathcal{V} \implique f(x)\geq m.\]
\end{itemize}
\end{proposition}

\begin{preuve}
Pour changer, faisons cette preuve avec $a=+\infty$ et $\ell \in \R$.
On pose $\epsilon=M-\ell >0$. Alors, il existe $m\in \R$ tel que $\forall x\geq m, \abs{f(x)-\ell}\leq \epsilon$, i.e. $f(x)\leq \ell+\epsilon=M$.
\end{preuve}

\begin{remarques}
\remarque En pratique, il conviendra d'expliciter les voisinages. Par
  exemple, si $f(x)$ tend vers $l<M$ lorsque $x$ tend vers $a\in\R$, il
  existe $\epsilon>0$ tel que
  \[\forall x\in\mathcal{D} \qsep \abs{x-a}\leq\epsilon \implique
    f(x)\leq M.\]
\remarque Si une fonction complexe $f$ admet une limite $l\in\C$
  non nulle en $+\infty$, il existe $m\in\R$ tel que
  \[\forall x\in\mathcal{D} \qsep x\geq m \implique f(x)\neq 0.\] 
\end{remarques}

\begin{theoreme}[utile=-3, nom={Théorème des gendarmes}]
Soit $f$, $g$ et $h:\mathcal{D}\to\R$ telles que
\[\forall x\in\mathcal{D} \qsep f(x)\leq g(x)\leq h(x).\]
On suppose que $f$ et $h$ admettent la même limite finie $l\in\R$ en $a\in\Rbar$.
Alors
\[g(x) \tendvers{x}{a} l.\]
\end{theoreme}

\begin{preuve}
Soit $(u_n)$ une suite d'éléments de $\mathcal{D}$ qui tend vers $a$. On a d'une part $\forall n \in \N, f\p{u_n}\leq g\p{u_n}\leq h\p{u_n}$ et d'autre part par caractérisation séquentielle de la limite, $(f(u_n))$ et $(h(u_n))$ tendent vers $\ell$. Donc $(g(u_n))$ tend vers $\ell$ par théorème des gendarmes sur les suites, et ceci étant valable pour n'importe quelle suite $(u_n)$ convergeant vers $a$, on en déduit que $g(x) \tendvers{x}{a} \ell$ par caractérisation séquentielle de la limite.
\end{preuve}


\begin{proposition}[utile=-3]
Soit $f$ et $g:\mathcal{D}\to\R$ telles que
\[\forall x\in\mathcal{D} \qsep f(x)\leq g(x)\]
et $a\in\Rbar$.
\begin{itemize}
\item Si $f(x)\tendvers{x}{a} +\infty$, alors $g(x)\tendvers{x}{a}+\infty$.
\item Si $g(x)\tendvers{x}{a} -\infty$, alors $f(x)\tendvers{x}{a}-\infty$.
\end{itemize}
\end{proposition}

\begin{preuve}
Par caractérisation séquentielle.
\end{preuve}

\begin{proposition}[utile=-3]
Soit $f:\mathcal{D}\to\K$ une fonction, $a\in\Rbar$, $l\in\K$ et $g$ une
fonction réelle positive telle que
\begin{itemize}
\item $\forall x\in\mathcal{D} \qsep \abs{f(x)-l}\leq g(x)$,
\item $g(x)\tendvers{x}{a} 0$.
\end{itemize}
Alors
\[f(x)\tendvers{x}{a} l.\]
\end{proposition}

\begin{preuve}
Par caractérisation séquentielle.
\end{preuve}

\begin{exoUnique}
\exo Déterminer la limite, si elle existe, de la fonction d'expression
  $\frac{x}{2+\sin\p{\frac{1}{x}}}$ en 0.
\begin{sol}
On le coince entre $x/3$ et $x$ si $x\geq 0$ et entre $x$ et $x/3$ si $x\leq 0$. Donc $\forall x \in \Rs, -\abs{x}\leq \frac{x}{2+\sin\p{\frac{1}{x}}} \leq \abs{x}$ + thm des gendarmes.
\end{sol}
% \exo On a
%   \[\integ{0}{1}{\frac{e^{-xt}}{1+t^2}}{t}\tendvers{x}{+\infty}0\]
\end{exoUnique}

\subsection{Limite à gauche, à droite}

\begin{definition}[utile=-3]

\begin{itemize}
\item Soit $f$ une fonction définie au voisinage à gauche de $a\in\R$ et
  $l\in\RbarouC$. On dit que $f$ admet $l$ pour \emph{limite à gauche} en $a$ lorsque
  la restriction de $f$ à $\mathcal{D}\cap\intero{-\infty}{a}$ admet $l$ pour
  limite en $a$. Si tel est le cas, on note
  \[f(x)\tendversgp{x}{a} l.\]
  La propriété \og tend vers $l$ lorsque $x$ tend vers $a$ par la gauche \fg est
  locale à gauche en $a$.
\item Soit $f$ une fonction définie au voisinage à droite de $a\in\R$ et
  $l\in\RbarouC$. On dit que $f$ admet $l$ pour \emph{limite à droite} en $a$ lorsque
  la restriction de $f$ à $\mathcal{D}\cap\intero{a}{+\infty}$ admet $l$ pour
  limite en $a$. Si tel est le cas, on note
  \[f(x)\tendversdp{x}{a} l.\]
  La propriété \og tend vers $l$ lorsque $x$ tend vers $a$ par la droite \fg est
  locale à droite en $a$.
\end{itemize}
\end{definition}

\begin{remarqueUnique}
\remarque On définit de même les notions de limite à gauche au sens large, de limite à
  droite et de limite épointée (pour $x\neq a$). On écrit alors respectivement
  \[f(x)\tendversg{x}{a} l, \qquad
    f(x)\tendversd{x}{a} l, \qquad
    f(x)\tendversp{x}{a} l.\]
\end{remarqueUnique}


\begin{proposition}[utile=-3]
Soit $f$ une fonction définie au voisinage de $a$ et $l\in\RbarouC$. Alors
$f(x)$ tend vers $l$ lorsque $x$ tend vers $a$ si et seulement si, les objets
ci-dessous susceptibles d'avoir un sens
\[\lim_{\substack{x\to a\\x<a}} f(x), \quad f(a) \et
  \lim_{\substack{x\to a\\x>a}} f(x)\]
existent et sont égaux à $l$.
\end{proposition}

\begin{preuve}
Sens gauche-droite immédiat.
Pour le sens droite-gauche, écrire les deux limites avec $\eta_1$ et $\eta_2$ puis prendre le $\min$.
\end{preuve}

\begin{remarques}
\remarque Par exemple, si $f$ est une fonction définie sur $\Rs$, et $l\in\Rbar$,
  alors $f(x)$ tend vers $l$
  lorsque $x$ tend vers 0 si et seulement si
  \[f(x)\tendversgp{x}{0}l \quad\text{et}\quad f(x)\tendversdp{x}{0} l.\]
  De même, si $f$ est une fonction définie sur $\RP$ et $l\in\R$, alors $f(x)$ tend vers $l$
  lorsque $x$ tend vers 0 si et seulement si
  \[f(0)=l \quad\text{et}\quad f(x)\tendversdp{x}{0} l.\]
\remarque On a bien entendu des théorèmes similaires faisant intervenir les limites au sens
  large. Par
  exemple, si $f$ est définie sur $\R$ et $a\in\R$, alors $f(x)$ tend vers $l$ lorsque $x$
  tend vers $a$ si et seulement si
  \[f(x)\tendversg{x}{a}l \quad\text{et}\quad f(x)\tendversdp{x}{a} l.\]
\end{remarques}

\begin{exoUnique}
\exo Soit $f$ la fonction définie sur $\R$ par
  \[\forall x\in\R \qsep f(x)\defeq
    \begin{cases}
    0 & \text{si $x\leq 0$}\\
    \e^{-1/x} & \text{si $x>0$}.
    \end{cases}\]
  Montrer que $f(x)\tendvers{x}{0} 0$.
\end{exoUnique}



\begin{theoreme}[nom={Théorème de la limite monotone}]
Soit $f:I\to\R$ une fonction croissante sur un intervalle $I$.
\begin{itemize}
\item Si $a\in I$ n'est pas une borne de $I$, $f$ admet une limite finie 
  à gauche (au sens strict) et une limite finie à droite (au sens strict) en $a$. De plus
  \[\lim_{\substack{x\to a\\x<a}} f(x) \leq f(a) \leq
    \lim_{\substack{x\to a\\x>a}} f(x).\]
\item Si $a$ est la borne supérieure de $I$, $f$ admet une limite en $a$.
  Cette limite est finie si $f$ est majorée, et est égale à $+\infty$ sinon.
\item Si $a$ est la borne inférieure de $I$, $f$ admet une limite en $a$.
  Cette limite est finie si $f$ est minorée, et est égale à $-\infty$ sinon.
\end{itemize}
\end{theoreme}

\begin{preuve}
Nous allons traiter seulement certains cas (plus simples) pour s'éviter des détails techniques superflus :
\begin{itemize}
\item[$\bullet$] Traitons le cas où $f:\R\mapsto \R$ est une fonction croissante et $a\in \R$. On veut montrer que $f$ admet une limite à gauche en $a$.
\begin{itemize}
\item Posons $A=\set{f(x), x<a}$. $A$ est non vide ($f(a-1)\in A$) et majorée par $f(a)$ par croissance. $A$ admet donc une borne supérieure que l'on note $\ell_1$. $f(a)$ étant un majorant, on a $\ell_1\leq f(a)$.
\item Soit $\epsilon>0$. Alors, il existe $y_0\in A$ tel que $y_0\geq \ell_1-\epsilon$, i.e. il existe $x_0\in ]-\infty;a[$ tel que $f(x_0)\geq \ell_1-\epsilon$. Posons alors $\eta=a-x_0>0$. Pour tout $x\in [a-\eta;a[, a-\eta=x_0\leq x$ donc $f(x_0)\leq f(x)$ par croissance de $f$ donc $\ell_1-\epsilon\leq f(x)$. Et $f(x)\in A$ donc $f(x)\leq \ell_1$. Finalement, pour tout $x\in [a-\eta;a[, \abs{f(x)-\ell_1}\leq \epsilon$.\\
\item $\ell_1$ est donc bien la limite à gauche de $f$ en $a$ et $\ell_1\leq f(a)$.
\item On montre de même qu'il existe une limite à droite en $a$ et qu'elle est supérieure à $f(a)$.
\end{itemize}
\item[$\bullet$] Traitons le cas où $f:\RPs\mapsto \R$ est une fonction croissante et où $a=0$. Montrons que $f$ admet une limite à droite en $0$ qui soit ou bien finie (si elle est minorée) ou bien $-\infty$.\\
On pose $A=\set{f(x), x\in \RPs}$.
\begin{itemize}
\item Si $A$ n'est pas minorée, montrons qu'alors la limite à droite vaut $-\infty$. Soit $M\in \R$. Puisque $A$ n'est pas minorée, il existe $x_0\in \RPs$ tel que $f(x_0)<M$. Mais alors, par croissance, pour tout $x\in ]0;x_0], f(x)(\leq f(x_0))\leq M$, ce qui prouve le résultat souhaité.
\item Si $A$ est minorée, comme $A$ est non vide, il existe une borne inférieure qu'on note $\ell$. Soit $\epsilon>0$, il existe $a_0\in A$ tel que $ a_0 \leq \ell+\epsilon$. Mais $\exists x_0\in \RPs$ tel que $a_0=f(x_0)$ et pour tout $x\in ]0;x_0]$ par croissance et parce-que $f(x)\in A$, on a $\ell \leq f(x) \leq \ell+\epsilon$ d'où $\abs{f(x)-\ell}\leq \epsilon$.
\end{itemize}
\end{itemize}
\end{preuve}

\begin{remarques}
\remarque Une proposition similaire existe pour les fonctions décroissantes.
\remarque Si $f:\intero{-\infty}{a}\to\R$ est une fonction croissante admettant une limite $l\in\R$ en $a\in\R\cup\ens{+\infty}$, alors
  \[\forall x\in\intero{-\infty}{a} \qsep f(x)\leq l.\]
  De plus, si $f$ est strictement croissante, alors
  \[\forall x\in\intero{-\infty}{a} \qsep f(x)< l.\]
\end{remarques}

\section{Continuité}

\subsection{Continuité ponctuelle}

\begin{definition}[utile=-3]
On dit qu'une fonction $f:\mathcal{D}\to\K$ est \emph{continue} en $x_0\in\mathcal{D}$ lorsque
\[f(x)\tendvers{x}{x_0} f\p{x_0}.\]
La propriété \og est continue en $x_0$ \fg est locale en $x_0$. On appelle
domaine de continuité de $f$ l'ensemble des $x_0\in\mathcal{D}$ en lesquels
$f$ est continue.
\end{definition}


\begin{remarques}
\remarque La continuité de $f$ en $x_0$ s'écrit 
  \[\forall \epsilon>0 \qsep \exists \eta>0 \qsep \forall x\in\mathcal{D}
    \qsep \abs{x-x_0}\leq\eta \implique \abs{f(x)-f\p{x_0}}\leq\epsilon.\]
\remarque Une fonction $f:\mathcal{D}\to\K$ est continue en $x_0\in\mathcal{D}$ si et seulement si
  elle admet une limite en $x_0$.
\begin{sol}
Sens gauche-droite évident. Si maintenant, il existe $\ell \in \Rbar$ tel que $f(x)\tendvers{x}{x_0} \ell$ alors $f$ étant définie en $x_0$, on obtient d'après la définition de la limite $\forall \epsilon>0, |f(x_0)-\ell|\leq \epsilon$ (étant donné que pour tout $\eta>0$, $|x_0-x_0|\leq \eta$) d'où $\ell=f(x_0)$.
\end{sol}
\remarque On dit qu'une fonction $f$ est \emph{continue à droite} en $x_0$ lorsque
  \[f(x)\tendversdp{x}{x_0} f\p{x_0}.\]
  De même, on définit la notion de \emph{continuité à gauche}. Une fonction est
  continue en $x_0$ si et seulement si elle est continue à droite et à gauche
  en $x_0$.
\remarque On dit qu'une fonction $f$ admet une \emph{discontinuité de première
  espèce} en $x_0$ lorsqu'elle admet des limites à droite et à gauche
  et lorsque l'une de ces limites est différente de $f(x_0)$.
  Par exemple, la fonction partie entière admet une discontinuité de première
  espèce en tout point $x_0\in\Z$.
% \remarque Toute fonction lipschitzienne est continue en tout point de son
%   domaine de définition.
%   \begin{sol}
%   Soit $k>0$. On rappelle qu'une fonction est dite $k$-lipschitzienne si $\forall (x,y)\in \mathcal{D}^2, \abs{f(x)-f(y)}\leq k\abs{x-y}$. Pour voir qu'elle est continue, pour $\epsilon>0$, il suffit de prendre $\eta=\frac{\epsilon}{k}$.
%   \end{sol}
% \remarque La fonction $f$ définie sur $\R$ par $f(x)\defeq\sin\p{1/x}$ si
%   $x\neq 0$ et $f\p{0}\defeq 0$ est discontinue en 0. Cette discontinuité n'est pas
%   une discontinuité de première espèce.
% \begin{sol}
% On parle alors de discontinuité essentielle ou discontinuité de deuxième espèce lorsque au moins une des deux limites à gauche et à droite n'existe pas ou est infinie.
% \end{sol}
\remarque Les fonctions valeur absolue, puissance (en particulier les puissances
  entières et les racines $n$-ièmes), $\ln$, $\exp$, les fonctions
  trigonométriques circulaires et hyperboliques, directes et réciproques sont
  continues en tout point de leur domaine de définition.
  \begin{sol}
La seule fonction "qui a un nom" qui admet des discontinuités (elles sont de première espèce) est la fonction partie entière.
\end{sol} 
\remarque Une fonction continue en $x_0$ est bornée au voisinage de $x_0$.
\end{remarques}

\begin{exoUnique}
\exo Soit $f:\R\to\R$ la fonction définie par
  \[\forall x\in\R\qsep f(x)\defeq\begin{cases}\sin\p{\frac{1}{x}} & \text{si $x\neq 0$}\\
    0 & \text{si $x=0$.}\end{cases}\]
  Montrer que $f$ est discontinue en 0 et que cette discontinuité n'est pas de première
  espèce.
\end{exoUnique}

\begin{definition}[utile=-3]
Soit $f:\mathcal{D}\to\K$ une fonction définie au voisinage d'un point $a\in\R$ n'appartenant pas
à $\mathcal{D}$. Lorsque $f(x)$ admet une limite finie $l\in\K$ lorsque $x$
tend vers $a$, on dit que $f$ est \emph{prolongeable par continuité} en $a$. La fonction
\[\dspappli{\hat{f}}{\mathcal{D} \cup\ens{a}}{\K}{x}{
  \begin{cases}
  f(x) & \text{si $x\neq a$}\\
  l      & \text{si $x=a$}  
  \end{cases}}\]
est alors appelée \emph{prolongement par continuité} de $f$ en $a$. C'est une fonction
continue en $a$.
\end{definition}
  

% \begin{applications}
% \application Soit $f:\RPs\to\R$ une fonction croissante telle que
%   $g:x\mapsto f(x)/x$ soit décroissante. Alors $f$ est continue sur $\RPs$.
% \end{applications}
\begin{proposition}[utile=-3]
Soit $f:\mathcal{D}\to\K$ une fonction et $x_0\in\mathcal{D}$. Alors $f$ est continue en $x_0$
si et seulement si pour toute suite $\p{u_n}$ d'éléments de $\mathcal{D}$
convergeant vers $x_0$
\[f\p{u_n} \tendvers{n}{+\infty} f\p{x_0}.\]
\end{proposition}

\begin{preuve}
C'est la caractérisation séquentielle de la limite.
\end{preuve}


\begin{remarques}
\remarque Soit $f$ une fonction continue sur l'intervalle $I$ et $\p{u_n}$
  une suite telle que
  \[\forall n\in\N \qsep u_{n+1}=f\p{u_n}\]
  Alors, si $\p{u_n}$ admet une limite $l\in\Rbar$, c'est une borne de
  $I$ ou un point fixe de $f$.
  \begin{sol}
  En effet, supposons que $\ell$ ne soit pas une borne de $I$.
  Alors $\ell \in I$. Puisque $f$ est continue en $I$, $f(u_n)\tendvers{n}{+\infty}f(\ell)$. Or $u_{n+1}\tendvers{n}{+\infty}\ell$. Par unicité de la limite $f(\ell)=\ell$.
  \end{sol}
\remarque Soit $f$ et $g$ deux fonctions continues ent tout point de $\R$. Si elles coïncident
  sur $\Q$, alors $f=g$.
  \begin{sol}
  Soit $x\in \R$, $\Q$ étant dense dans $\R$, il existe une suite $(u_n)$ de rationnels tendant vers $x$. Donc $\forall n\in \N, f(u_n)=g(u_n)$ et par passage à la limite $f(x)=g(x)$.
  \end{sol}
\remarque Cette proposition est utile pour prouver qu'une fonction $f$ n'est
  pas continue en $x_0$. Pour cela, il suffit de trouver une suite $\p{u_n}$
  convergeant vers $x_0$ telle que la suite de terme général $f\p{u_n}$ ait
  une limite différente de $f\p{x_0}$.
\end{remarques}

\begin{exos}
\exo Montrer que la fonction caractéristique de $\Q$ est discontinue en
  tout point.
  \begin{sol}
  Soit $x\in \Q$, $\R\setminus\Q$ étant dense dans $\R$, il existe une suite $(u_n)$ d'irrationnels tendant vers $x$. Alors $\forall n \in N, f(u_n)=0\tendvers{n}{+\infty}0\neq f(x)=1$, donc $f$ est discontinue en $x$.\\
  Soit $x\in \R\setminus\Q$, $Q$ étant dense dans $\R$, il existe une suite $(u_n)$ de rationnels tendant vers $x$. Alors $\forall n \in N, f(u_n)=1\tendvers{n}{+\infty}1\neq f(x)=0$, donc $f$ est discontinue en $x$.
  \end{sol}
\exo Quelles sont les fonctions $f:\R\to\R$, continues en tout point de $\R$, telles que
  \[\forall x,y\in\R \qsep f\p{x+y}=f(x)+f(y).\]
  \begin{sol}
  \begin{itemize}
  \item[$\bullet$] \underline{\textbf{Analyse :}} Soit $f$ continue de $\R$ dans $\R$ vérifiant $f(x+y)=f(x)+f(y), \forall (x,y)\in \R^2$.
  \begin{itemize}
  \item On montre d'abord par récurrence sur $n\in \N$ $H_n : "f(nx)=nf(x)"$.
  \item Soit $n\in \Z$ tel que $n\leq 0$. Alors $-n\in \N$. D'où $f(nx)+f(-nx)=f(0)=0$ donc $f(nx)=-f(-nx)=-(-nf(x))=nf(x)$. On a donc montré que $\forall n \in \Z, f(nx)=nf(x)$.
  \item Soit $r\in \Q$. Il existe $p\in \Z$ et $q\in \Ns$ tels que $r=p/q$. On a $pf(x)=f(px)=f(qrx)=qf(rx)$ d'où $f(rx)=p/qf(x)=rf(x)$.
  Finalement, $\forall r\in \Q, f(r)=rf(1)$.
  \item Soit $x\in \R$, il existe une suite $(r_n)$ de rationnels qui tend vers $x$ (par densité). Par continuité de $f$ en $x$, $r_nf(1)=f(r_n)\tendvers{n}{+\infty}f(x)$. Donc par unicité de la limite $f(x)=xf(1)$.
  \item En posant $\alpha=f(1)$ on a donc $f(x)=\alpha x$.
  \end{itemize}
  \item[$\bullet$] \underline{\textbf{Synthèse :}} Soit $\alpha \in \R$. On définit $f$ sur $\R$ par $f(x)=\alpha x$. Alors $f$ est continue sur $\R$ et $\forall (x,y)\in \R^2$, on a $f(x+y)=\alpha(x+y)=\alpha x +\alpha y=f(x)+f(y)$.
  \end{itemize}
  \end{sol}
\end{exos}

\begin{proposition}[utile=-3, nom={Théorèmes usuels}]
Soit $f$ et $g:\mathcal{D}\to\K$ deux fonctions continues en $x_0\in\mathcal{D}$. Alors
\begin{itemize}
\item Si $\lambda$, $\mu\in\K$, $\lambda f+\mu g$ est continue en $x_0$.
\item $fg$ est continue en $x_0$.
\item Si $g\p{x_0}\neq 0$, $g$ ne s'annule pas au voisinage de $x_0$ et
  $f/g$ est continue en $x_0$.
\end{itemize}
\end{proposition}

% \begin{remarques}
% \remarque Si $\mathcal{D}$ est une partie de $\R$, l'ensemble
%   $\mathcal{C}^0\p{\mathcal{D},\R}$ des fonctions continues de $\mathcal{D}$
%   dans $\R$ est une $\R$-algèbre.
% \end{remarques}

\begin{proposition}[utile=-3,nom={Théorèmes usuels}]
Soit $f:\mathcal{D}_f\to\R$ et $g:\mathcal{D}_g\to\K$ deux fonctions telles que
$g\circ f$ est défini au voisinage de $x_0\in\mathcal{D}_f$. Si $f$ est continue en
$x_0$ et $g$ est continue en $f\p{x_0}$, alors $g\circ f$ est continue en $x_0$.
\end{proposition}

\begin{preuve}
Les deux résultats précédents viennent directement des propriétés d'opérations sur les limites.
\end{preuve}

\begin{remarqueUnique}
% \remarque Les deux propositions précédentes sont regroupées sous la
%   dénomination de \og théorèmes usuels \fg.
\remarque La somme d'une fonction continue en $x_0$ et d'une fonction
  discontinue en $x_0$ est discontinue en $x_0$. Les autres propositions de ce
  type peuvent être fausses. Par exemple, si $f$ et $g$ sont les fonctions
  définies sur $\R$ par
  \[\forall x\in\R \qsep f(x)\defeq x \et g(x)\defeq
    \begin{cases}
    \sin\p{\frac{1}{x}} & \text{si $x\neq 0$}\\
    0 & \text{si $x=0$,}
    \end{cases}\]
  alors $f$ est continue en 0 et $g$ ne l'est pas. Pourtant $f\cdot g$ l'est.
  On retiendra que les réciproques des théorèmes usuels peuvent être
  fausses.
\end{remarqueUnique}

\begin{sol}
Pour voir que $fg$ est continue en $0$, $\abs{(fg)(x)}\leq |x|$.
\end{sol}

\begin{proposition}[utile=-3]
Soit $f:\mathcal{D}\to\K$ une fonction continue en $x_0\in\mathcal{D}$. Alors $\conj{f}$ et $\abs{f}$ sont
continues en $x_0$.  
\end{proposition}

%% Remarque
%% 1) Si f et g sont continues sup(f,g) et inf(f,g) le sont.

\begin{proposition}[utile=-3]
Soit $f:\mathcal{D}\to\C$ une fonction et $x_0\in\mathcal{D}$. Alors
\[\cro{\text{$f$ est continue en $x_0$}} \quad\ssi\quad
  \cro{\text{$\Re(f)$ et $\Im(f)$ sont continues en $x_0$}}.\]
\end{proposition}

\begin{remarqueUnique}
\remarque Si $f$ est une fonction à valeurs dans $\C$ continue en $x_0$, alors
  $\e^f$ est continue en $x_0$.
\end{remarqueUnique}
\begin{sol}
$\e^{f(x)}=\e^{\Re(f)(x)+\ii\Im(f(x))}=\cos(\Im(f(x)))\e^{\Re(f)(x)}+\ii\cos(\Im(f(x)))\e^{\Re(f)(x)}$ dont les parties réelle et imaginaire sont continues par TOC
\end{sol}

\begin{exos}
\exo Soit $f$ la fonction définie sur $\R$ par
  \[\forall x\in\R \qsep f(x)\defeq
    \begin{cases}
    \frac{\e^{\ii x}-1}{x} & \text{si $x\neq 0$}\\
    \ii & \text{si $x=0$.}
    \end{cases}\]
  Montrer que $f$ est continue en tout point de $\R$.
\end{exos}

\begin{sol}
Taux d'accroissement.
Sinon, pour voir la continuité en $0$, écrire pour $x\neq 0$, $$\frac{\e^{\ii x}-1}{x}=\frac{\cos(x)-1}{x}+i\frac{\sin(x)}{x}$$ donc $f(x)\tendversp{x}{0}f(0)$ d'où $f(x)\tendvers{x}{0}f(0)$ donc $f$ est continue en $0$.
\end{sol}

%% Exemple :
%% 1) Prolongement par continuité en 0 de sin x/x. On obtient le sinus cardinal.



\subsection{Continuité sur une partie}

\begin{definition}[utile=-3]
Soit $f:\dom\to\K$.
\begin{itemize}
\item On dit que $f$ est \emph{continue} lorsqu'elle est continue en tout point de
  $\dom$.
\item Si $A$ est une partie de $\dom$, on dit que $f$ est continue sur $A$ lorsque la
 restriction de $f$ à $A$ est continue.
\end{itemize}
\end{definition}

\begin{remarqueUnique}
\remarque Soit $f:\mathcal{D}\to\K$ et $A$ une partie de $\mathcal{D}$.
  Si $f$ est continue en tout point de $A$, alors $f$ est continue sur $A$. Cependant,
  la réciproque est fausse. En effet, si $f$ est la fonction de $\R$ dans $\R$ définie
  par
  \[\forall x\in\R\qsep f(x)\defeq\begin{cases}
    1 & \text{si $x\geq 0$}\\
    0 & \text{sinon}
  \end{cases}\]
  alors $f$ est continue sur $\RP$ mais n'est pas continue en 0.
\end{remarqueUnique}

\begin{definition}[utile=-3]
Soit $k\in\RP$. On dit qu'une fonction $f:\mathcal{D}\to\K$ est \emph{$k$-lipschitzienne} lorsque
\[\forall x,y\in \mathcal{D} \qsep \abs{f(x)-f(y)}\leq k\abs{x-y}.\] 
\end{definition}

\begin{exoUnique}
\exo Montrer que les fonctions \og valeur absolue\fg et \og sinus\fg sont
  1-lipschitziennes.
% \exo La fonction valeur absolue est $1$-lipschitzienne. En effet
%   \[\forall x,y\in\R \qsep \abs{\abs{x}-\abs{y}}\leq\abs{x-y}\]
% \exo La fonction $\sin$ est $1$-lipschitzienne.
\end{exoUnique}

% \begin{definition}
% On dit qu'une fonction $f:\mathcal{D}\to\RouC$ est \emph{lipschitzienne} lorsqu'il existe
% $M\geq 0$ tel que
% \[\forall x,y\in\mathcal{D}\qsep \abs{f(x)-f(y)}\leq M\abs{x-y}.\]
% \end{definition}

\begin{proposition}
Une fonction lipschitzienne est continue.
\end{proposition}

Dans la suite du cours d'analyse, un s'intéressera le plus souvent à des fonctions dont
le domaine de définition est une partie relativement simple de $\R$. Afin de formaliser cela,
une partie de $\R$ sera dite \emph{élémentaire} lorsque c'est une réunion d'un nombre fini
d'intervalles. Par exemple $\R$ et $\Rs=\intero{-\infty}{0}\cup\intero{0}{+\infty}$ sont des parties élémentaires
de $\R$ alors que $\Q$ n'en est pas une. Notons que cette définition est propre à ce cours.
Les parties élémentaires de $\R$ jouissent de nombreuses propriétés~:
$\emptyset$ et $\R$ sont des parties élémentaires,
une union finie de parties élémentaires est une partie élémentaire,
une intersection finie de parties élémentaires est une partie élémentaire,
le complémentaire d'une partie élémentaire est une partie élémentaire.
Autrement dit, tout ensemble construit à partir de parties élémentaires à l'aide d'un nombre
fini d'opérations est élémentaire.\\

On dit que deux intervalles $I$ et $J$ sont \emph{bien disjoints} lorsque $I\cup J$ n'est
pas un intervalle. En particulier deux intervalles bien disjoints sont disjoints. Mais
la réciproque est fausse comme le montre l'exemple de $I=\interf{0}{1}$ et
$J=\interof{1}{2}$ qui sont disjoints mais qui ne sont pas bien disjoints.
On montre
que si $A$ est une partie élémentaire de $\R$, à réordonnement près, il existe un unique
$n$-uplet $(I_1,\ldots,I_n)$ d'intervalles deux à deux biens disjoints tels que
$A=I_1\cup I_2\cup\cdots\cup I_n$. On dit que les $I_k$ sont les
\emph{composantes connexes} de $A$. Enfin, on dit que $a\in\Rbar$ est
une \emph{extrémité} de $A$ lorsque c'est l'extrémité d'un des $I_k$.\\

Si $\mathcal{D}$ est une partie élémentaire, une fonction $f:\mathcal{D}\to\K$ 
  est définie au voisinage de $a\in\Rbar$ si et seulement si $a$ est un élément
  ou une extrémité de $\mathcal{D}$.
  D'autre part, un élément $a\in\mathcal{D}$ est intérieur à
  $\mathcal{D}$ si et seulement si ce n'est pas
  une extrémité de $\mathcal{D}$.\\
  
\begin{exoUnique}
\exo Montrer que $\mathcal{D}\defeq\RP\setminus\ens{1}$ est une partie élémentaire. Si
  $f:\mathcal{D}\to\K$, déterminer
  \begin{itemize}
  \item L'ensemble des $a\in\Rbar$ pour lesquels $f$ est définie au voisinage de $a$.
  \item L'ensemble des $a\in\mathcal{D}$ intérieurs à $\mathcal{D}$.
  \end{itemize}
\end{exoUnique}
% \remarque Une partie élémentaire $\mathcal{D}$ est un ouvert si et seulement si
%   elle ne contient pas ses extrémités.
% \remarque Si $\mathcal{D}$ est une partie élémentaire de $\R$, les points intérieurs à
%   $\mathcal{D}$ sont les éléments de $\mathcal{D}$ qui n'en sont pas des extrémités.



\begin{proposition}
Soit $f:\dom\to\K$ une fonction définie sur un domaine \emph{élémentaire} et
$\dom=I_1\cup\cdots\cup I_n$ la décomposition de $\dom$ en composantes connexes. Alors $f$
est continue si et seulement si, pour tout $k\in\intere{1}{n}$, $f$ est continue
sur $I_k$.
\end{proposition}

\begin{proposition}
Soit $f:I\to\K$ une fonction définie sur un intervalle $I$ et $a\in I$.
\begin{itemize}
\item Si $f$ est continue sur $I\cap\intero{a}{+\infty}$, elle est continue en tout
  point de $I\cap\intero{a}{+\infty}$.
\item Si $f$ est continue sur $I\cap\interfo{a}{+\infty}$, elle est continue à droite
  en $a$ et en tout point de $I\cap\intero{a}{+\infty}$.
\end{itemize}
\end{proposition}


\subsection{Théorème des valeurs intermédiaires}

\begin{theoreme}[nom={Théorème des valeurs intermédiaires}]
Soit $f$ une fonction réelle continue sur $\interf{a}{b}$. Si
  $y_0\in\intergf{f(a)}{f(b)}$, il existe $x_0\in\interf{a}{b}$ tel
  que $f(x_0)=y_0$.
\end{theoreme}

\begin{preuve}
On définit la fonction $g$ sur $I$ par
\[\forall x\in I \qsep g(x)=f(x)-y_0\]
Alors, puisque $f$ est continue sur $I$, il en est de même pour $g$. Comme
de plus $y_0\in\intergf{f(a)}{f(b)}$, on en déduit que $g(a)$ et
$g(b)$ sont de signes distincts, c'est-à-dire que $g(a)g(b)\leq 0$.
Quitte à les échanger, on peut supposer que $a\leq b$. Nous souhaitons
montrer que $g$ s'annule en un point de $\interf{a}{b}$. Pour cela, nous
allons utiliser le procédé dit de dichotomie et construire deux suites
$\p{a_n}$ et $\p{b_n}$ telles que~:
\begin{itemize}
\item[$\bullet$] Pour tout $n\in\N$, $a_n\leq b_n$.
\item[$\bullet$] La suite $\p{a_n}$ est croissante et la suite $\p{b_n}$ est décroissante.
\item[$\bullet$] Pour tout $n\in\N$
  \[b_{n+1}-a_{n+1}=\frac{b_n-a_n}{2}\]
\item[$\bullet$] Pour tout $n\in\N$, $g\p{a_n}g\p{b_n}\leq 0$.
\end{itemize}
On construit ces deux suites de la manière suivante~:
\begin{itemize}
\item On pose $a_0=a$ et $b_0=b$. On a bien $a_0\leq b_0$ et
  $g\p{a_0}g\p{b_0}\leq 0$.
\item On suppose que les suites sont construites jusqu'au rang $n$. Construisons
  $a_{n+1}$ et $b_{n+1}$. On pose $c=\p{a_n+b_n}/2$. Alors
  $g\p{a_n}g(c)\leq 0$ ou $g(c)g\p{b_n}\leq 0$ (Autrement dit, soit
  $g\p{a_n}$ et $g(c)$ sont de signes distincts, soit $g(c)$ et $g\p{b_n}$
  sont de signes distincts). En effet, supposons que
  $g\p{a_n}g(c)> 0$ et $g(c)g\p{b_n}> 0$. Alors
  $g\p{a_n}g\p{b_n}g(c)^2>0$, donc $g\p{a_n}g\p{b_n}>0$ ce qui est absurde.
  \begin{itemize}
  \item Si $g\p{a_n}g(c)\leq 0$, on pose $a_{n+1}=a_n$ et $b_{n+1}=c$.
    On vérifie bien que $a_n\leq a_{n+1}$ et que $b_{n+1}\leq b_n$, que
    $a_{n+1}\leq b_{n+1}$, que $b_{n+1}-a_{n+1}=\p{b_n-a_n}/2$ et enfin
    que $g\p{a_{n+1}}g\p{b_{n+1}}\leq 0$.
  \item Sinon, $g(c)g\p{b_n}\leq 0$ et on pose $a_{n+1}=c$ et $b_{n+1}=b_n$.
  \end{itemize}
\end{itemize}
On a donc ainsi construit deux suites vérifiant les quatre assertions données
plus haut. Puisque
\[\forall n\in\N \qsep b_{n+1}-a_{n+1}=\frac{b_n-a_n}{2}\]
une récurrence immédiate nous donne
\[b_n-a_n=\frac{b-a}{2^n}\tendvers{n}{+\infty} 0\]
Comme de plus $a_n\leq b_n$ et que les suites $\p{a_n}$ et $\p{b_n}$ sont
respectivement croissantes et décroissantes, elles sont adjacentes.
Elles convergent donc
vers la même limite $x_0$. De plus $x_0\in\interf{a_0}{b_0}=\interf{a}{b}$.
Comme $g$ est continue en $x_0$, on en déduit que
\[g\p{a_n}g\p{b_n}\tendvers{n}{+\infty} g\p{x_0}^2\]
Or
\[\forall n\in\N \qsep g\p{a_n}g\p{b_n}\leq 0\]
Par passage à la limite $g\p{x_0}^2\leq 0$. On en déduit que $g\p{x_0}=0$
donc que $f\p{x_0}=y_0$.
\end{preuve}

\begin{remarques}
\remarque Une fonction réelle continue ne s'annulant pas sur un intervalle $I$ est
  de signe constant.
  \begin{sol}
  En effet, procédons par l'absurde et supposons qu'il existe
  $x_0\in I$ tel que $f\p{x_0}\geq 0$ et $x_1\in I$ tel que $f\p{x_1}\leq 0$.
  Puisque $f$ est continue sur $I$, d'après le théorème des valeurs intermédiaires,
  il existe $x\in I$ tel que $f(x)=0$. C'est absurde, donc $f$ est de signe
  constant sur $I$.    
  \end{sol}
\remarque Soit $f$ une fonction continue sur un intervalle $I$ telle que
  \[\forall x\in I \qsep \cro{f(x)=0 \ou f(x)=1}.\]
  Alors $f$ est constante. Plus généralement, si sur un intervalle, une fonction
  continue prend un nombre fini de valeurs, alors elle est constante.
\end{remarques}

\begin{exoUnique}
\exo Soit $f$ une fonction continue de $\interf{0}{1}$ dans $\interf{0}{1}$.
  Montrer que $f$ admet un point fixe.
\begin{sol}
$g(x)=f(x)-x$ puis $g(0)\geq 0$, $g(1)\leq 0$. Or $g$ continue donc d'après le TVI, il existe $x_0\in \interf{0}{1}$ tel que $g(x_0)=0$, d'où $f(x_0)=x_0$.
\end{sol}
% \exo Sur un intervalle, une fonction continue injective est strictement
%   monotone.
\end{exoUnique}

\begin{proposition}[utile=-3]
Soit $f$ une fonction réelle continue sur $\intero{a}{b}$ admettant
  respectivement pour limite $l_a$ et $l_b\in\Rbar$ en $a$ et $b\in\Rbar$. 
  Si $y_0\in\intergo{l_a}{l_b}$, il existe $x_0\in\intero{a}{b}$ tel que
  $f(x_0)=y_0$.
\end{proposition}

\begin{preuve}
Nous allons traiter le cas où $a=-\infty$, $b=+\infty$ et $l_a<l_b$, les autres
cas se traitant de manière similaire. Soit $y_0\in\intero{l_a}{l_b}$. Puisque
\[f(x)\tendvers{x}{-\infty} l_a<y_0\]
il existe $A\in\R$ tel que
\[\forall x\in\R \qsep x\leq A \implique f(x)\leq y_0\]
En particulier $f(A)\leq y_0$. De même, puisque
\[f(x)\tendvers{x}{+\infty} l_b>y_0\]
il existe $B\in\R$ tel que
\[\forall x\in\R \qsep x\geq B \implique f(x)\geq y_0\]
En particulier $f(B)\geq y_0$. Puisque $f$ est continue sur $\R$, on en
déduit, d'après le théorème des valeurs intermédiaires, qu'il existe
$x_0\in\intergf{A}{B}$ tel que $f\p{x_0}=y_0$.
\end{preuve}


\begin{exoUnique}
\exo Montrer que tout polynôme réel de degré impair admet au moins une racine réelle.
\end{exoUnique}

\begin{proposition}[utile=-3]
L'image d'un intervalle par une fonction réelle continue est un intervalle.
\end{proposition}

\begin{preuve}
Soit $f:I\mapsto\R$ une fonction continue sur un intervalle $I$. Montrons que $J=f(I)$ est un intervalle. Soit alors $y_0,y_1 \in J$ tels que $y_0\leq y_1$ et considérons $y\in [y_0,y_1]$. L'objectif est de montrer que $y\in J$.\\
Comme $J=f(I)$, il existe $x_0$ et $x_1$ dans $I$ tels que $f(x_0)=y_0$ et $f(x_1)=y_1$. Ainsi, $y\in \interf{f(x_0)}{f(y_0)}$ et $f$ est continue sur $I$, donc d'après le TVI, il existe $x\in "\interf{x_0}{x_1}"$ tel que $f(x)=y$. Or, $I$ est un intervalle donc $x\in I$ et finalement $y\in f(I)=J$. $J$ est donc bien un intervalle.
\end{preuve}

\begin{remarques}
\remarque Cette proposition est une reformulation du
  théorème des valeurs intermédiaires.
\remarque Il est possible que les intervalles $I$ et
  $f(I)$ ne soient pas de même nature (ouvert, fermé, ouvert à gauche et
  fermé à droite).
  Par exemple, si $f$ est la fonction définie sur
  $\R$ par \[\forall x\in\R \qsep f(x)\defeq \frac{1}{1+x^2}\]
  on a $f\p{\intero{-\infty}{+\infty}}=\interof{0}{1}$.
  \begin{sol}
  En dérivant $f$, on a $f$ croissante jusqu'en $0$ de valeurs $1$ et décroissante ensuite.
    \end{sol}
\end{remarques}


\begin{theoreme}[nom={Théorème de la bijection}]
\begin{itemize}
\item Soit $f$ une fonction continue, strictement croissante sur $[a,b]$. Alors elle réalise une bijection de $[a,b]$ sur \[f([a,b])=[f(a),f(b)].\]
\item Soit $a,b\in\Rbar$ et $f$ une fonction continue, strictement croissante sur $]a,b[$. On pose
\[l_a\defeq\lim_{x\to a} f(x) \et l_b\defeq\lim_{x\to b} f(x).\]
Alors $f$ réalise une bijection de $]a,b[$ sur \[f\p{\intero{a}{b}}=\intero{l_a}{l_b}.\]
\end{itemize}
\end{theoreme}

\begin{proposition}[utile=-3]
Soit $I$ un intervalle de $\R$ et $f$ une fonction réelle continue strictement
monotone sur $I$. Alors $f$ induit une bijection de l'intervalle $I$ sur
l'intervalle $J\defeq f(I)$ et sa bijection réciproque $f^{-1}:J\to I$ est
continue sur $J$.
\end{proposition}

\begin{preuve}
On suppose que $f$ est strictement croissante sur $I$, la démonstration dans le
cas où $f$ est strictement décroissante se faisant de même. Puisque $f$ est
continue sur l'intervalle $I$, d'après la proposition précédente, $J=f(I)$
est un intervalle. Puisque $f$ est strictement croissante, on en déduit
qu'elle est injective. Sa corestriction à $J$ (que l'on notera encore $f$)
est donc une bijection de $I$ dans $J$.\\
Montrons que $f^{-1}$ est continue. On se donne $y_0\in J$ et on souhaite
montrer que $f^{-1}$ est continue en $y_0$. Nous supposerons dans la suite que
$y_0$ est intérieur à $J$, la démonstration dans le cas où $y_0$ est une borne
de $J$ se faisant de même. Puisque $f$ est strictement croissante, il en est
de même pour $f^{-1}$, donc $x_0=f^{-1}\p{y_0}$ n'est pas une borne de $I$. On
se donne donc $\epsilon>0$ tel que
$\interf{x_0-\epsilon}{x_0+\epsilon}\subset I$. Par stricte croissance de $f$,
on a
\[f\p{x_0-\epsilon}<f\p{x_0}=y_0<f\p{x_0+\epsilon}\]
Il existe donc $\eta>0$ tel que $\interf{y_0-\eta}{y_0+\eta}\subset
\interf{f\p{x_0-\epsilon}}{f\p{x_0+\epsilon}}$. Soit
$y\in\interf{y_0-\eta}{y_0+\eta}$. Alors
\begin{eqnarray*}
y_0-\eta\leq y\leq y_0+\eta
&\donc& f\p{x_0-\epsilon}\leq y\leq f\p{x_0+\epsilon}\\
&\donc& x_0-\epsilon \leq f^{-1}(y) \leq x_0+\epsilon\\
&     & \text{par croissance de $f^{-1}$}\\
&\donc& \abs{f^{-1}(y)-x_0}\leq\epsilon\\
&\donc& \abs{f^{-1}(y)-f^{-1}\p{y_0}}\leq\epsilon
\end{eqnarray*}
Donc $f^{-1}$ est continue en $y_0$.
\end{preuve}

\begin{remarqueUnique}
\remarque La fonction $\sin$ est strictement croissante sur $\interf{-\pi/2}{\pi/2}$.
  Comme $\sin\p{-\pi/2}=-1$ et $\sin\p{\pi/2}=1$, elle réalise une bijection
  de $\interf{-\pi/2}{\pi/2}$ sur $\interf{-1}{1}$. Sa bijection réciproque,
  la fonction $\arcsin$ est donc continue sur $\interf{-1}{1}$.
\end{remarqueUnique}

\begin{exoUnique}
\exo Soit $f$ la fonction définie sur $\RP$ par
  \[\forall x\in\RP \qsep f(x)\defeq x \e^x\]
  Montrer que $f$ réalise une bijection continue de $\RP$ sur $\RP$, que
  $f^{-1}$ est continue et que \[f^{-1}(x)\tendvers{x}{+\infty}+\infty.\]
%   Alors $f$ est continue et strictement croissante sur $\interfo{-1}{+\infty}$.
%   Comme $f\p{-1}=-1/e$ et
%   \[f(x)\tendvers{x}{+\infty}+\infty\]
%   on en déduit que $f$ réalise une bijection de $\interfo{-1}{+\infty}$ sur
%   $\interfo{-1/e}{+\infty}$. Sa bijection réciproque, appelée fonction de Lambert,
%   est donc continue sur $\interfo{-1/e}{+\infty}$.
\end{exoUnique}

\begin{proposition}
Soit $f$ une fonction réelle, continue et injective sur un intervalle $I$. Alors $f$ est strictement monotone.
\end{proposition}

\begin{preuve}
On procède par l'absurde et on suppose qu'il existe $(x_1,y_1) \in I^2, x_1<y_1 \et f(x_1)\geq f(y_1)$ et $(x_2,y_2) \in I^2, x_2<y_2 \et f(x_2)\leq f(y_2)$.

Définissons alors $\phi : \interf{0}{1}
\mapsto \R$ définie par $$\phi(t)=f((1-t)x_1+tx_2)-f((1-t)y_1+ty_2).$$
$\phi$ est continue. $\phi(0)=f(x_1)-f(y_1)\geq 0$ et $\phi(1)=f(x_2)-f(y_2)\leq 0$ donc par le TVI, il existe $t\in [0,1]$ tel que $\phi(t)=0$. En posant $x_0=(1-t)x_1+tx_2$ et $y_0=(1-t)y_1+ty_2$, on a $x_0<y_0$ et $f(x_0)=f(y_0)$ ce qui contredit l'injectivité de $f$.
\end{preuve}

\subsection{Théorème de compacité}

\begin{definition}[utile=-3]
Soit $f$ une fonction réelle définie sur un ensemble non vide $X$. Si $f$ est
majorée sur $X$, $\ensim{f(x)}{x\in X}$ est une partie non vide majorée de
$\R$. Elle admet donc une borne supérieure notée
\[\sup_{x\in X} f(x).\]
On dit que cette borne est atteinte lorsqu'il existe $x_0\in X$ tel que
\[f\p{x_0}=\sup_{x\in X} f(x)\]
c'est-à-dire lorsque l'ensemble $\ensim{f(x)}{x\in X}$ admet un plus grand
élément; si tel est le cas, la borne supérieure est notée
\[\max_{x\in X} f(x).\]
\end{definition}

\begin{remarques}
\remarque On définit de même la notion de borne inférieure.
\remarque Soit $f$ la fonction définie sur $\R$  par
  \[\forall x\in\R \qsep f(x)\defeq x\p{1-x}.\]
  Alors $f$ est bornée et atteint ses bornes sur $\interf{0}{1}$.
  \[\sup_{x\in\interf{0}{1}} x\p{1-x}=f\p{\frac{1}{2}}=\frac{1}{4} \et
    \inf_{x\in\interf{0}{1}} x\p{1-x}=f\p{0}=f\p{1}=0.\]
\remarque Soit $f$ la fonction définie sur $\RPs$ par
  \[\forall x\in\RPs \qsep f(x)\defeq \frac{1}{x}.\]
  Alors $f$ n'est pas majorée sur $\RPs$. De plus, elle est minorée mais
  n'atteint pas sa borne inférieure sur $\RPs$.
  \[\inf_{x\in\RPs} \frac{1}{x}=0.\]
\end{remarques}

\begin{theoreme}[nom={Théorème de compacité}]
Sur un segment, une fonction réelle continue est bornée et atteint ses bornes.
\end{theoreme}

\begin{preuve}
Soit $f$ une fonction continue sur le segment $\interf{a}{b}$.
\begin{itemize}
\item Montrons que $f$ est bornée.\\
  Puisque $f$ est continue, il en est de même pour $\abs{f}$. Montrons que
  $\abs{f}$ est majorée. On raisonne par l'absurde et on suppose que
  \[\non\cro{\exists M\in\R \qsep \forall x\in\interf{a}{b} \qsep
           \abs{f(x)}\leq M}\]
   donc
  \[\forall M\in\R \qsep \exists x\in\interf{a}{b} \qsep \abs{f(x)}>M\]
   Étant donné $n\in\N$, on pose $M=n$. Il existe donc $x_n\in\interf{a}{b}$
   tel que $\abs{f\p{x_n}}\geq n$. On construit ainsi une suite
   $\p{x_n}$ d'éléments de $\interf{a}{b}$. Cette suite étant bornée, d'après
   le théorème de Bolzano-Weierstrass, il existe une extractrice $\phi$ et un
   réel $l$ tel que
   \[x_{\phi(n)}\tendvers{n}{+\infty} l\]
   Comme
   \[\forall n\in\N \qsep a\leq x_{\phi(n)}\leq b\]
   on en déduit que $l\in\interf{a}{b}$. En utilisant la continuité de
   $\abs{f}$ en $l$, il vient
   \[\abs{f\p{x_{\phi(n)}}}\tendvers{n}{+\infty} \abs{f(l)}\]
   C'est contradictoire avec le fait que
   \[\forall n\in\N \qsep \abs{f\p{x_{\phi(n)}}}\geq \phi(n)\geq n
     \tendvers{n}{+\infty} +\infty\]
   Donc $\abs{f}$ est majorée, donc $f$ est bornée.
\item Montrons que $f$ atteint ses bornes.\\
   Montrons d'abord que $f$ atteint sa borne supérieure. On pose
   \[M=\sup_{x\in\interf{a}{b}} f(x)\]
   Par définition de la borne supérieure
   \[\forall \epsilon>0 \qsep \exists x\in\interf{a}{b} \qsep
     f(x)\geq M-\epsilon\]
   Étant donné $n\in\N$, en posant $\epsilon=1/2^n>0$, il existe donc
   $x_n\in\interf{a}{b}$ tel que $f\p{x_n}\geq M-1/2^n$. 
   On construit ainsi une suite $\p{x_n}$ d'éléments de $\interf{a}{b}$. Cette
   suite étant bornée, d'après le théorème de Bolzano-Weierstrass, il existe
   une extractrice $\phi$ et un réel $l$ tel que
   \[x_{\phi(n)}\tendvers{n}{+\infty} l\]
   Comme
   \[\forall n\in\N \qsep a\leq x_{\phi(n)}\leq b\]
   on en déduit que $l\in\interf{a}{b}$. Comme de plus
   \[\forall n\in\N \qsep
     \underbrace{M-\frac{1}{2^{\phi(n)}}}_{\tendvers{n}{+\infty} M}
     \leq f\p{x_{\phi(n)}}\leq M\]
   on en déduit, d'après le théorème des gendarmes, que
   \[f\p{x_{\phi(n)}}\tendvers{n}{+\infty} M\]
   Or, par continuité de $f$ en $l$
   \[f\p{x_{\phi(n)}}\tendvers{n}{+\infty} f(l)\]
   donc $M=f(l)$. La borne supérieure de $f$ est donc atteinte.\\
   Comme $f$ est
   continue, $-f$ est continue donc atteint sa borne supérieure. Donc $f$
   atteint sa borne inférieure.
\end{itemize}
\end{preuve}

\begin{remarqueUnique}
\remarque Si $f:[a,b]\to\K$ est une fonction continue, on applique souvent ce théorème à la fonction $\abs{f}$. Il existe donc $M\in\RP$ tel que
\[\forall x\in[a,b]\qsep \abs{f(x)}\leq M.\]
\end{remarqueUnique}

\begin{exoUnique}
\exo Soit $f$ une fonction continue sur un segment $\interf{a}{b}$ telle
  que $\forall x\in\interf{a}{b} \qsep 0\leq f(x)<1$. Montrer que si
  $(u_n)$ est une suite d'éléments de $[a,b]$, alors
  \[f(u_n)^n\tendvers{n}{+\infty} 0.\]
  \begin{sol}
  $f$ étant continue sur $\interf{a}{b}$, elle est majorée et atteint ses bornes. Il existe donc $\alpha\in \interf{a}{b}$ tel que $\forall x \in \interf{a}{b}, 0\leq f(x)\leq f(\alpha)<1$. D'où le résultat en élevant à la puissance $n$ et d'après le théorème des gendarmes.
  \end{sol}
\end{exoUnique}


 
\begin{proposition}[utile=-3]
L'image d'un segment par une fonction réelle continue est un segment.
\end{proposition}

\begin{preuve}
Soit $f$ un fonction continue sur un segment $S$. Puisque $S$ est un
intervalle, d'après le théorème des valeurs intermédiaires, $f(S)$ est
un intervalle. D'après la proposition précédente, $f(S)$ est borné et contient
ses bornes. C'est donc un segment.
\end{preuve}

\subsection{Continuité uniforme}

\begin{definition}[utile=-3]
On dit qu'une fonction $f:\mathcal{D}\to\K$ est \emph{uniformément continue} lorsque
\[\forall \epsilon>0 \qsep \exists \eta>0 \qsep \forall x,y\in\mathcal{D}\qsep
  \abs{x-y}\leq\eta \implique \abs{f(x)-f(y)}\leq\epsilon.\]
\end{definition}

\begin{remarqueUnique}
\remarque Une fonction lipschitzienne est uniformément continue.
% \remarque Soit $\alpha\in\RPs$. On dit qu'une fonction $f$ est
%   $\alpha$-Hölderienne lorsqu'il existe $M\in\RP$ tel que
%   \[\forall x\in\mathcal{D} \qsep \abs{f(x)-f(y)}\leq M\abs{x-y}^\alpha\]
%   Une fonction Hölderienne est uniformément continue.
\end{remarqueUnique}

\begin{exoUnique}
\exo
  \begin{questions}
\question Montrer que
  \[\forall x,y\in\RP \qsep \abs{\sqrt{x}-\sqrt{y}}\leq\sqrt{\abs{x-y}}.\]
\question En déduire que la fonction $x\mapsto \sqrt{x}$, définie sur $\RP$,
   est uniformément continue, mais n'est pas lipschitzienne.
  \end{questions}
\end{exoUnique}

\begin{proposition}[utile=-3]
Si $f$ est uniformément continue, alors elle est continue.
\end{proposition}

\begin{remarqueUnique}
\remarque Soit $f$ une fonction continue. Alors
  \[\forall x\in\mathcal{D} \qsep \forall \epsilon>0 \qsep \exists \eta>0
    \qsep \forall y\in\mathcal{D} \qsep \abs{x-y}\leq\eta \implique
    \abs{f(x)-f(y)}\leq\epsilon.\]
  Les deux premiers quantificateurs étant de même nature, on peut les échanger,
  donc
  \[\forall \epsilon>0 \qsep \forall x\in\mathcal{D} \qsep \exists \eta>0
    \qsep \forall y\in\mathcal{D} \qsep \abs{x-y}\leq\eta \implique
    \abs{f(x)-f(y)}\leq\epsilon.\]
  Une fonction est donc uniformément continue lorsqu'on peut échanger les
  quantificateurs portant sur $x$ et $\eta$, c'est-à-dire lorsqu'il est
  possible de choisir $\eta$ indépendamment de $x$.
% \remarque Supposons que $f$ soit uniformément continue. Alors, si $\p{u_n}$
%   et $\p{v_n}$ sont deux suites d'éléments de $\mathcal{D}$ telles que
%   \[v_n-u_n\tendvers{n}{+\infty}0\]
%   on a $f\p{u_n}-f\p{v_n}\tendvers{n}{+\infty}0$. En effet, soit $\epsilon>0$.
%   Puisque $f$ est uniformément continue, il existe $\eta>0$ tel que
%   \[\forall x,y\in\mathcal{D}\qsep
%     \abs{x-y}\leq\eta \implique \abs{f(x)-f(y)}\leq\epsilon\]
%   Puisque $u_n-v_n$ tend vers 0, il existe $N\in\N$ tel que
%   \[\forall n\geq N \qsep \abs{u_n-v_n}\leq\eta\]
%   Donc, pour $n\geq N$, on a $\abs{f\p{u_n}-f\p{v_n}}\leq\epsilon$. En
%   conclusion
%   \[f\p{u_n}-f\p{v_n}\tendvers{n}{+\infty}0\]
%   En particulier, si on trouve deux suites $\p{u_n}$ et $\p{v_n}$ d'élements
%   de $\mathcal{D}$ telles que $u_n-v_n$ converge vers 0 et $f\p{u_n}-f\p{v_n}$
%   ne converge pas vers 0, on peut en conclure que $f$ n'est pas uniformément
%   continue.\\
%   On démontre ainsi que $x\mapsto x^2$ n'est pas uniformément continue.
%   En effet si $\p{u_n}$ et $\p{v_n}$ sont définies par
%   \[\forall n\in\Ns \qsep u_n=n+\frac{1}{n} \et v_n=n\]
%   alors $v_n-u_n=1/n$ tend vers 0, et
%   \[u_n^2-v_n^2=\p{n+\frac{1}{n}}^2-n^2=2+\frac{1}{n^2}\tendvers{n}{+\infty}
%     2\neq 0\]
%   Donc $x\mapsto x^2$ n'est pas uniformément continue.
\end{remarqueUnique}

\begin{exoUnique}
\exo Montrer que la fonction $f$ définie sur $\R$ par $f(x)\defeq x^2$
  n'est pas uniformément continue.
  \begin{sol}
  On souhaite montrer que
  \[\non\cro{\forall \epsilon>0 \qsep \exists \eta>0 \qsep
           \forall x,y\in\R\qsep
           \abs{x-y}\leq\eta \implique \abs{x^2-y^2}\leq\epsilon}\]
  c'est-à-dire que
  \[\exists \epsilon>0 \qsep \forall \eta>0 \qsep \exists x,y\in\R
    \qsep
    \begin{cases}
    \abs{x-y}\leq\eta \\
    \abs{x^2-y^2}>\epsilon
    \end{cases}\]
  On pose $\epsilon=1$. Soit $\eta>0$. Alors, il existe $n\in\Ns$ tel que
  $1/n\leq\eta$. On pose alors $x=n+1/n$ et $y=n$. Alors $\abs{x-y}=1/n\leq\eta$
  et $\abs{x^2-y^2}=2+1/n^2>1=\epsilon$. Donc $f$ est continue mais n'est pas
  uniformément continue.      
  \end{sol}
\end{exoUnique}

\begin{theoreme}[nom=Théorème de \nom{Heine}]
Sur un segment, une fonction continue est uniformément continue.
\end{theoreme}

\begin{preuve}
Soit $f$ une fonction continue sur $\interf{a}{b}$. Montrons qu'elle y est
uniformément continue. On raisonne par l'absurde et on suppose que
\[\non\cro{\forall \epsilon>0 \qsep \exists \eta>0 \qsep
           \forall x,y\in\interf{a}{b}\qsep
           \abs{x-y}\leq\eta \implique \abs{f(x)-f(y)}\leq\epsilon}\]
donc
\[\exists \epsilon>0 \qsep \forall \eta>0 \qsep \exists x,y\in\interf{a}{b}
  \qsep
  \begin{cases}
  \abs{x-y}\leq\eta \\
  \abs{f(x)-f(y)}>\epsilon
  \end{cases}\]
Soit un tel $\epsilon$. Étant donné $n\in\N$, en posant $\eta=1/2^n>0$, il
existe $x_n,y_n\in\interf{a}{b}$ tels que $\abs{x_n-y_n}\leq 1/2^n$
et $\abs{f\p{x_n}-f\p{y_n}}\geq \epsilon$. On construit ainsi deux suites
$\p{x_n}$ et $\p{y_n}$ d'éléments de $\interf{a}{b}$. Comme $\p{x_n}$ est
bornée, d'après le théorème de Bolzano-Weierstrass, il existe une extractrice
$\phi$ et un réel $l$ tel que
\[x_{\phi(n)}\tendvers{n}{+\infty} l\]
Comme
\[\forall n\in\N \qsep a\leq x_{\phi(n)}\leq b\]
on en déduit que $l\in\interf{a}{b}$. De plus
\begin{eqnarray*}
\forall n\in\N \qsep \abs{y_{\phi(n)}-l}
&=& \abs{y_{\phi(n)}-x_{\phi(n)}+x_{\phi(n)}-l}\\
&\leq& \abs{y_{\phi(n)}-x_{\phi(n)}}+\abs{x_{\phi(n)}-l}\\
&\leq& \frac{1}{2^{\phi(n)}}+\abs{x_{\phi(n)}-l}\tendvers{n}{+\infty} 0
\end{eqnarray*}
Donc 
\[y_{\phi(n)}\tendvers{n}{+\infty} l\]
Par continuité de $f$ en $l$, on en déduit que
\[\abs{f\p{x_{\phi(n)}}-f\p{y_{\phi(n)}}}\tendvers{n}{+\infty}
  \abs{f(l)-f(l)}=0\]
C'est contradictoire avec le fait que
\[\forall n\in\N \qsep \abs{f\p{x_{\phi(n)}}-f\p{y_{\phi(n)}}}\geq \epsilon\]
Donc $f$ est uniformément continue.
\end{preuve}
%END_BOOK

\end{document}



\section{Exercices}
\setcounter{numeroexercice}{1}
\documentclass{magnolia}

\magtex{tex_driver={pdftex}}
\magfiche{document_nom={Continuité, limites},
          auteur_nom={François Fayard},
          auteur_mail={fayard.prof@gmail.com}}
\magexos{exos_matiere={maths},
         exos_niveau={mpsi},
         exos_chapitre_numero={10},
         exos_theme={Continuité, limites}}
\magmisenpage{misenpage_presentation={tikzvelvia},
              misenpage_format={a4},
              misenpage_nbcolonnes={1},
              misenpage_sol={non}}
\maglieudiff{lieu_lycee={Aux Lazaristes},
             lieu_classe={MPSI 1},
             lieu_annee={2020--2021}}
\magprocess

\begin{document}

%BEGIN_BOOK
% Ce qui est affirmé sans preuve peut être nié sans preuve.
% Euclide
\magsection{Fonction numérique, topologie élémentaire}

\exercice{nom={Monotonie}}
Soit $f:\R\to\R$ une fonction telle que $f\circ f$ est croissante et
$f\circ f\circ f$ est strictement décroissante. Montrer que $f$ est
strictement décroissante.

\begin{sol}
Soit $a<b$. Procédons par l'absurde et supposons que $f(a)\leq f(b)$. Alors, par croissance de $f^2$, $f^3(a)\leq f^3(b)$. Mais par stricte décroissance de $f^3$, on a $f^3(a)> f^3(b)$, d'où la contradiction.
\end{sol}



\exercice{nom={Une fonction périodique étrange}}
Soit $f$ la fonction définie sur $\R$ par
\[\forall x\in\R \qsep
  f(x)\defeq
  \begin{cases}
  1 & \text{si $x\in\Q$}\\
  0 & \text{si $x\in\R\setminus\Q$.}  
  \end{cases}\]
Montrer que l'ensemble des périodes de $f$ est $\Q$.
\textit{On a donc construit une fonction périodique non constante qui n'admet pas
de plus petite période strictement positive.}

\begin{sol}
On montre que tout rationnel est une période par disjonction des cas. Puis si $T$ est une période, alors en particulier $f(0)=1=f(T)$ donc $T$ est rationnel.
\end{sol}


\magsubsection{Propriété locale}


\magsection{Limite}

\magsubsection{Définition, propriétés élémentaires}


\exercice{nom={Non existence d'une limite}}
Montrer que la fonction définie sur $\intero{0}{1}$ par
\[\forall x\in\intero{0}{1} \qsep f(x)\defeq\sin\p{\frac{1}{x-x^2}}\]
n'a pas de limite en 0.

\exercice{nom={Manipulation de limite}}
Soit $f$ une fonction définie sur $\R$ telle que
\[f(x)\tendvers{x}{0}0 \et \frac{f\p{2x}-f(x)}{x}\tendvers{x}{0}0.\]
En remarquant que
\[f(x)=\sum_{k=1}^n \cro{f\p{\frac{x}{2^{k-1}}}-f\p{\frac{x}{2^k}}} +
  f\p{\frac{x}{2^n}}\]
montrer que
\[\frac{f(x)}{x}\tendvers{x}{0}0.\]

\begin{sol}
Soit $\epsilon>0$. Il existe $\eta>0$ tel que $\forall x\in \RPs$, $$\abs{x}\leq \eta \Longrightarrow \abs{\frac{f\p{2x}-f(x)}{x}}\leq \frac{\epsilon}{2}.$$
Ainsi, d'après le résultat (à montrer) qu'on nous demande de remarquer, $\forall x \in \RPs$ tel que $\abs{x}\leq \eta$, on a:
$$\abs{\dfrac{f(x)}{x}}\leq \sum_{k=1}^n \frac{1}{2^k}\abs{\frac{\p{f\p{2\frac{x}{2^{k}}}-f\p{\frac{x}{2^k}}}}{\frac{x}{2^k}}} +
  \abs{\frac{f\p{\frac{x}{2^n}}}{x}}\leq \frac{\epsilon}{2}\p{1-\p{\frac{1}{2}}^n}+\abs{\frac{f\p{\frac{x}{2^n}}}{x}}\leq \frac{\epsilon}{2}+ \abs{\frac{f\p{\frac{x}{2^n}}}{x}}.$$
  Ceci étant vrai $\forall n \in \N$, on peut maintenant "regarder" cette inégalité à $x$ fixé et majorer le dernier terme par $\epsilon/2$ d'où le résultat souhaité.
\end{sol}


\magsubsection{Limite et ordre sur $\R$}
\magsubsection{Limite à gauche, limite à droite}

\exercice{nom={Existence et calculs de limites}}
Existence et calcul des limites des expressions suivantes
\[\ent{\frac{1}{x}} \ \text{à droite en 0}, \qquad x\cos\p{\frac{1}{x}} \ \text{en 0},\]
\[x\ent{\frac{1}{x}} \ \text{en 0}, \qquad
  \cos\p{\frac{1}{x}} \ \text{en 0}, \qquad
  \frac{x^x}{\ent{x}^{\ent{x}}} \ \text{en $+\infty$}.\]

\exercice{nom={Existence et calculs de limites}}
Existence et calcul des limites des expressions suivantes
\[\frac{6x^2+5x-4}{2x-1} \ \text{en $\frac{1}{2}$}, \qquad
  \frac{3}{x^3-1}-\frac{2}{x^2-1} \ \text{en 1},\]
\[x^n \e^{-1/x^2} \ \text{en 0 avec $n\in\Z$}, \qquad
%%  \frac{\p{1-\cos x}\ln\p{1+x^2}}{x^2\tan^2 x} \ \text{en 0}\]
%%\[\frac{\ln\p{\cos x}}{1-\cos\p{2x}} \ \text{en 0}  \qquad
  \frac{\ln\p{\ch\p{\alpha x}}}{\ln\p{\ch x}} \ \text{en $+\infty$ avec
  $\alpha\in\R$},\]
\[\frac{r \e^{\ii\alpha t}-1}{t} \ \text{en 0 avec $r\in\RPs$ et $\alpha\in\R$},
  \qquad \frac{\e^{2\ii t}\tan t}{t^2} \ \text{en 0}.\]
\begin{sol}
\[\frac{6x^2+5x-4}{2x-1}=3x+4\tendvers{x}{\frac{1}{2}}\frac{11}{2} \qquad
  \frac{3}{x^3-1}-\frac{2}{x^2-1}=-\frac{2x+1}{(x+1)(x^2+x+1)}\tendvers{x}{1}-\frac{1}{2}\]
\[x^n e^{-1/x^2}\tendvers{x}{0}0 \qquad
  \frac{\ln\p{\ch\p{\alpha x}}}{\ln\p{\ch x}}\tendvers{x}{+\infty}\abs{\alpha}\]
\[i\alpha \text{ si $r=1$, pas de limite sinon} \qquad \text{pas de limite}\]
\end{sol}





\magsection{Continuité}

\magsubsection{Continuité ponctuelle}

\exercice{nom={Une fonction discontinue en tout point}}
Soit $f$ la fonction définie sur $\R$ par
\[\forall x\in\R \qsep f(x)\defeq
  \begin{cases}
  1 & \text{si $x\in\Q$}\\
  0 & \text{si $x\in\R\setminus\Q$.}
  \end{cases}\]
Montrer que $f$ est discontinue en tout point.

\begin{sol}
  Soit $x\in \Q$, $\R\setminus\Q$ étant dense dans $\R$, il existe une suite $(u_n)$ d'irrationnels tendant vers $x$. Alors $\forall n \in N, f(u_n)=0\tendvers{n}{+\infty}0\neq f(x)=1$, donc $f$ est discontinue en $x$.\\
  Soit $x\in \R\setminus\Q$, $Q$ étant dense dans $\R$, il existe une suite $(u_n)$ de rationnels tendant vers $x$. Alors $\forall n \in N, f(u_n)=1\tendvers{n}{+\infty}1\neq f(x)=0$, donc $f$ est discontinue en $x$.
  \end{sol}

\exercice{nom={Une fonction continue en tout point de
$\R\setminus\Q$}}
Soit $f$ la fonction définie sur $\R$ par
\[\forall x\in\R \qsep f(x)\defeq
  \begin{cases}
  0 & \text{si $x\in\R\setminus\Q$}\\
  \frac{1}{q} & \text{si $x\in\Q$ et $x=\frac{p}{q}$ avec
    $\p{p,q}\in\Z\times\Ns$ et $\p{p,q}=1$.}
  \end{cases}\]
\begin{questions}
\question Montrer que $f$ est discontinue en tout $x\in\Q$.
\question Le but de cette question est de montrer que $f$ est continue en tout
  $x\in\R\setminus\Q$. Soit $x\in\R\setminus\Q$.
  \begin{questions}
  \question Soit $\p{p_n}$ une suite d'éléments de $\Z$ et $\p{q_n}$ une suite
    d'éléments de $\Ns$ telles que
    \[\frac{p_n}{q_n}\tendvers{n}{+\infty} x.\]
    Montrer que $\p{q_n}\tendvers{n}{+\infty}+\infty$.
  \question En déduire que $f$ est continue en $x$.
  \end{questions}
\end{questions}

\begin{sol}
\begin{questions}
\question Soit $x\in \Q$, $\R\setminus\Q$ étant dense dans $\R$, il existe une suite $(u_n)$ d'irrationnels tendant vers $x$. Alors $\forall n \in N, f(u_n)=0\tendvers{n}{+\infty}0\neq f(x)=1/q$, donc $f$ est discontinue en $x$.
\question 
  \begin{questions}
  \question Raisonnons par l'absurde. Cela signifie qu'il existe $m\in \R$ tel que $\forall N \in \N, \exists n\geq N$ tel que $q_n < m$.

Fixons un tel $m$ et construisons une extractrice $\phi$ telle que $\forall n \in \N$, $1\leq q_{\phi(n)}\leq m$.
\begin{itemize}
\item[$\bullet$] On pose $N=0$. Il existe $n\geq 0$ tel que $q_n<m$. On pose $n=\phi(0)$.
\item[$\bullet$] Soit $k\in \N$. On suppose construit $\phi(0)< \ldots< \phi(k)$ tels que $1\leq u_{\phi(k)}\leq m$.
On pose $N=\phi(k)+1$. Alors, il existe $n\geq N$ tel que $q_n<m$. On pose $\phi(k+1)=n$...
\end{itemize}
La suite $(q_{\phi(n)})$ est donc bien construite et bornée. Elle admet donc une sous-suite convergente, c'est-à-dire qu'il existe $\psi$ telle que $(q_{\phi(\psi(n))})$ converge vers $\ell \in \R$. Comme il s'agit d'une suite d'entiers convergentes, elle est constante apcr et donc $\ell \in \Ns$. Mais alors $\displaystyle p_{\phi(\psi(n))}=\frac{p_{\phi(\psi(n))}}{q_{\phi(\psi(n))}}q_{\phi(\psi(n))} \tendvers{n}{+\infty}{\ell x}=\ell'$ et pour les mêmes raisons $\ell' \in \Z$.
Ainsi, $\displaystyle \frac{p_{\phi(\psi(n))}}{q_{\phi(\psi(n))}} \tendvers{n}{+\infty} \frac{\ell'}{\ell} \in \Q$. Donc $x\in \Q$ par unicité de la limite, ce qui fournit une contradiction.
  \question En déduire que $f$ est continue en $x$.
  \end{questions}
\end{questions}
\end{sol}


\exercice{nom={Une équation fonctionnelle}}
Soit $f$ une application continue de $\R$ dans $\R$ telle que
\[\forall x,y\in\R \qsep f\p{x+y}=f(x)f(y).\]
\begin{questions}
\question On suppose qu'il existe $x_0\in\R$ tel que $f\p{x_0}=0$. Montrer que
  $f$ est la fonction nulle.
\enonce Dans la suite, on suppose que $f$ n'est pas la fonction nulle.
\question Montrer que
  \[\forall x\in\R \qsep f(x)>0.\]
\question En déduire qu'il existe $a\in\R$ tel que
  \[\forall x\in\R \qsep f(x)=\e^{ax}.\]
  \textit{On pourra penser à utiliser les résultats d'un exercice similaire vu en cours.}
\end{questions}

\begin{sol}
\begin{questions}
\question $\forall x \in \R$, $f(x)=f(x-x_0+x_0)=f(x-x_0)f(x_0)=0$ donc $f$ est la fonction nulle.
\question $\forall x\in\R \qsep f(x)=f\p{\dfrac{x}{2}}^2\geq 0$ et non nul par hypothèse, d'où le résultat.
\question Déjà, $f(0)=f(0)^2$ donc $f(0)=1$ car non nul. De plus, $\forall n \in \N, f(nx)=f(x)^n$ et $f(0)=f(nx)f(-nx)$ donc $f(-nx)=1/f(nx)=f(x)^{-n}$. 
Ainsi, pour tout rationnel $r=p/q$ avec $p\in \Z$ et $q\in \Ns$, on a :
$$f(1)^p=f(p)=f\p{q\frac{p}{q}}=f\p{\frac{p}{q}}^q$$ d'où $$f(r)=f(1)^r=\e^{r\ln(f(1))}.$$
On conclut grâce à la densité des rationnels.
\end{questions}
\end{sol}


% \exercice{nom={Une équation fonctionnelle}}
% Soit $f$ une fonction continue de $\R$ dans $\R$ telle que
% \[\forall x,y\in\R \qsep f\p{\frac{x+y}{2}}=\frac{f(x)+f(y)}{2}.\]
% \begin{questions}
% \question Montrer que
%   \[\mathcal{D}\defeq\ensim{\frac{p}{2^n}}{p\in\Z \et n\in\N}\]
%   est dense dans $\R$.
% \question Soit $a,b\in\R$ tels que $a<b$.
%   \begin{questions}
%   \question Montrer que
%     \[\forall x\in\mathcal{D}\cap\interf{0}{1} \qsep
%       f\p{\alpha\cdot a +\p{1-\alpha} \cdot b}=\alpha f(a)+
%       \p{1-\alpha} f(b).\]
%   \question En déduire qu'il existe $\alpha,\beta\in\R$ tels que
%     \[\forall x\in\interf{a}{b} \qsep f(x)=\alpha+\beta x.\]
%   \end{questions}
% \question Montrer que
%   \[\forall x\in\R \qsep f(x)=\alpha+\beta x.\]
% \end{questions}



% \exercice{nom={Une équation fonctionnelle}}
% Soit $f$ une fonction continue sur $\R$ telle que
% \[\forall x,y\in\R \qsep f\p{x+y}=f(x)+f(y).\]
% \begin{questions}
% \question Montrer que
%   \[\forall x\in\R \qsep \forall n\in\Z \qsep f\p{nx}=nf(x).\]
% \question Montrer que
%   \[\forall x\in\R \qsep \forall q\in\Q \qsep f\p{qx}=qf(x).\]
% \question En déduire qu'il existe $a\in\R$ tel que
%   \[\forall x\in\R \qsep f(x)=ax.\]
% % \question Le but de cette question est de montrer qu'il existe une fonction $g$
% %   de $\R$ dans $\R$ (qui sera donc non continue) telle que
% %   \[\forall x,y\in\R \qsep g\p{x+y}=g(x)+g(y)\]
% %   et telle que $g$ ne soit pas de la forme $g(x)=ax$.
% %   \begin{questions}
% %   \question On considère $\R$ comme $\Q$-espace vectoriel. En admettant que
% %     $A=\Q$ admet un supplémentaire $B$ dans $\R$ montrer qu'il existe un
% %     unique endomorphisme $g$ tel que~:
% %     \[\forall x\in A \quad g(x)=x \et \forall x\in B \quad g(x)=0\]
% %   \question Conclure.
% %   \end{questions}
% \end{questions}

% \begin{sol}
% Soit $f$ continue de $\R$ dans $\R$ vérifiant $f(x+y)=f(x)+f(y), \forall (x,y)\in \R^2$.
%   \begin{itemize}
%   \item On montre d'abord par récurrence sur $n\in \N$ $H_n : "f(nx)=nf(x)"$.
%   \item Soit $n\in \Z$ tel que $n\leq 0$. Alors $-n\in \N$. D'où $f(nx)+f(-nx)=f(0)=0$ donc $f(nx)=-f(-nx)=-(-nf(x))=nf(x)$. On a donc montré que $\forall n \in \Z, f(nx)=nf(x)$.
%   \item Soit $r\in \Q$. Il existe $p\in \Z$ et $q\in \Ns$ tels que $r=p/q$. On a $pf(x)=f(px)=f(qrx)=qf(rx)$ d'où $f(rx)p/qf(x)=rf(x)$.
%   Finalement, $\forall r\in \Q, f(r)=rf(1)$.
%   \item Soit $x\in \R$, il existe une suite $(r_n)$ de rationnels qui tend vers $x$ (par densité). Par continuité de $f$ en $x$, $r_nf(1)=f(r_n)\tendvers{n}{+\infty}f(x)$. Donc par unicité de la limite $f(x)=xf(1)$.
%   \item En posant $\alpha=f(1)$ on a donc $f(x)=\alpha x$.
%   \end{itemize}
% \end{sol}





\exercice{nom={Prolongement d'inégalités}}
\begin{questions}
\question Soit $f$ et $g$ deux fonctions de continues $\R$ dans $\R$ telles
  que
  \[\forall x\in\Q \qsep f(x)<g(x).\]
  \begin{questions}
  \question Montrer que
    \[\forall x\in\R \qsep f(x)\leq g(x).\]
  \question Montrer que l'on a pas nécessairement
    \[\forall x\in\R \qsep f(x)<g(x).\]
  \end{questions}
\question Soit $f$ une fonction continue de $\R$ dans $\R$ telle que
  \[\forall x,y\in\Q \qsep x<y \implique f(x)<f(y).\]
  Montrer que $f$ est strictement croissante, c'est-à-dire que
  \[\forall x,y\in\R \qsep x<y \implique f(x)<f(y).\]
\end{questions}

\begin{sol}
\begin{questions}
\question 
  \begin{questions}
  \question Suite de rationnels qui tend vers $x$ plus passage à la limite.
  \question $f=0$ et $g(x)=|x-\sqrt{2}|$.
  \end{questions}
\question Soit $x<y$. On peut trouver deux rationnels tels que $x<a<b<y$. Considérons alors $(x_n)$ et $(y_n)$ des suites de rationnels qui tendent vers $x$ et $y$. Pour $n$ assez grand, $x_n\leq a$ et donc $f(x_n)\leq f(a)$ et par passage à la limite $f(x)\leq f(a)$. De même on obtient $f(b)\leq f(y)$ et donc finalement :
$$f(x)\leq f(a)<f(b \leq f(y)$$ car $a$ et $b$ sont deux rationnels tels que $a<b$.
\end{questions}

\end{sol}

\magsubsection{Continuité sur une partie}
\magsubsection{Théorème des valeurs intermédiaires}

\exercice{nom={Équation fonctionnelle}}
Déterminer les fonctions $f$, continues de $\R$ dans $\R$, telles que
\[\forall x\in\R \qsep f(x)^2-2xf(x)-1=0.\]

\exercice{nom={Théorème de la corde raide}}
Soit $f$ une fonction continue de $\interf{0}{1}$ dans $\R$ telle que
$f\p{0}=f\p{1}$.
\begin{questions}
\question Montrer qu'il existe $x\in\interf{0}{\frac{1}{2}}$ tel que
  \[f\p{x+\frac{1}{2}}=f(x).\]
\question 
  Montrer que si un coureur parcourt 20~km en une
  une heure, il existe un intervalle de temps d'une demi-heure pendant lequel
  il a exactement parcouru 10~km.
\question Plus généralement, montrer que si $n\in\Ns$, il existe
  $x\in\interf{0}{1-\frac{1}{n}}$ tel que
  \[f\p{x+\frac{1}{n}}=f(x).\]
\question Si $\alpha\in\interf{0}{1}$, existe-t-il toujours
  $x\in\interf{0}{1-\alpha}$ tel que $f\p{x+\alpha}=f(x)$~?
\end{questions}

\begin{sol}
\begin{questions}
\question $g(x)=f(x+1/2)-f(x)$. $g(0)$ et $g(1)$ sont de signes opposées donc par continuité de $g$ et grâce au TVI, il existe ce $x$.
\question 
\begin{itemize}
\item[$\bullet$] \textbf{Première méthode :} On pose $f(x)=d(x)-20x$ et on applique la première question. 
\item[$\bullet$] \textbf{Deuxième méthode :}  $d(1)=20$ $d(0)=0$. On définit sur $[0,1/2], g(x)=d(x+1/2)-d(x)-10$ alors $g(0)=d(1/2)-10$ et $g(1/2)=10-d(1/2)$. Ils sont opposés...  
\end{itemize}
\question On considère la fonction $f_n$ définie sur $\interf{0}{1-1/n}$ par $$f_n(x)=f(x+1/n)-f(x).$$ On a alors $$0=f(1)-f(0)=\sum_{k=0}^{n-1}f_n\p{\frac{k}{n}}.$$ Comme les $f_n\p{\frac{k}{n}}$ ne peuvent pas être tous strictement positifs ou tous strictement négatifs, il existe $k\in \llbracket 0, n-2\rrbracket$ tel que $f_n\p{\frac{k}{n}}$ et $f_n\p{\frac{k+1}{n}}$ sont de signes opposés...
\question Définissons $f$ par $$f(x)=\cos\p{\frac{2\pi x}{\alpha}}-x\p{\cos\p{\frac{2\pi}{\alpha}}-1}.$$
Remarquons d'abord que $f(0)=f(1)=1$. Par ailleurs, s'il existait $x\in \interf{0}{1-\alpha}$ avec $f(x)=f(x+\alpha)$, on obtiendrait après simplifications $\cos\p{\frac{2\pi}{\alpha}}=1$ ce qui n'est pas possible si $1/\alpha$ n'est pas un entier.
\end{questions}
\end{sol}

\exercice{nom={Preuve du théorème des valeurs intermédiaires}}

Soit $f$ une fonction continue de $[a,b]$ dans $\R$. On suppose que $f(a)\leq 0$ et
$f(b)\geq 0$ et l'on souhaite montrer qu'il existe $c\in[a,b]$ tel que $f(c)=0$.
Comme c'est immédiat lorsque $f(b)=0$, on suppose que $f(b)>0$. On pose alors
\[X\defeq\enstq{x\in[a,b]}{f(x)\leq 0}.\]
\begin{questions}
\question Montrer que $X$ admet une borne supérieure que l'on note $c$, puis que
  $c\in[a,b]$.
\question Montrer que $f(c)\leq 0$.
\question Montrer que $c<b$, puis que $f(c)\geq 0$. Conclure.
\end{questions}

\exercice{nom={Monotonie des injections continues sur un intervalle}}
Soit $f:I\to\R$ une fonction continue et injective sur un intervalle $I$. Le but de cet
exercice est de montrer que $f$ est strictement monotone. On raisonne par l'absurde et
on suppose que $f$ ne l'est pas. On pose
\[A\defeq\ensim{(x,y)}{x,y\in I \quad\text{et}\quad x < y}.\]
\begin{questions}
\question Montrer qu'il existe $(x_0,y_0)\in A$ tel que $f(x_0)\leq f(y_0)$ et
  $(x_1,y_1)\in A$ tel que $f(x_1)\geq f(y_1)$.
\question Montrer soigneusement que
  \[\forall t\in[0,1]\qsep (t x_0+(1-t)x_1,t y_0+(1-t)y_1)\in A.\]
\question On définit la fonction $g$ de $[0,1]$ dans $\R$ par
  \[\forall  t\in[0,1]\qsep g(t)\defeq f(t x_0 +(1-t)x_1)-f(t y_0+(1-t)y_1).\]
  Montrer qu'il existe $t\in[0,1]$ tel que $g(t)=0$ et conclure.
\end{questions}

\magsubsection{Théorème de compacité}

\exercice{nom={Théorème des bornes atteintes}}
\begin{questions}
\question Montrer qu'une fonction continue sur $\R$ et périodique est bornée et
  atteint ses bornes.
\question Soit $a,b\in\R$ tels que $a<b$ et $f$ une fonction continue de $\interf{a}{b}$
  dans $\R$. Montrer que
  \[\sup_{x\in\interf{a}{b}} f(x) = \sup_{x\in\intero{a}{b}} f(x).\]
  On commencera par montrer que ces bornes supérieures existent.
\question Soit $f$ et $g$ deux fonctions continues définies sur $\interf{a}{b}$ à
  valeurs dans $\R$ telles que
  \[\forall x\in\interf{a}{b} \qsep 0<f(x)<g(x).\]
  Montrer qu'il existe $k\in\intero{0}{1}$ tel que
  \[\forall x\in\interf{a}{b} \qsep f(x)\leq kg(x).\]
\end{questions}

\begin{sol}
\begin{questions}
\question Soit $T > 0$ une période de $f$.
Sur $[0 ; T]$, $f$ est bornée par un certain $M$ car $f$ est continue sur un segment. Pour tout $x\in\R$, $x-nT \in [0,T]$ pour $n\ent{\frac{x}{T}}$ donc
$$|f(x)| = |f(x-nT)|\leq M.$$
Ainsi $f$ est bornée par $M$ sur $\R$.
\question 
\question Suffit d'appliquer le théorème de compacité à $f(x)/g(x)$. 
\end{questions}
\end{sol}


\exercice{nom={Théorème des bornes atteintes}}
Soit $f$ et $g$ deux fonctions continues de $\interf{a}{b}$ dans $\R$.
On définit $h:\R\to\R$ par
\[\forall t\in\R \qsep h(t)\defeq\sup_{x\in\interf{a}{b}}
  \p{f(x)+t g(x)}.\]
\begin{questions}
\question Montrer que $h$ est bien définie.
\question On pose 
\[M\defeq \sup_{x\in\interf{a}{b}} \abs{g(x)}.\]
Montrer que $h$ est $M$-lipschitzienne. En déduire qu'elle est continue.
\end{questions}

\begin{sol}
  $x\mapsto f(x)+tg(x)$ est définie et continue sur $[0,1]$ donc y est bornée et atteint ses bornes. Par suite, $\phi(t)$ est bien définie et plus précisément, il existe $x_t \in [0,1]$ tel que $\phi(t)=f(x_t)+tg(x_t)$. Puisque $g$ est continue sur $[0,1]$ elle y est bornée par un certain $M$. Pour tout $t,\tau \in [0,1]$, on a :
  $$\phi(t)-\phi(\tau)=f(x_t)+tg(x_t)-(f(x_\tau)+\tau g(x_\tau)).$$ Or, $$f(x_t)+\tau g(x_t)\leq f(x_\tau)+\tau g(x_\tau),$$ donc $$\phi(t)-\phi(\tau)\leq (t-\tau)g(x_t)\leq M|t-\tau|.$$
  De même, $$\phi(\tau)-\phi(t)\leq M|t-\tau|$$ et finalement $\phi$ est $M$-lipschitzienne.
\end{sol}

\exercice{nom={Exercice}}
Soit $f$ une fonction continue de $\R$ dans $\R$ admettant des limites finies
en $+\infty$ et $-\infty$.
\begin{questions}
\question Montrer que $f$ est bornée.
\question $f$ atteint-elle ses bornes~?
\end{questions}

\begin{sol}
\begin{questions}
\question Notons $\ell_1$ sa limite en $+\infty$ et $\ell_2$ en $-\infty$. Pour $\epsilon=1$, il existe $m\in \R$ tel que $\forall x\leq m, |f(x)-\ell_2|<1$ soit $|f(x)|\leq |\ell_2|+1$. Il existe aussi $M\in \R$ tel que $\forall x\geq M, |f(x)-\ell_1|<1$ soit $|f(x)|\leq |\ell_1|+1$.
\begin{itemize}
\item[$\bullet$] Si $m\geq M$, alors $\forall x \in \R$, $|f(x)|\leq \max(|\ell_1|+1,|\ell_2|+1)$.
\item[$\bullet$] Si $m<M$, $f$ est continue sur $[m,M]$ donc y est bornée par $A$. Alors, $\forall x \in \R$, $|f(x)|\leq \max(|\ell_1|+1,|\ell_2|+1,A)$.
\end{itemize}
\question Non, $\arctan$.
\end{questions}
\end{sol}

\magsubsection{Continuité uniforme}

\exercice{nom={Uniforme continuité}}
Soit $f$ une fonction uniformément continue de $\R$ dans $\R$. Montrer qu'il
existe $a,b\in\RP$ tels que
\[\forall x\in\R \qsep \abs{f(x)}\leq a+b\abs{x}.\]

% \exercice{nom={Uniforme continuité}}
% \begin{questions}
% \question Montrer que
%   \[\forall x,y\in\RP \qsep \abs{\sqrt{x}-\sqrt{y}}\leq\sqrt{\abs{x-y}}.\]
% \question En déduire que la fonction $x\mapsto \sqrt{x}$ définie sur $\RP$
%    est uniformément continue.
% \end{questions}

\exercice{nom={Fonctions Hölderiennes}}
Soit $\alpha>0$. On dit qu'une fonction $f$ est $\alpha$-Hölderienne lorsqu'il
existe $c\in\RP$ tel que
\[\forall x,y\in\R \qsep \abs{f(x)-f(y)}\leq c\abs{x-y}^{\alpha}.\]
\begin{questions}
\question Montrer que si $f$ est $\alpha$-Hölderienne avec $\alpha>1$, $f$ est
  constante.
\question Montrer que si $f$ est $\alpha$-Hölderienne, $f$ est uniformément
  continue.
\end{questions}


\exercice{nom={Généralisation du théorème de Heine}}
Soit $f$ une fonction continue de $\R$ dans $\R$ telle que~:
\[f(x) \tendvers{x}{-\infty} l_1 \et f(x) \tendvers{x}{+\infty} l_2\]
Le but de cet exercice est de montrer que $f$ est uniformément continue.
\begin{questions}
\question Soit $\epsilon>0$. Montrer qu'il existe $a,b\in\R$, avec $a\leq b$
  tels que
  \[\forall x,y \in\interof{-\infty}{a} \qsep \abs{f(x)-f(y)}\leq\epsilon
    \et
    \forall x,y \in\interfo{b}{+\infty} \qsep \abs{f(x)-f(y)}\leq\epsilon.\]
\question Montrer qu'il existe $\eta>0$ tel que
  \[\forall x,y\in\interf{a}{b} \qsep \abs{x-y}\leq\eta \implique
    \abs{f(x)-f(y)}\leq\epsilon.\]
\question Conclure.
\end{questions}
%END_BOOK

\end{document}

\chapter{Anneaux, corps, polynômes}
\setcounter{numeroexercicecours}{1}
\documentclass{magnolia}

\magtex{tex_driver={pdftex},
        tex_packages={xypic}}
\magfiche{document_nom={Cours sur les structures algebriques},
          auteur_nom={François Fayard},
          auteur_mail={fayard.prof@gmail.com}}
\magcours{cours_matiere={maths},
          cours_niveau={mpsi},
          cours_chapitre_numero={12},
          cours_chapitre={Anneaux, Corps, Polynômes}}
\magmisenpage{misenpage_presentation={tikzvelvia},
          misenpage_format={a4},
          misenpage_nbcolonnes={1},
          misenpage_preuve={non},
          misenpage_sol={non}}
\maglieudiff{}
\magprocess


\begin{document}

%BEGIN_BOOK
\magtoc


\section{Anneau, corps}
\subsection{Anneau}

\begin{definition}
Soit $\p{A,+}$ un groupe commutatif (d'élément neutre $0_A$) et $\times$ une
loi de composition interne sur $A$. On dit que $\p{A,+,\times}$ est un \emph{anneau}
lorsque
\begin{itemize}
\item $\times$ est associatif,
\item $\times$ admet un élément neutre $1_A$,
\item $\times$ est distributive par rapport à $+$
  \begin{eqnarray*}
  \forall a,b,c\in A, & & a\times\p{b+c}=a\times b+a\times c,\\    
                     & & \p{a+b}\times c=a\times c+b\times c.
  \end{eqnarray*}
\end{itemize}
Un élément $a\in A$ est dit \emph{inversible} lorsqu'il est inversible pour la loi
$\times$. Un anneau $\p{A,+,\times}$ est dit \emph{commutatif} lorsque $\times$ est
commutative.
\end{definition}

\begin{remarqueUnique}
\remarque Si $a,b\in A$, on dit que $a$ et $b$ \emph{commutent} lorsque $a\times b=b\times a$.
\end{remarqueUnique}

\begin{exemples}
\exemple $(\C,+,\times)$, $(\R,+,\times)$ et $(\Z,+,\times)$ sont des anneaux commutatifs.
\exemple Si $(A,+,\times)$ est un anneau et $X$ est un ensemble,
  l'ensemble $\mathcal{F}(X,A)$ des fonctions de $X$ dans $A$, muni des lois $+$ et $\times$ définies par
  \begin{eqnarray*}
  \forall f,g\in\mathcal{F}(X,A) \qsep \forall x\in X \qsep (f+g)(x)&\defeq&f(x)+g(x),\\
                                              (f\times g)(x)&\defeq&f(x)\times g(x)
  \end{eqnarray*}
  est un anneau dont l'élément neutre pour l'addition est la fonction $x\mapsto 0_A$ et
  l'élément neutre pour la multiplication est $x\mapsto 1_A$. Si $A$ est commutatif, alors
  $\mathcal{F}(X, A)$ l'est aussi. En particulier, $(\mathcal{F}(\R,\R),+,\times)$ et
  $(\R^\N,+,\times)$ sont des anneaux commutatifs.
\exemple Si $n\in\N$, alors $(\mat{n}{\R},+,\times)$ est un anneau dont l'élément neutre
  pour la multiplication est la matrice $I_n$. Il est non commutatif dès que
  $n\geq 2$. Les éléments inversibles de $\mat{n}{\R}$ sont les matrices inversibles. 
\exemple Si $E$ est un \Kev, alors $(\mathcal{L}(E),+,\circ)$ est un anneau dont l'élément
  neutre est ${\rm Id}$. En général, il n'est pas commutatif. Les éléments inversibles
  de $\mathcal{L}(E)$ sont les isomorphismes.
\exemple Soit $\mathbb{F}_2$ l'ensemble à deux éléments $\ens{\bar{0},\bar{1}}$. On
  définit sur $\mathbb{F}_2$ les lois $+$ et $\times$ par
  \begin{center}
  \begin{tabular}{|c|c|c|}
  \hline
  $+$\vphantom{\Large{A}} & $\bar{0}$ & $\bar{1}$\\
  \hline
  $\bar{0}$\vphantom{\Large{A}} & $\bar{0}$ & $\bar{1}$\\
  \hline
  $\bar{1}$\vphantom{\Large{A}} & $\bar{1}$ & $\bar{0}$\\
  \hline
  \end{tabular}
  \hspace{2cm}
  \begin{tabular}{|c|c|c|}
  \hline
  $\times$\vphantom{\Large{A}} & $\bar{0}$ & $\bar{1}$\\
  \hline
  $\bar{0}$\vphantom{\Large{A}} & $\bar{0}$ & $\bar{0}$\\
  \hline
  $\bar{1}$\vphantom{\Large{A}} & $\bar{0}$ & $\bar{1}$\\
  \hline
  \end{tabular}
  \end{center}
  Alors $(\mathbb{F}_2,+,\times)$ est un anneau commutatif.
\end{exemples}



%% Exemple :
%% 1) Z, R, C sont des anneaux
%% 2) Si A est un anneau et X est en ensemble non vide, F(X,A) est un anneau


\begin{proposition}
Soit $\p{A,+,\times}$ un anneau. Alors
\begin{eqnarray*}
\forall a\in A, & & 0_A \times a=0_A \quad\text{et}\quad a \times 0_A = 0_A\\
\forall a,b\in A, & & a\times\p{-b}=\p{-a}\times b=-\p{a\times b}\\
\forall a,b\in A \qsep \forall n\in\Z, & &
  \p{n\cdot a}\times b=a\times\p{n\cdot b}=n\cdot\p{a\times b}.
\end{eqnarray*}
\end{proposition}

\begin{preuve}
Clés : 
$$0\times a=(0+0)\times a=0\times a + 0\times a.$$
$$a(-b)+ab=a(-b+b)=0 \text{ donc } a(-b)=-(ab)$$
et $$(-a)b+ab=(-a+a)b=0 \text{ donc } (-a)b=-(ab)=a(-b).$$
\end{preuve}

\begin{remarqueUnique}
\remarque Soit $\p{A,+,\times}$ un anneau dans lequel $0_A=1_A$. Alors $A=\ens{0_A}$. Réciproquement, si $A$ est un ensemble contenant un unique élément muni des seules lois de composition interne $+$ et $\times$ que l'on peut définir sur cet ensemble, alors $A=\ens{0_A}$ et $(A,+,\times)$ est un anneau. On dit que cet anneau est l'anneau \emph{trivial}.
\end{remarqueUnique}

\begin{proposition}
Soit $\p{A,+,\times}$ un anneau et $a,b\in A$ tels que $a\times b=b\times a$.
Alors, pour tout $n\in\N$
\[\p{a+b}^n = \sum_{k=0}^n \binom{n}{k} \cdot \p{a^{n-k} \times b^k} \quad\et\quad
  a^n-b^n=\p{a-b}\times\cro{\sum_{k=0}^{n-1} a^{\p{n-1}-k}\times b^k}.\]
\end{proposition}

\begin{remarques}
\remarque Ces relations peuvent être fausses lorsque $a$ et $b$ ne commutent
  pas. Par exemple, si $a$ et $b$ sont deux éléments d'un anneau, alors
  \[(a+b)^2=a^2+2\cdot a\times b+b^2 \quad\ssi\quad a\times b=b\times a.\]
\remarque Remarquons que si $a\in A$, alors $a$ commute
  avec $1_A$, donc ces formules sont valables pour développer $(1_A+a)^n$ et
  factoriser $a^n-1_A$.
\end{remarques}

\begin{definition}
On dit qu'un élément $a\in A$ est nilpotent lorsqu'il existe $n\in\N$ tel que $a^n=0_A$.
\end{definition}

\begin{exoUnique}
\exo Montrer que si $x$ est nilpotent, alors $1_A-x$ est inversible.
\end{exoUnique}
\begin{sol}
On vérifie que son inverse à gauche et à droite est $\displaystyle\sum_{k=0}^{n-1}x^k$.
\end{sol}

\begin{definition}
Soit $\p{A,+,\times}$ un anneau. L'ensemble $U_A$ des éléments inversibles de $A$
est un groupe pour la multiplication.
\end{definition}

\begin{remarqueUnique}
\remarque L'ensemble $U_A$ des inversibles de $A$ est parfois noté $A^\times$. Il est important de ne pas
  confondre cet ensemble avec $A^*\defeq A\setminus\ens{0_A}$.
\end{remarqueUnique}

\begin{exemples}
\exemple Le groupe des inversibles de $\Z$ est $(\ens{-1,1},\times)$.
\exemple Le groupe des inversibles de $\mat{n}{\R}$ est $(\gl{n}{\R},\times)$. Le
  groupe des inversibles de $\mathcal{L}(E)$ est $({\rm GL}(E),\circ)$.
\end{exemples}

\begin{sol}
Exemple de $(\Z,+,.)$ pour qui ses éléments inversibles $\set{-1;1}$ forment un groupe pour la multiplication.
\end{sol}


\begin{definition}
On dit qu'un anneau $\p{A,+,\times}$ est \emph{intègre} lorsque
\begin{itemize}
\item $1_A\neq 0_A$
\item $\times$ est commutative
\item $\forall a,b\in A \qsep a\times b=0_A \implique \cro{a=0_A \ou b=0_A}.$
\end{itemize}
\end{definition}

\begin{remarqueUnique}
\remarque Si $(A,+,\times)$ est intègre, tout élément non nul $a\in A^*$ est régulier pour $\times$~:
    \[\forall b,c \in A\qsep a\times b=a\times c \quad\implique\quad b=c.\]
\end{remarqueUnique}

\begin{exoUnique}
\exo L'anneau $(\mathcal{F}(\R,\R),+,\times)$ est-il intègre~?
\end{exoUnique}

\begin{sol}
Non, évidemment. On peut prendre les indicatrices de deux parties disjointes. Le produit est nul sans que l'une des deux ne soient nulle.
\end{sol}

% \begin{remarqueUnique}
% \remarque 
Dans la suite, lorsqu'il n'y a pas de confusion possible,
  les éléments $0_A$ et $1_A$ seront le plus souvent notés $0$ et $1$.
% \end{remarqueUnique}

\begin{definition}
Soit $\p{A,+,\times}$ un anneau et $B$ une partie de $A$. On dit que $B$ est un
\emph{sous-anneau} de $A$ lorsque
\begin{itemize}
\item $0\in B \et 1\in B$
\item $\forall b_1,b_2\in B \qsep b_1+b_2\in B \qsep -b_1\in B \et b_1\times b_2\in B.$
\end{itemize}
Si tel est le cas $\p{B,+,\times}$ est un anneau.
\end{definition}

\begin{remarques}
\remarque Si $B$ est un sous-anneau de $(A,+,\times)$, $B$ est un
  sous-groupe de $(A,+)$.
\remarque Si $B$ est un sous-anneau de $\C$, alors $\Z\subset B$.
\end{remarques}

\begin{exoUnique}
\exo Montrer que $\Z[\ii]\defeq\ensim{a+\ii b}{a,b\in\Z}$ est un sous-anneau de $\C$.
\end{exoUnique}

\begin{definition}
Soit $\p{A,+,\times}$ et $\p{B,+,\times}$ deux anneaux. On dit qu'une
application $\phi$ de $A$ dans $B$ est un \emph{morphisme d'anneau} lorsque
\begin{eqnarray*}
\forall a_1,a_2\in A, & & \phi\p{a_1+ a_2}=\phi\p{a_1}+\phi\p{a_2}\\
\forall a_1,a_2\in A, & & \phi\p{a_1\times a_2}=\phi\p{a_1}\times\phi\p{a_2}\\
& & \phi\p{1_A}=1_B.
\end{eqnarray*}
\end{definition}

\begin{proposition}
Soit $\phi$ un morphisme d'anneau de $\p{A,+,\times}$ dans $\p{B,+,\times}$.
Alors
\begin{eqnarray*}
\forall a\in A \qsep \forall n\in\Z, & & \phi\p{n\cdot a}=n\cdot\phi(a)\\
\forall a\in A \qsep \forall n\in\N, & & \phi\p{a^n}=\cro{\phi(a)}^n.
\end{eqnarray*}
De plus, si $a\in A$ est inversible, il en est de même pour $\phi(a)$ et
\[\forall n\in\Z \qsep \phi\p{a^n}=\cro{\phi(a)}^n.\]
\end{proposition}

% \begin{remarques}
% \remarque Soit $\phi$ un morphisme d'anneau de $(A,+,\cdot)$ dans
%   $(B,+,\cdot)$.
% \end{remarques}

\begin{proposition}
\begin{itemize}
\item La composée de deux morphismes d'anneaux est un morphisme d'anneau.
\item La bijection réciproque d'un isomorphisme est un isomorphisme.
\end{itemize}
\end{proposition}

\begin{proposition}
Soit $\phi$ un isomorphisme de l'anneau $(A,+,\times)$ dans l'anneau $(B,+,\times)$. Alors
\[\forall x\in A\qsep x\in U_A \ssi \phi(x)\in U_B.\]
De plus $\phi$ induit un isomorphisme du groupe $(U_A,\times)$ dans le groupe $(U_B,\times)$.
\end{proposition}

\begin{definition}
On dit qu'une partie $\mathcal{I}$ d'un anneau commutatif $(A,+,\times)$ est un \emph{idéal} de $A$
lorsque
\begin{itemize}
\item $0\in\mathcal{I}$
\item $\forall x,y\in\mathcal{I}\qsep \forall a,b\in A\qsep ax+by\in\mathcal{I}$
\end{itemize}
\end{definition}

\begin{proposition}
Soit $(A,+,\times)$ un anneau commutatif et $x\in A$. Alors
\[x A\defeq\ensim{ax}{a\in A}\]
est un idéal de $A$. Un tel idéal est appelé idéal \emph{principal}.
\end{proposition}

\begin{remarqueUnique}
\remarque Soit $\mathcal{I}$ un idéal de $\Z$. Alors, il existe un unique $n\in\N$
  tel que $\mathcal{I}=n\Z$. Dans $\Z$, tout idéal est donc principal.
\end{remarqueUnique}

\subsection{Corps}

\begin{definition}
On dit qu'un anneau $\p{\K,+,\times}$ est un \emph{corps} lorsque
\begin{itemize}
\item $1_{\K}\neq 0_{\K}$
\item $\times$ est commutative
\item Tout élément non nul de $\K$ admet un inverse pour la loi $\times$.
\end{itemize}
\end{definition}

\begin{exemples}
\exemple Muni des lois usuelles d'addition et de multiplication, $\C$ est un corps.
\exemple $(\mathbb{F}_2,+,\times)$ est un corps.
\end{exemples}

\begin{remarqueUnique}
\remarque Si $\K$ est un corps, l'ensemble des
  inversibles de $\K$ est $\Ks\defeq\K\setminus\ens{0_{\K}}$. Autrement dit, $\K^{\times}=\Ks$.
\end{remarqueUnique}

\begin{proposition}
Un corps est intègre.  
\end{proposition}

\begin{exoUnique}
\exo Résoudre l'équation $x^2=1$ sur le corps $\K$.
\end{exoUnique}


\begin{definition}
Soit $\p{\KL,+,\times}$ un corps et $\K$ une partie de $\KL$. On dit que $\K$
est un \emph{sous-corps} de $\KL$ lorsque
\begin{itemize}
\item $\K$ est un sous-anneau de $\KL$
\item $\forall x\in\K\setminus\ens{0} \qsep x^{-1}\in\K$.
\end{itemize}
Si tel est le cas, $\p{\K,+,\times}$ est un corps.
\end{definition}

\begin{remarqueUnique}
\remarque $\Q$ et $\R$ sont des sous-corps de $\C$.
\remarque Si $\K$ est un sous-corps de $\C$, alors $\Q\subset \K$.    
\end{remarqueUnique}

\begin{definition}
Si $\p{\K,+,\times}$ et $\p{\KL,+,\times}$ sont deux corps, on appelle morphisme
de corps de $\K$ dans $\KL$ tout morphisme d'anneau pour les structures
sous-jacentes.
\end{definition}

\begin{remarques}
\remarque Si $\phi$ est un morphisme d'un sous-corps $\K$ de $\C$ dans un sous-corps
  $\KL$ de $\C$, alors
  \[\forall r\in\Q\qsep \forall x\in\K\qsep \phi(rx)=r\phi(x).\]
\remarque Si $\phi$ est un morphisme de corps, alors $\phi$ est injective.
\end{remarques}
\begin{sol}
On considère $\phi:\K\to \KL$ un morphisme de corps. En particulier, c'est un morphisme de groupes. Raisonnons par l'absurde et supposons qu'il existe $x\neq 0$ tel que $\phi(x)=0_{\KL}$. Comme $x$ est différent de $0$ et que $\K$ est un corps, alors $xx^{-1}=1_{\K}$ donc $$0_{\KL}=0_{\KL}\phi(x^{-1})=\phi(x)\phi(x^{-1})=\phi(xx^{-1})=\phi(1_{\K})=1_{\KL}$$ ce qui est absurde car $\KL$ est un corps et donc $0_{\KL}\neq 1_{\KL}$.
\end{sol}


\begin{exoUnique}
\exo Déterminer les morphismes de corps $\phi$ de $\C$ dans $\C$ tels
  que~: $\forall x\in\R \qsep \phi(x)=x$.
\end{exoUnique}

\begin{sol}
On va montrer que $\phi$ est ou bien l'identité ou bien la conjugaison. On raisonne par analyse-synthèse.
\begin{itemize}
\item[$\bullet$] \textbf{Analyse :} On a $\phi(i)^2=\phi(i^2)=\phi(-1)=-\phi(1)=-1$ donc $\phi(i)=i$ ou $\phi(i)=-i$. On montre alors que si $\phi(i)=i$ alors $\phi$ est l'identité et si $\phi(i)=-i$, c'est la conjugaison.
\item[$\bullet$] \textbf{Synthèse :} On vérifie que ce sont bien des morphismes de corps qui vérifie l'hypothèse.
\end{itemize}
\end{sol}

\begin{definition}
Soit $\K$ un corps. Alors, l'application
\[\dspappli{\phi}{\Z}{\K}{k}{k\cdot 1_{\K}}\]
est un morphisme du groupe $(\Z,+)$ dans $(\K,+)$. Il existe donc un unique $p\in\N$ tel que $\ker\phi=p\Z$.
L'entier $p$ est soit nul soit un nombre premier et est appelé \emph{caractéristique} de $\K$.
\end{definition}

\begin{remarques}
\remarque Les sous-corps de $\C$ sont de caractéristique nulle. Le corps $\mathbb{F}_2$ est de caractéristique 2.
\remarque Lorsque $\K$ est un sous-corps de $\C$, pour tout $k\in\Z$, on a $k\cdot 1_{\K}=k$. Si $\K$ est un corps
  quelconque, on confondra le plus souvent $k\cdot 1_{\K}$ et $k$.
\end{remarques}

\section{Espace vectoriel, algèbre}
\subsection{Espace vectoriel}


\begin{definition}[utile=-3]
Soit $\K$ un corps, $\p{E,+}$ un groupe commutatif d'élément neutre $0_E$ et
$\cdot$ une loi de composition externe.
\[\dspappli{\cdot}{\K\times E}{E}{\p{\lambda,x}}{\lambda\cdot x}\]
On dit que $\p{E,+,\cdot}$ est un \emph{\Kev} lorsque
\begin{eqnarray*}
\forall x,y\in E \qsep \forall \lambda\in\K, & &
  \lambda\cdot\p{x+y}=\lambda\cdot x+\lambda\cdot y\\
\forall x\in E \qsep \forall \lambda,\mu\in\K, & &
  \p{\lambda+\mu}\cdot x=\lambda\cdot x+\mu\cdot x\\
\forall x\in E \qsep \forall \lambda,\mu\in\K, & &
  \lambda\cdot\p{\mu\cdot x}=\p{\lambda\mu}\cdot x\\
\forall x\in E, & & 1\cdot x=x.
\end{eqnarray*}
Les éléments de $\K$ sont appelés \emph{scalaires}, ceux de $E$, \emph{vecteurs}.
\end{definition}

\begin{remarqueUnique}
\remarque L'ensemble du cours sur les espaces vectoriels reste valide pour un corps quelconque, excepté le paragraphe
  sur les symétries qui n'est vrai que pour les corps de caractéristique différente de 2.
\end{remarqueUnique}

\begin{proposition}[utile=1]
Soit $\p{E,+,\cdot}$ un \Lev et $\K$ un sous-corps de $\KL$. Alors
$\p{E,+,\cdot}$ est un \Kev. En particulier $\KL$ est un \Kev. 
\end{proposition}

\begin{remarques}
\remarque[utile=-2] $\C$ est un \Rev.
\remarque[utile=-1] Muni des lois usuelles, $\mathcal{F}(\R,\C)$ est un \Cev.
  Comme $\R$ est un sous-corps de $\C$, $\mathcal{F}(\R,\C)$ est aussi un \Rev.
\end{remarques}


\subsection{Algèbre}

\begin{definition}
On dit qu'un anneau $\p{A,+,\times}$ muni d'une loi de composition externe
$\cdot$ sur un corps $\K$ est une $\K$-algèbre lorsque
\begin{itemize}
\item $\p{A,+,\cdot}$ est un \Kev
\item $\times$ est compatible avec la loi de composition externe
  \[\forall x,y\in A \qsep \forall \lambda\in\K \qsep
    \p{\lambda\cdot x}\times y=x\times\p{\lambda\cdot y}=
     \lambda\cdot\p{x\times y}\]
\end{itemize}
On dit que l'algèbre $\p{A,+,\cdot,\times}$ est commutative lorsque $\times$
est commutatif.
\end{definition}

\begin{exemples}
\exemple $\K$ est une $\K$-algèbre.
\exemple Soit $X$ un ensemble. Alors $(\mathcal{F}(X,\K),+,\cdot,\times)$ est une $\K$-algèbre commutative. En particulier, l'ensemble des fonctions de $\R$ dans $\R$ est
une $\R$-algèbre commutative et l'ensemble des suites réelles est une $\R$-algèbre
commutative.
\end{exemples}
% Exemples :
% 1) La R-algèbre des fonctions de R dans R

\begin{definition}
Soit $\p{A,+,\cdot,\times}$ et $\p{B,+,\cdot,\times}$ deux $\K$-algèbre. On dit qu'une application $\phi$ de $A$ dans $B$ est un morphisme d'algèbre lorsque c'est un morphisme
d'anneau et une application linéaire, c'est-à-dire lorsque
\begin{eqnarray*}
  \forall x,y\in A \qsep \forall \lambda,\mu\in\K, & &
    \phi\p{\lambda x+\mu y}=\lambda \phi(x)+\mu \phi(y)\\
  \forall x,y\in A, & &
    \phi\p{xy}=\phi(x)\phi(y)\\
  & & \phi(1_A)=1_B.
  \end{eqnarray*}
\end{definition}


\begin{proposition}
Soit $E$ un \Kev. Alors $(\Endo{E},+,\cdot,\circ)$ est une $\K$-algèbre.
\end{proposition}

% \begin{remarqueUnique}
% \remarque Le groupe des inversibles de $\Endo{E}$ est $(\gl{}{E},\circ)$.
% \end{remarqueUnique}

\begin{proposition}
Soit $\K$ un corps et $n\in\N$. Alors $(\mat{n}{\K},+,\cdot,\times)$ est une $\K$-algèbre.
\end{proposition}

% \begin{remarqueUnique}
% \remarque Le groupe des inversibles de $\mat{n}{\K}$ est $(\gl{n}{\K},\times)$.
% \end{remarqueUnique}

\section{L'algèbre $\polyK$}

\subsection{Définition}
\begin{definition}
Soit $\K$ un corps. Alors il existe une unique algèbre commutative $\polyK$ ainsi
qu'un élément $X\in\polyK$, appelé \emph{indéterminée}, tels que
\begin{itemize}
\item Pour tout $P\in\polyK$, il existe $n\in\N$ et $a_0,\ldots,a_n\in\K$ tels
  que
  \[P=a_0+a_1 X+\cdots+a_n X^n\]
  où, par abus de notation, $a_0=a_0\cdot 1_{\polyK}=a_0 X^0$.
\item Pour tout $n\in\N$ et $a_0,\ldots,a_n\in\K$
  \[a_0+a_1X+\cdots+a_n X^n=0 \quad\implique\quad  a_0=\cdots=a_n=0.\]
\end{itemize}
On l'appelle \emph{algèbre des polynômes à coefficients dans $\K$}.
\end{definition}

\begin{remarqueUnique}
\remarque Soit $P\in\polyK$ et $a_0,\ldots,a_n\in\K$ tels que
  $P=a_0+a_1 X+\cdots+a_n X^n$.
  Il est d'usage de définir $a_k$ pour tout $k>n$ en posant $a_k\defeq 0$.
  Alors, quel que soit $q\geq n$
  \[P=\sum_{k=0}^q a_kX^k.\]
  D'autre part, si $b_0,\ldots,b_m\in\K$ sont tels que
  $P=b_0+b_1X+\cdots+b_m X^m$, en définissant $b_k\defeq 0$
  pour tout $k>m$, on a 
  \[\forall k\in\N\qsep a_k=b_k.\]
  Aurement dit, la suite $(a_k)$ est unique; on l'appelle famille des
  \emph{coefficients de $P$}.
% \remarque Soit $P\in\polyK$, $a_0,\ldots,a_n\in\K$ et $b_0,\ldots,b_m\in\K$ tels
%   que $P=a_0+a_1X+\cdots+a_n X^n$ et . Alors,
%   en prolongeant la définition des $a_k$ et $b_k$ comme dans la remarque précédente, on a
%   La suite $(a_k)_{k\in\N}$ est nulle à partir d'un certain rang et est appelée suite des 
%   Deux polynômes sont égaux si
%   et seulement si ils ont les même coefficients.
  % Si on prolonge
  % les définitions des suites $a$ et $b$ en posant $a_k=0$ pour $k>n$ et $b_k=0$
  % pour $k>m$, alors les suites $\p{a_k}$ et $\p{b_k}$ sont égales.
  % On dit que les $a_k$ sont les coefficients du polynôme $P$.
% \remarque Deux polynômes sont égaux si et seulement si ils ont les mêmes
%   coefficients.
% \remarque L'application de $\K$ dans $\polyK$ qui à $\lambda\in\K$ associe
%   $\lambda\cdot 1_{\polyK}\in\polyK$ est un morphisme injectif d'algèbre. Il
%   permet
%   d'identifier $\K$ à une sous-algèbre de $\polyK$. Par abus de langage, on
%   dit que $\K$ est une sous-algèbre de $\polyK$.
\end{remarqueUnique}

\begin{proposition}[nom={Produit de Cauchy}]
Soit $P,Q\in\polyK$ deux polynômes dont les coefficients sont respectivement
$(a_k)$ et $(b_k)$. Alors, si on note $(c_k)$ la famille des coefficients du
produit $PQ$, on a
\[\forall k\in\N\qsep c_k=\sum_{l=0}^k a_{k-l} b_l.\]
\end{proposition}

\begin{remarqueUnique}
\remarque On a donc
  \[\p{\sum_{k=0}^n a_k X^k}\p{\sum_{k=0}^m b_k X^k}=
    \sum_{k=0}^{n+m} \p{\sum_{l=0}^k a_{k-l} b_l}X^k\]
  si on utilise la convention $a_k\defeq 0$ pour $k>n$ et $b_k\defeq 0$
  pour $k>m$.
\end{remarqueUnique}

\begin{exoUnique}
\exo Soit $n\in\N$. Montrer que
  \[\sum_{k=0}^{n} \binom{n}{k}^2=\binom{2n}{n}\]
  en calculant $(1+X)^{2n}$ de deux manières différentes.
  \begin{sol}
  Il suffit de développer la relation $\p{1+X}^{2n}=\p{1+X}^n\p{1+X}^n$ avec
  le binôme de Newton.
    \end{sol}
  
\end{exoUnique}

\subsection{Substitution}

\begin{definition}
Soit $\mathcal{A}$ une $\K$-algèbre, $x\in\mathcal{A}$ et
$P=a_0+a_1X+\cdots+a_nX^n\in\polyK$. On définit $P(x)$ par
\[P(x)\defeq a_0 1_{\mathcal{A}}+a_1 x+\cdots+a_n x^n\in\mathcal{A}.\]
On dit que l'on a substitué l'élément $x\in\mathcal{A}$ à l'indéterminée $X$.
\end{definition}

\begin{remarques}
\remarque Si $\mathcal{A}$ une $\K$-algèbre et $x\in\mathcal{A}$,
  l'application $\phi$ de $\polyK$ dans $\mathcal{A}$ qui à $P$ associe
  $P(x)$ est un morphisme d'algèbre. Autrement dit
  \begin{eqnarray*}
  \forall P,Q\in\polyK \qsep \forall \lambda,\mu\in\K, & &
    \p{\lambda P+\mu Q}(x)=\lambda P(x)+\mu Q(x)\\
  \forall P,Q\in\polyK, & &
    \p{PQ}(x)=P(x)Q(x)\\
  & & 1_{\polyK}(x)=1_{\mathcal{A}}.
  \end{eqnarray*}
% \remarque Si $x\in\mathcal{A}$ et $n\in\Ns$, le calcul naïf de $x^n$ nécessite
%   $n-1$ multiplications dans $\mathcal{A}$. Si
%   $P=a_0+a_1X+\cdots+a_nX^n\in\polyK$, le calcul de $P(x)$ nécessite
%   donc $n\p{n-1}/2$ multiplications dans $\mathcal{A}$. Cependant, si on
%   écrit
%   \[P(x)=\p{\p{\cdots\p{\p{a_n x+a_{n-1}}x+a_{n-2}}x+\cdots+a_2}x+a_1}x+a_0\]
%   le calcul de $P(x)$ nécessite $n-1$ multiplications dans $\mathcal{A}$.
%   Cette méthode de calcul est connue sous le nom d'algorithme de Hörner.
\remarque On dit qu'un polynôme $P$ est un polynôme annulateur de
  $x\in\mathcal{A}$ lorsque $P(x)=0$. Par exemple, si $\K=\R$ et
  $\mathcal{A}=\C$, $P\defeq X^2+1$ est un polynôme annulateur de $\ii$. Si $E$
  est un \Kev et si $s\in\Endo{E}$ est une symétrie, alors $P\defeq X^2-1$ est un
  polynôme annulateur de $s$.
\remarque On dit qu'un élément $z\in\C$ est algébrique lorsqu'il existe un polynôme
  non nul $P\in\polyQ$ tel que $P(z)=0$. Par exemple $z_1=\p{1+\sqrt{5}}/2$
  est algébrique car $P_1\defeq X^2-X-1\in\polyQ$
  est un polynôme annulateur de $z_1$. De même,
  $\jj$ est algébrique car $P_2\defeq X^3-1\in\polyQ$
  est un polynôme annulateur de $\jj$. Lorsqu'on
  effectue des calculs avec un nombre algébrique $z$, il est souvent plus économe en
  calculs d'exploiter le fait que $P(z)=0$ plutôt que de remplacer $z$ par une
  expression parfois complexe. Par exemple, si $x\defeq\p{1+\sqrt{5}}/2$, en exploitant le fait que
  $x^2=x+1$, on a
  \[\p{\frac{1+\sqrt{5}}{2}}^3
    =x^3=x\cdot x^2=x\p{x+1}=x^2+x=2x+1=2+\sqrt{5}\]
  Si on souhaite calculer $1/x$, on exploite le fait que $x^2-x-1=0$, ce qui donne $x\p{x-1}=1$, puis $1/x=x-1$. Donc
  \[\frac{1}{\p{\frac{1+\sqrt{5}}{2}}}=\frac{-1+\sqrt{5}}{2}.\]

\remarque On dit qu'un élément de $\C$ est \emph{transcendant} lorsqu'il n'est pas
  algébrique. On peut montrer, mais c'est difficile, que $\e$ et $\pi$ sont
  transcendants.
\end{remarques}

\begin{exoUnique}
\exo Montrer que $1+\sqrt{7}$ et $\sqrt{2}+\sqrt{5}$ sont algébriques.
  \begin{sol}
  On pose $x=1+\sqrt{7}$ et on part de $(x-1)^2=7$. De même avec $x=\sqrt{2}+\sqrt{5}$, on l'élève au carré puis on élève $x^2-7$ au carré.
  Les polynômes minimaux sont $X^2-2X-6$ et $X^4-14X^2+9$.
  \end{sol}
\exo Soit $D\defeq\diag{\lambda_1,\ldots,\lambda_n}\in\mat{n}{\K}$ et $P\in\polyK$. Montrer que
  $P(D)=\diag{P\p{\lambda_1},\ldots,P\p{\lambda_n}}$.  
  \begin{sol}
  L'expliquer et ajouter qu'alors $P=\prod_{k=1}^n(X-\lambda_k)$ est un polynôme annulateur de $D$.
  \end{sol}
\end{exoUnique}

%% Algorithme de Hörner
%%
%% - exemple : calcul de 1+2x-x^2+3x^3 : 3 multiplications entre réels
%%             1+x(2+x(-1+x*3))        : 2 multiplications entre réels
%% - implémentation en Maple
%%   Si P est une liste contenant les coefficients
%%
%%   horner:=proc(P,x)
%%     if P:=[] then
%%       return(0);
%%     else
%%       return(P[1]+x*horner(P[2..-1],x));
%%     end if;
%%   end proc;

\begin{definition}
Soit $P,Q\in\polyK$. On définit le polynôme $P\circ Q$ par
\[P\circ Q\defeq P(Q).\]  
\end{definition}

\begin{remarqueUnique}
\remarque Si $P\in\polyK$, $P(X)=P$. Un polynôme peut donc indifféremment
  être noté $P$ ou $P(X)$.
\end{remarqueUnique}

% \begin{definition}
% Soit $\mathcal{A}$ une $\K$-algèbre et $x\in\mathcal{A}$. Alors,
% il existe une plus petite sous-algèbre de $\mathcal{A}$ contenant $x$; on
% l'appelle algèbre engendrée par $x$ et on la note $\K[x]$. De plus~:
% \[\K[x]=\enstq{P(x)}{P\in\polyK}\]
% En particulier $\K[x]$ est commutative.
% \end{definition}

% \begin{definition}
% Soit $x\in\mathcal{A}$ et $P\in\polyK$. On dit que $P$ est un polynôme
% annulateur de $x$ lorsque $P(x)=0$.
% \end{definition}

\begin{definition}
Soit $P\in\polyK$. On dit que
\begin{itemize}
\item $P$ est \emph{pair} lorsque $P\p{-X}=P(X)$.
\item $P$ est \emph{impair} lorsque $P\p{-X}=-P(X)$.
\end{itemize}
\end{definition}

\begin{proposition}
Soit $\K$ un corps qui n'est pas de caractéristique 2 et $P\in\polyK$. Alors
\begin{itemize}
\item $P$ est pair si et seulement si ses coefficients d'indices impairs sont
  nuls.
\item $P$ est impair si et seulement si ses coefficients d'indices pairs sont
  nuls.
\end{itemize}
\end{proposition}

% \begin{remarques}
% \remarque Cette proposition est fausse lorsque $\K$ est un corps où
%   $1_\K+1_\K=0_\K$
%   (on parle de corps de caractéristique 2) ce qui ne sera jamais le cas en
%   {\sc MPSI}, les corps considérés étant des sous-corps de $\C$.
% \end{remarques}

\subsection{Degré d'un polynôme}

\begin{definition}
Soit $P\in\polyK$. On définit le \emph{degré} de $P$ que l'on note $\deg P$ par
\begin{itemize}
\item Si $P=0$, on pose $\deg P\defeq -\infty$.
\item Sinon, il existe $n\in\N$ et $a_0,\ldots,a_n\in\K$ tels que
  \[P=a_0+a_1X+\cdots+a_n X^n \et a_n\neq 0.\]
  De plus $n$ et les $a_0,\ldots,a_n$ sont uniques; on pose alors $\deg P\defeq n$.
  Le coefficient $a_n$ est appelé \emph{coefficient dominant} de $P$.
\end{itemize}
\end{definition}

\begin{preuve}
L'unicité de $n$ et des $a_i$ provient directement du second axiome de la première définition du cours de $\polyK$, ce qu'on peut démontrer.
\end{preuve}

\begin{remarques}
\remarque Si $P\in\polyK$ est non nul, son coefficient dominant est parfois noté
  ${\rm cd}(P)$.
\remarque Un polynôme $P\in\polyK$ est de degré inférieur ou égal à $n\in\N$
  si et seulement si il existe $a_0,\ldots,a_n\in\K$ tels que
  \[P=\sum_{k=0}^n a_k X^k.\]
\remarque On dit qu'un polynôme $P$ est constant lorsqu'il existe $\lambda\in\K$ tel
  que $P=\lambda$, c'est-à-dire lorsque son degré est inférieur ou égal à 0.
\end{remarques}

\begin{proposition}
Soit $P,Q\in\polyK$ et $n\in\N$.
\begin{itemize}
\item Soit $\lambda,\mu\in\K$. Si $\deg P\leq n$ et $\deg Q\leq n$, alors
  \[\deg\p{\lambda P+\mu Q}\leq n.\]
\item Soit $\lambda\in\Ks$ et $\mu\in\K$. Si $\deg P=n$ et $\deg Q<n$, alors
  \[\deg\p{\lambda P+\mu Q}=n \quad\et\quad {\rm cd}(\lambda P+\mu Q)=\lambda{\rm cd}(P).\]
\end{itemize}
\end{proposition}

\begin{remarqueUnique}
\remarque Lorsque $P$ et $Q$ sont des polynômes de degré $n$, il est possible que
  $P+Q$ soit de degré strictement inférieur à $n$. Par exemple $P\defeq X+1$ et
  $Q\defeq -X$ sont de degré 1 mais $P+Q=1$ est de degré 0.
\end{remarqueUnique}

\begin{exoUnique}
\exo Soit $P\in\polyC$.
  Calculer le degré de $P\p{X+1}-P(X)$ en fonction de celui de $P$.
  \begin{sol}
  est égal à $\deg(P)-1$ si $\deg P\geq 1$ et à $-\infty$ sinon.  
  \end{sol}
\end{exoUnique}

\begin{definition}
Soit $n\in\N$. On note $\polyK[n]$ l'espace vectoriel des polynômes de degré inférieur ou égal à $n$.
\end{definition}

\begin{remarqueUnique}
\remarque Si $n\geq 1$, $\polyK[n]$ n'est pas stable par produit. En effet,
  $X^n\in\polyK[n]$ mais $X^{2n}=X^{n}\cdot X^n\not\in\polyK[n]$.
\end{remarqueUnique}

% \begin{proposition}
% Soit $n\in\N$. Alors $1,X,\ldots,X^n$ est une base de $\polyK[n]$ appelée base
% canonique. En particulier, $\polyK[n]$ est de dimension finie et~:
% \[\dim \polyK[n]=n+1\]
% \end{proposition}

% \begin{proposition}
% Soit $P_0,\ldots,P_n\in\polyK[n]$ une famille de polynômes tels que~:
% \[\forall k\in\intere{0}{n} \quad \deg P_k=k\]
% Alors $P_0,\ldots,P_n$ est une base de $\polyK[n]$.
% \end{proposition}

\begin{proposition}
Soit $P,Q\in\polyK$. Alors
\[\deg\p{PQ}=\deg P+\deg Q.\]
De plus, si $P$ et $Q$ sont non nuls, alors ${\rm cd}(PQ)={\rm cd}(P){\rm cd}(Q)$.
\end{proposition}

\begin{preuve}
L'égalité de la proposition tient compte de la convention $-\infty+(-\infty)=-\infty$ et $-\infty+n=-\infty$.
On évacue tout d'abord le cas où $P$ ou $Q$ est nul.
Ensuite, on peut poser $p=\deg P$, $q=\deg Q$ et écrire $$P=\sum_{n\in \N}a_nX_n \et  Q=\sum_{n\in \N}b_nX_n \et PQ=\sum_{n\in \N}c_nX_n.$$
Soit $n\geq p+q+1$, on a :
$$c_n=\sum_{k=0}^pa_k\underbrace{b_{n-k}}_{=0 (n-k\geq q+1)}+\sum_{k=p+1}^n\underbrace{a_k}_{=0}b_{n-k}$$
et de plus on montre que $c_{p+q}$ est non nul en séparant la somme entre les termes avant $p-1$ et après $p+1$ qui sont tous nuls et $a_pb_q\neq 0$.
\end{preuve}

\begin{remarques}
\remarque Si $P\in\polyK$ est non nul et si $n\in\N$, alors $\deg\p{P^n}=n\deg P$.
\remarque Si $P\in\polyK$ et $Q\in\polyK$ n'est pas constant,
  alors $\deg\p{P\circ Q}=\deg(P)\deg(Q)$.
\end{remarques}

\begin{proposition}
$\polyK$ est une algèbre intègre
\[\forall P,Q\in\polyK \qsep PQ=0 \quad\implique\quad \cro{P=0 \ou Q=0}.\]
\end{proposition}

\begin{proposition}
Les éléments inversibles de $\polyK$ sont les polynômes de degré 0, c'est-à-dire
les polynômes constants non nuls.
\end{proposition}

\begin{definition}
On dit qu'un polynôme non nul $U$ est unitaire lorsque son coefficient dominant est
égal à 1. Tout polynôme $P$ non nul s'écrit de manière unique sous la forme
$P=\lambda P_u$ où $\lambda\in\Ks$ et $P_u$ est unitaire. Lorsque $P=0$, on pose
par convention $P_u\defeq 0$.
\end{definition}


\begin{definition}[nom={Division euclidienne}]
  Soit $A,B\in\polyK$ avec $B\neq 0$. Alors, il existe un unique couple
  $\p{Q,R}\in\polyK^2$ tel que
  \[A=QB+R \et \deg R<\deg B.\]
  $Q$ est appelé \emph{quotient} de la division euclidienne de $A$ par $B$, $R$ son
  \emph{reste}.
  \end{definition}
  
  \begin{preuve}
  Pour l'existence, on fait une preuve effective fournissant un algorithme :
  
  Fixons $B\in \polyK$, $B\neq 0$. On définit $H_n : " \forall A \in \polyK[n], \exists (Q,R)\in \polyK^2 \text{ tq } A=BQ+R \et \deg R<\deg B".$.
  \begin{itemize}
  \item[$\bullet$] $H_0$ : Soit $A\in \polyK[0]=\K$. Si $\deg B=0$, $A=(AB^{-1})B+0$ convient. Sinon, $A=0B+A$ convient.
  \item[$\bullet$] Hérédité : On suppose $\deg A=n+1$, sinon, c'est bon d'après l'HR. Si $\deg A<\deg B$, $A=0B+A$ convient.
  \'Ecrivons alors $$A=\sum_{k=0}^{n+1}a_k X^k \et B=\sum_{k=0}^{d}b_k X^k$$ avec $a_{n+1}\neq 0$, $b_d\neq 0$ et $d\leq n+1$.
  
  Posons $A_1=A-\dfrac{a_{n+1}}{b_d}X^{n+1-d}B$. $A_1\in \polyK[n]$ donc d'après l'HR, $A_1=BQ_1+R$ et donc $$A=B\p{Q_1+\dfrac{a_{n+1}}{b_d}X^{n+1-d}}+R.$$
  \end{itemize}
  \end{preuve}
  
  \begin{remarques}
  \remarque Si $B$ est un polynôme annulateur non nul de $x$ et $A\in\polyK$.
    Alors $A(x)=R(x)$ où $R$ est le reste de la division euclidienne de $A$
    par $B$. En effet
    \[A(x)=Q(x)\underbrace{B(x)}_{=0}+R(x)\]
  \exo Il est parfois utile de calculer le reste de la division euclidienne
    de $A$ par $B$ sans calculer son quotient.
      Par exemple, si $A\defeq X^n$
      et $B\defeq \p{X-1}\p{X-2}$, le reste $R$ de la division euclidienne de $A$ par $B$
      est de degré inférieur ou égal à 1 donc il existe $a,b\in\R$ tels que
      $R=aX+b$. Comme $A=QB+R$, on en déduit que
      $A\p{1}=Q\p{1}B\p{1}+R\p{1}$. Comme $B\p{1}=0$, on a $A\p{1}=R\p{1}$. De
      même $A\p{2}=R\p{2}$. Donc
      \[\syslin{a&+b&=&1\hfill\cr
               2a&+b&=&2^n.\hfill}\]
      On en déduit que $a=2^n-1$ et $b=2-2^n$. Donc $R=\p{2^n-1}X+\p{2-2^n}$.
      Cette méthode fonctionne dès que le polynôme $B$, de degré $n$, admet $n$
      racines deux à deux distinctes.
  \end{remarques}
  
  \begin{exos}
  \exo Calculer $x^5+x^4-1$ où $x\defeq(1+\sqrt{5})/2$.
    \begin{sol}
    On a $X^5+X^4-1=(X^3+2X^2+3X+5)(X^2-X-1)+(8X+4)$. En particulier,
    si $x=(1+\sqrt{5})/2$
    \[\p{\frac{1+\sqrt{5}}{2}}^5+\p{\frac{1+\sqrt{5}}{2}}^4-1=x^5+x^4-1=8x+4=
      8+4\sqrt{5}\]  
    \end{sol}
  % \exo Montrer que le polynôme $P\defeq X^3+pX+q\in\polyR$ admet 3 racines réelles
  %   deux à deux distinctes si et seulement si $4p^3+27q^2<0$.
  %   \begin{sol}
  %   On a déjà forcément une racine réelle. Ensuite (faire un dessin), il nous faut regarder les annulations de la dérivée. Pour qu'il y ait trois racines réelles distinctes, il nous faut deux racines pour la dérivée $x_0$ et $x_1$ telles que $P(x_1)P(x_0)<0$.
    
  %   $P'=3X^2+p$ ne s'annule pas si $p>0$ (dans ce cas, la fonction est strictement croissante et n'admet donc pas trois racines distinctes). Si $p\leq 0$, alors $P'$ s'annule en $x_0$ et $x_1$ vérifiant de plus $x_0+x_1=0$ et $x_0x_1=\frac{p}{3}$. On cherche à calculer $P(x_1)P(x_0)$. On peut faire la DE de $P$ par $P'$ puis on trouve $P(x_1)P(x_0)=\frac{1}{27}\p{4p^3+27q^2}$.
  %   \end{sol}
      \exo On pose
      \[A\defeq
      \begin{pmatrix}
      4 & -1\\
      2 & 1
      \end{pmatrix}\]
      Vérifier que $A^2-5A+6I_2=0$ puis calculer $A^n$ pour tout $n\in\N$.
      \begin{sol}
      Après avoir vérifié l'égalité, on fait la DE de $X^n$ par $X^2-5X+6=(X-3)(X-2)$ où seul le reste nous intéresse.
      On trouve
      \[A^n=
      \begin{pmatrix}
      -2^n+2\cdot 3^n & 2^n-3^n\\
      -2\cdot 2^n+2\cdot 3^n & 2\cdot 2^n-3^n
      \end{pmatrix}\]
      \end{sol}
  \end{exos}

\begin{proposition}
Soit $\mathcal{I}$ un idéal de $\polyK$. Alors, il existe un unique polynôme unitaire
ou nul $P$ tel que $\mathcal{I}=P\polyK$.
\end{proposition}

\begin{remarqueUnique}
\remarque En particulier, dans $\polyK$, tout idéal est principal.
\end{remarqueUnique}

\subsection{Racines, fonctions polynôme}

\begin{definition}
Soit $P\in\polyK$. On appelle racine de $P$ tout élément $\alpha\in\K$ tel que
$P\p{\alpha}=0$.
\end{definition}

\begin{remarques}
\remarque Les polynômes de degré 1 admettent une unique racine.
\remarque D'après le théorème des valeurs intermédiaires, tout polynôme
  réel de degré impair admet (au moins) une racine réelle.
\remarque La notion de racine dépend du corps considéré. En effet, si on le
  considère comme élément de $\polyC$, les racines de \mbox{$(X^2-2)(X^2+1)$} sont
  $\sqrt{2},-\sqrt{2},\ii,-\ii$. Considéré comme élément de $\polyR$,
  ses racines sont $\sqrt{2},-\sqrt{2}$. Enfin il n'a aucune racine si
  on le considère comme un élément de $\polyQ$.
\remarque Si $\K$ est un sous-corps de $\KL$, $P\in\polyK$
  et $\alpha\in\KL$, on dit que $\alpha$ est une racine de $P$ sur $\KL$
  lorsque $P\p{\alpha}=0$.
\end{remarques}


\begin{proposition}
Soit $P\in\polyK$. Si $\alpha\in\K$ est une racine de $P$, alors il existe $Q\in\polyK$ tel
que
\[P=(X-\alpha)Q.\]
\end{proposition}

\begin{remarqueUnique}
\remarque Cette factorisation se calcule en pratique en effectuant la
  division euclidienne de $P$ par $X-\alpha$.
\end{remarqueUnique}

\begin{proposition}
Si $n\in\N$, tout polynôme de degré $n$ admet au plus $n$ racines.
\end{proposition}

\begin{preuve}
Se démontre par récurrence sur $n$ avec la proposition comme prédicat. Pour l'hérédité, ou bien il n'y a pas de racines auquel cas c'est tout bon, ou bien il en existe une $\alpha$ et on peut alors écrire $P=P(X)-P(\alpha)$ qui en écrivant avec les coefficients va faire apparaitre des $X^k-\alpha^k$ qu'on peut tous factoriser par $X-\alpha$. On a donc écrit $P=(X-\alpha)Q$. Reste à conclure.
\end{preuve}

\begin{remarques}
\remarque On en déduit qu'un polynôme de degré inférieur ou égal à $n$
  admettant $n+1$ racines deux à deux distinctes est nul. De même, si deux
  polynômes de degrés inférieurs ou égaux à $n$ prennent la même valeur
  en $n+1$ points deux à deux distincts, alors ils sont égaux.
\remarque Un polynôme admettant une infinité de racines est donc nul. De même,
  deux polynômes prenant la même valeur sur un ensemble infini sont
  égaux.
\end{remarques}

\begin{exos}
\exo Montrer que les polynômes de $\polyC$ tels que $P(X)=P(X+1)$ sont les
  polynômes constants.
  \begin{sol}
  On peut utiliser l'exercice d'avant en utilisant les degrés. Ou bien on montre que $P-P(0)$ admet une  infinité de racines.
  \end{sol}
\exo Montrer qu'il n'existe pas de polynôme $P\in\polyC$ tel que, pour tout
  $z\in\C$, $P(z)=\conj{z}$.
  \begin{sol}
  Tous les réels sont racines de $P(X)-X$.
  \end{sol}
% \exo On se donne $n+1$ éléments de $\K$ deux à deux distincts
%   $x_0,x_1,\ldots,x_n$ et $y_0,\ldots,y_n\in\K$. Montrer qu'il existe un unique
%   polynôme $P$ de degré inférieur ou égal à $n$ tel que
%   \[\forall k\in\intere{0}{n} \quad P\p{x_k}=y_k\]
%   On dit que $P$ est le polynôme interpolateur de Lagrange associé aux
%   familles $\p{x_k}$ et $\p{y_k}$.
%   \begin{sol}
%   On a l'unicité en en prenant deux et en montrant que le polynôme de la différence admet trop de racines. 
%   Pour l'existence, on peut poser $\forall k \in \llbracket 0,n \rrbracket$ $$L_k=\prod_{\underset{i\neq k}{i=0}}^n\p{\frac{X-x_i}{x_k-x_i}}.$$
%   On peut le faire sentir car on a besoin de $L_k(x_j)=\delta_{jk}$.
%   \end{sol}
\end{exos}

\begin{proposition}[nom={Polynôme interpolateur de \nom{Lagrange}}]
Soit $x_0,\ldots,x_n\in\K$, $n+1$ éléments deux à deux distincts et $y_0,\ldots,y_n\in\K$. Alors,
il existe un unique polynôme $P$ de degré inférieur ou égal à $n$ tel que
\[\forall i\in\intere{0}{n} \qsep P(x_i)=y_i.\]
On l'appelle \emph{polynôme interpolateur de \nom{Lagrange}} associé aux familles $\p{x_0,\ldots,x_{n}}$ et $\p{y_0,\ldots,y_{n}}$.
\end{proposition}

\begin{remarqueUnique}
\remarque  Pour tout $i\in\intere{0}{n}$, on note
$L_i$ le polynôme défini par
\[L_i\defeq \prod_{\substack{k=0\\k\neq i}}^{n} \frac{X-x_k}{x_i-x_k}.\]
Les polynômes $L_i$ sont appelés \emph{polynômes de Lagrange}.
Le polynôme interpolateur de \nom{Lagrange} associé aux familles $\p{x_0,\ldots,x_{n}}$ et $\p{y_0,\ldots,y_{n}}$ est donné par
\[P=\sum_{i=0}^{n} y_i L_i.\]
\end{remarqueUnique}

% \begin{proposition}
% Soit $x_0,\ldots,x_{n}\in\K$, $n+1$ éléments deux à deux distincts.
% \end{proposition}

\begin{preuve}
On a l'unicité en en prenant deux et en montrant que le polynôme de la différence admet trop de racines. 
  Pour l'existence, on peut poser $\forall k \in \llbracket 0,n \rrbracket$ $$L_k=\prod_{\underset{i\neq k}{i=0}}^n\p{\frac{X-x_i}{x_k-x_i}}.$$
  On peut le faire sentir car on a besoin de $L_k(x_j)=\delta_{jk}$.
  \end{preuve}

% \begin{remarqueUnique}
% \remarque Soit $x_0,\ldots,x_{n}\in\K$ deux à deux distincts, $y_0,\ldots,y_{n}\in\K$ et $P\in\polyK$. Alors
%   \[\forall i\in\intere{0}{n} \qsep P(x_k)=y_k\]
%   si et seulement si il existe un polynôme $Q\in\polyK$ tel que
%   \[P=\sum_{i=0}^{n} y_i L_i + Q\prod_{k=0}^{n} \p{X-x_k}.\]
% \end{remarqueUnique}


\begin{definition}
On dit qu'une application $f:\K\to\K$ est une fonction polynôme lorsqu'il existe
$P\in\polyK$ tel que
\[\forall x\in\K \qsep f(x)=P(x).\]
% L'ensemble des fonctions polynomiales est une sous-algèbre de
% $\mathcal{F}\p{\K,\K}$.
\end{definition}

\begin{proposition}
Si $\K$ est infini,
l'application de l'algèbre $\polyK$ dans l'algèbre $\mathcal{F}(\K,\K)$, qui au polynôme $P$ associe la fonction polynôme
$\tilde{P}$, est injective.  
\end{proposition}

\begin{preuve}
Si $P$ et $Q$ ont la même image, alors $P-Q$ admet une infinité (car $\K$ est infini) de racines donc $P=Q$.
\end{preuve} 

\begin{remarques}
% \remarque Lorsque $\K$ est un sous-corps de $\C$, l'algèbre des fonctions polynomiales hérite donc des propriétés de
%   $\polyK$. En particulier, elle est intègre, contrairement à l'algèbre
%   $\mathcal{F}\p{\K,\K}$.
\remarque Cette proposition permet, lorsque $\K$ est infini, d'identifier polynômes et
  fonctions polynôme. C'est pourquoi de rares énoncés se
  permettent de les confondre, identification
  que nous ne ferons que lorsque l'énoncé le demande explicitement.
\remarque Cette proposition est fausse lorsque le corps $\K$ est fini. En effet,
  si $\K=\ens{a_1,\ldots,a_n}$, le polynôme
  \[P\defeq\prod_{k=1}^n \p{X-a_k}\]
  est non nul car $\deg P=n$, mais la fonction polynôme associée est nulle.
\end{remarques}

\subsection{Polynôme dérivé}

\begin{definition}
Soit $P=a_0+a_1 X+\cdots+a_n X^n\in\polyK$. On définit le polynôme dérivé de
$P$ par
\begin{eqnarray*}
P' &\defeq& a_1 + 2a_2 X+\cdots+na_nX^{n-1}\\
   &=& \sum_{k=1}^n ka_k X^{k-1}.
\end{eqnarray*}
\end{definition}

\begin{remarqueUnique}
\remarque Dans le cas où $\K=\R$, la fonction polynôme associée à $P'$ est
  la dérivée (comme définie dans le cours d'analyse) de la fonction
  polynôme associée à $P$.
\end{remarqueUnique}

\begin{proposition}
Soit $P,Q\in\polyK$ et $\lambda,\mu\in\K$. Alors
\[\p{\lambda P+\mu Q}'=\lambda P'+\mu Q' \qquad
  \p{PQ}'=P'Q+PQ' \et \p{P\circ Q}'=Q'\p{P'\circ Q}.\]
\end{proposition}

\begin{remarqueUnique}
\remarque On peut utiliser le polynôme dérivé pour calculer la reste de la division
  euclidienne de $A$ par $B$ lorsque $B$ possède des racines multiples.
  Par exemple, si $A\defeq X^n$
  et $B\defeq \p{X-1}^2$, le reste $R$ de la division euclidienne de $A$ par $B$
  est de degré inférieur ou égal à 1 donc il existe $a,b\in\R$ tels que
  $R=aX+b$. Comme plus haut, $A\p{1}=R\p{1}$. En dérivant la relation
  $A=QB+R$, on obtient $A'=B'Q+BQ'+R'$. Puisque 1 est racine de $B$ et de $B'$, on
  en déduit que $A'\p{1}=R'\p{1}$. Donc
  \[\syslin{a&+b&=&1\hfill\cr
            a&  &=&n.\hfill}\]
  On en déduit que $a=n$ et $b=1-n$, donc $R=nX+\p{1-n}$.
\end{remarqueUnique}

\begin{definition}
Soit $P\in\polyK$. On définit par récurrence la dérivée $n$-ième de $P$ par
\begin{itemize}
\item $P^{(0)}\defeq P$
\item $\forall n\in\N \qsep P^{(n+1)}\defeq \cro{P^{(n)}}'$.
\end{itemize}
\end{definition}

\begin{remarqueUnique}
\remarque Soit $n\in\N$. Alors
  \[\forall k\in\N \qsep \p{X^n}^{(k)}=
    \begin{cases}
    \frac{n!}{\p{n-k}!}X^{n-k} &\text{si $k\leq n$}\\
    0 & \text{sinon.}
    \end{cases}\]
  En particulier, si $P=a_0+a_1X+\cdots+a_nX^n\in\polyK$, alors
  \[\forall k\in\N \qsep P^{(k)}\p{0}=k! a_k.\]
\end{remarqueUnique}

\begin{proposition}
Soit $P,Q\in\polyK$ et $n\in\N$
\begin{itemize}
\item Soit $\lambda,\mu\in\K$. Alors
  \[\p{\lambda P+\mu Q}^{(n)}=\lambda P^{(n)}+\mu Q^{(n)}.\]
\item On a
  \[\p{PQ}^{(n)}=\sum_{k=0}^n \binom{n}{k}P^{(n-k)}Q^{(k)}.\]
  Cette formule est appelée formule de \nom{Leibniz}.
\end{itemize}
\end{proposition}

\begin{exoUnique}
\exo Calculer $(X^2P)^{(n)}$ en fonction des dérivées successives de $P$.
\end{exoUnique}

Dans la suite du cours, on suppose que $\K$ est un corps de caractéristique nulle.

\begin{proposition}
Soit $P\in\polyK$. Alors
\[\deg P'=
  \begin{cases}
  \deg(P)-1 & \text{si $\deg P\geq 1$,}\\
  -\infty & \text{sinon.}
  \end{cases}\]
\end{proposition}

\begin{preuve}
Notons qu'on se place ici dans un corps de caractéristique nulle. Par exemple, dans $\Z/p\Z$, $P=X^p$ est de degré $p$ mais son polynôme dérivé est nul.
\end{preuve}

\begin{remarques}
\remarque $P'=0$ si et seulement si $P$ est constant.
\remarque Pour tout $P\in\polyK$, $\deg P'\leq\deg(P)-1$.
\remarque Soit $P\in\polyK$ et $n\in\N$. Alors 
\[\deg P^{(n)}=
  \begin{cases}
  \deg(P)-n & \text{si $\deg P\geq n$,}\\
  -\infty & \text{sinon.}
  \end{cases}\]
  En particulier, $\deg P^{(n)}\leq\deg(P)-n$.
\end{remarques}

\begin{proposition}[nom={\nom{Formule de Taylor}}]
Soit $P\in\polyK$ un polynôme de degré inférieur ou égal à $n$ et $\alpha\in\K$. Alors
\[P=\sum_{k=0}^n \frac{P^{(k)}\p{\alpha}}{k!}\p{X-\alpha}^k.\]  
\end{proposition}


\begin{preuve}
\begin{eqnarray*}
P&=&\sum_{k=0}^n a_kX^k\\
&=& \sum_{k=0}^n a_k\p{X-\alpha+\alpha}^k\\
&=&\sum_{k=0}^n a_k\sum_{i=0}^k\binom{k}{i}(X-\alpha)^i\alpha^{k-i}\\
&=&\sum_{i=0}^n\p{\sum_{k=i}^n \binom{k}{i}a_k\alpha^{k-i}}(X-\alpha)^i\\
&=&\sum_{i=0}^n\frac{1}{i!}\underbrace{\p{\sum_{k=i}^n \frac{k!}{(k-i)!}a_k\alpha^{k-i}}}_{=P^{(i)}(\alpha)}(X-\alpha)^i
\end{eqnarray*}
\end{preuve}
% \begin{remarques}
% \remarque Si $\K$ est un corps pour lequel il existe $n\in\Ns$ tel que
%   \[\underbrace{1_\K+1_\K+\cdots+1_\K}_{\text{$n$ fois}}=0_\K\]
%   (on parle de corps de caractéristique non nulle), les deux 
%   propositions précédentes deviennent fausses. Cependant, ce ne sera pas le cas
%   en {\sc MPSI}, les corps considérés étant des sous-corps de $\C$.  
% \end{remarques}

%END_BOOK

\end{document}
\section{Exercices}
\setcounter{numeroexercice}{1}
\documentclass{magnolia}

\magtex{tex_driver={pdftex}}
\magfiche{document_nom={Corps, polynomes},
          auteur_nom={Emmanuel Roblet, François Fayard},
          auteur_mail={fayard.prof@gmail.com}}
\magexos{exos_matiere={maths},
         exos_niveau={mpsi},
         exos_chapitre_numero={12},
         exos_theme={Corps, polynomes}}
\magmisenpage{}
\maglieudiff{}
\magprocess


\begin{document}
%BEGIN_BOOK

\magsection{Anneau, corps}

\magsubsection{Anneau}

\exercice{nom={Anneau de Boole}}
Soit $E$ un ensemble. Pour tout $A,B\in\mathcal{P}(E)$, on définit $A\Delta B$ par
\[A\Delta B\defeq(A\setminus B)\cup(B\setminus A).\]
Pour tout $A\in\mathcal{P}(E)$, on définit la fonction caractéristique de $A$ comme
l'application $\mathds{1}_A:A\to\mathbb{F}_2$ définie par
\[\forall x\in E \qsep \mathds{1}_A(x)\defeq
\begin{cases}
\bar{1} & \text{si $x\in A$}\\
\bar{0} & \text{sinon.}
\end{cases}\]
\begin{questions}
\question Montrer que pour tout $A,B\in\mathcal{P}(E)$,
  $\mathds{1}_{A\Delta B}=\mathds{1}_A + \mathds{1}_B$ et
  $\mathds{1}_{A\cap B}=\mathds{1}_A \mathds{1}_B$.
\question En déduire que $(\mathcal{P}(E),\Delta,\cap)$ est un anneau commutatif.
\question Montrer que cet anneau est intègre si et seulement si $E$ possède un unique
  élément.
\end{questions}

\exercice{nom={Fonction définie sur un anneau}}
Soit $(A,+,\times)$ un anneau et $f$ une application de $A$ dans $\RP$ telle
que
\begin{itemize}
\item $\forall x\in A \qsep f(x)=0 \ssi x=0$.
\item $\forall x,y\in A \qsep f(xy)=f(x)f(y)$.
\item $\forall x,y\in A \qsep f(x+y)\leq \max\p{f(x),f(y)}$.
\end{itemize}
Montrer que $\enstq{x\in A}{f(x)\leq 1}$ est un sous-anneau de~$A$.

\magsubsection{Corps}

\exercice{nom={Exemple de corps}}
On définit sur $\R$ deux lois $\oplus$ et $\otimes$ par
\begin{eqnarray*}
\forall x,y\in\R\qsep x\oplus y&\defeq&x+y-1\>,\\
\forall x,y\in\R\qsep x\otimes y&\defeq&x+y-xy\>.
\end{eqnarray*}
Montrer que $(\R,\oplus,\otimes)$ est un corps.

\exercice{nom={Morphisme de corps}}
Montrer qu'un morphisme de corps est toujours injectif. Trouver un morphisme
d'anneaux non injectif.

\begin{sol}
Soient $K_1, K_2$ des corps et $f$ un morphisme de $K_1$ dans $K_2$. Comme $f$ est un morphisme de corps, $f(1)=1$. Supposons $f(a)=f(b)$. On a alors $f(a-b)=0$. Si jamais on avait $a \neq b$, on aurait $a-b$ inversible (notons $c$ son inverse) et par suite $1=f(1)=f(c(a-b))=f(c)f(a-b)=f(c)\times 0=0$, ce qui n'est pas possible (dans un corps, $0$ et $1$ sont distincts). On a donc nécessairement $a=b$. Cela signifie bien que $f$ est injective.
\end{sol}

\exercice{nom={Extension quadratique}}
Soit $\alpha\in\N$ tel que $\sqrt{\alpha}\not\in\Q$. On pose 
\[\Q(\sqrt{\alpha})\defeq\ensim{a+b\sqrt{\alpha}}{(a,b)\in\Q^2}\]
\begin{questions}
\question Soit $x\in\Q\p{\sqrt{\alpha}}$. Montrer qu'il existe un unique
  couple $\p{a,b}\in\Q^2$ tel que $x=a+b\sqrt{\alpha}$.
\question Montrer que $\Q(\sqrt{\alpha})$ est un sous-corps de $(\R,+,\times)$.
% \enonce Dans la suite, on suppose que $\alpha=2$ et on note $\K$ le corps
%   $\Q(\sqrt{2}) $.
\question Pour $x=a+b\sqrt{\alpha}\in\Q(\sqrt{\alpha})$, on pose
  $\overline{x}=a-b\sqrt{\alpha}$; on l'appelle le conjugué de $x$. Montrer que
  l'application  $x\mapsto \overline{x}$ est bien définie et est un
  automorphisme du corps $\Q(\sqrt{\alpha})$.
\question Montrer que l'automorphisme construit à la question précédente est
  le seul automorphisme non trivial de $\Q(\sqrt{\alpha})$.
% \question Pour $x\in\K$, on pose $N(x)=x\overline{x}$. Pour tous  $x$ et $y$
%   dans~$\K$, montrer que 
%   \begin{itemize}
%   \item $N(x)=0$ si et seulement si $x=0$ ;
%   \item $N(xy)=N(x)N(y)$.
%   \end{itemize}
%   Que dire de l'application $N$ de $(\Ks,\times)$ dans $(\Qs,\times)$~?
%   Calculer $N(1/x)$ lorsque $x$ est un élément non nul de~$\K$.
% \question Soit $A=\enstq{a+b\sqrt{2}}{(a,b)\in\Z^2}$. Montrer que $A$ est un
%   sous-anneau intègre du corps $\K$. Montrer que $x\in A$ admet un inverse dans
%   $A$ si et seulement si $|N(x)|=1$ ; que vaut alors l'inverse de $x$~?
%    Donner un exemple non trivial d'élément de~$A$ inversible dans $A$.
\end{questions}

% \exercice{nom={Théorème de Wilson}}
% Soit $p$ un nombre premier.
% \begin{questions}
% \question Montrer que l'application $\phi$ de $(\Z/p\Z)^*$ dans $(\Z/p\Z)^*$
%   qui à $x$ associe $1/x$ est une bijection. Quels sont les points fixes de
%   cette bijection~?
% \question En déduire que $(p-1)!\equiv -1\ [p]$.
% \end{questions}

% \begin{sol}
% \begin{questions}
% \question C'est une bijection d'inverse elle-même. Elle a pour point fixe $1$ et $-1$.
% \question On a un nombre pair d'éléments dans ce produit. En regroupant chaque élément avec son inverse, il reste $1$ et $-1$ dont le produit fait bien $-1$.
% \end{questions}
% \end{sol}

\magsection{Espace vectoriel, algèbre}
\magsubsection{Espace vectoriel}
\magsubsection{Algèbre}


\magsection{L'algèbre $\polyK$}

\magsubsection{Définition}

\exercice{nom={Produit}}
On note $\mathcal{A}$ l'ensemble des polynômes $P\in\polyR$ tels qu'il existe
$a_0,\ldots,a_n\in\RP$ tels que
\[P=\sum_{k=0}^n (-1)^k a_k X^k.\]
Montrer que $\mathcal{A}$ est stable par produit.


\magsubsection{Substitution}
\magsubsection{Degré d'un polynôme}

\exercice{nom={Équations sur $\polyR$}}
Déterminer l'ensemble des $P\in\polyR$ tels que
\[P(X^2)=(X^2+1)P(X).\]
On commencera par effectuer une analyse et un cherchera des informations sur
le degré de $P$.

% On trouve P = lambda(X^2 - 1)

\exercice{nom={Division euclidienne}}
Soit $n\in\N$. Calculer le reste de la division euclidienne
\begin{questions}
\question de $X^n(X+1)^2$ par $(X-1)(X-2)$.
\question de $X^n$ par $(X-1)^2(X+1)$.
% \question de $X^{2n}$ par $(X^2+1)^n$.
\question de $(X+1)^{2n+1} - X^{2n+1}$ par $X^2+X+1$.
\end{questions}

\magsubsection{Racines, fonctions polynôme}

\exercice{nom={Exercice}}
\begin{questions}
\question Déterminer tous les polynômes $P\in\polyR$ tels que
  \[\forall n\in\N\qsep P(n)=n^2.\]
\question Déterminer tous les polynômes $P\in\polyR$ tels que
  \[\forall n\in\N\qsep P(n)=n^2+(-1)^n.\]
\end{questions}


\exercice{nom={Polynômes de \nom{Tchebychev}}}
On définit la suite de polynômes $\p{T_n}$ par
\[T_0\defeq 1, \quad T_1\defeq X \et \forall n\in\N \qsep T_{n+2}\defeq 2XT_{n+1}-T_n.\]
\begin{questions}
\question Calculer les polynômes $T_n$ pour $n\in\intere{0}{5}$.
\question Calculer le degré de $T_n$ et son coefficient dominant.
\question Montrer que pour tout $n\in\N$, $T_n$ est l'unique polynôme
  tel que
  \[\forall x\in\R \qsep T_n\p{\cos x}=\cos\p{nx}.\]
\question En déduire les racines de $T_n$.
\question En dérivant deux fois la relation obtenue dans la question 3, montrer que
  \[\forall n\in\N\qsep (X^2-1)T_n''+XT_n'-n^2T_n=0.\]
\question Déterminer une expression explicite de $T_n$ en exploitant la formule de
  Moivre.
\end{questions}
\begin{sol}
$\quad$
\begin{questions}
\question On trouve $P_0=1$, $P_1=X$, $P_2=2X^2-1$, $P_3=4X^3-3X$,
  $P_4=8X^4-8X^2+1$, $P_5=16X^5-20X^3+5X$.
\question Faire une récurrence double sur $n$.
\question Ce sont les
  \[\cos\p{\frac{\pi}{2n}\p{1+2k}}\]
  pour $k\in\intere{0}{n-1}$. Ce sont des racines et il faut démontrer par
  récurrence sur $n$ que le degré de $P_n$ est $n$.
\end{questions}
\end{sol}

\exercice{nom={Exercice}}
On note $\mathcal{A}$ l'ensemble des fonctions $f\in\mathcal{F}(\R,\R)$ telles qu'il
existe $\lambda_0,\ldots,\lambda_n\in\R$ tels que
\[\forall x\in\R\qsep f(x)=\sum_{k=0}^n \lambda_k \e^{kx}.\]
\begin{questions}
\question Montrer que $\mathcal{A}$ est une sous-algèbre de $\mathcal{F}(\R,\R)$.
\question Montrer que $\mathcal{A}$ est intègre.
\question Déterminer $U_{\mathcal{A}}$. 
\end{questions}


\exercice{nom={Polynômes de Lagrange}}
Soit $n\geq 2$, $x_1,\ldots,x_n\in\K$ deux à deux distincts et $L_1,\ldots,L_n$ les
polynôme de Lagrange associés. Simplifier les sommes
\[\sum_{i=1}^n L_i \quad\text{et}\quad \sum_{i=1}^n x_i L_i.\]

\exercice{nom={Polynômes de Lagrange}}
On note $L_0,\ldots,L_n$ les polynômes de Lagrange de $0,\ldots,n$.
\begin{questions}
\question Pour tout $k\in\intere{0}{n}$, exprimer le coefficient dominant de $L_k$ au
  moyen de factorielles.
\question Exprimer de deux manières différentes l'unique polynôme $P\in\polyR$ de degré
  inférieur ou égal à $n$ tel que
  \[\forall k\in\intere{0}{n}\qsep P(k)=k^n.\]
\question En déduire une simplification de
  \[\sum_{k=0}^n \binom{n}{k}(-1)^{n-k} k^n.\]
\end{questions}

% Cela donne n!

\exercice{nom={Exercice}}
Soit $P\in\polyC$.
\begin{questions}
\question Donner une condition nécessaire et suffisante sur les coefficients de $P$
  pour que \[\forall x\in\R\qsep P(x)\in\R.\]
\question Donner une condition nécessaire et suffisante sur les coefficients de $P$
  pour que \[\forall x\in\Q\qsep P(x)\in\Q.\]
  \emph{On pourra utiliser les polynômes de Lagrange.}
\end{questions}


\magsubsection{Polynôme dérivé}

\exercice{nom={Exercice}}
Soit $n\in\N$ et $k\in\intere{0}{n}$. On pose
\[S\defeq\sum_{i=k}^n \binom{i}{k} \quad\text{et}\quad
  P\defeq\sum_{i=k}^n (X+1)^i.\]
\begin{questions}
\question Exprimer $S$ en fonction de $P^{(k)}(0)$.
\question Simplifier $(X+1)P - P$, puis dériver $k+1$ fois la
  relation obtenue.
\question En déduire une expression simple de $S$.
\end{questions}

\exercice{nom={Équations sur $\polyR$}}
\begin{questions}
\question Déterminer les polynômes $P\in\polyR$ tels que
  \[P\p{2X}=P'(X)P''(X).\]
\question Déterminer les polynômes $P\in\polyR$ tels que
  \[X\p{X+1}P''+\p{X+2}P'-P=0.\]
\end{questions}
\begin{sol}
$\quad$
\begin{questions}
\question Regarder le degré du polynôme. On trouve qu'il est inférieur ou égal
  à 3. Finalement, les solutions sont $0$ et $\frac{4}{9}X^3$.
\question Regarder le coefficient dominant. Cela implique que le degré du
  polynôme est inférieur ou égal à 1. On trouve finalement les solutions
  suivantes~: les $\lambda(X+2)$ avec $\lambda\in\R$.
\end{questions}
\end{sol}



\exercice{nom={Coefficients binomiaux}}
On donne un entier $n\geq 1$.
\begin{questions}
\question Pour $a,b\in\R$, calculer la dérivée $n$-ième du polynôme
  $P=\p{X-a}^n\p{X-b}^n$.
\question En déduire une expression simplifiée de la somme
  \[S=\sum_{k=0}^n \binom{n}{k}^2\]
\end{questions}


\exercice{nom={Primitive}}
Soit $P\in\polyC$ un polynôme de degré inférieur ou égal à $n$. Montrer que
\[Q\defeq \sum_{k=0}^n \frac{(-1)^k P^{(k)}(X)}{(k+1)!}X^{k+1}\]
est l'unique primitive de $P$ qui s'annule en 0.

% \exercice{nom={Calcul de degré}} %%Tout-en-un (2ème édition), p.749
% Soit $P$ un polynôme de $\polyR$. Déterminer le degré du polynôme
% $$P(X+1)-P(X)$$
% en fonction du degré de $P$.













% \exercice{nom={Une base de $\polyK[n]$}}
% Soit $n\in\N$. On définit pour tout $k\in\intere{0}{n}$~:
% \[P_k=X^k\p{1-X}^{n-k}\]
% Montrer que $P_0,\ldots,P_n$ est une base de $\polyK[n]$.



% \exercice{nom={Polynômes de Hilbert}}
% Soit $n\in\N$. Pour tout $k\in\intere{0}{n}$, on définit le polynôme $P_k$ par~:
% \[P_k=\frac{1}{k!}\prod_{l=0}^{k-1} \p{X-l}\]
% \begin{questions}
% \question Montrer que la famille $P_0,\ldots,P_n$ est une base de $\polyR[n]$.
% \question On définit l'application $\Delta$ par~:
%   \[\dspappli{\Delta}{\polyR[n]}{\polyR[n]}{P}{P\p{X+1}-P(X)}\]
%   \begin{questions}
%   \question Montrer que $\Delta$ est linéaire.
%   \question Pour tout $k\in\intere{0}{n}$, calculer $\Delta\p{P_k}$.
%   \question En déduire $\ker \Delta$ et $\im \Delta$.
%   \end{questions}
% \question
%   \begin{questions}
%   \question Montrer que~:
%     \[\forall P\in\polyR[n] \quad P=\sum_{k=0}^n \p{\Delta^k P}\p{0}P_k\]
%   \question Montrer que~:
%     \[\forall k\in\intere{0}{n} \quad \forall l\in\Z \quad P_k(l)\in\Z\]
%   \question En déduire que si $P\in\polyR[n]$~:
%     \[\cro{\forall l\in\Z \quad P(l)\in\Z} \quad \ssi\]
%     \[\cro{\text{Les coordonnées de $P$ dans la base $P_0,\ldots,P_n$ sont
%       entières.}}\]
%   \end{questions}
% \end{questions}







% \exercice{nom={Bases de $\polyR[n]$}}
% \begin{questions}
% \question Soit $n\in\N$ et $P_0,\ldots,P_n$ $n+1$ polynômes tels que :
%   $$\forall k\in\intere{0}{n} \quad \deg P_k=k$$
%   Montrer que la famille $P_0,\ldots,P_n$ est une base de $\polyR[n]$.
% \question Soit $\alpha\in\R$. Montrer que la famille
%   $\p{(X-\alpha)^n}_{n\in\N}$ est une base de $\polyR$.
% \end{questions}


% \exercice{nom={Inversion d'une matrice}}
% Soit $n\in\N$. On considère l'application $\phi$ qui au polynôme $P$
% associe le polynôme $P(X+1)$.
% \begin{questions}
% \question Montrer que $\phi$ est un isomorphisme de $\polyC[n]$ et déterminer
%   $\phi^{-1}$.
% \question Expliciter la matrice de $\phi$ dans la base canonique de
%   $\polyC[n]$.
% \question Montrer que cette matrice est inversible et calculer son inverse.
% \end{questions}


%END_BOOK


\end{document}



\chapter{Dénombrement}
\setcounter{numeroexercicecours}{1}
\documentclass{magnolia}


\magtex{tex_driver={pdftex},
        tex_packages={xypic}}
\magfiche{document_nom={Cours sur le dénombrement},
          auteur_nom={François Fayard},
          auteur_mail={fayard.prof@gmail.com}}
\magcours{cours_matiere={maths},
          cours_niveau={mpsi},
          cours_chapitre_numero={10},
          cours_chapitre={Dénombrement}}
          \magmisenpage{misenpage_presentation={tikzvelvia},
          misenpage_format={a4},
          misenpage_nbcolonnes={1},
          misenpage_preuve={non},
          misenpage_sol={non}}
\maglieudiff{}
\magprocess


\begin{document}

%BEGIN_BOOK
\hidemesometimes{
\hfill\includegraphics[width=0.45\textwidth]{../../Commun/Images/maths-cours-gaston.png}}

\magtoc

\section{Cardinal}

\subsection{Équipotence}

\begin{definition}
Soit $A$ et $B$ deux ensembles. On dit que $A$ est \emph{équipotent} à $B$
lorsqu'il existe une bijection de $A$ dans $B$.
\end{definition}

\begin{proposition}
La relation \og est équipotent à \fg est une relation d'équivalence.
\end{proposition}

\begin{remarques}
\remarque Une fois que nous aurons défini le cardinal d'un ensemble fini, nous verrons que deux ensembles
  finis sont équipotents si et seulement si ils ont le même nombre d'éléments.
\remarque Il est possible qu'un ensemble soit équipotent à l'une de ses parties strictes; ces ensembles sont infinis.
  Par exemple l'application $f$ de $\N$ dans $\Ns$ qui à $n$ associe $n+1$ est une bijection, ce qui montre que $\N$
  est équipotent à $\Ns$. Pourtant $\Ns$ est une partie stricte de $\N$. De même, l'application $f$
  de $\interf{0}{1}$ dans $\interf{0}{2}$ qui à $x$ associe $2x$ est une bijection. Pourtant
  $\interf{0}{1}$ est une partie stricte de $\interf{0}{2}$.
\remarque Il existe des ensembles infinis qui ne sont pas équipotents. Par exemple, on peut montrer que,
  quel que soit l'ensemble $X$, les ensembles $X$ et $\mathcal{P}(X)$ ne sont pas équipotents. En particulier
  $\N$ et $\mathcal{P}(\N)$ ne sont pas équipotents. Il existe donc des ensembles infinis qui ont
  \og plus d'éléments \fg que d'autres.
\remarque On dit qu'un ensemble est \emph{dénombrable} lorsqu'il est équipotent à $\N$. On peut montrer
  que $\Z$, $\N^n$ (pour $n\in\Ns$) et $\Q$ sont dénombrables. On peut montrer cependant que $\R$ n'est
  pas dénombrable.  On dit qu'un ensemble est \emph{au plus dénombrable} lorsqu'il est fini ou dénombrable. 
\end{remarques}

\subsection{Ensemble fini, cardinal}

On rappelle que pour tout $n\in\N$, on a $\intere{1}{n}\defeq\enstq{k\in\N}{1\leq k\leq n}$.
En particulier $\intere{1}{0}=\emptyset$, $\intere{1}{1}=\ens{1}$,
$\intere{1}{2}=\ens{1, 2}$, $\intere{1}{3}=\ens{1, 2, 3}$, etc. 
\vspace{2ex}
\begin{definition}
On dit qu'un ensemble $A$ est \emph{fini} lorsqu'il existe $n\in\N$ tel que $A$ est
équipotent à $\intere{1}{n}$. On dit qu'il est \emph{infini} dans le cas contraire.
\end{definition}

\begin{definition}
Soit $A$ un ensemble fini. Alors il existe un unique $n\in\N$ tel que $A$ est
équipotent à $\intere{1}{n}$. On l'appelle \emph{cardinal} de $A$ et on le note
$\card(A)$ ou $|A|$.
\end{definition}

\begin{preuve}
On commence par l'admettre sur le dessin. Ensuite~:
\begin{itemize}
\item Existence : par définition
\item Unicité : On dit que A est de cardinal n lorsqu'il existe une bijection
  entre [1,n] et A.
  On va montrer par récurrence sur n que [1,n] ne peut être de
  cardinal n et de cardial p avec n<p.
  \begin{itemize}
  \item $H_0$ : A ne peut être vide et non vide
  \item $H_n \implique H_(n+1)$ : Supposons que A soit de cardinal n+1 et de
     cardinal p avec
     n+1<p. Si x est un élément de A $A\setminus\ens{x}$ est de cardinal n et de
     cardinal p-1 avec n<p-1. C'est absurde
  \end{itemize}
\end{itemize}
\end{preuve}

\begin{remarques}
  \remarque Un ensemble $A$ est fini de cardinal $n\in\N$ si et seulement si il existe une
    bijection de $\interefo{0}{n}$ dans $A$.
  \remarque Soit $a,b\in\Z$ tels que $a\leq b + 1$. L'ensemble $\intere{a}{b}$ est fini de
  cardinal $b-a+1$.
\end{remarques}

\begin{proposition}
Soit $A$ un ensemble fini et $B$ un ensemble. Alors, $A$ et $B$ sont équipotents si et
seulement si $B$ est fini et $\card(A)=\card(B)$.  
\end{proposition}


\begin{exos}
\exo Dénombrer les couples $(a,b) \in \N^2$ tels que $3a + b = 833$.
\exo On a utilisé 6921 chiffres (les caractères d'imprimerie) pour numéroter les
  pages d'un dictionnaire. Combien de pages ce dictionnaire contient-il~?
  Chaque page est numérotée une seule fois, la première portant le numéro 1.
  \begin{sol}
  La réponse est 2007.
  \end{sol}
\end{exos}

\begin{definition}
Soit $A$ une partie de $\N$.
\begin{itemize}
\item Si $A$ est fini, il est l'image d'une unique application
  strictement croissante de $\intere{1}{\card(A)}$ dans $\N$.
\item Sinon, $A$ est infini et il est l'image d'une unique application strictement
  croissante de $\N$ dans $\N$.
\end{itemize}
Une telle application est appelée une \emph{énumération} de $A$.
\end{definition}

\begin{preuve}
Utiliser l'algorithme de construction
\end{preuve}

% \begin{remarqueUnique}
% \remarque Si $\phi$ est une application strictement croissante de $\N$ dans
%   $\N$, alors, pour tout $k\in\N$, $\phi(k) \geq k$.
% \end{remarqueUnique}

\begin{proposition}
Une partie de $\N$ est finie si et seulement si elle est majorée.
\end{proposition}

\begin{preuve}
Si elle est finie, elle est majorée.
Réciproquement, si elle est majorée, elle est incluse dans $[1,n]$ donc il ne
peut y avoir d'application strictement croissante de $N$ dans $A$.  
\end{preuve}

% \begin{proposition}
% Toute partie non vide majorée de $\N$ admet un plus grand élément.
% \end{proposition}

% \begin{exoUnique}
% \exo Soit $P$ un prédicat sur $\N$ tel que
%   \begin{itemize}
%   \item $P(n)$ est vrai pour une infinité de $n$.
%   \item Pour tout $n\in\Ns$, si $P(n)$ est vrai, alors $P(n-1)$ est vrai.
%   \end{itemize}
%   Montrer que $P(n)$ est vrai pour tout $n\in\N$.
% \end{exoUnique}

\begin{proposition}
Soit $E$ un ensemble fini et $A$ une partie de $E$. Alors
\begin{itemize}
\item $A$ est un ensemble fini et $\card(A)\leq\card(E)$.
\item $A=E$ si et seulement si $\card(A)=\card(E)$.
\end{itemize}
\end{proposition}

\begin{preuve}
Admis (voir Tissier, Acx, Desnoux page 33) :
1) On sa ramène à la démonstration du fait que toute partie de [1,n] est
   finie de cardinal inférieur à n. C'est est évident car elle est majorée.
    et on utilise l'application strictement croissante.
\end{preuve}


\begin{proposition}
Soit $E$ et $F$ deux ensembles. Alors
\begin{itemize}
\item Si $F$ est fini, il existe une injection de $E$ dans $F$ si et seulement
  si $E$ est fini et $\card (E)\leq\card (F)$.
\item Si $E$ est fini et $F$ est non vide, il existe une surjection de $E$ dans $F$ si et seulement
  si $F$ est fini et $\card (F)\leq\card (E)$.
\item Si l'un des ensembles est fini, il existe une bijection de $E$ dans $F$ si
  et seulement si l'autre est fini et $\card(E)=\card(F)$.
\end{itemize}
\end{proposition}

\begin{proposition}[nom={Principe des tiroirs}]
Soit $E$ et $F$ deux ensembles finis tels que $\card(F)<\card(E)$ et $f$ une application de $E$ dans $F$. Alors, il existe
$x_1,x_2\in E$ tels que $x_1\neq x_2$ et $f(x_1)=f(x_2)$.
\end{proposition}

% \begin{remarqueUnique}
% % \remarque Si $A$ et $B$ sont finis, il existe une injection de $A$ dans $B$ si
% %   et seulement si $\card A=\card B$.
% \remarque Soit $E$ et $F$ deux ensembles finis et $f\in\mathcal{F}(E,F)$. Si $\card(F)< \card(E)$, alors il existe deux éléments distincts de
%   $E$ qui ont même image par $f$. Cet énoncé est connu sous le nom de \og
%   principe des tiroirs \fg ou \og pigeonhole principle \fg.
% \end{remarqueUnique}

\begin{exos}
\exo Soit $n \geq 2$. En supposant que la relation \og est ami avec \fg est
  symétrique, montrer que dans une assemblée de $n$ personnes, il y en a au moins deux qui
  ont le même nombre d'amis.
% \exo De combien de personnes a-t-on besoin pour pouvoir affirmer avec certitude que $2$ sont nées le même jour de l'année ? De même avec $3$.
% \exo Montrer qu'il existe des entiers $a,b,c$ non tous nuls de valeur absolue strictement inférieure à $10^{6}$ tels que $|a+b \sqrt{2}+ c \sqrt{3}| < 10^{-11}$.
\exo Soit $x_1,\ldots,x_{n+1}\in\interf{0}{1}$. Montrer qu'il existe
  $i,j\in\intere{1}{n+1}$ tels que $i\neq j$ et
  \[\abs{x_i-x_j}\leq\frac{1}{n}.\]
% A mettre dans le cours sur les réels
%

\end{exos}

\begin{proposition}
Soit $E$ un ensemble fini, $F$ un ensemble et $f$ une application de $E$ dans
$F$. Alors
\begin{itemize}
\item $f(E)$ est un ensemble fini et $\card\p{f(E)} \leq \card(E)$
\item $\card\p{f(E)}=\card(E)$ si et seulement si $f$ est injective.
\end{itemize}
Si de plus $F$ est un ensemble fini.
\begin{itemize}
\item $\card\p{f(E)}\leq\card(F)$
\item $\card\p{f(E)}=\card(F)$ si et seulement si $f$ est surjective.
\end{itemize}
\end{proposition}





\begin{proposition}
Soit $E$ et $F$ deux ensembles finis et $f$ une application de $E$ dans $F$.
\begin{itemize}
\item Si $f$ est injective et $\card(E)=\card(F)$, alors $f$ est bijective.
\item Si $f$ est surjective et $\card(E)=\card(F)$, alors $f$ est bijective.
\end{itemize}
Autrement dit, si $\card(E)=\card(F)$, alors
\[\text{$f$ est injective} \quad\ssi\quad \text{$f$ est bijective} \quad\ssi\quad
  \text{$f$ est surjective}.\]
\end{proposition}

\section{Dénombrement}

\subsection{Dénombrement élémentaire}

\begin{proposition}
Soit $E$ un ensemble fini et $A$, $B$ deux parties disjointes de $E$, c'est-à-dire telles
que $A\cap B=\emptyset$. Alors
\[\card\p{A\cup B}=\card (A)+\card(B).\]
\end{proposition}

\begin{remarques}
\remarque Le \og ou exclusif \fg se traduit donc par un $ + $ en dénombrement.
\remarque Si $A$ et $B$ sont deux parties disjointes, leur réunion est parfois notée $A\sqcup B$.
\end{remarques}

\begin{proposition}
Soit $E$ un ensemble fini.
\begin{itemize}
\item Si $A$ est une partie de $E$
  \[\card\p{\bar{A}}=\card(E)-\card(A).\]
\item Si $A$ et $B$ sont deux parties de $E$
  \[\card\p{A\cup B}=\card(A)+\card(B)-\card\p{A\cap B}.\]
\item Si $(A_1,\ldots,A_n)$ est une partition de $E$, alors
  \[\card(E)=\card\p{A_1}+\cdots+\card\p{A_n}.\]
\end{itemize}
\end{proposition}

%% Preuve :
%% 1) On pourrait le faire
%% 2) On pourrait le faire : commencer par le cas ou A et B sont disjoints
%% 3) Faire une récurrence sur n


\begin{remarqueUnique}
\remarque L'avantage du passage au complémentaire est qu'il permet de prendre la négation
  de la propriété qui définit l'ensemble. Cela donne parfois une phrase plus simple à
  manipuler et donc un ensemble plus simple à dénombrer.
\end{remarqueUnique}

\begin{exos}
\exo Quel est le nombre d'entiers entre $1$ et $100$ qui ne sont pas divisibles par $3$~?
\exo Dénombrer \[A\defeq\enstq{n \in \intere{1}{100}}{2 |n \text{ ou } 3 |n}.\]
\exo Quel est le nombre de rythmes de $n$ temps que l'on peut composer uniquement avec
  des noires (1~temps) et des blanches (2~temps)~?
\end{exos}

\begin{proposition}[nom={Formule du crible}]
Soit $A_1,\ldots,A_n$ des parties d'un même ensemble fini $E$. Alors
  \[\card\p{\bigcup_{i=1}^n A_i}=\sum_{k=1}^n (-1)^{k+1}
    \sum_{1\leq i_1<\cdots<i_k\leq n}
    \card\p{A_{i_1}\cap\cdots\cap A_{i_k}}.\]
\end{proposition}

\begin{remarqueUnique}
\remarque   Par exemple, pour $n=3$, la formule du crible s'écrit
\begin{eqnarray*}
\card\p{A_1\cup A_2\cup A_3}&=&\card(A_1)+\card(A_2)+\card(A_3)\\
  & & - \cro{\card(A_2\cap A_3)+\card(A_1\cap A_3)+\card(A_1\cap A_2)}\\
& &+\card(A_1\cap A_2\cap A_3).
\end{eqnarray*}
\end{remarqueUnique}

\begin{proposition}
\begin{itemize}
\item Si $A$ et $B$ sont deux ensembles finis, alors $A\times B$ est fini et
  \[\card\p{A\times B}=\card(A)\card(B).\]
  Plus généralement, si $A_1,\ldots,A_n$ sont des ensembles finis,
  $A_1\times\cdots\times A_n$ est fini et
  \[\card(A_1\times \cdots\times A_n)=\prod_{k=1}^n \card(A_k).\]
\item Si $A$ est un ensemble fini et $n\in\N$, alors $A^n$ est fini et
  \[\card\p{A^n}=\card(A)^n.\]
\end{itemize}
\end{proposition}


\begin{remarqueUnique}
\remarque La formule $\card\p{A_1 \times \dots \times A_n}=\card(A_1)\cdots\card(A_n)$ dit
  simplement que lorsque l'on doit faire une succession \og et \fg de choix indépendants,
  on dénombre ces choix et on les multiplie.
\end{remarqueUnique}

\begin{exos}
\exo Montrer que dans un village de 700 personnes, deux au moins ont les mêmes initiales.
%   \begin{sol}
%   Il y en a $26^2$.
%   \end{sol}
% \exo On lance 5 dés. Combien de résultats possibles~?
\exo Combien de menus différents peut-on faire avec 4 entrées, 6 plats et 2 desserts~?
\exo Quel est le nombre de mots de 4 lettres contenant au moins un \og e\fg ? 
% \exo Quel est le nombre de mots de 5 lettres contenant un seul \og f\fg, un seul \og z\fg et pas de \og t\fg ? 
\end{exos}
  

\begin{proposition}
\begin{itemize}
\item Soit $E$ et $F$ deux ensembles finis. Alors $\mathcal{F}\p{E,F}=F^E$ est fini et
  \[\card\p{\mathcal{F}\p{E,F}}=\card\p{F^E}=\card(F)^{\card(E)}.\]
\item Soit $E$ un ensemble fini. Alors $\mathcal{P}(E)$ est fini et
  \[\card\p{\mathcal{P}(E)}=2^{\card(E)}.\]
\end{itemize}
\end{proposition}

\begin{preuve}
\begin{itemize}
\item Utiliser le 3 de la question précédente
\item si $E$ est de cardinal $E$, construire une bijection de $F(E,F)$ dans $F^n$
\item Bijection entre $P(E)$ est les applications de $E$ dans $\{0,1\}$
\end{itemize}
\end{preuve}

\begin{exoUnique}
\exo Quel est le nombre de possibilités de répartir $p$ boules distinctes
  dans $n$ urnes distinctes~?
  \begin{sol}
  Il y en a $n^p$.
  \end{sol}
\exo Une urne contient $n$ boules distinctes. On effectue $p$ tirages
  successifs avec remise (c'est-à-dire que l'on remet la boule dans l'urne
  après chaque tirage). Combien y-a-t-il de possibilités~?
  \begin{sol}
  Il y en a $n^p$.
  \end{sol}
\exo De combien de manières peut-on descendre $n+1$ marches (donc $n$
  \og paliers \fg), en en sautant éventuellement certaines~?
  \begin{sol}
  C'est l'ensemble des paliers qui seront sautés. Donc $2^n$.
  \end{sol}
\end{exoUnique}

\begin{proposition}[nom={Lemme des bergers}]
Soit $E$ et $F$ deux ensembles finis et $f$ une application de $E$ dans $F$. On
suppose qu'il existe $p\in\N$ tel que
\[\forall y\in F \qsep \card\p{f^{-1}\p{\ens{y}}}=p.\]
Alors $\card(E)=p\card(F)$.
\end{proposition}

\begin{remarqueUnique}
\remarque Soit $E$ l'ensemble des pattes des moutons foulant un pré, $F$
  l'ensemble des moutons du pré et $f$ l'application de $E$ dans $F$ qui à chaque
  patte associe son propriétaire. Comme chaque mouton a 4 pattes
  \[\forall m\in F\qsep \card\p{f^{-1}\p{\ens{m}}}=4.\] 
  On en déduit que le nombre de pattes foulant le pré est égal à quatre fois
  le nombre de moutons. C'est de cet exemple que la proposition précédente
  tire son nom de \og lemme des bergers \fg.
\end{remarqueUnique}

% \begin{exoUnique}
% \exo Soit $E$ un ensemble de cardinal $n$. Combien existe-t-il de couples
%   $(a,b)\in E^2$ avec $a \neq b$~? Combien existe-t-il de parties de $E$ de cardinal $2$~?
% \end{exoUnique}
  
\subsection{Arrangement, combinaison}

% On commence par compter le nombre de manières de choisir $p$ éléments parmi $n$ lorsque
% l'ordre compte.

\begin{definition}[nom={$p$-listes}]
Soit $E$ un ensemble et $p\in\N$. On appelle \emph{$p$-liste} d'éléments de $E$ tout
$p$-uplet $(a_1,\ldots,a_p)\in E^p$.
\end{definition}

\begin{exoUnique}
\exo Si $E=\ens{1,2,3}$, donner les $2$-listes d'éléments de $E$.
\end{exoUnique}

\begin{remarques}
\remarque Les $p$-listes d'éléments de $E$ sont les fonctions de $\intere{1}{p}$ dans $E$.
\remarque Choisir une $p$-liste, c'est choisir $p$ éléments de $E$ en tenant
  compte de  l'\emph{ordre} et en autorisant les \emph{répétitions}.
\end{remarques}

\begin{proposition}[nom={Nombre de listes}]
Soit $E$ un ensemble de cardinal $n$ et $p\in\N$. Alors, il existe
\[n^p\]
$p$-listes d'éléments de $E$.
\end{proposition}

% On compte maintenant le nombre de manières de choisir $p$ éléments \emph{distincts} parmi $n$ où l'\emph{ordre compte}.

\begin{definition}[nom={$p$-arrangements}]
Soit $E$ un ensemble et $p \in \N$. On appelle \emph{$p$-arrangement} d'éléments de $E$
toute $p$-liste $(a_1,\dots,a_p)\in E^p$ telle que
\[\forall i,j\in\intere{1}{p}\qsep i\neq j\quad\implique\quad a_i \neq a_j.\]
\end{definition}

\begin{exoUnique}
\exo Si $E=\ens{1,2,3}$, donner les $2$-arrangements d'éléments de $E$.
\end{exoUnique}

\begin{remarques}
\remarque Les $p$-arrangements d'éléments de $E$ sont les fonctions injectives $\intere{1}{p}$ dans $E$.
\remarque Choisir un $p$-arrangement, c'est choisir $p$ éléments de $E$ en tenant
  compte de  l'\emph{ordre} et en n'autorisant pas les \emph{répétitions}.
% \remarque C'est aussi le nombre de possibilités de choisir de manière ordonnée $p$
%   objets parmi $n$ objets.
\end{remarques}

\begin{proposition}[nom={Nombre d'arrangements}]
Soit $E$ un ensemble de cardinal $n$ et $p\in\N$. Alors il existe
\[A_n^p\defeq
  \begin{cases}
  \dsp\frac{n!}{\p{n-p}!} & \text{si $p\leq n$}\\
  0 & \text{si $p>n$}
  \end{cases}\]
  $p$-arrangements d'éléments de $E$.
\end{proposition}

\begin{preuve}
Récurrence sur $k$. On démontre facilement en utilisant le principe des bergers
que $A(n,k+1)=n.A(n-1,k)$.
\end{preuve}

\begin{remarqueUnique}
\remarque Si $n\in\N$ et $p\in\intere{0}{n}$, alors
  \[A_n^p = \underbrace{n\p{n-1}\cdots\p{n-(p-1)}}_{\text{$p$ termes}}.\]
\end{remarqueUnique}

\begin{exos}
\exo On répartit $p$ boules distinctes dans $n$ urnes distinctes. Quel est
  le nombre de répartitions pour lesquelles aucune urne ne contient plus d'une
  boule~?
  \begin{sol}
  C'est $A_n^p$.
  \end{sol}
\exo Une urne contient $n$ boules distinctes. On effectue $p$ tirages
  successifs sans remise. Combien y-a-t-il de possibilités~?
  \begin{sol}
  C'est $A_n^p$.
  \end{sol}
\end{exos}

\begin{proposition}
\begin{itemize}
\item Il existe $A_n^p$ injections d'un ensemble à $p\in\N$ éléments dans un ensemble à $n\in\N$ éléments.
\item Il existe $n!$ bijections d'un ensemble à $n\in\N$ éléments dans un ensemble à $n$ éléments.
\end{itemize}
\end{proposition}

\begin{preuve}
Prendre $k=n$ dans la proposition précédente.
\end{preuve}

\begin{remarqueUnique}
\remarque En particulier, si $E$ est un ensemble à $n$ éléments, il existe $n$ bijections de $E$ dans $E$.
  De telles applications sont appelées des \emph{permutations} de $E$.
\end{remarqueUnique}

\begin{exoUnique}
\exo Combien d'anagrammes peut-on former avec les mots \og maths \fg, \og chimie \fg et
  \og anagramme \fg~?
% \exo $2n$ personnes doivent se placer autour d'une table ronde. Combien de dispositions
%   sont possibles~? On suppose que ces $2n$ personnes sont constituées de $n$ hommes et $n$
%   femmes. Combien de dispositions sont possibles si l'on veut respecter l'alternance
%   homme/femme~?
%   \begin{sol}
%   Si on numérote les places, il y a $(2n)!$ possibilités. Si on ne les numérote pas, il y a
%   $2n$ rotations possibles, donc $(2n-1)!$ possibilités. Si on considère en plus les
%   symétries par rapport à une droite, et $n\geq 2$, il y a $(2n-1)!/2$ possibilités.\\
%   Si on numérote les places et on place une femme sur la première place, il y a $n!^2$
%   possibilités. Sans la contrainte de la première place, il y a $2 n!^2$. Si on ne
%   numérote plus les place, il y a $(n!)^2/(2n)$ possibilités.
%   \end{sol}
\end{exoUnique}

\begin{definition}[nom={$p$-combinaisons}]
Soit $E$ un ensemble et $p\in\N$. On appelle \emph{$p$-combinaison} d'éléments de $E$ toute
partie de $E$ à $p$ éléments.
\end{definition}

\begin{remarqueUnique}
\remarque Choisir une $p$-combinaison, c'est choisir $p$ éléments de $E$ sans tenir compte
  de l'\emph{ordre} et en n'autorisant pas les \emph{répétitions}.
\end{remarqueUnique}

\begin{proposition}[nom={Nombre de combinaisons}]
Soit $E$ un ensemble de cardinal $n$ et $p\in\N$. Alors il existe
\[C_n^p\defeq
    \begin{cases}
    \dsp\frac{n!}{p!\p{n-p}!} & \text{si $p\leq n$}\\
    0 & {\text{si $p>n$}}
    \end{cases}\]
$p$-combinaisons d'éléments de $E$.
\end{proposition}


\begin{preuve}
  Utiliser le principe des Bergers avec l'application phi qui a une injection
  de $\intere{1}{k}$ dans $E$ associe son image.
  \end{preuve}

\begin{remarques}
\remarque Si $n\in\N$ et $p\in\intere{0}{n}$, alors
  \[C_n^p = \frac{\overbrace{n\p{n-1}\cdots\p{n-(p-1)}}^{\text{$p$ termes}}}{p!}.\]
\remarque Le nombre de combinaisons $C_n^p$ est aussi noté $\binom{n}{p}$. Nous utiliserons par la suite
  cette notation qui est la notation internationale.
\end{remarques}
\vspace{1ex}
\begin{exos}
\exo Une urne contient $6$ boules numérotées de $1$ à $6$. Dénombrer les tirages
  possibles si on tire $3$ boules
  \begin{itemize}
  \question successivement et avec remise.
  \question successivement et sans remise.
  \question simultanément.
  \end{itemize}
\exo Un code de coffre-fort est composé de 6 chiffres entre $0$ et $9$ dont l'ordre
  compte. Dénombrer 
  \begin{itemize}
  \item tous les codes possibles.
  \item les codes dont tous les chiffres sont distincts.
  \item les codes ne contenant pas $0$.
  \item les codes contenant au plus deux fois le chiffre $1$.
  \item les codes contenant autant de chiffres pairs que de chiffres impairs.
  \item les codes contenant la succession $123$ quelque-part.
  \item les codes strictement croissants.
  \end{itemize}
% \exo On dispose de $n\in\Ns$ joueurs et l'on souhaite constituer une équipe de $p\in\Ns$ joueurs dont l'un sera désigné capitaine.
%   Dénombrer de deux fa\c{c}ons les équipes possibles et en d\'eduire
%   \[\binom{n}{p}=\frac{n}{p}\binom{n-1}{p-1}.\]
  % \exo Montrer que si $q \leq n, p$ alors \[{n+p \choose q}=
  % \sum\limits_{k = 0}^q {n \choose q-k} {p \choose k}.\]
  %   Cette identité est appelée identité de Vandermonde.
% À mettre en exo
% \exo Soit $n\in\Ns$. Combien y a-t-il de parties de $\intere{1}{n}$ ne contenant pas
%   deux entiers consécutifs~?
\exo On lance 5 dés. Combien y a-t-il de possibilités pour faire l'une des
  figures suivantes~?
  \begin{itemize}
  \item Un Yam's~: les 5 dés ont la même valeur.
  \item Une suite~: $(1,2,3,4,5)$ ou $(2,3,4,5,6)$.
  \item Un Full~: 3 dés identiques et les 2 autres identiques mais différents des premiers.
  \end{itemize}
\exo Quel est le nombre de manières de répartir $p$ boules indiscernables
dans $n$ urnes distinctes sachant qu'on ne peut pas mettre plus d'une boule par urne~?
  \begin{sol}
  C'est $\binom{n}{p}$.
  \end{sol}  
\exo Déterminer le nombre de solutions de l'équation $x_1+\cdots+x_n=p$ dans
  $\ens{0,1}^n$.
  % , puis dans $\Ns^n$ (on commencera par considérer l'inéquation
  % $x_1+\cdots+x_n\leq p$)~?
  \begin{sol}
  Pour la première question, il y a $\binom{n}{p}$ solutions. Si on associe
  à une solution de $x_1+\cdots+x_n\leq p$ le $n$-uplet
  $(x_1,x_1+x_2,\ldots,x_1+\cdots+x_n)$ on remarque que le nombre de solution
  de l'inéquation est le nombre de fonctions strictement croissantes de
  $\intere{1}{n}$ dans $\intere{1}{p}$. C'est donc $\binom{p}{n}$. Le nombre
  de solutions cherchées est donc
  $\binom{p}{n}-\binom{p-1}{n}=\binom{p-1}{n-1}$.
  \end{sol}
\exo Combien peut-on former de mots contenant $p$ fois
  la lettre O et $q$ fois la lettre I~?
  \begin{sol}
  C'est $\binom{n}{p}$.
  \end{sol}
\exo Quel est le nombre de manières de répartir $p$ boules indiscernables
  dans $n$ urnes distinctes~?
  \begin{sol}
  On peut modéliser cela par \og $..||...|.|..$ \fg. On est ramené un problème
  d'anagramme. On trouve donc $\binom{n+p-1}{p}$.
  \end{sol}
\exo Déterminer le nombre de solutions de l'équation $x_1+\cdots+x_n=p$ dans
  $\N^n$.
\exo Combien y a-t-il d'applications strictement
  croissantes de $\intere{1}{p}$ dans $\intere{1}{n}$~?
  \begin{sol}
  C'est $\binom{n}{p}$.
  \end{sol}
\exo Combien y a-t-il d'applications croissantes de $\intere{1}{p}$ dans $\intere{1}{n}$~?
  \begin{sol}
  Une application croissante est déterminée par le nombre d'antécédents de
  chaque élément de $\intere{1}{n}$. Le cardinal recherché est
  $\binom{n+p-1}{p}$.
  \end{sol}
\exo Soit $n\in\Ns$. Combien y a-t-il de surjections de $\intere{1}{n+1}$
  dans $\intere{1}{n}$~?
\end{exos}

\begin{proposition}
\begin{eqnarray*}
\forall n\in\N \qsep \forall p\in\intere{0}{n}, & &
  \binom{n}{p}=\binom{n}{n-p}\\
\forall n,p\in\N, & & \binom{n}{p}+\binom{n}{p+1}=\binom{n+1}{p+1}\\
\forall n,p\in\Ns, & & \binom{n}{p}=\frac{n}{p}\binom{n-1}{p-1}.
\end{eqnarray*}
\end{proposition}

\begin{preuve}
Utiliser des considérations ensemblistes
\end{preuve}

\begin{remarqueUnique}
\remarque La dernière formule est parfois appelée \og formule du capitaine \fg ou formule du
  \og comité président \fg.
\end{remarqueUnique}

\begin{proposition}
\begin{eqnarray*}
\forall n\in\N, & & \sum_{p=0}^n \binom{n}{p}=2^n\\
\forall a,b,n\in\N, & & \sum_{p=0}^n \binom{a}{p}\binom{b}{n-p}=\binom{a+b}{n}. 
\end{eqnarray*}
\end{proposition}

\begin{remarqueUnique}
\remarque La seconde formule est appelée formule de \nom{Vandermonde}.
\end{remarqueUnique}

\begin{proposition}[nom={Binôme de \nom{Newton}}]
\[\forall a,b\in\C\qsep \forall n\in\N\qsep \p{a+b}^n =\sum_{p=0}^n \binom{n}{p} a^{n-p} b^p.\]
\end{proposition}


% \subsection{Quelques exercices}

% \begin{exos}
% \exo Soit $k \in \N^*$. Combien existe-t-il de manières de colorier les
% arêtes d'un carré avec $k$ couleurs sans que deux arêtes
% adjacentes aient la même couleur ?
% \exo On lance $5$ dés. Quelle est la probabilité de faire l'une des figures suivantes :
% \begin{enumerate}
% \item Yam's : les $5$ dés ont la même valeur.
% \item Suite : $(1,2,3,4,5)$ ou $(2,3,4,5,6)$.
% \item Full : $3$ dés identiques et les $2$ autres identiques mais différents des premiers.
% \end{enumerate}
% \exo Dans le jeu des tours de Hanoï avec $3$ piquets et $n$ disques, de combien de manières peut-on disposer les $n$ disques ?
% \exo Calculer le nombre de parties de $\intere{1}{n}$ ne contenant pas deux entiers
% consécutifs.
% \exo Soit $n \in \N^*$. Combien existe-t-il de surjections de $I_{n+1}$ dans $I_n$ ?
% \exo Soit $n$ et $p \in \N^*$. Dénombrer les applications de $\intere{1}{p}$ dans $\intere{1}{n}$ :
% \begin{enumerate}
% \item strictement croissantes,
% \item croissantes,
% \item monotones.
% \end{enumerate}
% \exo 
% Soient $E$ et $F$ deux ensembles de cardinaux respectifs $n$ et $p$. On note
% $S_n^p$ le nombre de surjections de $E$ dans $F$.
% \begin{enumerate}
% \item Calculer $S_n^1$, $S_n^n$ et $S_n^p$ pour $p>n$.
% \item On suppose $p \leq n$ et on considère $a$ un élément de $E$. En observant qu'une surjection de $E$ dans $F$ réalise ou ne réalise pas
% une surjection de $E \setminus \{a\}$ dans $F$, montrer que
% $$S_n^p=p \left( S_{n-1}^{p-1} + S_{n-1}^p \right)$$
% \item En déduire \[S_n^p=\sum\limits_{k=0}^p (-1)^{p-k} {p \choose k} k^n.\]
% \end{enumerate}
% \end{exos}
%END_BOOK

% \begin{proposition}
% On a la formule du binôme de Newton~:
% \[\forall z_1,z_2\in\C \quad \forall n\in\N \quad
%   \p{z_1+z_2}^n=\sum_{k=0}^n \binom{n}{k} z_1^{n-k} z_2^k\]  
% \end{proposition}

% \begin{exos}
% \exo Soit $n\in\N$. Calculer
%   \[\sum_{\substack{k=0\\k\equiv 0\ [2]}}^n \binom{n}{k} \et
%     \sum_{\substack{k=0\\k\equiv 1\ [2]}}^n \binom {n}{k}\]
% \end{exos}
%END_BOOK

% \subsection{Problèmes classiques de dénombrement}

% \subsubsection{Applications de $\intere{1}{p}$ dans $\intere{1}{n}$}

% \subsubsection{Injections de $\intere{1}{p}$ dans $\intere{1}{n}$}

% \subsubsection{Parties à $p$ éléments d'un ensemble à $n$ éléments}


\end{document}

%%%%%%%%%%%%%%%%%%%%%%%%%%%%%%%%%%%%%%%%%%%%%%%%%%%%%%%%%%%%%%%%%%%%%%%%%%%%%%%%%%%
%%%%%%%%%%%%%%%%%%%%%%%%%%%%%%%%%%%%%%%%%%%%%%%%%%%%%%%%%%%%%%%%%%%%%%%%%%%%%%%%%%%

% \begin{itemize}
% \item {\bf Applications strictement croissantes de $\intere{1}{p}$ dans
%   $\intere{1}{n}$}~:\\
%   Étant donnés $n,p\in\Ns$, on cherche le cardinal de l'ensemble $E$ des fonctions
%   strictement croissantes de $\intere{1}{p}$ dans $\intere{1}{n}$. De telles
%   applications étant injectives, $\card E=0$ dès que $p>n$. On suppose donc dans
%   la suite que $p\leq n$.On remarque alors que l'application (où
%   $\mathcal{P}_p\p{\intere{1}{n}}$ désigne l'ensemble des parties de
%   $\intere{1}{n}$ à $p$ éléments)~:
%   \[\dspappli{\phi}{E}{\mathcal{P}_p\p{\intere{1}{n}}}{f}{\im f}\]
%   est une bijection. En effet~:
%   \begin{itemize}
%   \item $\phi$ est à valeurs dans $\mathcal{P}_p\p{\intere{1}{n}}$~:\\
%     En effet, soit $f$ une application strictement croissante de $\intere{1}{p}$
%     dans $\intere{1}{n}$. Alors $f$ est injective, donc
%     $\card\p{\im f}=\card \p{f\p{\intere{1}{p}}}=p$.
%   \item $\phi$ est injective~:\\
%     En effet, soit $f_1$ et $f_2$ deux applications strictement croissantes de
%     $\intere{1}{p}$ dans $\intere{1}{n}$ telles que $\im f_1=\im f_2$. Alors $f_1$
%     et $f_2$ sont des énumérations de $Y=\im f_1=\im f_2$, donc $f_1=f_2$.
%   \item $\phi$ est surjective~:\\
%     En effet, soit $Y$ une partie à $p$ éléments de $\intere{1}{n}$. Alors, il
%     existe une énumération de $Y$, c'est-à-dire une application strictement 
%     croissante $f$ de $\intere{1}{p}$ dans $\intere{1}{n}$ dont l'image est $Y$.
%     Alors $f\in E$ et $\phi(f)=Y$.
%   \end{itemize}
%   $\phi$ est bijection, donc~:
%   \[\card E=\card\p{\mathcal{P}_p\p{\intere{1}{n}}}=\binom{n}{p}\]
% \end{itemize}


% \end{document}



\section{Exercices}
\setcounter{numeroexercice}{1}
\documentclass{magnolia}

\magtex{tex_driver={pdftex}}
\magfiche{document_nom={Exercices sur le dénombrement},
          auteur_nom={François Fayard},
          auteur_mail={fayard.prof@gmail.com}}
\magexos{exos_matiere={maths},
         exos_niveau={mpsi},
         exos_chapitre_numero={10},
         exos_theme={Dénombrement}}
\magmisenpage{}
\maglieudiff{}
\magprocess

\begin{document}
%BEGIN_BOOK
\magsection{Cardinal}
\magsubsection{Équipotence}

\exercice{nom={Non dénombrabilité de $[0,1]$}}
Le but de cet exercice est de démontrer que $\interf{0}{1}$ n'est pas dénombrable,
c'est-à-dire qu'il n'existe pas de surjection de $\N$ dans $\interf{0}{1}$. On raisonne
par l'absurde et on suppose qu'il existe une suite $(u_n)$ d'éléments de $[0,1]$ telle que
\[\forall x\in[0,1]\qsep \exists n\in\N\qsep u_n=x.\]
\begin{questions}
\question Construire deux suites $(a_n)$ et $(b_n)$ de réels telles que
  \begin{itemize}
  \item $a_0= 0 \et b_0=1.$
  \item $\forall n\in\N \qsep a_n\leq a_{n+1}< b_{n+1}\leq b_n.$
  \item $\forall n\in\N \qsep b_{n+1}-a_{n+1}=\frac{1}{3}\p{b_n-a_n}.$
  \item $\forall n\in\N \qsep u_n \not\in \interf{a_{n+1}}{b_{n+1}}.$
  \end{itemize}
\question Montrer que les suites $(a_n)$ et $(b_n)$ convergent vers une limite commune 
  $l\in[0,1]$, puis que
  \[\forall n\in\N \qsep u_n\neq l.\]
  Conclure.
\end{questions}


\begin{sol}
Tout roule bien. En particulier, on fait de la trichotomie, car avec de la dichotomie, $u_n$ pourrait être le point milieu à l'étape $n$ de la construction. Pour la fin, fixons $n_0\in \N$, $u_{n_0}\notin \interf{a_{n_0}}{b_{n_0}}$ donc $u_{n_0}$ n'est pas dans l'intersection donc n'appartient pas à $\set{\ell}$. Ainsi, $u_{n_0}\neq \ell$ et ce quelque soit $n_0$.\\

Pour conclure, pour toute suite à valeurs dans $[0,1]$, c'est-à-dire toute fonction de $\N$ dans $[0,1]$, on a construit un élément jamais atteint. Donc elle n'est pas surjective.
\end{sol}

\magsubsection{Ensemble fini, cardinal}

\exercice{nom={Ensemble d'entiers}}
Étant donnés 51 entiers compris entre 1 et 100, montrer qu'il en existe toujours 2 consécutifs.

\exercice{nom={Théorème d'approximation de \nom{Dirichlet}}}
\begin{questions}
\question Soit $x\in\R$ et $n\in\Ns$. Pour tout $k\in\intere{0}{n}$, on pose
  $\delta_k\defeq kx-\ent{kx}$. En appliquant le principe des tiroirs aux
  réels $\delta_k$, montrer le théorème d'approximation de \nom{Dirichlet}
  \[\exists p\in\Z\qsep \exists q\in\intere{1}{n}\qsep \abs{x-\frac{p}{q}}<\frac{1}{nq}.\]
\question Soit $x\in\R\setminus\Q$.
  \begin{questions}
  \question Montrer qu'il existe une infinité de $(p,q)\in\Z\times\Ns$ tels
    que
    \[\abs{x-\frac{p}{q}}<\frac{1}{q^2}.\]
  \question Montrer qu'il existe une infinité de $p\in\Z$ pour lesquels
    \[\exists q\in\Ns\qsep \abs{x-\frac{p}{q}}<\frac{1}{q^2}.\]
  \end{questions}
\question On admet que $\pi$ est irrationnel. Dans ces conditions
  $\sin n\neq 0$ pour tout $n\in\Ns$ et on peut poser
  \[\forall n\in\Ns\qsep u_n\defeq\frac{1}{n\sin n}.\]
  On souhaite montrer que la suite $(u_n)$ n'admet pas de limite. On raisonne
  par l'absurde et on suppose que $(u_n)$ admet une limite
  $l\in\Rbar$.
  \begin{questions}
  \question Montrer que $l=0$.
  \question Obtenir une contradiction en appliquant les résultats de
    la question 2.b au réel $\pi$.
  \end{questions}

\end{questions}
% Soit $x\in\R$ et $q\geq 2$. Montrer qu'il existe $a\in\Z$ et
%   $b\in\intere{1}{q-1}$ tels que
%   \[\abs{bx-a}\leq \frac{1}{q}.\]
%   \begin{sol}
%   En effet, on considère 0, 1 et
%   les $\{bx\}$ pour $b\in\intere{1}{q-1}$ où $\{y\}$ désigne la distance de $y$
%   à sa partie entière. D'après la remarque précédente, il
%   existe deux éléments deux à deux distincts dont la distance est inférieure à
%   $1/q$ ce qui permet de conclure. (ENS Ulm/Lyon 1995) 
%   \end{sol}

\exercice{nom={7 nombres réels}}
Soit sept nombres réels $x_1,\ldots,x_7$. Montrer qu'il existe deux indices $i$
et $j$ distincts tels que 
\[0<\frac{x_i-x_j}{1+x_ix_j}<\frac{1}{\sqrt{3}}\>\cdot\]
\emph{Rappelez vous que la trigonométrie se cache même aux endroits où on ne l'attend pas}.
% {\it On pourra introduire, pour tout $k\in\intere{1}{7}$, $\theta_k$ tel que
% $\tan \theta_k=x_k$.}

\magsection{Dénombrement}
\magsubsection{Dénombrement élémentaire}


\exercice{nom={Couples dans le plan}}
Combien y a-t-il de couples $(i,j)$
\begin{questions}
\question dans $\intere{1}{n}^2$ pour lesquels $i+j=n$~?
\question dans $\intere{1}{n}^2$ pour lesquels $i<j$~?
\question dans $\intere{1}{n}\times\intere{1}{2n}$ pour lesquels $i<j$~?
\question dans $\intere{1}{n}^2$ pour lesquels $\abs{i-j}\leq 1$~?
\end{questions}

\exercice{nom={Autour du crible}}
Une certaine ville compte $17\ 500$ actifs dans sa population. À l'issue
d'un recensement, on a obtenu les informations suivantes sur ces $17\ 500$
actifs~:
\begin{itemize}
\item 4 actifs sur 7 sont des femmes et 6 d'entre elles sur 10 ont voté
  aux dernières municipales.
\item 3 actifs sur 5 ont voté aux dernières municipales et $40 \%$ de ces
  personnes sont au chômage.
\item Le chômage touche 1 actif sur 4 et $60 \%$ des demandeurs d'emploi
  sont des femmes.
\item $60 \%$ des femmes au chômage ont voté aux dernières municipales.
\end{itemize}
Combien d'hommes qui ne sont pas au chômage sont restés chez eux le jour des
élections municipales~?



\magsubsection{Arrangement, combinaison} 



\exercice{nom={Anagrammes}}
Dénombrer les anagrammes des mots suivants
\begin{center}
{\sc couvert, couteau, fourchette}.
\end{center}
\begin{sol}
Pour couvert, il y en a 7!. Pour couteau, il y en a 7!/2!. Pour fourchette, il
y en a 10!/(2! 2!).
\end{sol}

\exercice{nom={Le livreur}}
Un livreur doit distribuer des colis à cinq personnes A, B, C, D, E. Combien y
a-t-il de trajets possibles~? S'il souhaite livrer A avant B et C, combien
y a-t-il de trajets possibles~?
\begin{sol}
Il y a 5! trajets possibles. Si on livre $A$ avant $B$ et $C$, il y a
$2\times 4\times 5$ possibilités.
\end{sol}

\exercice{nom={Surjections d'un ensemble fini dans un autre}}
Quels que soient $n,p\in\N$, on note $S_{p,n}$ le nombre de surjections d'un
ensemble à $p$ éléments dans un ensemble à $n$ éléments.
\begin{questions}
\question On suppose que $p\leq n$. Que vaut $S_{p,n}$ ?
\question Calculer $S_{n+1,n}$ et $S_{p,2}$.
\question Montrer que
  \[\forall p,n \in \N \qsep n^p=\sum_{i=0}^n \binom{n}{i}S_{p,i}.\]
\end{questions}
\begin{sol}
$\quad$
\begin{questions}
\question Si $p<n$ il n'y a pas de surjection, donc $S_{p,n}=0$.
\question On a $S_{n+1,n}=\binom{n+1}{2}n!$. Si $p\leq 1$, alors $S_{p,n}=0$.
  Sinon, il faut choisir les antécédents d'un des élément. Il faut donc
  prendre une partie non vide de l'ensemble de départ différente de l'ensemble
  en entier. Il y a donc $2^p-2$ possibilités.
\question Considérer l'ensemble des fonctions d'un ensemble à $p$ éléments dans
  un ensemble à $n$ éléments. Faire ensuite une partition selon le cardinal de
  l'image de $f$.
\end{questions}
\end{sol}

\exercice{nom={Partitions d'un ensemble}}
Quels que soient $n,p\in\N$, on note $P_{n,p}$ le nombre de partitions d'un ensemble
de cardinal $np$ en $n$ parties à $p$ éléments. Montrer que
\[\forall n,p\in\N\qsep P_{n+1,p}=\frac{1}{n+1}\binom{(n+1)p}{p}P_{n,p}.\]
En déduire $P_{n,p}$.

\exercice{nom={Tours de Hanoï}}
Le jeu des tours de Hanoï se compose de trois tiges sur lesquelles on peut
empiler $n$ disques deux à deux distincts ($n\ge1$). Initialement, les $n$
disques sont empilés sur la première tige, par ordre décroissant de
taille, du bas vers le haut. Le but du jeu est de transporter la tour
complète sur une autre tige par une suite de mouvements consistant à
déplacer un disque à la fois, et en respectant les deux règles suivantes :
\begin{itemize}
\item on ne peut ôter d'une tige que le disque se trouvant au sommet de la
  pile ;
\item on ne peut empiler un disque sur une tige que si elle est vide ou bien si
  l'on pose le disque en question sur un autre plus grand.
\end{itemize}

\bigskip
\setlength{\unitlength}{2pt}
\begin{center}
\begin{picture}(130,15)
\put(0,0){\line(1,0){130}}
\put(5,2){\line(1,0){30}}
\put(7,4){\line(1,0){26}}
\put(9,6){\line(1,0){22}}
\put(11,8){\line(1,0){18}}
\put(13,10){\line(1,0){14}}
\put(5,0){\line(0,1){2}}
\put(35,0){\line(0,1){2}}
\put(7,2){\line(0,1){2}}
\put(33,2){\line(0,1){2}}
\put(9,4){\line(0,1){2}}
\put(31,4){\line(0,1){2}}
\put(11,6){\line(0,1){2}}
\put(29,6){\line(0,1){2}}
\put(13,8){\line(0,1){2}}
\put(27,8){\line(0,1){2}}
\put(19,10){\rule{2mm}{5mm}}
\put(59,0){\rule{2mm}{15mm}}
\put(99,0){\rule{2mm}{15mm}}
\end{picture}
\end{center}

\bigskip
\noindent
Notons $a_n$ le nombre minimal de mouvements nécessaires au transport de
la tour initiale de $n$ disques.
\begin{questions}
\question Montrer que $a_1=1$ et $a_2=3$. 
\question Établir une relation entre $a_{n+1}$ et $a_n$ pour tout $n\ge 1$.
\question Montrer que $a_n=2^n-1$ pour tout $n\ge 1$. 
\end{questions}





\exercice{nom={Exercice}}
Sur une étagère, on range les $n$ tomes d'une encyclopédie. Combien y a-t-il de manières de les ranger tout en étant sûr que le tome 1 et le tome 2 sont côte à côte et dans cet ordre ?


\exercice{nom={Exercice}}%http://images.math.cnrs.fr/Juin-2016-3e-defi.html
De combien de façons différentes peut-on ranger les nombres $1,2,\dots,n$ si l'on veut que le produit de deux nombres voisins soit toujours pair ?



\exercice{nom={Exercice}}%Monier-Haberer-Lardon Livre d'exos MPSI
Une urne contient 10 boules numérotées de $1$ à $10$. La boule 1 est jaune,
les boules 2 et 3 sont bleues, les boules 4,5,6 sont rouges et les boules 7,8,9,10
sont vertes. On tire dans l'urne successivement et avec remise 5 boules. Le résultat
est donc la liste ordonnée des cinq numéros des boules tirées. Déterminer le nombre de
résultats
\begin{enumerate}
\item en tout,
\item pour lesquels les cinq boules sont toutes de la même couleur,
\item pour lesquels les quatre couleurs apparaissent parmi les cinq boules,
\item pour lesquels la boule numéro 8 a été tirée et exactement deux des boules tirées sont rouges.
\end{enumerate}
\begin{sol}
\begin{questions}
\question Il y a $10^5$ tirages.
\question Il y a $1+2^5+3^5+4^5=1300$ tirages.
\question 4 couleurs différents. Si il y a deux jaunes, cela fait $2\times 3\times 4\times 5!/2=1440$.
  Si il y a deux bleues, cela fait $3\times 4\times 5!+2\times 3\times 4\times 5!/2=2880$.
  Si il y a deux rouges, 4320. Si il y a deux verts, 5760. Au total, cela donne $14400$.
\question Si on tire 3 fois 8 et 2 rouges, on a $(3 parmi 5)3^2=90$. Si on tire 2 fois 8 et 2 rouges,
  $(2 parmi 5)(2 parmi 3)3^2 6^1=1620$. Si on tire une fois 8 et 2 rouges $9720$.
  Cela fait au total 11430.\\
  Autre méthode~: On choisit où placer les rouges. Il y a 2 parmi 5 possibilités. Ensuite, on choisit
  les rouges, il y a $3^2$ possibilités. Il faut ensuite trouver les 3-listes avec au moins un 8. On peut
  utiliser le complémentaire. On trouve
  $(2 parmi 5) 3^2(7^3-6^3)$.
\end{questions}
\end{sol}




\exercice{nom={Exercice}}%Monier-Haberer-Lardon Livre d'exos MPSI
On dispose de trois urnes notées $A,B,C$ et de six boules. On répartit les six boules dans
les trois urnes (chaque urne peut contenir de $0$ à $6$ boules). Une répartition est une
liste ordonnée de trois nombres indiquant le nombre de boules contenues dans les urnes
$A,B,C$. Par exemple, la répartition $(2,4,0)$ indique que l'urne $A$ contient $2$ boules,
l'urne $B$ en contient $4$ et l'urne $C$ est vide.
Déterminer le nombre de répartitions
\begin{enumerate}
\item en tout,
\item telles que l'urne $A$ est vide,
\item telles que l'urne $A$ est la seule urne vide,
\item telles qu'une urne et une seulement est vide,
\item telles qu'aucune urne est vide,
\item telles qu'au moins une urne est vide.
\end{enumerate}
\begin{sol}
\begin{questions}
\question En utilisant la représentation $OOOIOOOI$, on trouve $\choose{8}{2}=28$ possibilités.
\question En utilisant la représentation $OOOIOOO$, on trouve $\choose{7}{1}=7$ possibilités.
\question On enlève les deux possibilités $B$ vide et $C$ vide. On a $7-2=5$ possibilités.
\question On a $3\times 5=15$ possibilités.
\question Il y a 10 possibilités.
\question Complémantaire du cas précédent. $28-10=18$.
\end{questions}
\end{sol}



\exercice{nom={Exercice}}
On considère un quadrillage de $n$ lignes et $m$ colonnes. On part de la case en haut à gauche pour arriver à la case en bas à droite. Les seuls mouvements possibles sont de se déplacer d'une case à droite ou d'une case en bas. 
Combien existe-t-il de chemins ?


%\exercice{nom={Exercice}}
%Sur un échiquier (donc 8 lignes et 8 colonnes), il ne reste plus que les 2 tours blanches et les deux noires. Combien de positions possibles pour que les tours ne se menacent pas entre-elles ?
%


\exercice{nom={Exercice}}%vincent
Un domino est un rectangle constitué de deux carrés, chacun comportant entre 0 et 6 points. 
% \begin{center}
% \includegraphics[width=5cm]{domino.jpg}
% \end{center}
\begin{enumerate}
\item Combien existe-t-il de dominos ?
\item Combien de paires peut-on former avec des dominos ayant un nombre en commun~?
\end{enumerate}
\begin{sol}
\begin{questions}
\question Si on ordonne les chiffres des dominos, il y en a $7\times 7=49$. Si on n'ordonne pas, il y en a $(49-7)/2+7=28$.
\question On choisit le chiffre en commun. Puis on choisit 2 chiffres parmi 7. On trouve 147.
\end{questions}
\end{sol}


\exercice{nom={Exercice}}
Soit $E$ un ensemble de cardinal $2n$. On appelle partition de $E$ en paires tout ensemble 
$\{ \{a_1,b_1 \}, \dots, \{ a_n, b_n \} \}$ avec pour tout $i \in \intere{1}{n}$, 
$a_i \neq b_i$ et $\p{\{a_1,b_1 \}, \dots, \{ a_n, b_n \}}$ est une partition de $E$.
Dénombrer les  partitions de $E$ en paires.


\exercice{nom={Exercice}}
Soit $E$ une partie de cardinal $n\in\Ns$. Dénombrer de deux manières différentes les couples
$(a,A)\in E\times\mathcal{P}(E)$ tels que $a\notin A$. En déduire
\[\forall n\in\Ns \qsep \sum_{k=1}^n k \binom{n}{k}=n 2^{n-1}.\]
\begin{sol}
On choisit d'abord $a\in A$. On a $n$ possibilités. On choisit ensuite une partie dans ce qui reste. On a $2^{n-1}$ possibilités.

\end{sol}

\exercice{nom={Compter les matrices}}
Combien existe-t-il de matrices de $\mat{q,p}{\R}$ dont les entrées sont
\og 0 \fg ou \og 1 \fg et
\begin{questions}
\question dont chaque ligne contient exactement un coefficient \og 1 \fg ~?
\question dont chaque ligne contient exactement deux coefficients \og 1 \fg ~?
\question dont chaque ligne et chaque colonne contiennent exactement un
  coefficient \og 1 \fg (on suppose ici $q=p$)~?
\end{questions}

\exercice{nom={Le crible}}
Soit $n\in\Ns$.
\begin{questions}
\question Montrer que
  \[\sum_{k=0}^n (-1)^{n-k}\binom{n}{k}k^n = n!.\]
  On pourra appliquer la formule du crible aux ensembles
  $\intere{1}{n}\setminus\ens{k}$ pour $k\in\intere{1}{n}$.
\question Simplifier de même
  \[\sum_{k=0}^n (-1)^{n-k}\binom{n}{k}k^p\]
  pour tout $p\in\intere{0}{n}$.
\end{questions}

%\exercice{nom={Exercice}}
%On dit qu'une partie $A$ de $\intere{1}{n}$ vérifie la propriété $(\star)$ ssi pour tout $a \in A$, $a+1$ et $a+2 \notin A$. On note $u_n$ le nombre de parties de $\intere{1}{n}$ vérifiant la propriété $(\star)$.
%
%Déterminer une relation de récurrence vérifiée par la suite $(u_n)$. En déduire la forme d'un équivalent de $u_n$.
%

\exercice{nom={Exercice}}
Pour tout $n\in\N$, on note $u_n$ le nombre d'applications $f$ de $\intere{1}{n}$ dans
$\intere{1}{n}$ telles que $f\circ f=\id$. Déterminer une relation de récurrence vérifiée
par la suite $(u_n)$.


%END_BOOK

\end{document}


















\chapter{Dimension finie}
\setcounter{numeroexercicecours}{1}
\documentclass{magnolia}

\magtex{tex_driver={pdftex},
        tex_packages={xypic}}
\magfiche{document_nom={Cours sur la dimension finie},
          auteur_nom={François Fayard},
          auteur_mail={fayard.prof@gmail.com}}
\magcours{cours_matiere={maths},
          cours_niveau={mpsi},
          cours_chapitre_numero={13},
          cours_chapitre={Dimension finie}}
\magmisenpage{misenpage_presentation={tikzvelvia},
          misenpage_format={a4},
          misenpage_nbcolonnes={1},
          misenpage_preuve={non},
          misenpage_sol={non}}
\maglieudiff{}
\magprocess

\begin{document}

%BEGIN_BOOK
\magtoc

\section{Famille libre, famille génératrice, base}
\subsection{Famille libre}

\begin{definition}[utile=-3]
On dit qu'une famille $(x_1,\ldots,x_p)\in E^p$ est \emph{libre} lorsque quels que
soient $\lambda_1,\ldots,\lambda_p\in\K$
\[\lambda_1 x_1+\cdots+\lambda_p x_p=0 \quad\implique\quad
  \lambda_1=\cdots=\lambda_p=0.\]
Sinon, on dit qu'elle est \emph{liée}. Dans ce cas, toute relation du type
$\lambda_1 x_1+\cdots+\lambda_p x_p=0$ où les $\lambda_k$ ne sont pas tous nuls
est appelée \emph{relation de liaison}.
\end{definition}

\begin{sol}
Faire le lien tout de suite entre liberté et identification
\end{sol}

\begin{exos}
\exo Donner une condition nécessaire et suffisante sur
  $\lambda\in\R$ pour que la famille $(\p{1,2,0}, \p{2,1,1}, \p{3,3,\lambda})$
  soit une famille libre de $\R^3$.
  \begin{sol} On résout $a(1,2,0)+b(2,1,1)+c(3,3,\lambda)=(0,0,0)$ à l'aide du pivot de Gauss.
  Si $\lambda\neq 1$, le système est de Cramer (après résolution avec la méthode du pivot de Gauss). Sinon, c'est équivalent à : il existe $t\in \R$ tel que $a+2b+3c=0$, $b+c=0$ et $c=t$ d'où $a=-t$, $b=-t$ et $c=t$ donc il existe une solution non nulle (même une infinité).
  \end{sol}
\exo Montrer que sur $\mathcal{F}\p{\R,\R}$, la famille $\sin,\cos$
  est libre. Que dire si on lui adjoint la fonction
  $x\mapsto\sin(x+\pi/4)$~?
  \begin{sol}
  $\forall x \in \R$, $\sin\p{x+\pi/4}=\sin(x)\cos(\pi/4)+\cos(x)\sin(\pi/4)$ ce qui donne une relation entre les fonctions.
  \end{sol} 
% \exo Si $a_0,\ldots,a_n\in\K$ sont $n+1$ scalaires deux à deux distincts,
%   La famille $P_0,\ldots,P_n\in\polyK$ définie par
%   \[\forall k\in\intere{0}{n} \quad P_k=\prod_{\substack{i=0\\i\neq k}}^n
%     \frac{X-a_i}{a_k-a_i}\]
%   est libre.
\exo Sur $\mathcal{F}\p{\R,\R}$, considéré comme un \Rev, montrer que si
  $\alpha_1,\cdots,\alpha_p\in\R$ sont tels que $\alpha_1<\cdots<\alpha_p$, la famille des fonctions d'expressions
  $\e^{\alpha_1 x},\e^{\alpha_2 x},\ldots,\e^{\alpha_p x}$ est libre.
  \begin{sol}
  Pour faire plusieurs méthodes, je rajoute l'exo avec les $\alpha_i$ réels deux à deux distincts.
  Pour les complexes, on le prouve par récurrence sur $p$. Pour l'hérédité, on dérive l'égalité et on enlève $\alpha_{p+1}$ fois la combinaison linéaire nulle à ce qu'on vient d'obtenir, ce qui permet d'annuler le terme le plus haut et donc d'utiliser l'HR sur ce qui reste. On conclut grâce aux $\alpha_i$ distincts deux à deux.
  \end{sol}
\exo Montrer que la famille $1,\sqrt{2},\sqrt{3}$ est libre dans $\R$
  considéré comme un $\Q$-espace vectoriel.
  \begin{sol}
  On élève $a+b\sqrt{2}=-c\sqrt{3}$ au carré :
  $a^2+2b^2+ab\sqrt{2}=3c^2$ donc si $ab\neq 0$, $\sqrt{2}\in \Q$ donc $a=0$ ou $b=0$. Ensuite, selon le cas, on recommence le même genre de méthodes.
  \end{sol}
\end{exos}

\begin{remarques}
\remarque Lorsque $(x_1,\ldots,x_n)$ est libre, on dit aussi que $x_1,\ldots,x_n$ sont
  \emph{linéairement indépendants}.
\remarque Si $(x_1,\ldots,x_p)$ est une famille libre d'éléments de $E$ alors,
  quels que soient $\lambda_1,\ldots,\lambda_p,\mu_1,\ldots,\mu_p\in\K$ tels que
  \[\lambda_1 x_1+\cdots+\lambda_p x_p=\mu_1 x_1+\cdots+\mu_p x_p\]
  on a $\lambda_1=\mu_1,\ldots,\lambda_p=\mu_p$. La liberté d'une famille permet
  donc \og d'identifier \fg les coefficients d'une combinaison linéaire.
\remarque[utile=-3] Une famille composée d'un unique vecteur est libre si et seulement si
  ce vecteur est non nul.
\remarque Une famille $(x,y)\in E^2$ formée de deux vecteurs est liée si et
  seulement si ces vecteurs sont \emph{colinéaires}, c'est-à-dire si et seulement si il existe
  $\lambda\in\K$ tel que $x=\lambda y$ ou $y=\lambda x$. Plus généralement, une famille est
  liée si et seulement si il existe un vecteur qui est combinaison linéaire des
  autres.  
\remarque[utile=-3] Une famille libre reste libre lorsqu'on effectue une permutation
  de ses vecteurs ou lorsqu'on retire certains de ses vecteurs. En particulier, une famille
  libre ne contient ni vecteur nul, ni doublon, ni vecteurs colinéaires.
\remarque Si $(x_1,\ldots,x_p)$ est
une famille libre et si $x\in E\setminus\vect\p{x_1,...,x_p}$, alors la famille
$(x_1,\ldots,x_p,x)$ est libre.
\remarque[utile=-1] Sur $\C$, considéré comme un $\R$-espace vectoriel, la famille
  $(1,\ii)$ est libre. Cependant, si $\C$ est considéré comme un \Cev, elle est liée. La notion
  de liberté dépend donc du corps considéré.
% \remarque[utile=-3] Sur $\C$, si $z\in\C\setminus\R$, la famille $(1,z)$ est $\R$-libre.
% \begin{sol}
% Si $a,b\in \R$, $a+bz=0$ conduit à $a=-bz \in \C\setminus\R$ si $b\neq 0$.
% \end{sol}
% \remarque L'image d'une famille liée par une application linéaire est une famille
%   liée. En particulier, si l'image d'une famille par une application
%   linéaire est libre, alors cette famille est libre.
% \remarque Pour tout $n\in\N$, la famille $1,X,\ldots,X^n$ est une famille
%   libre de $\polyK$.
\end{remarques}

\begin{proposition}
L'image d'une famille libre par une application linéaire injective est libre.
\end{proposition}

\subsection{Famille génératrice}

\begin{definition}[utile=-3]
On dit qu'une famille $(x_1,\ldots,x_p)\in E^p$ est \emph{génératrice} de $E$ lorsque, quel que soit
$x\in E$, il existe $\lambda_1,\ldots,\lambda_p\in\K$ tels que
\[x=\lambda_1 x_1+\cdots+\lambda_p x_p.\]
Autrement dit, la famille $(x_1,\ldots,x_p)$ est génératrice si et seulement si
$\vect\p{x_1,\ldots,x_p}=E$.
\end{definition}

\begin{remarques}
% \remarque Afin de montrer que $(x_1,\ldots,x_p)$ est une famille génératrice de $E$
  % il ne faut pas oublier de montrer que $x_1,\ldots,x_p$ sont des éléments de $E$.
\remarque Une famille génératrice reste génératrice lorsqu'on effectue une
  permutation de ses vecteurs ou lorsqu'on lui rajoute d'autres vecteurs.
\remarque Si $(x_1,\ldots,x_p)$ est une famille génératrice et si
  $x_p\in\vect\p{x_1,\ldots,x_{p-1}}$, alors la famille $(x_1,\ldots,x_{p-1})$ est
  génératrice.
\end{remarques}


\begin{exos}
  \exo Soit $\lambda\in\R$. Donner une condition nécessaire et suffisante sur
    $\lambda$ pour que la famille $((1,2,0), (2,1,1), (3,3,\lambda))$
    soit une famille génératrice de $\R^3$.
    \begin{sol}
    Il faut $\lambda\neq 1$.
    \end{sol}
  \exo Soit $\K$ un corps de caractéristique nulle, $n\in\N$ et $\alpha\in\K$. Montrer que la famille
    $(1,X-\alpha,(X-\alpha)^2,\ldots,(X-\alpha)^n)$ est une famille génératrice de $\polyK[n]$.
  \end{exos}


\begin{proposition}[utile=-2]
Soit $f$ et $g$ deux applications linéaires de $E$ dans $F$. On suppose que
$(x_1,\ldots,x_p)$ est une famille génératrice de $E$ telle que
\[\forall k\in\intere{1}{p} \qsep f\p{x_k}=g\p{x_k}.\]
Alors $f=g$.
\end{proposition}
  
\begin{exoUnique}
\exo Montrer la formule de \nom{Simpson}.
  \[\forall P\in\polyR[3] \qsep \integ{0}{1}{P(t)}{t}=
    \frac{P\p{0}+4P\p{1/2}+P\p{1}}{6}.\]
    \begin{sol}
    Les termes de gauche et de droite sont les images d'applications linéaires qui sont égales en $1, X, X^2$ et $X^3$ d'où le résultat.
    \end{sol}
\end{exoUnique}

\begin{proposition}[utile=-2]
Soit $f$ une application linéaire de $E$ dans $F$. On suppose que chaque
élément d'une famille génératrice $(y_1,\ldots,y_p)$ de $F$ admet
au moins un antécédent par $f$. Alors $f$ est surjective.
\end{proposition}

\begin{proposition}[utile=-2]
Soit $f$ une application linéaire de $E$ dans $F$. Alors l'image
$(f\p{x_1},\ldots,f\p{x_p})$ d'une famille génératrice $(x_1,\ldots,x_p)$ de $E$ est
une famille génératrice de $\im f$.
En particulier, si $f$ est surjective, l'image par $f$
d'une famille génératrice de $E$ est une famille génératrice de $F$.
\end{proposition}

\subsection{Base}

\begin{definition}[utile=-3]
On dit qu'une famille $(e_1,\ldots,e_n)\in E^n$ est une \emph{base} de $E$ lorsque,
quel que soit $x\in E$, il existe un unique $n$-uplet $\p{\lambda_1,\ldots,\lambda_n}\in\K^n$ tel que
\[x=\lambda_1 e_1+\cdots+\lambda_n e_n.\]
Autrement dit, la famille $(e_1,\ldots,e_n)$ est une base de $E$ si et seulement si elle est libre et génératrice.
\end{definition}

\begin{remarques}
\remarque Si $\mathcal{B}\defeq(e_1,\ldots,e_n)$ est une base de $E$ et $x\in E$, les
  $\lambda_1,\ldots,\lambda_n\in\K$ tels que $x=\lambda_1 e_1+\cdots+\lambda_n e_n$ sont appelées
  \emph{coordonnées} de $x$ relativement à la base $\mathcal{B}$.
\remarque[utile=-3] Une base reste une base lorsqu'on effectue une permutation de ses
  vecteurs.
\end{remarques}

\begin{definition}[utile=-3]
\begin{itemize}
\item Si $n\in\N$, la famille $(e_1,\ldots,e_n)$ de $\K^n$ définie par
  \[\forall k\in\intere{1}{n} \qsep e_k\defeq \p{0,\ldots,0,1,0,\ldots,0}\]
  est une base de $\K^n$, appelée \emph{base canonique}.
\item Si $n\in\N$, la famille $(1,X,\ldots,X^n)$ est une base de $\polyK[n]$.
  On l'appelle \emph{base canonique} de $\polyK[n]$.
\item Si $p,q\in\N$, la famille $(E_{i,j})_{1\leq i\leq q, 1\leq j\leq p}$ des
  matrices élémentaires est une base de $\mat{q,p}{\K}$ appelée
  \emph{base canonique}.
\end{itemize}
\end{definition}

% \begin{proposition}
% Soit $\mathcal{F}\defeq(e_1,\ldots,e_n)$ une famille d'éléments de $E$. Alors, les assertions suivantes
% sont équivalentes~:
% \begin{itemize}
% \item $\mathcal{F}$ est une base de $E$.
% \item $\mathcal{F}$ est \emph{libre maximale}, c'est-à-dire qu'elle est libre et que
%   l'ajout de tout élément la rend liée.
% \item $\mathcal{F}$ est \emph{génératrice minimale}, c'est-à-dire qu'elle est
%   génératrice et que le retrait de tout élément lui fait perdre son caractère
%   générateur.
% \end{itemize}
% \end{proposition}

\begin{proposition}[utile=3]
Soit $E$ et $F$ deux \Kevs tels que $E$ admette une base $(e_1,\ldots,e_n)$.
Alors, quel que soit $(y_1,\ldots,y_n)\in F^n$, il existe une unique application
linéaire $f$ de $E$ dans $F$ telle que
\[\forall k\in\intere{1}{n} \qsep f\p{e_k}=y_k.\]
De plus
\begin{itemize}
\item $f$ est injective si et seulement si $(y_1,\ldots,y_n)$ est libre.
\item $f$ est surjective si et seulement si $(y_1,\ldots,y_n)$ est génératrice dans $F$.
\item $f$ est un isomorphisme si et seulement si $(y_1,\ldots,y_n)$ est une base de $F$.
\end{itemize}
\end{proposition}

\begin{preuve}
Pour la première partie du théorème, on démontre l'unicité à l'aide de la proposition 1.2.
Pour l'existence, on définit alors $f$ en tout $x\in E$ en le décomposant sur $(e_1,\ldots,e_n)$. On montre ensuite que cette application est bien linéaire.\\


Pour la seconde partie du théorème :
On a déjà montré les implications gauche droite des deux premiers points. 
\begin{itemize}
\item[$\bullet$] On suppose que $(y_1,\ldots,y_n)$ est libre. Montrons que $f$ est injective. On prend $x\in \ker f$ et on démontre qu'il est nul en  le décomposant sur la base $(e_1,\ldots,e_n)$.
\item[$\bullet$] On suppose maintenant que $(y_1,\ldots,y_n)$ est génératrice et on démontre que $f$ est surjective.
\item[$\bullet$] On obtient le troisième point en combinant les deux premiers.
\end{itemize}
\end{preuve}

% \begin{remarques}
% \remarque[utile=3] En reprenant les notations du théorème,
%   $f$ est injective si et seulement si $y_1,\ldots,y_n$
%   est libre. De même, $f$ est surjective si et seulement si $y_1,\ldots,y_n$ est
%   une famille génératrice de $F$. Enfin, $f$ est bijective si et seulement si
%   $y_1,\ldots,y_n$ est une base de $F$.
% \remarque Soit $p,q\in\Ns$. Un application $f$ de $\K^p$ dans $\K^q$ est
%   linéaire si et seulement si il existe une famille de scalaires
%   $\p{a_{i,j}}_{1\leq i\leq q,1\leq j\leq p}$ telle que, pour tout
%   $\p{x_1,\ldots,x_p}\in\K^p$
%   \[f\p{x_1,\ldots,x_p}=\p{\sum_{k=1}^p a_{1,k} x_k,\ldots,
%                            \sum_{k=1}^p a_{q,k} x_k}\]
%   En particulier, si $n\in\Ns$, une application $f$ de $\K^n$ dans $\K$ est une
%   forme linéaire si et seulement si il existe $a_1,\ldots,a_n\in\K$ tels
%   que, pour tout $\p{x_1,\ldots,x_n}\in\K^n$
%   \[f\p{x_1,\ldots,x_n}=a_1 x_1+\cdots+a_n x_n\]
% \end{remarques}

\begin{exoUnique}
\exo On pose $E\defeq\R^3$. Montrer qu'il existe $u\in\mathcal{L}(E)$ tel que
  $u^3=\id$ et $u\neq \id$.
  \begin{sol}
  Suffit de bien la définir sur la base canonique en prenant par exemple $u(e_1)=e_2$, $u(e_2)=e_3$ et $u(e_3)=e_1$.
  \end{sol}
\end{exoUnique}

\begin{proposition}
Soit $A$ et $B$ deux sous-espaces vectoriels de $E$, $\mathcal{A}\defeq(a_1,\ldots,a_p)$ une famille d'éléments de $A$ et
$\mathcal{B}\defeq(b_1,\ldots,b_q)$ une famille d'éléments de $B$. On pose $\mathcal{E}\defeq(a_1,\ldots,a_p,b_1,\ldots,b_q)$.
\begin{itemize}
\item Si $A$ et $B$ sont en somme directe et $\mathcal{A}$ et $\mathcal{B}$ sont libres, alors $\mathcal{E}$ est libre.
\item Si $A+B=E$ et $\mathcal{A}$ et $\mathcal{B}$ sont respectivement génératrices de $A$ et $B$, alors
  $\mathcal{E}$ est génératrice de $E$.
\item Si $A\oplus B=E$ et $\mathcal{A}$ et $\mathcal{B}$ sont respectivement des bases de $A$ et $B$, alors
  $\mathcal{E}$ est une base de $E$. On dit qu'une telle base est \emph{adaptée} à la
  décomposition $E=A\oplus B$.
\end{itemize}
\end{proposition}

\begin{definition}
Soit $E$ un \Kev admettant une base $\mathcal{B}\defeq(e_1,\ldots,e_n)$. Pour tout $k\in\intere{1}{n}$,
on définit $e_k^\star\in E^\star$ comme étant l'application de $E$ dans $\K$ qui à $x$ associe la
$k$-ième coordonnée de $x$ relativement à la base $\mathcal{B}$. Les applications $e_k^\star$ sont
appelées \emph{applications coordonnées} de la base $\mathcal{B}$.
\end{definition}

\begin{remarques}
\remarque Si $\mathcal{B}\defeq(e_1,\ldots,e_n)$ est une base de $E$ et $x\in E$,
  alors $(e_1^\star(x),\ldots,e_n^\star(x))\in\K^n$ est la famille des coordonnées de $x$
  relativement à la base $\mathcal{B}$. Autrement dit
  \[\forall x\in E\qsep x=\sum_{k=1}^n e_k^{\star}(x)e_k.\]
\remarque En particulier, si $\mathcal{B}\defeq(e_1,\ldots,e_n)$ est la base canonique
  de $\K^n$ et $x\defeq(x_1,\ldots,x_n)\in\K^n$, alors
  \[\forall k\in\intere{1}{n}\qsep e_k^\star(x)=x_k.\]
\end{remarques}

\begin{definition}
Soit $E$ un \Kev admettant une base $\mathcal{B}\defeq(e_1,\ldots,e_n)$. Alors $(e_1^\star,\ldots,e_n^\star)$
est une base de $E^\star$, appelée \emph{base duale} de $\mathcal{B}$.
\end{definition}

\begin{remarqueUnique}
\remarque Dans $E\defeq \K^n$, si $\mathcal{B}\defeq(e_1,\ldots,e_n)$ est la base canonique
  de $E$ et $\phi\in E^*$, il existe $a_1,\ldots,a_n\in\K$ tels que
  \[\phi=a_1 e_1^*+\cdots+a_n e_n^*.\]
  Pour tout $x=(x_1,\ldots,x_n)\in\K^n$, on a donc $\phi(x)=a_1 x_1+\cdots+a_n x_n$.
\end{remarqueUnique}

\subsection{Cas des familles infinies}

\begin{definition}
  On dit qu'une famille $(\lambda_i)_{i\in I}$ d'éléments de $\K$ est \emph{presque nulle} lorsqu'il existe une partie finie
  $J\defeq\ens{i_1,\ldots,i_n}$ de $I$ telle que
  \[\forall i\in I\setminus J\qsep \lambda_i=0.\]
  On note $\K^{(I)}$ l'ensemble des familles presque nulles d'éléments de $\K$ indexées par $I$.
\end{definition}

\begin{remarques}
\remarque Si $I$ est fini, toute famille d'éléments de $\K$ indexée par $I$ est presque
  nulle. Cette notion n'a donc d'intérêt que lorsque $I$ est infini.
\remarque Si $(\lambda_i)_{i\in I}$ est une famille d'éléments de $\K$, on appelle
  \emph{support} de cette famille l'ensemble
  \[J\defeq\enstq{i\in I}{\lambda_i \neq 0}.\]
  Les familles presque nulles sont donc les familles à support fini.
\remarque Une suite $(\lambda_n)_{n\in\N}$ est presque nulle si et seulement si elle est
  nulle à partir d'un certain rang. En particulier, la famille des coefficients d'un
  polynôme est presque nulle.
\end{remarques}

\begin{definition}
Soit $(x_i)_{i\in I}$ est une famille d'éléments de $E$ et $(\lambda_i)_{i\in I}$ une
famille presque nulle d'éléments de $\K$. Soit $J\defeq\ens{i_1,\ldots,i_n}$ une partie
de $I$ en dehors de laquelle $\lambda_i$ est nul. On définit alors
\[\sum_{i\in I} \lambda_i x_i \defeq \sum_{k=1}^n \lambda_{i_k} x_{i_k}.\]
\end{definition}

\begin{remarques}
\remarque Si $P\in\polyK$ et $(a_k)_{k\in\N}$ est la famille de ses coefficients, alors
  \[P=\sum_{k\in\N} a_k X^k.\]
\remarque Soit $(\lambda_i)_{i\in I}$ et $(\mu_i)_{i\in I}$ deux familles presque nulles d'éléments de $\K$. On pose
  \[x\defeq\sum_{i\in I} \lambda_i x_i \quad\et\quad y\defeq\sum_{i\in I} \mu_i x_i.\]
  Alors, quels que soient $\lambda$ et $\mu\in\K$, $(\lambda\lambda_i+\mu\mu_i)_{i\in I}$ est une famille
  presque nulle et
  \[\lambda x+\mu y=\sum_{i\in I} (\lambda \lambda_i+\mu \mu_i)x_i.\]
\end{remarques}


\begin{proposition}
Soit $(x_i)_{i\in I}$ une famille d'éléments de $E$. Alors
\[\vect\ensim{x_i}{i\in I}=\ensim{\sum_{i\in I} \lambda_i x_i}{(\lambda_i)_{i\in I} \in \K^{(I)}}.\]
\end{proposition}

\begin{definition}[utile=-3]
On dit qu'une famille $(x_i)_{i\in I}$ d'éléments de $E$ est \emph{libre} lorsque, quelle que soit la famille
presque nulle $(\lambda_i)_{i\in I}$
\[\sum_{i\in I} \lambda_i x_i=0 \quad\implique\quad \forall i\in I\qsep \lambda_i=0.\]
Sinon, on dit qu'elle est \emph{liée}. Dans ce cas, toute relation du type
$\sum_{i\in I} \lambda_i x_i=0$ où les $\lambda_i$ ne sont pas tous nuls est appelée \emph{relation de liaison}.
\end{definition}

\begin{remarques}
\remarque On dit qu'une partie $A$ de $E$ est libre lorsque la famille de ses éléments est
  libre.
\remarque Une famille $(x_i)_{i\in I}$ d'éléments de $E$ est libre si et seulement si
  toute sous-famille finie de $(x_i)_{i\in I}$ est libre, c'est-à-dire si et seulement si
  pour tout $i_1,\ldots,i_n\in I$ deux à deux distincts et $\lambda_1,\ldots,\lambda_n\in\K$
  tels que
  \[\lambda_1 x_{i_1}+\cdots+\lambda_n x_{i_n}=0\]
  on a $\lambda_1=\cdots=\lambda_n=0$.
\remarque Une famille $(x_i)_{i\in\N}$ est libre si et seulement si pour tout $n\in\N$
  et $\lambda_0,\ldots,\lambda_n\in\K$ tels que
  \[\lambda_0 x_0+\cdots+\lambda_n x_n=0\]
  on a $\lambda_0=\cdots=\lambda_n=0$.
\end{remarques}

\begin{exoUnique}
\exo Soit $(P_n)_{n\in\N}$ une famille d'éléments de $\polyK$ telle que
\[\forall n\in\N\qsep \deg(P_n)=n.\]
Montrer qu'elle est libre.
\end{exoUnique}

\begin{sol}
On peut prendre comme exemple la famille $(x\mapsto |x-a|)_{a\in \R}$.
\end{sol}

\begin{definition}[utile=-3]
On dit qu'une famille $(x_i)_{i\in I}$ d'éléments de $E$ est \emph{génératrice} de $E$ lorsque, quel que soit
$x\in E$, il existe une famille presque nulle $(\lambda_i)_{i\in I}$ telle que
\[x=\sum_{i\in I} \lambda_i x_i.\]
Autrement dit, la famille $(x_i)_{i\in I}$ est génératrice si et seulement si $\vect\ensim{x_i}{i\in I}=E$.
\end{definition}

\begin{remarques}
\remarque On dit qu'une partie $A$ de $E$ est génératrice lorsque la famille de ses éléments est génératrice.
\remarque Une famille $(x_i)_{i\in I}$ d'éléments de $E$ est génératrice de $E$ si et
seulement si pour tout $x\in E$, il existe $i_1,\ldots,i_n\in I$ et $\lambda_1,\ldots,\lambda_n\in\K$ tels que
\[x=\sum_{k=1}^n \lambda_k x_{i_k}.\]
\remarque Comme dans le cas des familles finies, si $f,g\in\mathcal{L}(E,F)$ coïncident sur une famille génératrice de
$E$, alors $f=g$. De même, si $f\in\mathcal{L}(E,F)$ et que chaque élément d'une famille génératrice de $F$ admet un antécédent
par $f$, alors $f$ est surjective.
\remarque Enfin, l'image par une application linéaire $f\in\mathcal{L}(E, F)$ d'une famille génératrice de $E$ est une famille
génératrice de $\im f$.
\end{remarques}

\begin{definition}[utile=-3]
On dit qu'une famille $(e_i)_{i\in I}$ d'éléments de $E$ est une \emph{base} de $E$ lorsque, quel que soit $x\in E$, il existe une
unique famille presque nulle $(\lambda_i)_{i\in I}$ telle que
\[x=\sum_{i\in I} \lambda_i e_i.\]
Autrement dit, la famille $(e_i)_{i\in I}$ est une base de $E$ si et seulement si elle est libre et génératrice.
\end{definition}

\begin{remarques}
\remarque Si $(e_i)_{i\in I}$ est une base de $E$ et $x\in E$, la
  famille $(\lambda_i)_{i\in I}$ telle que
  $x=\sum_{i\in I} \lambda_i e_i$ est appelée famille des
  \emph{coordonnées} de $x$ dans la base $(e_i)_{i\in I}$.
% \remarque Comme pour les familles finies, une famille d'éléments de $E$ est une base
%   si et seulement si elle est libre maximale ou génératrice minimale.
\remarque Si $(e_i)_{i\in I}$ est une base de $E$ et $(y_i)_{i\in I}$ est une famille
d'éléments de $F$, alors il existe une unique application linéaire $f$ de $E$ dans $F$ telle que
\[\forall i\in I\qsep f(e_i)=y_i.\]
De plus, $f$ est injective si et seulement si $(y_i)_{i\in I}$ est libre, $f$ est surjective si et seulement si
$(y_i)_{i\in I}$ est génératrice de $F$ et $f$ est un isomorphisme si et seulement si $(y_i)_{i\in I}$ est une base de $F$.
\remarque La proposition sur les sommes de sous-espaces vectoriels
  et les familles s'énonce de manière similaire avec des familles
  infinies.
\end{remarques}

% \begin{exempleUnique}
% \exemple La famille $(X^k)_{k\in\N}$ est une base de $\polyK$.
% \end{exempleUnique}

\begin{definition}
La famille $(X^n)_{n\in\N}$ est une base de $\polyK$, appelée \emph{base canonique}.
\end{definition}

\begin{remarqueUnique}
\remarque Si $P\in\polyK$ alors la famille 
  $(a_k)_{k\in\N}$ des coefficients de $P$ est la famille des coordonnées de  $P$ dans
  la base canonique.
\end{remarqueUnique}

\section{Dimension}

\subsection{Espace vectoriel de dimension finie}
\begin{definition}[utile=-3]
On dit qu'un \Kev $E$ est de \emph{dimension finie} lorsqu'il admet une famille
génératrice finie. Dans le cas contraire, on dit que $E$ est de \emph{dimension
infinie}.
\end{definition}

\begin{remarqueUnique}
\remarque $\K^n$ et $\polyK[n]$ sont de dimension finie. $\polyK$ est
  de dimension infinie.
\end{remarqueUnique}


\begin{proposition}[utile=2,nom={Lemme de \nom{Steinitz}}]
Soit $E$ un \Kev de dimension finie et $(g_1,\ldots,g_q)$ une famille génératrice de $E$.
\begin{itemize}
\item Toute famille libre de $E$ possède au plus $q$ éléments.
\item Il est possible de compléter toute famille libre $(e_1,\ldots,e_p)$ en une base
$(e_1,\ldots,e_p,g_{i_{p+1}},\ldots,g_{i_n})$ de $E$ où $g_{i_{p+1}},
\ldots,g_{i_n}$ sont des éléments de la famille $(g_1,\ldots,g_q)$.
\end{itemize}

\end{proposition}

\begin{preuve}
Soit $\mathcal{F}$ l'ensemble des familles libres de $E$ formées des éléments
$(e_1,\ldots,e_m)$ et de certains des éléments de $(g_1,\ldots,g_p)$. Cet ensemble
est non vide car il contient la famille $(e_1,\ldots,e_m)$. De plus les familles
de $\mathcal{F}$ sont toutes de cardinal inférieur à $m+p$. Il existe
donc une famille de $\mathcal{F}$ de cardinal maximal. Notons
$(e_1,\ldots,e_m,g_{i_{m+1}},\ldots,g_{i_n})$ une telle famille et montrons que c'est
une base de $E$
\begin{itemize}
\item Cette famille est libre par définition de $\mathcal{F}$.
\item Montrons qu'elle est génératrice.
  \begin{itemize}
  \item Montrons d'abord que~:
    \[\forall k\in\intere{1}{p} \quad
      g_k\in\vect(e_1,\ldots,e_m,g_{i_{m+1}},\ldots,g_{i_n})\]
    Soit $k\in\intere{1}{p}$. Si $g_k$ est un des $g_{i_{m+1}},\ldots,g_{i_n}$, le
    résultat est évident. \\
    Sinon, la famille
    $(e_1,\ldots,e_m,g_{i_{m+1}},\ldots,g_{i_n},g_k)$ possède $n+1$ éléments. Puisque
    les familles de $\mathcal{F}$ ont toutes au plus $n$ éléments, on en déduit
    que cette famille n'est pas dans $\mathcal{F}$. Elle est donc liée. Il existe
    donc $\lambda_1,\ldots,\lambda_{n+1} \in\K$ non tous nuls tels que~:
    \[\lambda_1 e_1+\cdots+\lambda_m e_m + \lambda_{m+1} g_{i_{m+1}}+\cdots+
      \lambda_n g_{i_n} + \lambda_{n+1} g_k=0\]
    Montrons que $\lambda_{n+1}\neq 0$. On raisonne par l'absurde et on suppose que
    $\lambda_{n+1}=0$. Alors~:
    \[\lambda_1 e_1+\cdots+\lambda_m e_m + \lambda_{m+1} g_{i_{m+1}}+\cdots+
      \lambda_n g_{i_n}=0\]
    Puisque la famille $(e_1,\ldots,e_m,g_{i_{m+1}},\ldots,g_{i_n})$ est libre, on en
    déduit que $\lambda_1=\cdots=\lambda_n=0$. Donc les $\lambda_k$
    ($1\leq k\leq n+1$) sont tous nuls. C'est absurde, donc $\lambda_{n+1}\neq 0$.
    On en déduit donc que~:
    \[g_k=\p{-\frac{\lambda_1}{\lambda_{n+1}}} e_1+\cdots+
          \p{-\frac{\lambda_m}{\lambda_{n+1}}} e_m +
          \p{-\frac{\lambda_{m+1}}{\lambda_{n+1}}} g_{i_{m+1}}+\cdots+
          \p{-\frac{\lambda_n}{\lambda_{n+1}}} g_{i_n}\]
    Donc $g_k\in\vect\ens{e_1,\ldots,e_m,g_{i_{m+1}},\ldots,g_{i_n}}$
  \item Montrons désormais que la famille $(e_1,\ldots,e_m,g_{i_{m+1}},\ldots,g_{i_n})$
    est génératrice, c'est-à-dire que~:
    \[E=\vect(e_1,\ldots,e_m,g_{i_{m+1}},\ldots,g_{i_n})\]
    Puisque
    \[\forall k\in\intere{1}{p} \quad
      g_k\in\vect(e_1,\ldots,e_m,g_{i_{m+1}},\ldots,g_{i_n})\]
    on en déduit que~:
    \[\vect(g_1,\ldots,g_p) \subset
      \vect(e_1,\ldots,e_m,g_{i_{m+1}},\ldots,g_{i_n})\]
    Or la famille $(g_1,\ldots,g_p)$ est génératrice donc
    $\vect(g_1,\ldots,g_p)=E$. Donc~:
    \[E\subset  \vect(e_1,\ldots,e_m,g_{i_{m+1}},\ldots,g_{i_n})\]
    L'inclusion réciproque étant immédiate, on en déduit que~:
    \[E= \vect(e_1,\ldots,e_m,g_{i_{m+1}},\ldots,g_{i_n})\]
    et donc que $(e_1,\ldots,e_m,g_{i_{m+1}},\ldots,g_{i_n})$ est génératrice.
  \end{itemize}
\end{itemize}
En conclusion $(e_1,\ldots,e_m,g_{i_{m+1}},\ldots,g_{i_n})$ est une base de $E$.  
\end{preuve}
\begin{preuve}
Montrons ce résultat par récurrence sur $n$~:
\begin{center}
$\mathcal{H}_n$ : \flqq\ 
\parbox[t]{0.75\linewidth}{%
Soit $E$ un \Kev et $(g_1,\ldots,g_n)$ une famille génératrice de $E$. Alors
toute famille d'éléments de $E$ possédant plus de $n+1$ éléments est liée.
\frqq}
\end{center}
\begin{itemize}
\item $\mathcal{H}_0$ est vraie.\\
  En effet, soit $E$ un \Kev engendré par une famille vide. Alors $E=\ens{0}$.
  Montrons que toute famille qui possède au moins 1 élément est liée. En effet
  si $x_1\in E$ alors $x_1=0$ donc $1\cdot x_1=0$ est une relation de liaison.
\item $\mathcal{H}_n\implique\mathcal{H}_{n+1}$.\\
  On suppose que $\mathcal{H}_n$ est vraie. Montrons que $\mathcal{H}_{n+1}$ est
 vraie. En effet, soit $E$ un \Kev engendré par une famille $(g_1,\ldots,g_{n+1})$.
  Montrons que toute famille qui possède au moins $n+2$ éléments est liée. Soit
  $(x_1,\ldots,x_{n+2})$ une famille de $n+2$ éléments. Puisque
  $(g_1,\ldots,g_{n+1})$ est génératrice il existe une famille de scalaires
  $\p{\lambda_{i,j}}$ avec $i\in\intere{1}{n+2}$ et $j\in\intere{1}{n+1}$ telle
  que~:
  \[\syslin{x_1&=&\lambda_{1,1}g_1&+\lambda_{1,2}g_2&+\cdots&
                  +\lambda_{1,n+1}g_{n+1}\cr
            x_2&=&\lambda_{2,1}g_1&+\lambda_{2,2}g_2&+\cdots&
                  +\lambda_{2,n+1}g_{n+1}\cr
               &\hfill\vdots\hfill&&&&\cr
            x_{n+2}&=&\lambda_{n+2,1}g_1&+\lambda_{n+2,2}g_2&+\cdots&
                  +\lambda_{n+2,n+1}g_{n+1}}\]
  Deux cas se présentent~:
  \begin{itemize}
  \item Si $\forall i\in\intere{1}{n+2} \quad \lambda_{i,1}=0$.\\
    Dans ce cas la famille $(x_1,\ldots,x_{n+2})$ est une famille de $n+2$
    vecteurs de l'espace $F=\vect\p{g_2,\ldots,g_{n+1}}$ engendré par $n$
    éléments. Puisque $\mathcal{H}_n$ est vraie, on en déduit que cette famille
    est liée.
  \item Sinon, il existe $i\in\intere{1}{n+2}$ tel que $\lambda_{i,1}\neq 0$.
    Quitte à changer l'ordre des $x_i$ on peut supposer que
    $\lambda_{1,1}\neq 0$. Alors, en utilisant $\lambda_{1,1}$ comme pivot~:
    \[L_2\gets L_2-\frac{\lambda_{2,1}}{\lambda_{1,1}}L_1 \quad \ldots \quad
      L_{n+2}\gets L_{n+2}-\frac{\lambda_{n+2,1}}{\lambda_{1,1}}L_1\]
    \[\syslin{x_1&=&\lambda_{1,1}g_1&+\lambda_{1,2}g_2&+\cdots&
                  +\lambda_{1,n+1}g_{n+1}\cr
              x_2-\frac{\lambda_{2,1}}{\lambda_{1,1}}x_1
                &=&&\p{?}g_2&+\cdots&+\p{?}g_{n+1}\cr
               &\hfill\vdots\hfill&&&&\cr
              x_{n+2}-\frac{\lambda_{n+2,1}}{\lambda_{1,1}}x_1
                &=&&\p{?}g_2&+\cdots&+\p{?}g_{n+1}}\]
    où les $\p{?}$ sont des scalaires dont la valeur exacte est inutile. On en
    déduit que la famille de $n+1$ vecteurs~:
    \[x_2-\frac{\lambda_{2,1}}{\lambda_{1,1}}x_1,\ldots,
      x_{n+2}-\frac{\lambda_{n+2,1}}{\lambda_{1,1}}x_1\]
    appartient à l'espace vectoriel $F=\vect\p{g_2,\ldots,g_{n+1}}$ engendré
    par $n$ vecteurs. Puisque $\mathcal{H}_n$ est vraie, on en déduit que cette
    famille est liée. Il existe donc $\lambda_2,\ldots,\lambda_{n+2}\in\K$ non
    tous nuls tels que~:
    \[\lambda_2\p{x_2-\frac{\lambda_{2,1}}{\lambda_{1,1}}x_1}+\cdots+
      \lambda_{n+2}\p{x_{n+2}-\frac{\lambda_{n+2,1}}{\lambda_{1,1}}x_1}=0\]
    Donc~:
    \[\p{?}x_1+\lambda_2 x_2+\cdots+\lambda_{n+2}x_{n+2}=0\]
    qui est une relation de liaison puisque les $\lambda_k$ ne sont pas tous
    nuls. En conclusion, la famille $(x_1,\ldots,x_{n+2})$ est liée.
  \end{itemize}
   Donc $\mathcal{H}_{n+1}$ est vraie.
\end{itemize}
Par récurrence sur $n$ on en déduit donc que $\mathcal{H}_n$ est vraie pour
tout $n\in\N$.
\end{preuve}

\begin{remarqueUnique}
\remarque[utile=-2] Pour montrer qu'un espace vectoriel est de dimension infinie, il
  suffit donc de trouver des familles libres possédant autant d'éléments que l'on
  souhaite.
\end{remarqueUnique}

\begin{exoUnique}
\exo \`A l'aide d'éléments de la famille $(1, X, X^2)$, compléter la famille
  $(X^2-1,X^2+1)$ en une base de $\polyR[2]$.
\begin{sol}
Il suffit de rajouter un élément de la base canonique (par exemple) qui n'est pas dans son $\vect$. $X$ fonctionne. On vérifie que cela forme une base.
\end{sol}
\end{exoUnique}



\begin{theoreme}[nom={Théorèmes de la base incomplète et de la base extraite}]
Soit $E$ un \Kev de dimension finie.
\begin{itemize}
\item Toute famille libre se complète en une base finie de $E$.
\item Il est possible d'extraire une base finie de toute famille génératrice de $E$.
\end{itemize}
\end{theoreme}

\begin{preuve}
Comme $E$ est de dimension finie, il existe une famille génératrice (définition) donc on peut appliquer le théorème de la base incomplète.

Pour le deuxième point, on complète à partir de la famille vide.
\end{preuve}

% \begin{remarques}
% \remarque Algorithme de construction d'une base~:
% \begin{itemize}
% \item On commence par une famille $e_1,\ldots,e_m$
% \item Si cette famille est génératrice, on arrête. C'est une base.
% \item Sinon, on y ajoute un vecteur quelconque de $g_1,\ldots,g_p$ qui n'est pas
%   dans $\vect(F)$. On retourne à l'étape 2.
% \end{itemize}
% \end{remarques}


\begin{proposition}[utile=-3]
Tout espace vectoriel de dimension finie admet une base finie.
\end{proposition}

\begin{preuve}
C'est une autre lecture du deuxième point de la proposition précédente.
\end{preuve}

\begin{remarques}
\remarque[utile=-3] On en déduit qu'un espace vectoriel est de dimension finie si et seulement si il
  admet une base finie.
\remarque On peut montrer que tout espace vectoriel admet une base mais ce théorème est
  peu utile, hors programme, difficile à montrer et fait appel à l'axiome du choix.
\end{remarques}

\subsection{Dimension d'un espace vectoriel}


\begin{definition}[utile=-3]
Soit $E$ un \Kev de dimension finie. Alors toutes les bases de $E$ sont finies et ont le
même nombre d'éléments; on appelle \emph{dimension} de $E$ et on note $\dim E$ cet
entier.
\end{definition}
\begin{preuve}
Soit $E$ un \Kev de dimension finie, $\mathcal{B}_1=(e_1,\ldots,e_{n_1})$ et
$\mathcal{B}_2=(f_1,\ldots,f_{n_2})$ deux bases de $E$. Montrons que
$n_1=n_2$.\\
Puisque $E$ est engendré par $\mathcal{B}_2$ et que $\mathcal{B}_1$ est une
famille libre de $E$ le théorème précédent montre que $n_1\leq n_2$. De même on
montre que $n_2\leq n_1$, donc $n_1=n_2$.
\end{preuve}

\begin{remarques}
\remarque[utile=-2] Un espace vectoriel $E$ est de dimension nulle si et seulement
  si $E=\ens{0}$. On appelle \emph{droite vectorielle} tout espace vectoriel de dimension 1,
  c'est-à-dire tout espace vectoriel $E$ tel qu'il existe $x\in E\setminus\ens{0}$ tel
  que $E=\K x$. On appelle \emph{plan vectoriel} tout espace vectoriel de dimension 2.
\remarque[utile=-1] Considéré comme un \Rev, $\C$ est de dimension 2. Cependant, $\C$
  est de dimension 1 lorsqu'on le considère comme un \Cev. La dimension
  est donc une notion qui dépend du corps.
\end{remarques}

\begin{exos}
\exo Dans $\R^3$, donner une base du sous-espace vectoriel d'équation $x+y+z=0$.
\exo Dans $\polyR[2]$, donner une base puis la dimension du sous-espace vectoriel
  \[F=\enstq{P\in\polyR[2]}{\integ{0}{1}{P(t)}{t}=0}.\]
  
  \begin{sol}
  $\vect(X-1/2,X^2-1/3)$ puis on montre que c'est une famille libre.
  \end{sol}
\exo Si $E$ est un \Kev de dimension $n$, montrer qu'il existe $u\in\mathcal{L}(E)$
  tel que $u^n=0$ et $u^{n-1}\neq 0$.
  \begin{sol}
  On définit $u$ par $u(e_k)=e_{k-1}$ si $k\geq 2$ et $u(e_1)=0$.
  \end{sol}
\end{exos}

\begin{proposition}[utile=-3]
\begin{itemize}
\item Si $n\in\N$, $\K^n$ est un \Kev de dimension $n$.
\item Si $n\in\N$, $\polyK[n]$ est un \Kev de dimension $n+1$.
\item Si $p,q\in\N$, $\mat{q,p}{\K}$ est un \Kev de dimension $pq$. 
\end{itemize}
\end{proposition}

\begin{remarques}
\remarque En particulier, $\mat{n}{\K}$ est de dimension $n^2$.
\remarque $\mathcal{T}_n(\K)$ et $\mathcal{S}_n\p{\K}$ sont de dimension $n(n+1)/2$. $\mathcal{A}_n\p{\K}$ est de dimension $n(n-1)/2$.
\end{remarques}

\begin{proposition}[utile=-3]
Soit $E$ un \Kev de dimension $n$ et $(x_1,\ldots,x_p)$ une famille de $p$
éléments de $E$.
\begin{itemize}
\item Si $(x_1,\ldots,x_p)$ est libre, alors $p\leq n$.
\item Si $(x_1,\ldots,x_p)$ est génératrice, alors $p\geq n$.
\end{itemize}
\end{proposition}
\begin{preuve}
Soit $\mathcal{B}=(e_1,\ldots,e_n)$ une base de $E$ et $(x_1,\ldots,x_p)$ une
famille de $E$. Si $(x_1,\ldots,x_p)$ est libre, alors $p\leq n$ car
$\mathcal{B}$ est génératrice. Si $(x_1,\ldots,x_p)$ est génératrice, alors
$n\leq p$ car $\mathcal{B}$ est libre.
\end{preuve}

\begin{exoUnique}
\exo Soit $E$ un \Kev de dimension $n$ et $f$ un endomorphisme nilpotent
  de $E$. En notant $m$ le plus petit entier tel que $f^{m}=0$, montrer
  qu'il existe $x\in E$ tel que la famille $(x,f(x),\ldots,f^{m-1}(x))$ est
  libre. Que peut-on en déduire sur $m$~?
  
  \begin{sol}
  On prend $x\in E$ tel que $f^{m_0-1}(x)\neq 0$. Montrons que la famille $(x,f(x),\ldots,f^{m_0-1}(x))$ est libre. On considère une CL nulle de cette famille et on la suppose par l'absurde non triviale. On prend le plus petit indice tel que $\lambda_i\neq 0$ et on compose par la bonne puissance de $f$ qui annule tous les autres et laisse tout seul l'indice $i_0$ pour aboutir à une contradiction.
  \end{sol}
\end{exoUnique}

\begin{proposition}[utile=2]
Soit $E$ et $F$ deux \Kevs. On suppose que $E$ est de dimension finie. Alors
$E$ et $F$ sont isomorphes si et seulement si $F$ est de dimension finie et
$\dim E=\dim F$.
En particulier, si $E$ est de dimension $n$, $E$ est isomorphe à $\K^n$.
\end{proposition}
\begin{preuve}
$\Longrightarrow$
Supposons que $E$ et $F$ soient isomorphes et montrons que $F$ est de dimension
finie et que $\dim F=\dim E$. Soit $(e_1,\ldots,e_n)$ une base de $E$ et $f$
un isomorphisme de $E$ dans $F$. Alors la famille $(f\p{e_1},\ldots,f\p{e_n})$
est une base de $F$. En particulier $F$ est de dimension finie et
$\dim F=n=\dim E$.\\
$\Longleftarrow$
Réciproquement, supposons que $F$ soit de dimension finie et que
$\dim F=\dim E$. Soit $(e_1,\ldots,e_n)$ une base de $E$ et $(f_1,\ldots,f_n)$ une
base de $F$. Alors, il existe une unique application linéaire $f$ de $E$ dans
$F$ telle que~:
\[\forall k\in\intere{1}{n} \quad f\p{e_k}=f_k\]
Puisque $(f\p{e_1},\ldots,f\p{e_n})$ est une base de $F$, on en déduit que $f$
est un isomorphisme.
\end{preuve}

\begin{exos}
\exo Soit $a_0,\ldots,a_{p-1}\in\K$ et $E$ l'ensemble des suites $(u_n)$ à valeurs dans
$\K$ telles que
\[\forall n\in\N\qsep u_{n+p}=a_{p-1}u_{n+p-1}+\cdots+a_0 u_n.\]
Montrer que
\[\dspappli{\phi}{E}{\K^p}{(u_n)}{(u_0,\ldots,u_{p-1})}\]
est un isomorphisme. En déduire que $E$ est de dimension finie et que
$\dim E=p$.
\exo Soit $a,b,c\in\R$ avec $a\neq 0$ et $E$ l'ensemble des solutions de l'équation
  différentielle
  \[\forall t\in\R\qsep ay''(t)+by'(t)+cy(t)=0.\]
  Montrer que si $t_0\in\R$, alors
  \[\dspappli{\phi}{E}{\R^2}{y}{(y(t_0),y'(t_0))}\]
  est un isomorphisme. En déduire que $E$ est un \Rev de dimension 2.
\end{exos}

\subsection{Existence et unicité en dimension finie}

\begin{proposition}[utile=3]
Soit $E$ un \Kev de dimension $n$. Alors
\begin{itemize}
\item Toute famille libre de $E$ comportant $n$ éléments est une base de $E$.
\item Toute famille génératrice de $E$ comportant $n$ éléments est une base de
  $E$.
\end{itemize}
Autrement dit, si la famille $\mathcal{F}$ est composée de $n$ éléments,
\[\text{$\mathcal{F}$ est libre} \quad\ssi\quad \text{$\mathcal{F}$ est une base}
  \quad\ssi\quad \text{$\mathcal{F}$ est génératrice.}\]
\end{proposition}
\begin{preuve}
Soit $(e_1,\ldots,e_n)$ une famille libre à $n$ éléments de $E$. D'après le
théorème de la base incomplète, on peut la compléter en une base
$\mathcal{B}=(e_1,\ldots,e_n,\ldots,e_{n+k})$ de $E$. Puisque $\dim E=n$ on en
déduit que $n+k=n$ donc $k=0$. En conclusion $\mathcal{B}=(e_1,\ldots,e_n)$ donc
$(e_1,\ldots,e_n)$ est une base de $E$.\\
Soit $(e_1,\ldots,e_n)$ une famille génératrice à $n$ éléments de $E$. En partant
de la famille vide, on peut d'après le théorème de la base à incomplète la
compléter avec des éléments de la famille génératrice $(e_1,\ldots,e_n)$ pour
en faire une base de $E$. Comme cette base a $n$ éléments, on en déduit que
c'est la famille $(e_1,\ldots,e_n)$.
\end{preuve}

\begin{proposition}
Soit $\mathcal{B}\defeq (P_0,\ldots,P_n)$ une famille de polynômes de degrés
\emph{échelonnés}, c'est-à-dire telle que
\[\forall k\in\intere{0}{n} \qsep \deg P_k=k.\]
Alors $\mathcal{B}$ est une base de $\polyK[n]$.
\end{proposition}

% \begin{remarqueUnique}
% \remarque[utile=3] On dit qu'une famille de polynômes $(P_0,\ldots,P_n)$ est
%   \emph{échelonnée} lorsque
%   Une telle famille est une base de $\polyK[n]$.
% %  En particulier
% %   si $\alpha\in\K$, la famille $1$, $(X-\alpha)$,
% %   $(X-\alpha)^2,\ldots,(X-\alpha)^n$ est une base de $\polyK[n]$.
% \end{remarqueUnique}

\begin{exos}
% \exo La famille $\p{0,1,1}$, $\p{1,0,1}$, $\p{1,1,0}$ est libre. C'est donc
%   une base de $\R^3$.
\exo Montrer que l'application de $\polyC[n+1]$ dans $\polyC[n]$ qui à $P$
  associe $P\p{X+1}-P(X)$ est surjective.
  \begin{sol}
  On montre que $\phi(X^k)$ est de degré $k-1$ donc on a $(\phi(X),\ldots,\phi(X^{n+1}))$ qui forme une famille de degrés échelonnés de $\polyC[n]$ donc ça en est une base, en particulier c'est générateur, d'où le résultat.
  \end{sol}
% \exo La famille $P_0,\ldots,P_n$ définie par $P_k=X^{n-k}(1-X)^k$ est une
%   base de $\polyK[n]$.
%\exo Base de Lagrange.
% \exo Soit $n\in\Ns$ et $E$ l'espace des fonctions affines par morceaux sur
%   $\interf{k/n}{(k+1)/n}$. Alors~:
%   \begin{itemize}
%   \item $E$ est de dimension $n+1$~: Construire l'isomorphisme
%     $f\mapsto \p{f(0),\ldots,f(1)}$.
%   \item La famille : $x\mapsto\abs{x-k/n}$ est libre
%   \item Comme elle comporte $n+1$ éléments, c'est une base de $E$.
%   \end{itemize}
\exo Quels sont les sous-espaces vectoriels de $\polyR$ stables par
  dérivation~?
  \begin{sol}
  $\polyK$ et $\polyK[n]$. On fait une analyse-synthèse. On introduit l'ensemble des degrés d'éléments de $E\setminus\set{0}$ et on étudie le cas où cet ensemble est majoré ou non.
  \end{sol}
\end{exos}

\begin{proposition}[utile=3]
Soit $E$ et $F$ deux \Kevs de dimension finie et $f$ une application linéaire
de $E$ dans $F$.
\begin{itemize}
\item Si $f$ est injective et $\dim E=\dim F$, alors $f$ est un isomorphisme.
\item Si $f$ est surjective et $\dim E=\dim F$, alors $f$ est un isomorphisme.
\end{itemize}
Autrement dit, si $\dim E=\dim F$
\[\text{$f$ est injective} \quad\ssi\quad \text{$f$ est bijective}
   \quad\ssi\quad \text{$f$ est surjective.}\]
\end{proposition}
\begin{preuve}
Supposons que $f$ soit injective et que $\dim E=\dim F=n$. Soit
$(e_1,\ldots,e_n)$ une base de $E$. Puisque $f$ est injective on en déduit que
$(f\p{e_1},\ldots,f\p{e_n})$ est libre dans $F$. Or $\dim F=n$ donc cette
famille est une base de $F$. Donc $f$ est un isomorphisme.\\
Supposons que $f$ soit surjective et que $\dim E=\dim F=n$. Soit
$(e_1,\ldots,e_n)$ une base de $E$. Puisque $f$ est surjective
$(f\p{e_1},\ldots,f\p{e_n})$ est génératrice dans $F$. Puisque $\dim F=n$ on
en déduit que c'est une base de $F$. Donc $f$ est un isomorphisme.
\end{preuve}

\begin{remarques}
\remarque[utile=-3] Si $f$ est un endomorphisme d'un \Kev de dimension
  finie $E$, pour montrer que c'est un automorphisme, il suffit de montrer
  qu'il est injectif (ou surjectif).
\remarque[utile=-1] Ce théorème est faux si $E$ et $F$ ne sont pas de même dimension
  ou si ils sont de dimension infinie. Par exemple, les applications linéaires
  \[\dspappli{f}{\R}{\R^2}{x}{\p{x,0}} \et
    \dspappli{g}{\polyK}{\polyK}{P}{XP}\]
  sont injectives mais pas surjectives. De même, les applications linéaires
  \[\dspappli{f}{\R^2}{\R}{\p{x,y}}{x+y} \et
    \dspappli{f}{\polyK}{\polyK}{P}{P'}\]
  sont surjectives mais ne sont pas injectives.
\end{remarques}

\begin{proposition}
Soit $E$ un \Kev de dimension finie et $f\in\Endo{E}$. Alors
\begin{itemize}
\item $f$ est inversible si et seulement si $f$ est inversible à gauche, c'est-à-dire
  si et seulement si il existe $g\in\mathcal{L}(E)$ tel que $g\circ f=\id_E$. Si tel
  est le cas, $g=f^{-1}$ et en particulier $f\circ g=\id_E$.
\item $f$ est inversible si et seulement si $f$ est inversible à droite, c'est-à-dire
  si et seulement si il existe $g\in\mathcal{L}(E)$ tel que $f\circ g=\id_E$. Si tel
  est le cas, $g=f^{-1}$ et en particulier $g\circ f=\id_E$.
\end{itemize}
\end{proposition}


\subsection{Dimension d'un sous-espace vectoriel}
\begin{proposition}[utile=3]
Soit $E$ un \Kev de dimension finie et $A$ un sous-espace vectoriel de $E$.
Alors
\begin{itemize}
\item $A$ est de dimension finie et $\dim A\leq \dim E$.
\item $A=E$ si et seulement si $\dim A=\dim E$.
\end{itemize}
\end{proposition}
\begin{preuve}
Soit $E$ un \Kev de dimension finie $n$ et $A$ un sous-espace vectoriel de $E$.
Montrons que $A$ est de dimension finie. Soit $\mathcal{A}$ l'ensemble des
familles libres de $A$. Puisque toute famille libre de $A$ est une famille
libre de $E$ ces familles ont au plus $n$ éléments. Il existe donc une
famille libre de $A$ de cardinal maximum. Soit $(e_1,\ldots,e_m)$ une telle
famille. Montrons que c'est une base de $A$.
\begin{itemize}
\item Cette famille est libre par définition.
\item Montrons que cette famille est génératrice. Soit $x$ un élément de $A$.
  Puisque $(e_1,\ldots,e_m,x)$ possède $m+1$ éléments, elle n'est pas libre.
  Il existe donc $\lambda_1,\ldots,\lambda_m,\lambda\in\K$ non tous nuls tels
  que~:
  \[\lambda_1 e_1+\cdots+\lambda_m e_m +\lambda x=0\]
  Alors $\lambda\neq 0$. En effet si tel était le cas on aurait
  $\lambda_1 e_1+\cdots+\lambda_m e_m=0$ donc tous les $\lambda_k$ seraient
  nuls car la famille $(e_1,\ldots,e_m)$ est libre. Ce serait absurde car
  la famille $(\lambda_1,\ldots,\lambda_m,\lambda)$ comporte au moins un élément
  non nul. Donc $\lambda\neq 0$ et~:
  \[x=\p{-\frac{\lambda_1}{\lambda}}e_1+\cdots+
      \p{-\frac{\lambda_m}{\lambda}}e_m\]
  En conclusion, la famille $(e_1,\ldots,e_m)$ est génératrice.
\end{itemize}
Donc la famille $(e_1,\ldots,e_m)$ est une base de $A$. $A$ est donc de dimension
finie. De plus $m\leq n$ par définition de $(e_1,\ldots,e_m)$. Donc $A$ est de
dimension finie et $\dim A\leq \dim E$.\\
Montrons que $A=E$ si et seulement si $\dim A=\dim E$. Si $E=A$ alors il est
évident que $E$ et $A$ sont de même dimension. Réciproquement supposons que
$E$ et $A$ soient de même dimension $n$. Soit $(e_1,\ldots,e_n)$ une base de
$A$. Puisque cette famille est libre dans $E$ et possède $n$ éléments, c'est
une base de $E$. Donc $E=A$.
\end{preuve}

\begin{remarqueUnique}
\remarque Si $E$ est un \Kev de dimension $n$ et $A$ un sous-espace vectoriel de $E$
  de dimension $p$, toute base $\mathcal{B}\defeq(e_1,\ldots,e_n)$ de $E$
  telle que $(e_1,\ldots,e_p)$ est une base de $A$ est appelée \emph{base adaptée}
  au sous-espace vectoriel $A$.
\end{remarqueUnique}

\begin{exoUnique}
\exo Soit $E$ un \Kev de dimension finie, $u\in\gl{}{E}$ et $F$ un
  sous-espace vectoriel de $E$. Montrer que si $F$ est stable par $u$,
  alors $u(F)=F$ et $F$ est stable par $u^{-1}$.

\end{exoUnique}

  \begin{sol}
  On sait que $\forall x \in F$, $u(x)\in F$ donc $u(F)\subset F$. Réciproquement, on pose $\dspappli{\phi}{F}{F}{x}{u(x)}$. $\phi$ est bien défini grâce à la stabilité. Elle est linéaire. On veut montrer qu'elle est surjective. Comme $F$ est de dimension finie, c'est équivalent de démontrer que $F$ est injective. On prend $x\in \ker\phi$. Alors $\phi(x)=u(x)=0$. Donc $x\in \ker u=\set{0}$ car $u$ est un isomorphisme. Ainsi, $u$ est bien injective donc surjective. Donc $\forall y \in F$, il existe $x\in F$ tel que $u(x)=\phi(x)=y$.\\
  Montrons que $F$ est stable par $u^{-1}$. Soit $y\in F$, il existe $x\in F$ tel que $u(x)=y$ d'après la première partie. Ainsi, $u^{-1}(y)=u^{-1}(u(x))=x \in F$ ce qui montre bien la stabilité.
  \end{sol}



\begin{proposition}[utile=2]
Dans un \Kev de dimension finie, tout sous-espace vectoriel admet un
supplémentaire.
\end{proposition}
\begin{preuve}
Puisque $A$ est un sous-espace vectoriel de $E$, il est de dimension finie.
Il admet donc une base $(e_1,\ldots,e_m)$. Cette famille est libre dans $E$, donc
d'après le théorème de la base incomplète on en déduit qu'on peut la
compléter en une base $(e_1,\ldots,e_m,\ldots,e_{m+k})$ de $E$. On pose
$B=\vect\p{e_{m+1},\ldots,e_{m+k}}$. Montrons que $B$ est un supplémentaire de
$A$ dans $E$.
\begin{itemize}
\item La somme est directe, c'est-à-dire que $A\cap B=\ens{0}$.\\
  En effet, soit $x\in A\cap B$. Puisque $x\in A$, il existe
  $\lambda_1,\ldots,\lambda_m\in\K$ tels que $x=\lambda_1 e_1+\cdots+
  \lambda_m e_m$. Comme de plus $x\in B$, il existe $\lambda_{m+1},\ldots,
  \lambda{m+k} e_{m+k}$ tels que $x=\lambda_{m+1} e_{m+1}+\cdots+
  \lambda_{m+k}e_{m+k}$. Donc~:
  \[\lambda_1 e_1+\cdots+\lambda_m e_m -\lambda_{m+1} e_{m+1}-\cdots-
    \lambda_{m+k}e_{m+k}=0\]
  Puisque la famille $(e_1,\ldots,e_{m+k})$ est libre, on en déduit que tous les
  $\lambda_k$ sont nuls. En conclusion $x=0$. Donc la somme est directe.
\item $E=A+B$.\\
  En effet, soit $x\in E$. Alors, puisque $e_1,\ldots,e_{m+k}$ est génératrice
  dans $E$, il existe $\lambda_1,\ldots,\lambda_{m+k}$ tels que~:
  \[x=\underbrace{\lambda_1 e_1+\cdots+\lambda_m e_m}_{\in A} +
      \underbrace{\lambda_{m+1} e_{m+1}+\cdots+\lambda_{m+k} e_{m+k}}_{\in B}\]
  Donc $x\in A+B$.
\end{itemize}
En conclusion $E=A\oplus B$, donc $B$ est un supplémentaire de $F$ dans $E$.
\end{preuve}

\begin{exoUnique}
\exo Trouver un supplémentaire dans $\R^3$ de l'espace vectoriel engendré par $\p{1,1,1}$.
  \begin{sol}
  On complète $\p{1,1,1}$ en une base avec des vecteurs d'une base (la canonique par exemple).
  \end{sol}
% \exo Soit $E$ un \Kev de dimension finie et $f,g\in\Endo{E}$. Montrer
%   qu'il existe $h\in\Endo{E}$ tel que $g=h\circ f$ si et seulement si
%   $\ker f\subset \ker g$. 
\end{exoUnique}

\begin{remarqueUnique}
% \remarque Si $E$ est un \Kev de dimension $n$, $A$ un sous-espace vectoriel de $E$
%   de dimension $p$ et $B$ un supplémentaire de $A$ de dimension $q$ (où $p+q=n$),
%   toute base
%   $\mathcal{B}\defeq(a_1,\ldots,a_p,b_1,\ldots,b_q)$ de $E$ telle que
%   $(a_1,\ldots,a_p)$ est une base de $A$ et $(b_1,\ldots,b_q)$ est une base de $B$ est
%   appelée \emph{base adaptée} à la décomposition $E=A\oplus B$.
\remarque Si $E$ est de dimension infinie, on peut montrer que tout sous-espace vectoriel
  de $E$ admet un supplémentaire, mais ce théorème est peu utile, hors programme, difficile à
  démontrer et fait appel à l'axiome du choix.
\end{remarqueUnique}

\subsection{Notion de rang}

\begin{definition}[utile=-3]
Soit $f$ une application linéaire de $E$ dans $F$. Lorsque $\im f$ est de
dimension finie, on appelle \emph{rang} de $f$ et on note $\rg f$ la dimension de
$\im f$.
\end{definition}

\begin{remarques}
\remarque[utile=-2] Si $F$ est de dimension finie, $\im f$ est de dimension finie et
  $\rg f\leq \dim F$. De plus, cette inégalité est une égalité si et seulement
  si $f$ est surjective.
\remarque Si $E$ est de dimension finie, $\im f$ est
  de dimension finie et $\rg f\leq \dim E$. De plus, cette inégalité est une
  égalité si et seulement si $f$ est injective.
\remarque Si $A$ est un sous-espace
  vectoriel de $E$, alors $\dim f(A)\leq \dim A$.
  \begin{sol}
  La deuxième partie de la remarque est moins triviale. Si $E$ est de dimension finie, $E$ admet une base $(e_1,\ldots,e_n)$ et $(f(e_1),\ldots,f(e_n))$ est une famille génératrice de $\im f$. Donc $\im f$ est de dimension finie et $\rg f\leq n=\dim E$.
  \end{sol}
\remarque[utile=-3] Si $u\in\lin{E}{F}$, alors $\rg u=0$ si et seulement si $u=0$.
\end{remarques}

\begin{exoUnique}
\exo Calculer le rang de l'application de $\polyK[n]$ dans lui-même
  qui à $P$ associe $P\p{X+1}-P(X)$.
\end{exoUnique}

  \begin{sol}
  On a déjà vu que $\phi$ est à valeurs dans $\polyK[n-1]$ et que l'image de $(X,X^2,\ldots,X^n)$ formait une famille échelonnée donc était libre. Ainsi, $\im \phi\subset \polyK[n-1]$ et $\dim( \im \phi)\geq n$ donc c'est une égalité et par égalité des dimensions + inclusion, les deux espaces sont égaux. Ainsi, le rang vaut $n$.
  \end{sol}


\begin{proposition}[utile=2]
Soit $E$ et $F$ deux \Kevs de même dimension $n$ et $f$ une application
linéaire de $E$ dans $F$. Alors $f$ est un isomorphisme si et seulement si
$\rg f=n$.
\end{proposition}
\begin{preuve}
Supposons que $f$ soit un isomorphisme. Alors $f$ est surjective donc
$\im f=F$. Donc $\rg f=\dim\im f=\dim F=n$.\\
Réciproquement, supposons que $\rg f=n$. Alors puisque $\dim F=n$, on en
déduit que $\im f=F$. Donc $f$ est surjective. Puisque $E$ et $F$ ont même
dimension, on en déduit que $f$ est un isomorphisme.
\end{preuve}

\begin{proposition}
\begin{itemize}
\item Soit $E, F, G$ des $\K$-espaces vectoriels de dimension finie, $f\in\lin{E}{F}$ et
  $g\in\lin{F}{G}$. Alors
  \[\rg(g\circ f)\leq\rg(g) \quad\et\quad \rg(g\circ f)\leq\rg(f).\]
\item On ne change pas le rang d'une application linéaire si on la compose à gauche ou à
  droite par un isomorphisme.
\end{itemize}
\end{proposition}

\begin{definition}[utile=-3]
On appelle \emph{rang} d'une famille $(x_1,\ldots,x_p)\in E^p$ et on note
$\rg\p{x_1,\ldots,x_p}$ la dimension de $\vect\p{x_1,\ldots,x_p}$.
\end{definition}

\begin{remarques}
\remarque Si $(x_1,\ldots,x_p)$ est une famille d'éléments d'un \Kev $E$ de
  dimension $n$, alors son rang est inférieur ou égal à $n$ et à $p$.
\remarque 
  Si $u\in\mathcal{L}(E,F)$, alors
  $\rg(u(x_1),\ldots,u(x_p))\leq\rg(x_1,\ldots,x_p)$. De plus, si $u$ est injective, alors
  \[\rg(u(x_1),\ldots,u(x_p))=\rg(x_1,\ldots,x_p).\]
\end{remarques}

\begin{sol}
On pose $G=\vect(x_1,\ldots,x_p)$ et $H=\vect(f(x_1),\ldots,f(x_p))$ et on définit $\dspappli{\phi}{G}{H}{x}{f(x)}$. $\phi$ est bien définie et est linéaire. Grâce à l'injectivité de $f$, $\phi$ est injective également. On montre ensuite que $g$ est surjective. Donc $G$ et $H$ ont même dimension.
\end{sol}
  


\section{Calcul de dimension et de rang, hyperplan}

\subsection{Somme de deux sous-espaces vectoriels}

% \begin{proposition}
% Soit $E$ un \Kev et $A$ et $B$ deux sous-espaces vectoriels de dimensions
% finies, supplémentaires dans $E$. Si $a_1,\ldots,a_{n_A}$ est une base de $A$ et
% $b_1,\ldots,b_{n_B}$ est une base de $B$, alors $a_1,\ldots,a_{n_A},b_1,\ldots,
% b_{n_B}$ est une base de $E$. Un telle base est dite adaptée à la décomposition
% $E=A \oplus B$.
% \end{proposition}

\begin{proposition}[utile=2]
Soit $E$ un \Kev de dimension finie et $A$, $B$ deux sous-espaces
vectoriels supplémentaires de $E$. Alors
\[\dim E=\dim A+\dim B.\]
\end{proposition}

\begin{preuve}
On concatène les bases.
\end{preuve}

% \begin{remarques}
% \remarque[utile=-2] Plus généralement, on a démontré que si $E$ est un \Kev admettant deux
%   sous-espaces vectoriels $A$ et $B$ supplémentaires de dimensions finies, alors
%   $E$ est de dimension finie et $\dim E=\dim A+\dim B$.
% \end{remarques}

\begin{proposition}[nom={Formule de \nom{Grassmann}}]
Soit $E$ un \Kev de dimension finie et $A$, $B$ deux sous-espaces vectoriels
de $E$. Alors
\[\dim\p{A+B}=\dim A+\dim B-\dim\p{A\cap B}.\]
\end{proposition}

\begin{preuve}
On pose $F=A+B$.
On pose $B'$ un supplémentaire de $A\cap B$ dans $B$. Montrons que $A\oplus B'=F$.
\begin{itemize}
\item[$\bullet$] Soit $x\in A\cap B'$. Comme $B'\subset B$, $x\in A\cap B\cap B'=\set{0}$ d'après la définition de $B'$. Donc $A\cap B'=\set{0}$ (l'autre inclusion étant triviale).
\item[$\bullet$] Montrons que $A+B'=F$. On a bien $A+B'\subset F$. Maintenant, considérons $x\in F=A+B$. $x$ s'écrit $x=a+b=a+(a'+b')=(a+a')+b' \in A+B'$.
\end{itemize}
Donc $F=A\oplus B'$. Donc $\dim(A+B)=\dim F=\dim A+\dim B'$. Or, $B=(A\cap B)\oplus B'$ donc $\dim B=\dim(A\cap B)+\dim B'$ donc $$\dim\p{A+B}=\dim A+\dim B-\dim\p{A\cap B}.$$
\end{preuve}

\begin{remarqueUnique}
\remarque En particulier, $\dim(A+B)\leq\dim A+\dim B$ et cette inégalité est
  une égalité si et seulement si $A$ et $B$ sont en somme directe.
\end{remarqueUnique}
\begin{exos}
\exo Dans $\R^3$, montrer que l'intersection de deux plans vectoriels est soit un
  plan soit une droite.
  \begin{sol}
  $\dim(A+B)=4-\dim(A\cap B)$. Or $A\subset A+B\subset R^3$ donc $\dim(A+B)\in \set{2,3}$, ce qui donne un plan ou une droite pour $A\cap B$.
  \end{sol}
\end{exos}

\begin{proposition}[utile=3]
Soit $E$ un \Kev de dimension finie et $A$ et $B$ deux sous-espaces vectoriels
de $E$.
\begin{itemize}
\item Si $A\cap B=\ens{0}$ et $\dim E=\dim A+\dim B$, alors $E=A\oplus B$.
\item Si $E=A+B$ et $\dim E=\dim A+\dim B$, alors $E=A\oplus B$.
\end{itemize}
Autrement dit, si $\dim E=\dim A+\dim B$
\[\text{$A$ et $B$ sont en somme directe} \quad\ssi\quad
  E=A\oplus B \quad\ssi\quad E=A+B.\]

\end{proposition}

\begin{preuve}
Immédiat pour les deux points avec l'application de la formule de Grassmann.
\end{preuve}


\subsection{Produit d'espaces vectoriels, espace $\lin{E}{F}$}

% \begin{proposition}
% Soit $E$ et $F$ deux \Kevs de dimensions finies. Soit
% $\mathcal{B}_E=x_1,\ldots,x_{n_E}$ une base de $E$ et
% $\mathcal{B}_F=y_1,\ldots,y_{n_F}$ une base de $F$. Alors
% $\p{x_1,0_F},\ldots,\p{x_{n_E},0_F},\p{0_E,y_1},\ldots,\p{0_E,y_{n_F}}$ est une
% base
% de $E\times F$.
% \end{proposition}

\begin{proposition}[utile=-2]
Soit $E$ et $F$ deux \Kevs de dimension finie. Alors $E\times F$ est de
dimension finie et
\[\dim\p{E\times F}=\dim E+\dim F.\]
En particulier, si $n\in\N$, $E^n$ est de dimension finie et
\[\dim\p{E^n}=n\dim E.\]
\end{proposition}

\begin{preuve}
Avec $(e_1,\ldots,e_n)$ une base de $E$ et $(f_1,\ldots,f_p)$ une base de $F$, on montre que \\$((e_1,0),\ldots,(e_n,0),(0,f_1),\ldots,(0,f_p))$ est une base de $E\times F$ en montrant son caractère libre et générateur.
\end{preuve}

\begin{proposition}[utile=-3]
Soit $E$ et $F$ deux \Kevs de dimension finie. Alors $\lin{E}{F}$ est de
dimension finie et
\[\dim\lin{E}{F}=\dim E\cdot\dim F.\]
En particulier, $\Endo{E}$ est de dimension finie et
\[\dim\Endo{E}=\p{\dim E}^2.\]
\end{proposition}

\begin{preuve}
Soit $(e_1,\ldots,e_n)$ une base de $E$. On pose $$\dspappli{\phi}{\lin{E}{F}}{F^n}{f}{(f(e_1),\ldots,f(e_n)}.$$
On montre que $\phi$ est linéaire. De plus, c'est une bijection car pour tout $(y_1,\ldots,y_n)\in F^n$, il existe un unique $f\in \lin{E}{F}$ tel que $\forall k \in \intere{1}{n}, f(e_k)=y_k$, i.e $\phi(f)=(y_1,\ldots,y_n)$.
Ainsi, $$\dim(\lin{E}{F})=\dim(F^n)=n\dim(F)=\dim(E)\dim(F).$$
\end{preuve}

\begin{proposition}[utile=-1]
Si $E$ est un \Kev de dimension finie, $E^\star$ est de dimension finie et
\[\dim E^\star=\dim E.\]
\end{proposition}


% \begin{definition}
% Soit $E$ un \Kev de dimension $n$ et $\mathcal{B}=e_1,\ldots,e_n$ une base de
% $E$. Alors, pour tout $j\in\intere{1}{n}$, il existe une unique forme linéaire
% $e_j^\star$ telle que~:
% \[\forall k\in\intere{1}{n} \quad e_j^\star\p{e_k}=\delta_{j,k}=
%   \begin{cases}
%   0 &\text{si $k\neq j$}\\
%   1 &\text{si $k=j$}
%   \end{cases}\]
% La forme linéaire $e_j^\star$ associe à tout élément $x\in E$ sa $j$-ième
% coordonnée relativement à la base $\mathcal{B}$. De plus, la famille
% $e_1^\star,\ldots,e_n^\star$ est une base de $E^\star$ notée $\mathcal{B}^\star$
% et appelée base duale de la base $\mathcal{B}$.
% \end{definition}


\subsection{Théorème du rang}

% \begin{proposition}
% Soit $E$, $F$ et $G$ trois \Kevs de dimensions finies.
% \begin{itemize}
% \item Soit $g$ une application linéaire surjective de $E$ dans $F$ et
%   $f$ une application linéaire de $F$ dans $G$. Alors~:
%   \[\rg\p{f\circ g}=\rg f\]
% \item Soit $g$ une application linéaire injective de $F$ dans $G$ et $f$
%   une application linéaire de $E$ dans $F$. Alors~:
%   \[\rg\p{g \circ f}=\rg f\]
% \end{itemize}
% En particulier, on ne change pas le rang d'une application linéaire si on la
% compose à droite ou à gauche par un isomorphisme.
% \end{proposition}

% \begin{proposition}[utile=1]
% On ne change pas le rang d'une application linéaire si on la compose à droite ou à
% gauche par un isomorphisme.
% \end{proposition}


% \begin{preuve}
% Soit $f\in \lin{E}{F}$ avec $E$ et $F$ de dimension finie.
% \begin{itemize}
% \item[$\bullet$] Soit $G$ un $\Kev$ de dimension finie et $g\in \lin{F}{G}$ un isomorphisme. Montrons que $\rg(g\circ f)=\rg (f)$. On pose $$\dspappli{\phi}{\im f}{\im(g\circ f)}{y}{g(y)}.$$
% Tout d'abord, $\phi$ est bien définie : soit $y\in \im f$. Alors, il existe $x\in E$ tel que $y=f(x)$, alors $\phi(y)=g(y)=g(f(x))=(g\circ f)(x)\in \im(g\circ f).$ De plus, $\phi$ est linéaire (évident), on montre qu'elle est injective et surjective.
% Ainsi, c'est un isomorphisme, on a donc égalité entre les dimensions de l'espace de départ et d'arrivée, ce qui correspond au résultat souhaité.

% Remarque : seul l'injectivité de $g$ a servi.
% \item[$\bullet$] Pour la droite, on montre directement que $\im(f\circ u)=\im(f)$ avec $u$ un isomorphisme (on peut remarquer qu'on utilise en fait seulement le caractère surjectif).
% \end{itemize}

% \end{preuve}


% \begin{remarqueUnique}
% \remarque[utile=-1] Plus généralement, on a démontré qu'on ne changeait pas le rang d'une
%   application linéaire en la composant par la droite par une application surjective
%   ou en la composant par la gauche par une application injective.  
% \end{remarqueUnique}

\begin{theoreme}[utile=3, nom=Théorème du rang]
Soit $E$ un \Kev de dimension finie, $F$ un \Kev et $f$ une application
linéaire de $E$ dans $F$. Alors $\im f$ est de dimension finie et
\[\dim E=\dim\ker f+\dim\im f.\]
Autrement dit
\[\rg f=\dim E-\dim\ker f.\]
\end{theoreme}
\begin{preuve}
Puisque $E$ est de dimension finie, $\ker f$ admet un supplémentaire $G$ dans
$E$. Soit $\phi$ l'application~:
\[\dspappli{\phi}{G}{\im f}{x}{f(x)}\]
Alors~:
\begin{itemize}
\item $\phi$ est bien à valeurs dans $\im f$.
\item $\phi$ est {\it injective}.\\
  En effet, montrons que $\ker\phi=\ens{0}$. Soit
  $x\in\ker\phi$. Alors $x\in G$ et $\phi(x)=0$. Donc $f(x)=0$,
  c'est-à-dire $x\in\ker f$. Donc $x\in\ker f\cap G$. Or $G$ est un
  supplémentaire de $\ker f$ dans $E$ donc $x=0$.
\item $\phi$ est {\it surjective}.\\
  En effet, si on se donne $y\in\im f$, il existe
  $x\in E$ tel que $y=f(x)$. Puisque $E=\ker f+G$, il existe $a\in\ker f$ et
  $b\in G$ tels que $x=a+b$. Alors
  \[f(x)=f\p{a+b}=\underbrace{f(a)}_{=0}+f(b)=f(b)\]
  Donc $y=f(b)=\phi(b)$. Donc $y\in\im\phi$.
\end{itemize}
En conclusion $\phi$ est un {\it isomorphisme} de $G$ dans $\im f$. Puisque $G$
est de dimension finie, il en est de même pour $\im f$ et $\dim\im f=\dim G$.
Or, puisque $E=\ker f\oplus G$, on a $\dim E=\dim \ker f+\dim G$. Donc~:
\[\dim E=\dim\ker f+\dim\im f\]
Puisque par définition $\rg f=\dim\im f$, il vient~:
\[\rg f=\dim E-\dim\ker f\]
\end{preuve}

\begin{remarqueUnique}
% \application Si $u_1,\ldots,u_p$ est une famille de vecteurs de $\K^n$, alors
%   l'application~:
%   \[\dspappli{f}{\K^p}{\K^n}{\p{\lambda_1,\ldots,\lambda_p}}{\lambda_1 u_1+
%     \cdots+\lambda_p u_p}\]
%   est linéaire et $\im f=\vect\p{u_1,\ldots,u_p}$. On en déduit donc que
%   $\rg\p{u_1,\ldots,u_p}=\rg f$. Le théorème du rang nous donne~:
%   \[\rg f=\dim \K^p-\dim\ker f=p-\dim\ker f\]
%   On est donc ramené à calculer $\ker f$, c'est-à-dire résoudre un système
%   linéaire à $n$ équations et $p$ inconnues.\\
%   Par exemple, calculons le rang de la famille $u_1$,$u_2$,$u_3$,$u_4$ avec~:
%   \[u_1=\p{1,2,3} \quad u_2=\p{4,5,6} \quad u_3=\p{7,8,9} \quad
%     u_4=\p{-2,-1,0}\]
%   Soient $\lambda_1,\ldots,\lambda_4\in\R$. Alors~:
%   \begin{eqnarray*}
%   &    &\p{\lambda_1,\lambda_2,\lambda_3,\lambda_4}\in\ker f\\
%   &\ssi& \lambda_1 u_1+\lambda_2 u_2+\lambda_3 u_3+\lambda_4 u_4=0\\
%   &\ssi& \syslin{\lambda_1&+4\lambda_2&+7\lambda_3&-2\lambda_4&=&0\cr
%                  2\lambda_1&+5\lambda_2&+8\lambda_3&-\lambda_4&=&0\cr
%                  3\lambda_1&+6\lambda_2&+9\lambda_3& &=&0}\\
%   &\ssi& \syslin{\lambda_1&+4\lambda_2&+7\lambda_3&-2\lambda_4&=&0&\cr
%    &-3\lambda_2&-6\lambda_3&+3\lambda_4&=&0&\quad L_2\gets L_2-2L_1\cr
%    &-6\lambda_2&-12\lambda_3&+6\lambda_4 &=&0&\quad L_3\gets L_3-3L_1}\\
%   &\ssi& \syslin{\lambda_1&+4\lambda_2&+7\lambda_3&-2\lambda_4&=&0&\cr
%    &\lambda_2&+2\lambda_3&-\lambda_4&=&0&\quad L_2\gets L_2/\p{-3}\cr
%    &\lambda_2&+2\lambda_3&-\lambda_4 &=&0&\quad L_3\gets L_3/\p{-6}}\\
%   &\ssi& \syslin{\lambda_1&+4\lambda_2&+7\lambda_3&-2\lambda_4&=&0\cr
%                             &\lambda_2&+2\lambda_3&-\lambda_4&=&0}\\
%   &\ssi& \exists t_1,t_2\in\R \quad 
%          \syslin{\lambda_1&+4\lambda_2&+7\lambda_3&-2\lambda_4&=&0\cr
%                              &\lambda_2&+2\lambda_3&-\lambda_4&=&0\cr
%                              &         &\lambda_3&            &=&t_1\cr
%                              &         &         &\lambda_4   &=&t_2}\\
%   &\ssi& \exists t_1,t_2\in\R \quad
%          \syslin{\lambda_1&=&t_1&-2t_2 \cr
%                  \lambda_2&=&-2t_1&+t_3 \cr
%                  \lambda_3&=&t_1&    \cr
%                  \lambda_4&=&   &t_2 }\\
%   &\ssi& \exists t_1,t_2\in\R \quad \p{\lambda_1,\lambda_2,\lambda_3,\lambda_4}
%          =t_1\p{1,-2,1,0}+t_2\p{-2,1,0,1}
%   \end{eqnarray*}
%   Remarquons que les vecteurs $e_1=\p{1,-2,1,0}$ et $e_2=\p{-2,1,0,1}$ forment
%   une famille libre. En effet si $\mu_1$ et $\mu_2$ sont deux réels tels que
%   $\mu_1 e_1+\mu_2 e_2=0$, alors la troisième composante de cette égalité
%   (respectivement la quatrième) nous donne $\mu_1=0$ (respectivement
%   $\mu_2=0$). Donc $e_1,e_2$ est une base de $\ker f$, donc $\ker f$ est
%   de dimension 2. En conclusion~:
%   \[\rg\p{u_1,u_2,u_3,u_4}=\dim \R^4-\dim\ker f=4-2=2\]
%   En analysant cette méthode, il est possible de montrer que la famille
%   $e_1,\ldots,e_k$ est toujours libre (la méthode utilisée dans cet exemple se
%   généralise aisément). C'est donc une base de $\ker f$,
%   donc $\dim \ker f=k$ et $\rg\p{u_1,\ldots,u_p}=p-k$. Puisque $k$ est
%   le nombre d'inconnues $p$ moins le nombre d'équations restantes $r$ lorsque
%   le système linéaire d'origine est triangularisé, on en déduit que
%   $\rg\p{u_1,\ldots,u_p}=p-\p{p-r}=r$. On aurait donc pu s'arrêter à cet
%   instant et dire que le rang du système était $r$.\\
%   Nous reviendrons sur cette méthode lors du calcul du rang d'une matrice par
%   la méthode du pivot de Gauss.
\remarque Le théorème du rang permet de retrouver le fait que si $f$ est une
  application linéaire d'un \Kev $E$ de dimension finie dans un \Kev $F$ de
  même dimension, alors il suffit de montrer que $f$ est injective ou 
  surjective pour montrer que $f$ est un isomorphisme.
%   En effet, si $f$ est injective $\ker f=\ens{0}$ donc $\dim\ker f=0$. D'après
%   le théorème du rang, $\rg f=\dim E=\dim F$. Donc $\im f$ est de même
%   dimension que $F$ donc $\im f=F$. En conclusion $f$ est surjective. C'est
%   donc un isomorphisme.\\
%   De la même manière, si $f$ est surjective $\im f=F$ donc $\rg f=n$. D'après
%   le théorème du rang $\dim E=\dim\ker f+\rg f$ donc $n=\dim\ker f+n$
%   donc $\dim\ker f=0$. On en déduit que $\ker f=\ens{0}$ puis que $f$ est
%   injective. En conclusion, $f$ est un isomorphisme.
\end{remarqueUnique}


\begin{exos}
\exo Soit $n\in\N$ et
  \[\dspappli{\phi}{\polyR[n]}{\polyR[n]}{P}{P(X+1)-P(X)}\]
  Déterminer le noyau de $\phi$, son rang, puis son image.
\exo Soit $E$ et $F$ deux \Kevs de dimension finie, $f\in\lin{E}{F}$ et $A$
  un sous-espace vectoriel de $E$. Montrer que
  \[\dim\cro{f(A)}=\dim A-\dim\p{A\cap\ker f}\]
  \begin{sol}
  Il suffit de l'appliquer à la restriction de $f$ à $A$ au départ.
  \end{sol}
%%\exo Si $E$ est de dimension finie et $f\in\Endo{E}$, alors
%  \[\dim\p{\im f}=\dim\p{\im f^2}+\dim\p{\ker f\cap\im f}\]
%   En effet, soit $\phi$ l'application~:
%   \[\dspappli{\phi}{\im f}{\im f^2}{x}{f(x)}\]
%   Alors on montre aisément que $\phi$ est bien à valeurs dans $\im f^2$,
%   qu'elle est surjective et que $\ker\phi=\ker f\cap\im f$. En appliquant le
%   théorème du rang à $\phi$, il vient~:
%   \[\dim \im f=\dim\p{\ker f\cap\im f}+\dim\im f^2\]
\exo Soit $E$ un \Kev de dimension finie et $g\in\Endo{E}$. Calculer le rang de
  \[\dspappli{\phi}{\Endo{E}}{\Endo{E}}{f}{g\circ f}\]
  \begin{sol}
  Déjà, on montre que $\phi$ est linéaire. Il nous faut déterminer $\dim(\ker\phi))$ puis appliquer le théorème du rang. 
  $$\ker \phi=\set{f\in \Endo{E} \text{ tel que } \im(f)\subset \ker g}.$$ Pour avoir sa dimension, on définit :
  $$\dspappli{\psi}{\lin{E}{\ker g}}{\ker \phi}{f}{\dspappli{\psi(f)}{E}{E}{x}{f(x)}}.$$
  On montre que $\psi$ est bien définie, linéaire, injective, surjective. On en déduit donc $$\dim E\times\dim(\ker g)=\dim(\lin{E}{\ker g})=\dim(\ker\phi)$$ et ainsi $$\rg(\phi)=(\dim E)^2-\dim E\times\dim(\ker g)=\dim E \times \rg(g).$$
  \end{sol}
\exo Soit $E,F$ et $G$ des espaces vectoriels de dimension finie,
  $f\in\lin{F}{E}$ et $g\in\lin{G}{E}$. Montrer qu'il existe une application
  linéaire $h$ de $F$ dans $G$ telle que $f=g\circ h$ si et seulement si
  $\im f\subset\im g$.  
  \begin{sol}
  Gauche-droite ok. Droite-gauche :
  Soit $G'$ un supplémentaire de $\ker g$ dans $G$. On définit :
  $$\dspappli{\phi}{G'}{\im g}{x}{g(x)}$$ alors $\phi$ est un isomorphisme (cf. preuve précédente et remarque). On pose alors $$\dspappli{h}{F}{G}{x}{\phi^{-1}(f(x))}.$$ On montre que $h$ est bien définie, linéaire et vérifie $f=g\circ h$.
  \end{sol}
\end{exos}

\subsection{Hyperplan}

% \begin{definition}
% Soit $E$ un \Kev. On appelle \emph{hyperplan} de $E$ tout noyau d'une forme linéaire non
% nulle.
% \end{definition}

% \begin{proposition}
% Soit $E$ un \Kev.
% \begin{itemize}
% \item Si $H$ est un hyperplan de $E$ et $D$ est une droite vectorielle non contenue dans
%   $H$, alors
%   \[E=H\oplus D.\]
% \item Si $D$ est une droite vectorielle, tout supplémentaire de $D$ est un hyperplan.
% \end{itemize}
% \end{proposition}

% \begin{proposition}[utile=1]
% Soit $H$ un hyperplan de $E$ et $\phi_0$ une forme linéaire telle que $H=\ker \phi_0$.
% Alors l'ensemble des formes linéaires de $E$ dont le noyau est $H$ est
% \[\Ks\phi_0=\ensim{\lambda\phi_0}{\lambda\in\K^*}.\]  
% \end{proposition}

% \begin{preuve}
% Soit $E$ un \Kev de dimension finie, $H$ un hyperplan de $E$ et $\phi_0$
% une forme linéaire telle que $H=\ker \phi_0$.

% Montrons que $H=\ker\phi\Longleftrightarrow \exists \lambda\in \Ks \text{ tel que } \phi=\lambda\phi_0$.
% Le sens droite-gauche est aisée. Réciproquement, soit $\phi\in E^\star$ telle que $H=\ker\phi$. Considérons $(e_1,\ldots,e_{n-1})$ une base de $H$, qu'on complète avec $e_n$ en une base de $E$. Comme $e_n\notin H$, $\phi_0(e_n)\neq 0$. On pose $\displaystyle \lambda=\frac{\phi(e_n)}{\phi_0(e_n)}$. On montre que $\phi=\lambda\phi_0$ en montrant qu'elles coïncident sur une base (la base $(e_1,\ldots,e_n)$).
% \end{preuve}


\begin{proposition}[utile=1]
Soit $E$ un \Kev de dimension $n$. Un sous-espace vectoriel $H$ de $E$ est
un hyperplan si et seulement si $\dim H=n-1$.
\end{proposition}

\begin{preuve}
Pour montrer l'équivalence entre les deux définitions. La remontée vient du théorème du rang.
Pour la descente, on prend un supplémentaire de $H$ qui est donc de dimension $1$ et on définit la forme linéaire qui à $x$ associe la coordonnée sur une base de ce supplémentaire. On vérifie que la fonction ainsi définie est bien linéaire et que son noyau est bien $H$.
\end{preuve}


\begin{exos}
\exo Dans $\R^3$, donner la dimension du sous-espace vectoriel d'équation $3x+2y-z=0$.
\exo Dans $\polyR[n]$, donner la dimension du sous-espace vectoriel
  \[H\defeq\enstq{P\in\polyR[n]}{\integ{0}{1}{P(t)}{t}=0}.\]
\exo Soit $E$ un \Kev de dimension finie et $\phi_1$, $\phi_2$ deux formes
  linéaires sur $E$. Calculer le rang de
  \[\dspappli{f}{E}{\K^2}{x}{\p{\phi_1(x),\phi_2(x)}}\]
  en fonction de $\phi_1$ et $\phi_2$.
\exo Quelle est la dimension de l'intersection de deux hyperplans~?
\begin{sol}
Soient $H_1$ et $H_2$ deux hyperplans de $E$. On a $\dim(H_1\cap H_2)=\dim H_1+ \dim H_2-\dim(H_1+H_2)\geq 2n-2-n=n-2$. De plus, comme $H_1\cap H_2 \subset H_1$, on a $\dim(H_1\cap H_2)\leq n-1$ donc la dimension vaut $n-1$ ou $n-2$. Supposons que ce soit $n-1$, alors par inclusion et égalité des dimensions, on a $H_1\cap H_2=H_1=H_2$. Réciproquement, si $H_1=H_2$, on a bien la dimension de l'intersection qui vaut $n-1$. Et on a alors la dimension qui vaut $n-2$ si et seulement s'ils sont distincts.
\end{sol}
\end{exos}

\begin{proposition}[utile=1]
Soit $E$ un \Kev de dimension $n$ et $\mathcal{B}\defeq(e_1,\ldots,e_n)$ une base de
$E$. 
\begin{itemize}
\item Si $a_1,\ldots,a_n\in\K$ ne sont pas tous nuls, l'ensemble
  $H$ d'équation
  \[a_1 e_1^\star(x)+\cdots+a_n e_n^\star(x)=0\]
  est un hyperplan de $E$.
\item Réciproquement, si $H$ est un hyperplan de $E$, il existe
  $a_1,\ldots,a_n\in\K$ non tous nuls tels que
  \[a_1 e_1^\star(x)+\cdots+a_n e_n^\star(x)=0\]
  est une équation de $H$. De plus, si $b_1,\ldots,b_n\in\K$, alors
  \[b_1 e_1^\star(x)+\cdots+b_n e_n^\star(x)=0\]
  est une équation de $H$ si et seulement si il existe $\alpha\in\K^*$ tel que
  \[\forall k\in\intere{1}{n} \qsep b_k=\alpha a_k.\]
\end{itemize}
\end{proposition}

\begin{preuve}
En considérant $(e_1,\ldots,e_n)$ une base de $E$. on définit $$\dspappli{\phi}{E}{\K}{x=\sum_{i=1}^nx_k e_k}{\sum_{i=1}^n\lambda_k x_k}.$$
On montre que $\phi$ est linéaire et non nulle (pour la non nullité, comme les $\lambda_i$ sont non tous nuls, on a $\phi(e_{i_0})=\lambda_{i_0}\neq 0$ pour un certain $i_0$).


Réciproquement, soit $H$ un hyperplan de $E$. Alors $H$ est le noyau d'une forme linéaire non nulle $\phi$. En posant $\lambda_k=\phi(e_k)$, ils sont non tous nuls (sinon, $\phi=0$). Et si $x=\displaystyle \sum_{k=1}^nx_ke_k$, on a $x\in H\Longleftrightarrow \phi(x)=0\Longleftrightarrow \displaystyle \sum_{k=1}^n\lambda_k x_k=0$, ce qui fournit bien une équation de $H$.


Soit $\mu_1,\ldots, mu_n \in \K$. On définit $$\dspappli{\psi}{E}{\K}{x=\sum_{i=1}^nx_k e_k}{\sum_{i=1}^n\mu_k x_k}.$$ L'ensemble d'équation $\sum_{i=1}^n\mu_k x_k=0$ est $\ker \psi$.
Et \begin{eqnarray*}
H=\ker \psi &\Longleftrightarrow & \exists \alpha\in \Ks, \psi=\lambda \phi\\
&\Longleftrightarrow & \exists \alpha\in \Ks, \forall k \in \intere{1}{n}, \psi(e_k)=\alpha\phi(e_k)\\
&\Longleftrightarrow & \exists \alpha\in \Ks, \forall k \in \intere{1}{n}, u_k=\alpha\lambda_k.
\end{eqnarray*}
\end{preuve}

\begin{proposition}[utile=1]
Soit $E$ un \Kev de dimension $n$.
\begin{itemize}
\item L'intersection de $p$ hyperplans est de dimension supérieure ou égale à $n-p$.
\item Si $p\in \intere{0}{n}$, tout sous-espace vectoriel de dimension $n-p$ est l'intersection de $p$ hyperplans.
\end{itemize}
\end{proposition}

\begin{preuve}
\begin{itemize}
\item Montrons par récurrence sur $p$ que l'intersection de $p$ hyperplans est de dimension supérieure ou égale à $n-p$.
La propriété est vraie pour $p=1$. Supposons la vraie pour $p$. D'après la formule de Grassmann $\dim(H_1\cap\ldots\cap H_{p+1})=\dim(H_1\cap\ldots\cap H_{p})+\dim H_{p+1}-\dim(H_1\cap\ldots\cap H_{p}+H_{p+1})\geq \dim(H_1\cap\ldots\cap H_{p})+n-1-n\underbrace{\geq}_{HR}n-p-1=n-(p+1).$
\item Montrons par récurrence sur $p\in \intere{1}{n}$ que tout sous-espace vectoriel de dimension $n-p$ est l'intersection de $p$ hyperplans. Vrai pour $p=1$. Supposons la propriété vraie pour $p\in \intere{1}{n-1}$ et montrons-là pour $p+1$. Soit $F$ un sous-espace vectoriel de dimension $n-(p+1)$. $F\neq E$ donc il existe $a\in E\setminus F$. Posons $D=\vect(a)$ et $G=F+D$. $a$ n'appartenant pas à $F$, $F$ et $D$ sont en somme directe et $G$ est donc de dimension $n-p$. $G$ est donc d'après l'hypothèse de récurrence l'intersection de $p$ hyperplans $H_1,\ldots, H_p$. Soit $G'$ un supplémentaire de $G$ dans $E$. On a $\dim(G)=p$. Posons $H_{p+1}=F+G'$. Cette somme est directe car si $x\in F\cap G'$, $x\in G\cap G'=\set{0}$. Ainsi, $\dim(H_{p+1})=\dim F+\dim G'=n-(p+1)+p=n-1$, donc $H_{p+1}$ est un hyperplan. Montrons pour finir que $F=G\cap H_{p+1}$, ainsi, nous aurons $F=H_1\cap\ldots\cap H_p\cap H_{p+1}$, ce qui correspond au résultat souhaité.

$F$ est contenu dans $G$ et dans $H_{p+1}$ par construction de ces deux sous-espaces. Ainsi, $F\subset G\cap H_{p+1}$. Réciproquement, soit $x\in G\cap H_{p+1}$. Comme $x\in H_{p+1}=F+G'$, il s'écrit $x=x_1+x_2$ avec $(x_1,x_2)\in F\times G'$. $x$ et $x_1$ sont deux éléments de $G$ donc $x-x_1=x_2\in G$. Mais $x_2\in G'$. Donc $x_2\in G\cap G'=\set{0}$, d'où $x_2=0$ et finalement, $x=x_1\in F$, ce qui achève de démontrer $G\cap H_{p+1}\subset F$.
\end{itemize}

\end{preuve}

\section{Sous-espace affine}


Soit $E$ un \Kev. Nous avons vu que les éléments de $E$ sont naturellement des vecteurs.
Dans cette partie, nous allons voir les éléments de $E$ comme des points. Afin de marquer
cette différence de point de vue, nous utiliserons dans cette partie les lettres
$a$, $b$, $c$, $d$ du début de l'alphabet pour désigner les éléments de $E$ que nous
considérerons comme des points et les lettres $x$, $y$, $z$ de la fin de l'alphabet pour
désigner les éléments de $E$ que nous considérerons comme des vecteurs.
\vspace{2ex}
\begin{definition}
Soit $a,b\in E$. On définit le vecteur $\ve{ab}\in E$ par
\[\ve{ab}\defeq b-a.\]
\end{definition}

\begin{proposition}
\begin{itemize}
\item Quel que soit $a\in E$, l'application
  \[\dspappli{\phi}{E}{E}{b}{\ve{ab}}\]
  est bijective.
\item Quels que soient $a, b, c\in E$
  \[\ve{ab}+\ve{bc}=\ve{ac}.\]
\end{itemize}
\end{proposition}

\begin{remarqueUnique}
\remarque Soit $a\in E$ et $x\in E$. Alors $b\defeq a+x$ est l'unique élément de $E$ tel
  que $\ve{ab}=x$.
\end{remarqueUnique}

\begin{definition}
Soit $x\in E$. On appelle translation de vecteur $x$ l'application
\[\dspappli{\tau_x}{E}{E}{a}{a+x.}\]
\end{definition}

\begin{remarques}
\remarque Si $x,y\in E$, alors $\tau_x \circ \tau_y=\tau_{x+y}$. 
\remarque Si $x\in E$, l'application $\tau_x$ est linéaire si et seulement si $x=0$.
  Dans ce cas, $\tau_x=\id$.
\end{remarques}



\begin{definition}
On dit qu'une partie $\mathcal{F}$ de $E$ est un \emph{sous-espace affine} de $E$
lorsqu'il existe $a\in E$ et un sous-espace vectoriel $F$ de $E$ tel que
\[\mathcal{F}=\ensim{a+x}{x\in F}.\]
On écrit alors $\mathcal{F}=a+F$. De plus $a\in\mathcal{F}$ et
\[F=\ensim{\ve{bc}}{b, c\in\mathcal{F}}.\]
On dit que $F$ est la \emph{direction} de $\mathcal{F}$.
\end{definition}

\begin{remarques}
\remarque En particulier, un sous-espace affine est non vide.
\remarque Si $\mathcal{F}$ est un sous-espace affine de direction $F$, alors, quel que
  soit $a\in\mathcal{F}$, $\mathcal{F}=a+F$.
\remarque Un sous-espace vectoriel est un sous-espace affine. Réciproquement, un sous-espace
  affine est un sous-espace vectoriel si et seulement si il contient 0.
\end{remarques}

\begin{exos}
\exo Déterminer l'ensemble des solutions dans $\R^3$ de
  \[\syslin{x&+2y&+3z&=1&\cr
           2x&+3y&+4z&=1&\cr
           3x&+4y&+5z&=1&}\]
  et l'exprimer comme un sous-espace affine de $\R^3$.
\begin{sol}
$$\mathcal{S}=SP+\vect((1,-2,1)).$$ C'est une droite affine de $\R^3$.
\end{sol}
\exo Résoudre sur $\R$ l'équation différentielle suivante et exprimer son ensemble de
  solutions comme un sous-espace affine de $\mathcal{D}(\R,\R)$.
  \[y'+2y=x^2-2x+3.\]

\begin{sol}
$y(x)=ce^{-2x}+(2x^2-6x+9)/4$
\end{sol}
\end{exos}

\begin{proposition}
Soit $u\in\lin{E}{F}$ et $b\in F$. Alors l'ensemble
\[\mathcal{S}\defeq\enstq{a\in E}{u(a)=b}\]
de solutions de l'équation $u(a)=b$ est
\begin{itemize}
\item Soit vide. 
\item Soit un sous-espace affine de $E$ de direction $\ker u$.
\end{itemize}
\end{proposition}

\begin{remarqueUnique}
\remarque En particulier, si $a_0\in E$ est une solution de l'équation $u(a)=b$, l'ensemble
  des solutions de cette équation est
  \[a_0 + \ker u = \ensim{a_0+x}{x\in\ker u}.\]
\end{remarqueUnique}

\begin{proposition}
Soit $(\mathcal{F}_i)_{i\in I}$ une famille de sous-espaces affines de $E$ et
$(F_i)_{i\in I}$ la famille de leurs directions respectives. On pose
\[\mathcal{F}\defeq\bigcap_{i\in I} \mathcal{F}_i \quad\et\quad F\defeq\bigcap_{i\in I} F_i.\]
Alors
\begin{itemize}
\item Soit $\mathcal{F}$ est vide. Dans ce cas, ce n'est pas un sous-espace affine.
\item Soit $\mathcal{F}$ est non vide. Dans ce cas, c'est un sous-espace affine de
  $E$ de direction $F$.
\end{itemize}
\end{proposition}

% \begin{exoUnique}
% \exo Soit $\mathcal{A}$ et $\mathcal{B}$ deus sous-espaces affines de $E$ de directions
%   respectives $A$ et $B$. On suppose que $A$ et $B$ sont supplémentaires dans $E$. Montrer
%   que $\mathcal{A}\cap\mathcal{B}$ contient excatement un point.
% \end{exoUnique}


  
% \begin{definition}
% Soit $\mathcal{F}_1$ et $\mathcal{F}_2$ deux sous-espaces affines de
% $E$ de directions respectives $F_1$ et $F_2$. On dit que
% \begin{itemize}
% \item $\mathcal{F}_1$ est parallèle à $\mathcal{F}_2$ lorsque
%   $F_1\subset F_2$.
% \item $\mathcal{F}_1$ et $\mathcal{F}_2$ sont parallèles lorsque
%   $F_1=F_2$.
% \end{itemize}
% \end{definition}

% \begin{remarqueUnique}
% \item Une droite peut être parallèle à un plan, mais pas l'inverse.
% \end{remarqueUnique}

% \begin{proposition}
% Soit $\mathcal{F}_1$ et $\mathcal{F}_2$ deux sous-espaces affines de
% $\mathcal{E}$.
% \begin{itemize}
% \item Si $\mathcal{F}_1$ est parallèle à $\mathcal{F}_2$, alors
%   $\mathcal{F}_1\cap\mathcal{F}_2=\emptyset$ ou
%   $\mathcal{F}_1\subset\mathcal{F}_2$.
% \item Si $\mathcal{F}_1$ et $\mathcal{F}_2$ sont parallèles
%   $\mathcal{F}_1\cap\mathcal{F}_2=\emptyset$ ou 
%   $\mathcal{F}_1=\mathcal{F}_2$.
% \end{itemize}
% \end{proposition}

% \begin{preuve}
% \begin{itemize}
% \item Si $\mathcal{F}_1$ est parallèle à $\mathcal{F}_2$ alors ou bien 
%   $\mathcal{F}_1\cap\mathcal{F}_2=\emptyset$ ou bien il existe $A \in \mathcal{F}_1\cap\mathcal{F}_2$ auquel cas $\mathcal{F}_1=A+\ve{\mathcal{F}_1} \subset A+\ve{\mathcal{F}_2}=\mathcal{F}_2$.
% \item D'après le point 1 et la définition d'être tous deux parallèles.
% \end{itemize}
% \end{preuve}

% \begin{proposition}
% Soit $\mathcal{F}_1$ et $\mathcal{F}_2$ deux sous-espaces affines de
% $\mathcal{E}$. Alors soit $\mathcal{F}_1\cap\mathcal{F}_2=\emptyset$, soit 
% $\mathcal{F}_1\cap\mathcal{F}_2$ est un sous-espace affine de direction
% $\ve{\mathcal{F}_1}\cap \ve{\mathcal{F}_2}$.
% \end{proposition}

% \begin{preuve}
% Supposons $\mathcal{F}_1\cap\mathcal{F}_2\neq \emptyset$. On peut alors fixer $A$ dans l'intersection. Montrons que $$\mathcal{F}_1\cap\mathcal{F}_2=A+\ve{\mathcal{F}_1}\cap\ve{\mathcal{F}_2}.$$
% On raisonne par double inclusion :
% \begin{itemize}
% \item[$\bullet$] Soit $B\in \mathcal{F}_1\cap\mathcal{F}_2$, comme il est dans chaque sous espace-affine, on peut écrire $$B=A+\ve{u} \text{ avec } \ve{u} \in \ve{\mathcal{F}_1} \et B=A+\ve{v} \text{ avec } \ve{v} \in \ve{\mathcal{F}_2}.$$
% Donc $\ve{u}=\ve{v}$ et $B$ s'écrit donc comme $A$+ un élément de $\ve{\mathcal{F}_1}\cap\ve{\mathcal{F}_2}$.
% \item[$\bullet$] La réciproque se voit encore plus clairement à la main.
% \end{itemize}
% \end{preuve}


% \begin{proposition}
% Soit $\mathcal{F}$ et $\mathcal{G}$ deux sous-espaces affines de $E$ de directions
% respectives $F$ et $G$. Si $E=F\oplus G$, alors $\mathcal{F}\cap\mathcal{G}$ est réduit à
% un unique point.
% \end{proposition}

% \begin{preuve}
% Commençons par montrer que $\mathcal{F}\cap\mathcal{G}$ est non vide. On considère $A\in \mathcal{F}$ et $B\in \mathcal{G}$. $\ve{AB}$ en tant qu'élément de $E$ s'écrit $\ve{u}+\ve{v}$ avec $(\ve{u},\ve{v})\in \ve{\mathcal{F}}\times\ve{\mathcal{G}}$.\\
% Posons alors $C=A+\ve{u} \in \mathcal{F}$. On a aussi $C=B-\ve{AB}+\ve{u}=B-\ve{v} \in \ve{\mathcal{G}}$. Ainsi, $\mathcal{F}\cap\mathcal{G}\neq \emptyset$.\\

% Mais alors, d'après la proposition précédente, $\mathcal{F}\cap\mathcal{G}$ est donc un sous-espace affine de $E$, de direction $\ve{\mathcal{F}}\cap\ve{\mathcal{G}}={\ve{0}}$ donc il s'agit d'un singleton.
% \end{preuve}
%END_BOOK

\end{document}


\section{Exercices}
\setcounter{numeroexercice}{1}
\documentclass{magnolia}

\magtex{tex_driver={pdftex}}
\magfiche{document_nom={Exercices sur les espaces vectoriels de dimension finie},
          auteur_nom={François Fayard},
          auteur_mail={fayard.prof@gmail.com}}
\magexos{exos_matiere={maths},
         exos_niveau={mpsi},
         exos_chapitre_numero={13},
         exos_theme={Espaces vectoriels de dimension finie}}
\magmisenpage{}
\maglieudiff{}
\magprocess

\begin{document}
%BEGIN_BOOK
\magsection{Famille libre, famille génératrice, base}
\magsubsection{Famille libre}

\exercice{nom={Familles de fonctions}}
\begin{questions}
\question
Soit $f_1,f_2,f_3,f_4\in\mathcal{F}(\R,\R)$ les fonctions définies par
\[\forall x\in\R\qsep f_1(x)\defeq\sin x\qsep f_2(x)\defeq \cos x\qsep f_3(x)\defeq x\sin x\qsep f_4(x)\defeq x\cos x.\]
Montrer que la famille $(f_1,f_2,f_3,f_4)$ est libre.
\question Soit $f_1,f_2,f_3\in\mathcal{F}(\R,\R)$ les fonctions définies par
\[\forall x\in\R\qsep f_1(x)\defeq \e^x\qsep f_2(x)\defeq \e^{2x}\qsep f_3(x)\defeq \e^{x^2}.\]
Montrer que la famille $(f_1,f_2,f_3)$ est libre.
\end{questions}


\exercice{nom={Familles de fonctions}}
Montrer que les familles de fonctions suivantes sont libres.
\begin{questions}
% \question La famille $f_0,f_1,\ldots,f_n$ définie par~:
%   \[\forall k\in\intere{0}{n} \quad \forall x\in \R \quad f_k(x)=x^k\]
\question La famille de fonctions $(f_0,\ldots,f_n)$ définies sur $\R$ par
  \[\forall k\in\intere{0}{n} \qsep \forall x\in\R \qsep f_k(x)\defeq x^k.\]
% \question $f_n(x)=x^n$, pour $A=\RPs$ et $n\in\Z$.
\question La famille $(f_{1},\ldots,f_{n})$
  où $\alpha_1,\ldots,\alpha_n\in\RP$ sont deux à deux distincts et
  \[\forall k\in\intere{1}{n} \qsep \forall x\in\R \qsep
    f_k(x)\defeq \cos\p{\alpha_k x}.\]
\question La famille de fonctions $(f_0,\ldots,f_n)$ définies par
  \[\forall k\in\intere{0}{n} \qsep \forall x\in\R \qsep f_k(x)\defeq \sin^k x.\]
% \question $f_\alpha(x)=e^{\alpha x}$, pour $A=\R$ et $\alpha\in\R$.
\question La famille $(f_1,\ldots,f_n)$ où
  $\alpha_1,\ldots,\alpha_n\in\R$ sont deux à deux distincts et
  \[\forall k\in\intere{1}{n} \qsep \forall x\in\R \qsep
    f_{k}(x)\defeq \e^{\ii \alpha_k x}.\]
% \question $f_\alpha(x)=e^{\alpha x}$, pour $A=[1,2]$ et $\alpha\in\C$.
\question La famille $(f_{1},\ldots,f_{n})$ où
  $\alpha_1,\ldots,\alpha_n\in\R$ sont deux à deux distincts et
  \[\forall k\in\intere{1}{n} \qsep \forall x\in\R \qsep
    f_{k}(x)\defeq\abs{x-\alpha_k}.\]
\end{questions}
\begin{sol}
$\quad$
\begin{questions}
\question   
\question dériver deux fois.
\question Introduire le polynôme
  \[P=\sum_{k=0}^n X^k\]
  Alors $\forall x\in\R \quad P(\sin x)=0$ donc $P$ admet une infinité de
  racines donc $P=0$, donc tous les $\lambda_k$ sont nuls.
\question dériver.
\question Remarquer que $f_{\alpha_k}$ n'est pas dérivable en $\alpha_k$.
\end{questions}
\end{sol}


\exercice{nom={Modification d'une famille libre}}
Soit $E$ un \Kev et $(x_1,\ldots,x_n)$ une famille libre de $n$ vecteurs. On se
donne $n$ scalaires $\lambda_1,\ldots,\lambda_n\in\K$ et on pose
\[y\defeq\sum_{k=1}^n \lambda_k x_k.\]
Pour tout $i\in\intere{1}{n}$, on pose $y_i\defeq x_i+y$ et l'on souhaite montrer que
$(y_1,\ldots,y_n)$ est libre si et seulement si
\[\sum_{k=1}^n \lambda_k \neq -1.\]
\begin{questions}
\question Soit $\mu_1,\ldots,\mu_n\in\K$. Montrer que
  \[\sum_{i=1}^n \mu_i y_i =\sum_{i=1}^n \p{\mu_i + \lambda_i \sum_{k=1}^n \mu_k}x_i.\]
\question Dans le cas où $\sum_{k=1}^n \lambda_k= -1$, déterminer une relation de liaison
  pour la famille $(y_1,\ldots,y_n)$.
\question Montrer que si $(y_1,\ldots,y_n)$ est libre, alors
  $\sum_{k=1}^n \lambda_k\neq -1$, puis conclure.

% On suppose que $\sum_{k=1}^n \lambda_k \neq -1$ et on se donne
%   $\mu_1,\ldots,m_u\in\K$ tels que $\mu_1 y_1+\cdots+\mu_n y_n
%   ,
%   alors la famille $(y_1,\ldots,y_n)$ est libre.
% \question Réciproquement, montrer que si la famille $(y_1,\ldots,y_n)$ est libre, alors
%  .
\end{questions}
% Donner une condition nécessaire
% et suffisante sur les $\lambda_k$ pour que la famille $y_1,\ldots,y_n$ soit libre.
\begin{sol}
La condition nécessaire et suffisante est
\[\sum_{k=1}^n \lambda_k\neq -1\]
On voit facilement que si cette somme est différente de $-1$, la famille est
libre (il suffit d'écrire les équations obtenues en utilisant le fait que la
famille des $x_k$ est libre et de sommer toutes les équations). Si la somme
est égale à $-1$
\[\sum_{k=1}^n \lambda_k y_k=0\]
est une relation de liaison.
\end{sol}



% \exercice{nom={Espace des fonctions trigonométriques}}
% Soit $E$ le $\R$-espace vectoriel des fonctions de $\R$ dans $\R$. Montrer que
% l'espace vectoriel engendré par les fonctions $x\rightarrow \cos^n(x)$ pour
% $n\in\N$ et celui engendré par les fonctions $x\rightarrow \cos(nx)$ pour
% $n\in\N$ coïncident.



\magsubsection{Famille Génératrice}

\exercice{nom={Exercice}}
Soit $E$ un \Kev et $u_1,\ldots,u_n\in E$. Pour tout $k\in\intere{1}{n}$, on pose
\[v_k\defeq u_1+\cdots+u_k.\]
\begin{questions}
\question Montrer que la famille $(u_1,\ldots,u_n)$ est libre si et seulement si
  la famille $(v_1,\ldots,v_n)$ l'est.
\question Montrer que la famille $(u_1,\ldots,u_n)$ engendre $E$ si et seulement si
  la famille $(v_1,\ldots,v_n)$ engendre $E$.
\end{questions}



\magsubsection{Base}

\magsubsection{Cas des familles infinies}

\exercice{nom={Polynômes}}
Pour tout $n\in\N$, on pose
\[P_n\defeq\prod_{k=0}^{n-1} (X-k).\]
Montrer que la famille $(P_n)_{n\in\N}$ est libre.

\exercice{nom={Famille des fonction polynôme-exponentielle}}
Pour tout entier naturel $n$ et tout nombre complexe $\alpha$, on définit la
fonction $f_{n,\alpha}$ sur $\R$ par
\[\forall x \in \R \qsep f_{n,\alpha}(x)\defeq x^n \e^{\alpha x}.\]
Le but de cet exercice est de montrer que la famille
$(f_{n,\alpha})_{(n,\alpha) \in\N\times\C}$ est libre.
\begin{questions}
\question Montrer que la famille $(f_{n,0})_{n\in\N}$ est libre.
\question Pour $\alpha\in\Cs$, résoudre l'équation différentielle~:
  $\forall x \in \R \qsep y'(x)+\alpha y(x)=0.$
\question Montrer par récurrence sur un majorant de
  $\deg{P_1}+\dots+\deg{P_n}$, que si $\alpha_1,\ldots,\alpha_n$ sont $n$
  complexes deux à deux distincts et si
  \[\forall x\in\R \qsep P_1(x)\e^{\alpha_1 x}+\ldots+P_n(x)\e^{\alpha_n x}=0\]
  alors $P_1=P_2=\ldots=P_n=0$.
\question Conclure.
\end{questions}

\exercice{nom={$\R$ comme $\Q$-ev}}
On considère $\R$ comme $\Q$-espace vectoriel et on note $\mathcal{P}$ l'ensemble des
nombres premiers. Montrer que la famille $(\ln p)_{p\in\mathcal{P}}$ est libre.



\magsection{Dimension}

\magsubsection{Espace vectoriel de dimension finie}


\exercice{nom={Centre de $\Endo{E}$}}
Soit $E$ un \Kev de dimension finie. Le but de cet exercice est de montrer que
le centre de $\Endo{E}$, c'est-à-dire l'ensemble des endomorphismes qui
commutent avec tous les endomorphismes est l'ensemble des homothéties.
\begin{questions}
\question Montrer que les homothéties sont dans le centre de $\Endo{E}$.
\question Soit $u$ un endomorphisme de $E$ tel que
  \[\forall x\in E \qsep \exists \lambda\in\K \qsep  u(x)=\lambda x.\]
  Montrer que $u$ est une homothétie.
\question Soit $u$ un endomorphisme de $E$ qui commute avec toutes les
  applications linéaires.
  \begin{questions}
  \question Soit $x\in E$. En considérant une symétrie par rapport à $\K x$,
    montrer qu'il existe $\lambda\in\K$ tel que $u(x)=\lambda x$.
  \question Conclure.
  \end{questions}
\end{questions}

\magsubsection{Dimension d'un espace vectoriel}
\magsubsection{Existence et unicité en dimension finie}

\exercice{nom={Changement de base}}
Dans $\R^3$, on considère les trois vecteurs $u\defeq(1,1,-1)$, $v\defeq(-1,1,1)$ et
$w\defeq(1,-1,1)$.
\begin{questions}
\question Montrer que $(u,v,w)$ forme une base de $\R^3$.
\question Donner les coordonnées de $(2,1,3)$ dans cette base.
\end{questions}
\begin{sol}
$\quad$
\begin{questions}
\question Il suffit de montrer que la famille est libre.
\question On trouve pour coordonnées $(3/2,2,5/2)$.
\end{questions}
\end{sol}

\exercice{nom={Base}}
\begin{questions}
\question Soit $\alpha\in\R$. Donner une condition nécessaire
  et suffisante sur $\alpha$ pour que la famille
  \[((2,0,\alpha),(2,\alpha,2),(\alpha,0,2))\]
  soit une base de $\R^3$.
\question La famille
  \[((1, 0, 2, 1), (0, 1, 1, 2), (2, 0, 1, 1), (2, 1, 0, 1))\]
  est-elle une base de $\R^4$~?
\end{questions}



\exercice{nom={Le Laplacien discret}}
Soit $\phi$ l'application de $\polyC$ dans $\polyC$ définie par
\[\forall P\in\polyC \qsep \phi(P)\defeq P\p{X+1}+P\p{X-1}-2P(X).\]
\begin{questions}
\question Montrer que $\phi$ est linéaire.
\question Calculer $\deg\cro{\phi(P)}$ en fonction de $\deg P$.
\question Quel est le noyau de $\phi$~?
\question Montrer que $\phi$ est surjective.
\end{questions}
\begin{sol}
$\quad$
\begin{questions}
\question
\question
\question Soit $Q\in\polyK$ et $n\in\N$ tel que $\deg Q\leq n$. On
  introduit~:
  \[\dspappli{\psi}{\polyK[n+2]}{\polyK[n]}{P}{\phi(P)}\]
  et on montre que $\psi$ est surjective. Pour cela, on montre que
  $(\psi(X^2),\ldots,\psi(X^{n+2}))$ est libre. Il suffit de remarquer
  que si
  \[\sum_{k=2}^{n+2} \lambda_k \psi(X^k)=0\]
  alors
  \[\sum_{k=2}^{n+2} \lambda_k X^k\in\ker\phi\]
  donc tous les $\lambda_k$ sont nuls.
\end{questions}
\end{sol}

\exercice{nom={Formes linéaires et trace sur $\mat{n}{\K}$}}
\begin{questions}
\question On définit l'application $\psi$ de $\mat{n}{\K}$ dans $\mat{n}{\K}^\star$,
  qui à la matrice $A\in\mat{n}{\K}$ associe l'application $\phi_A$ de
  $\mat{n}{\K}$ dans $\K$ définie par
  \[\forall X\in\mat{n}{\K} \qsep \phi_A(X)\defeq\tr(AX).\]
  Montrer que $\psi$ est bien à valeurs dans $\mat{n}{\K}^\star$ puis qu'elle est
  linéaire.
\question Montrer que $\psi$ est un isomorphisme. En déduire que si $\phi$ est
  une forme linéaire sur $\mat{n}{\K}$, il existe une et une seule matrice
  $A\in\mat{n}{\K}$ telle que
  \[\forall X\in\mat{n}{\K} \qsep \phi(X)=\tr(AX).\] 
\question Montrer que si $\phi$ est une forme linéaire sur $\mat{n}{\K}$, alors
\[\forall X,Y\in\mat{n}{\K} \qsep \phi(XY)=\phi(YX)\]
si et seulement si il existe $\lambda\in\K$ tel que
\[\forall X\in\mat{n}{\K}\qsep \phi(X)=\lambda \tr(X).\]
\end{questions}

\exercice{nom={Matrice triangulaire supérieure par blocs}}
Soit $n\in\N$, $p_1,p_2\in\N$ tels que $n=p_1+p_2$ et $A\in\mat{n}{\K}$ une matrice
triangulaire supérieure par blocs
\[A=
  \begin{pmatrix}
  A_{1,1} & A_{1,2}\\
  0 & A_{2,2}
\end{pmatrix} \qquad \text{où $A_{i,j}\in\mat{p_i,p_j}{\K}$.}\]
Montrer que $A$ est inversible si et seulement si $A_{1,1}$ et $A_{2,2}$ le
sont. Si tel est le cas, montrer que
\[A^{-1}=
  \begin{pmatrix}
  A_{1,1}^{-1} & \star \\
  0 & A_{2,2}^{-1}
\end{pmatrix}\]


\magsubsection{Dimension d'un sous-espace vectoriel}

\magsubsection{Notion de rang}

\exercice{nom={Rang et composition}}
Montrer qu'on ne change pas le rang d'une application linéaire en la composant par
la droite par une application surjective ou par la gauche par une application injective. 

\magsection{Calcul de dimension et de rang, hyperplan}
\magsubsection{Somme de deux sous-espaces vectoriels}

\exercice{nom={Dimension}}
Soit $A$ et $B$ deux sous-espaces vectoriels de dimension $3$ de $\R^5$.
Montrer que $A\cap B\not=\{0\}$.
\begin{sol}
Utiliser Grassman.
\end{sol}

\exercice{nom={Supplémentaire}}
Soit $\lambda\in\R$. À quelle condition nécessaire et suffisante
sur $\lambda$ les sous-espaces vectoriels
\[\vect((\lambda,\lambda,1))\quad\text{et}\quad
  \vect((1,\lambda,1),(2,1,1))\]
sont-ils supplémentaires dans $\R^3$~?

\exercice{nom={Supplémentaire}}
Soit $x_1,\ldots,x_n\in\R$ deux à deux distincts. On pose
\[F\defeq\enstq{f\in\mathcal{F}(\R,\R)}{\forall k\in\intere{1}{n}\qsep f(x_k)=0}.\]
Montrer que $F$ est un sous-espace vectoriel de $\mathcal{F}(\R,\R)$
et en déterminer un supplémentaire dans $\mathcal{F}(\R,\R)$. \emph{On pourra utiliser
les polynômes de Lagrange}.

\exercice{nom={Supplémentaire commun}}
Soit $E$ un \Kev de dimension finie $n$, et $A$ et $B$ deux sous-espaces
vectoriels de $E$ de même dimension $r$. Le but de cet exercice est de montrer
que $A$ et $B$ ont un supplémentaire commun dans $E$.
\begin{questions}
\question Montrer le résultat lorsque $n=r+1$.
\question Plus généralement, montrer que si le résultat est vrai si
  $\dim A=\dim B=r+1$, alors il est vrai si $\dim A=\dim B=r$.
\question Conclure.
\end{questions}

\magsubsection{Produit de deux sous-espaces vectoriels, $\mathcal{L}(E,F)$}
\magsubsection{Théorème du rang}



\exercice{nom={Noyau et image en somme directe}}
Soit $E$ un \Kev de dimension finie et $f$ un endomorphisme de $E$. Montrer
que
\[E=\im f \oplus \ker f  \quad\Longleftrightarrow\quad
 \im f = \im f^2.\]
Cette équivalence est-elle vraie en dimension infinie~?


 

\exercice{nom={Rang d'une somme d'applications linéaires}}
Soit $E$ et $F$ deux \Kev de dimension finie, $f$ et $g$ deux applications
linéaires de $E$ dans $F$.
\begin{questions}
\question Montrer que
  \[\abs{\rg f-\rg g} \leq \rg (f+g) \leq \rg f+\rg g.\]
\question On suppose dans cette question que $E=F$. Montrer que si
  $f\circ g=0$ et $f+g$ est inversible, alors $\rg f+\rg g=\dim E$.
\end{questions}

\exercice{nom={Dimension du noyau et composition}}
Soit $E$, $F$ et $G$ trois \Kevs de dimensions finies, $f\in\lin{E}{F}$ et
$g\in\lin{F}{G}$. En considérant la restriction de $f$ à
$\ker g\circ f$, montrer que
\[\dim\p{\ker g\circ f} \leq \dim\p{\ker g} + \dim\p{\ker f}.\]
\begin{sol}
Utiliser la théorème du rang avec~:
\[\dspappli{\phi}{\ker(g\circ f)}{F}{x}{f(x)}\]
\end{sol}

\exercice{nom={Endomorphisme $f$ tel que $f^2=0$}}
Soit $E$ un \Kev de dimension 4 et $f\in\Endo{E}$ tel que $f^2=0$. Montrer
que $\rg f\leq 2$. 
\begin{sol}
Utiliser le fait que $\im f\subset\ker f$. On utilise ensuite le théorème
du rang. 
\end{sol}

\exercice{nom={Factorisation}}
Soit $E$ un \Kev de dimension finie et $f,g\in\Endo{E}$. Montrer qu'il existe
$h\in\Endo{E}$ tel que $f=h\circ g$ si et seulement si $\ker g\subset \ker f$.

\exercice{nom={Noyaux et images itérés}}
Soit $E$ un \Kev de dimension finie et $f\in\Endo{E}$. On définit, pour tout
entier $n\in\N$
\[K_n\defeq\ker f^n \quad \text{et} \quad I_n\defeq\im f^n.\]
\begin{questions}
\question 
  \begin{questions}
  \question Montrer que la suite $\p{K_n}$ est croissante au sens de l'inclusion et
    que la suite $\p{I_n}$ est décroissante au sens de l'inclusion.
  \question 
    \begin{questions}
    \question Montrer qu'il existe $n\in\N$ tel que $K_{n}=K_{n+1}$. Dans la
      suite, on note $n_0$ le plus petit entier tel que $K_{n_0}=K_{n_0+1}$.
    \question Montrer que
      \[\forall n \geq n_0 \qsep K_n=K_{n_0}.\]
    \end{questions}
  \question 
    \begin{questions}
    \question Montrer qu'il existe $n\in\N$ tel que $I_{n}=I_{n+1}$. Dans la
      suite, on note $n_1$ le plus petit entier tel que $I_{n_1}=I_{n_1+1}$.
    \question Montrer que
      \[\forall n \geq n_1 \qsep I_n=I_{n_1}.\]
    \end{questions}
  \question Montrer que $n_0=n_1$. Dans la suite, on note $r$ cette valeur commune.
  \end{questions}
\question 
  \begin{questions}
  \question Montrer que $I_r \oplus K_r = E$.
  \question Montrer que $I_r$ et $K_r$ sont stables par $f$, puis que la
    restriction de $f$ à $K_r$ est nilpotente et que la restriction de $f$ à
    $I_r$ est un automorphisme.
  \end{questions}
\question 
  \begin{questions}
  \question Montrer que la suite de terme général $\dim I_{n} - \dim I_{n+1}$
    est décroissante. {\it On pourra pour cela considérer l'application
    $\phi$ de $I_n$ dans $I_{n+1}$ qui à $x$ associe $f(x)$.}
  \question Énoncer et démontrer un résultat semblable pour la suite $\p{K_n}$. 
  \end{questions}
\end{questions}
\magsubsection{Hyperplan}

\exercice{nom={Intersection d'hyperplans}}
Soit $E$ un \Kev de dimension finie $n\geq 2$ et $H_1$ et $H_2$ deux hyperplans
distincts de $E$. Calculer $\dim(H_1 \cap H_2)$.


\exercice{nom={Bidual}}
Soit $E$ un \Kev de dimension finie. On note $E^{\star\star}$ le bidual de $E$,
c'est-à-dire le dual du dual de $E$.
\begin{questions}
\question Montrer que $E$ et $E^{\star\star}$ sont isomorphes.
\question Soit $\phi$ l'application de $E$ dans $E^{\star\star}$ qui à $x$ associe
  l'application $\phi_x$ de $E^\star$ dans $\K$ qui à $\psi$ associe
  $\psi(x)$.
  \begin{questions}
  \question Montrer que pour tout $x\in E$, $\phi_x$ est une forme linéaire
    sur $E^\star$.
  \question Montrer que $\phi$ est un isomorphisme.
  \end{questions}
\end{questions}


\exercice{nom={Hyperplan de $\polyK$}}
Soit $\alpha\in\K$. Montrer que
\[H\defeq\enstq{P\in\polyK}{P(\alpha)=0}\]
est un hyperplan et en déterminer une base.


% \subsection{Endomorphisme nilpotent}
% Soit $E$ un \Kev de dimension $n$ et $f$ un endomorphisme nilpotent de $E$,
% c'est-à-dire tel qu'il existe un entier $m$ tel que $f^m=0$.
% \begin{questions}
% \question En notant $m_0$ le plus petit entier tel que $f^{m_0}=0$, montrer
%   qu'il existe $x\in E$ tel que la famille $x,f(x),\ldots,f^{m_0-1}(x)$ soit
%   libre.
% \question Que peut-on en déduire sur $m_0$ ?
% \end{questions}



% \subsection{Rang d'une famille de vecteurs}
% Déterminer le rang des familles suivantes :
% \begin{questions}
% \question $x_1=(1,0,1),x_2=(-1,-2,3),x_3=(-1,-1,1)$ dans $\R^3$.
% \question $x_1=(1,0,1,0),x_2=(1,1,0,0),x_3=(-1,-1,1,0),x_4=(0,0,2,0)$ dans
%   $\R^4$.
% \question $x_1=(3,2,1,0),x_2=(2,3,4,5),x_3=(0,1,2,3),x_4=(1,2,1,2)$ dans
%   $\R^4$.
% \question $\phi_1,\phi_2,\phi_3,\phi_4$ dans $\mathcal{L}(\R^4,\R)$ avec pour
%   $x,y,z,t\in\R$ :
%   \begin{itemize}
%   \item $\phi_1(x,y,z,t)=x+z$
%   \item $\phi_2(x,y,z,t)=-x+2y$
%   \item $\phi_3(x,y,z,t)=x+y-z+t$
%   \item $\phi_4(x,y,z,t)=y+t$
%   \end{itemize}
% \end{questions}






% \magsection{Calcul de dimension et de rang, hyperplan}








%END_BOOK
\end{document}



\end{document}
